\documentclass[11pt]{article}
\usepackage[margin=1.1in]{geometry}
\usepackage{amsmath,amssymb,amsthm,mathtools}
\usepackage{booktabs}
\usepackage[protrusion=true,expansion=false]{microtype}
\usepackage{xcolor}
\usepackage[colorlinks=true,linkcolor=blue!50!black,citecolor=blue!50!black,urlcolor=blue!50!black]{hyperref}
\usepackage{enumitem}

\newtheorem{theorem}{Theorem}[section]
\newtheorem{proposition}[theorem]{Proposition}
\newtheorem{lemma}[theorem]{Lemma}
\newtheorem{corollary}[theorem]{Corollary}
\newtheorem{conjecture}[theorem]{Conjecture}
\newtheorem{hypothesis}[theorem]{Hypothesis}
\theoremstyle{definition}
\newtheorem{definition}[theorem]{Definition}
\newtheorem{remark}[theorem]{Remark}
\newtheorem{observation}[theorem]{Observation}
\newtheorem{question}[theorem]{Question}
\newtheorem{problem}[theorem]{Problem}

\DeclareMathOperator{\per}{per}
\DeclareMathOperator{\sgn}{sgn}
\newcommand{\dtv}{d_{\mathrm{TV}}}
\newcommand{\N}{\mathbb{N}}
\newcommand{\R}{\mathbb{R}}
\newcommand{\E}{\mathbb{E}}
\renewcommand{\Pr}{\mathbb{P}}
\newcommand{\DD}{\mathcal{D}}
\newcommand{\PP}{P}
\newcommand{\QQ}{Q}
\newcommand{\CC}{C}
\newcommand{\spec}{\operatorname{spec}}

\title{The second row of a random Latin square:\\
towards a precise form of Cameron's conjecture\\[4pt]
\large Working notes}
\author{}
\date{August 28, 2026}

\begin{document}
\maketitle

\begin{abstract}
\noindent
Let $L$ be a uniformly random Latin square of order $n$ whose first row is the
identity, and let $\sigma$ be its second row, a uniformly ``Latin-tilted''
random derangement.  Cameron conjectured
[\emph{A niggling problem}, blog post, 24 Jan 2015] that the law of $\sigma$
tends \emph{very rapidly} to the uniform distribution on derangements.  These
notes make the conjecture quantitative and assemble exact, structural, and
numerical evidence for the following refined picture.  Write $N_n(\sigma)$ for
the number of completions of the $2\times n$ rectangle
$\binom{\mathrm{id}}{\sigma}$, which depends only on the cycle type
$\lambda$ of $\sigma$; set $\CC_n(\lambda)=N_n(\sigma)/(n-2)!$.
(i)~\textbf{Refined law (conjectural):} $\log \CC_n(\lambda)$ is, up to a
constant, \emph{additive} over the cycles of $\lambda$:
$\log \CC_n(\lambda)= c_0(n)+\sum_{\ell\ge 2}h_n(\ell)\,\lambda_\ell + \text{(small)}$,
with $h_n(2)= c_2\,n^{-3}(1+o(1))$, $c_2\approx 1.3$, and
$h_n(\ell)\asymp n^{1-2\ell}$; consequently
$\max_{\sigma,\tau}N_n(\sigma)/N_n(\tau)= 1+\Theta(n^{-2})$ and
$\dtv \asymp n^{-3}$ --- polynomially, not exponentially, fast.
(ii)~\textbf{Exact layer mechanism:} for the first extension row the analogue
of the refined law is a theorem.  With $M_n$ the m\'enage numbers, the number
of third rows compatible with $\binom{\mathrm{id}}{\sigma}$, $\sigma$ of type
$(\ell,n-\ell)$, equals $M_n+M_{n-2\ell}$ exactly --- an identity which turns
out to be an instance of the classical theory of discordant permutations of
Touchard and Riordan --- giving a per-$\ell$-cycle effect
$(1+o(1))\,n^{-2\ell}$; exact computations of the fourth row show the same
effect again, supporting the mechanism
``$n$ layers $\times$ $n^{-4}$ per layer $= n^{-3}$''.  A \emph{pair
inclusion--exclusion identity} reduces the fourth layer entirely to
punctured Touchard--Riordan counts and extends the layer table to $n=18$;
an \emph{exact switch calculus} (comparing adjacent cycle-type
representatives whose boards differ by one $2\times2$ switch, via
Wronskian identities for path matching polynomials) proves the
fourth-layer law: $r_4=r_3(1+O(1/n))$; the calculus transfers layer by
layer --- a triple inclusion--exclusion identity gives
$r_5=r_3(1+O(1/n))$ as well, with exact confirmation at $n\le9$ --- so
that each of the first three
extension layers contributes the same $(1+o(1))n^{-4}$ per $2$-cycle.  The same toolkit
yields as by-products exact joint intercalate statistics of $3\times n$
rectangles: the three pair counts are asymptotically independent
$\mathrm{Poisson}(1/2)$, with
$\E[X_{12}]-\tfrac12 = \tfrac{M_{n-4}}{2M_n}(1+O(n^{-2}))$ (the third row
tilts the second by exactly the layer coefficient) and pairwise covariances
$(\tfrac12+o(1))n^{-2}$.
(iii)~\textbf{An exact identity and a reduction:} the switching multigraph of
Cavenagh--Greenhill--Wanless yields an \emph{exact} identity expressing
adjacent completion ratios through a single structural probability: for a
split of a part $\alpha+\beta$ into $\alpha\neq\beta$,
$\CC_n(\lambda)/\CC_n(\mu) = \tfrac12\,(1+\bar q)$ where $\bar q\in[0,2]$ is
twice the probability that an explicit ``post-flip'' column pair is a
$B$-pair.  A row-swap involution forces the two constituent events of
$\bar q$ to have exactly equal probability.  A \emph{marked-pair
identity} then eliminates the joining measure and the post-flip events
altogether: $2\CC_n(\lambda)/\CC_n(\mu)-1$ equals the exact $B$:$A$
odds of a single marked column pair $\{j,\sigma^\alpha(j)\}$ in a
single uniform square of $S(\lambda)$ --- an identity with a short
self-contained proof, valid uniformly in $\alpha,\beta$ (including the
previously anomalous $\alpha=\beta$ case) and verified exhaustively
for $n\le7$.  Modulo an equidistribution
hypothesis for this single bit --- verified exhaustively for $n\le 6$,
exactly tabulated for $n\le 11$, and tested by Monte Carlo beyond ---
the weak (total-variation) form of Cameron's conjecture follows, and with it
the statement that almost all loops have multiplication group $S_n$.
A transfer programme for the Kwan--Petrova--Sawhney re-randomisation
machinery, whose quantifiers match the required $o(1/\log n)$ budget,
is laid out for this bit.
\end{abstract}

\tableofcontents

%======================================================================
\section{The conjecture and its context}\label{sec:intro}

\subsection{Setup and notation}

A \emph{Latin square} of order $n$ is an $n\times n$ array over the symbol set
$[n]=\{1,\dots,n\}$ in which every symbol occurs exactly once in each row and
each column.  Rows are permutations of $[n]$ (row $r$ maps $j\mapsto L(r,j)$
after identifying columns with symbols).  For two rows $r,s$ the permutation
$\sigma_{r,s}$ defined by $\sigma_{r,s}(L(r,j))=L(s,j)$ is a
\emph{derangement} of $[n]$.  We write $\DD_n$ for the set of derangements,
$D_n=|\DD_n|$.

Throughout, $\sigma$ denotes the second row of a uniformly random Latin
square whose first row is the identity; equivalently
$\sigma=\sigma_{1,2}$ under the convention row~1 $=\mathrm{id}$.  For a
derangement $\sigma$ let
\[
N_n(\sigma)\;=\;\#\bigl\{\text{Latin squares with rows } 1,2
 = (\mathrm{id},\sigma)\bigr\},
\qquad
\Pr(\text{second row}=\sigma)\;=\;\frac{N_n(\sigma)}{\sum_{\tau\in\DD_n}N_n(\tau)}.
\]
Cycle types of derangements are partitions
$\lambda=(2^{\lambda_2}3^{\lambda_3}\cdots)\vdash n$ with no part~$1$; we
write $\kappa(\lambda)=\sum_\ell \lambda_\ell$ for the number of parts and
$\gamma(\lambda)=n!/\prod_\ell \lambda_\ell!\,\ell^{\lambda_\ell}$ for the
number of derangements of type $\lambda$.

\begin{lemma}[type dependence and within-type uniformity]\label{lem:type}
$N_n(\sigma)$ depends only on the cycle type of $\sigma$.  Consequently,
conditionally on its cycle type, the second row is uniform on the
derangements of that type.
\end{lemma}

\begin{proof}
For $\pi\in S_n$, relabelling symbols by $\pi$ and then columns by $\pi^{-1}$
is a bijection of Latin squares sending rows
$(\mathrm{id},\sigma)\mapsto(\mathrm{id},\pi\sigma\pi^{-1})$.
\end{proof}

Following \cite{CGW08} we set, for $\lambda\vdash n$ with no part $1$,
\[
\CC_n(\lambda)\;=\;\frac{N_n(\sigma)}{(n-2)!}\qquad(\sigma\text{ of type
}\lambda),
\]
(the completions of a $2\times n$ rectangle fall into classes of size
$(n-2)!$ under permutation of the remaining rows), and define two probability
measures on cycle types:
\[
\PP_n(\lambda)=\frac{\gamma(\lambda)\CC_n(\lambda)}
     {\sum_{\mu}\gamma(\mu)\CC_n(\mu)},
\qquad
\QQ_n(\lambda)=\frac{\gamma(\lambda)}{D_n}.
\]
$\PP_n$ is the law of the cycle type of $\sigma$; $\QQ_n$ is the law of the
cycle type of a uniform derangement.  By Lemma~\ref{lem:type}, for any set
$S\subseteq \DD_n$ of derangements,
\begin{equation}\label{eq:setTV}
\bigl|\Pr(\sigma\in S)-|S|/D_n\bigr|\;\le\;\dtv(\PP_n,\QQ_n),
\end{equation}
so total variation at the level of cycle types controls total variation at
the level of derangements; in fact the two distances are \emph{equal},
because the two laws of $\sigma$ share the same conditional (uniform)
distribution within each type.

\subsection{The conjecture}

Cameron's blog post \cite{CamBlog15} states:

\begin{quote}\itshape
The probability distribution on derangements defined as above (where the
probability of a derangement is the probability that a random Latin square
with first row the identity has that derangement as its second row) tends
very rapidly to the uniform distribution as $n$ tends to infinity.
\end{quote}

The weak form was formalised by Cavenagh, Greenhill and Wanless as

\begin{conjecture}[{\cite[Conj.~6.1]{CGW08}}]\label{conj:cgw}
$\dtv(\PP_n,\QQ_n)=o(1)$ as $n\to\infty$.
\end{conjecture}

Even Conjecture~\ref{conj:cgw} is open.  The purpose of these notes is
(a)~to propose a sharp quantitative form of the full conjecture, based on new
exact and numerical analysis; (b)~to isolate exact structure (a
``layer'' theorem for the third row, and an exact switching identity) that
explains the observed rates; and (c)~to give a reduction of
Conjecture~\ref{conj:cgw} to a single equidistribution hypothesis about an
explicit local event, together with extensive verification.

\subsection{Known results}\label{sec:known}

We use the following throughout; $[n]$-asymptotics are as $n\to\infty$.

\begin{itemize}[leftmargin=2em]
\item \textbf{CGW switching bounds \cite{CGW08}.}  For partitions
$\lambda,\mu$ adjacent in the split/join graph (i.e.\ $\mu$ refines
$\lambda$ by splitting one part in two),
$\tfrac12\gamma(\lambda)/\gamma(\mu)\le |S(\lambda)|/|S(\mu)| \le
\tfrac32\gamma(\lambda)/\gamma(\mu)$, where $S(\lambda)$ is the set of Latin
squares whose first two rows have row-cycle type $\lambda$.  Iterating:
$\PP_n(\lambda)$ agrees with $\QQ_n(\lambda)$ within a factor
exponential in $\kappa(\lambda)$, whence, e.g., almost all cycle types have
$\PP_n/\QQ_n\in[n^{-3/2},n^{3/2}]$, the single-cycle probability is
$2n^{-2}\le \PP_n((n))\le 2n^{-2/3}$, and
$\Pr(\sigma\text{ even})\in[\tfrac14-o(1),\tfrac34+o(1)]$.
They also computed $\CC_n$ \emph{exactly} for $n\le 11$
\cite[Tables 1--2]{CGW08}; those tables are the data source for
\S\ref{sec:refined}.
\item \textbf{Intercalates.}  An \emph{intercalate} is a $2\times2$ Latin
subsquare; the number of intercalates inside rows $1,2$ equals
$\lambda_2(\sigma)$ exactly.  McKay--Wanless \cite{McKW99} and, much more
precisely, Kwan--Sah--Sawhney--Simkin \cite{KSSS22} determined the global
intercalate count: whp $\mathbf N=(1+o(1))n^2/4$, with large-deviation rates
$\exp(-\Theta(n^2))$ (lower tail) and $\exp(-\Theta(n^{4/3}\log n))$ (upper
tail).
\item \textbf{Parities.}  Kwan--Petrova--Sawhney \cite{KPS25} proved that the
vector of row parities of a random Latin square converges in total variation
to $n$ independent fair coins (resolving a different conjecture of Cameron),
introducing \emph{stable intercalate switchings} and a quantitative
comparison with the triangle-removal process.  Their results do not imply
Conjecture~\ref{conj:cgw}, but their machinery is discussed in
\S\ref{sec:kps}.
\item \textbf{Loops motivation.}  Cameron \cite{Cam92} proved that almost all
quasigroups have multiplication group $S_n$ or $A_n$;
H\"aggkvist--Janssen \cite{HJ96} ruled out $A_n$ (the proportion of all-even
Latin squares is $O(c^n)$, $c<1$); {\L}uczak--Pyber \cite{LP93} proved that
almost no permutation lies in a transitive subgroup other than $S_n,A_n$.
Conjecture~\ref{conj:cgw} would transfer this to \emph{loops} (reduced Latin
squares) via \eqref{eq:setTV}: almost all loops would have multiplication
group $S_n$, hence trivial character theory
\cite[Lemma 6.2]{CGW08}.
\end{itemize}

\begin{remark}[reduced squares: the conditioning is invisible]
\label{rem:reduced}
A loop on $[n]$ with identity $1$ is a \emph{reduced} Latin square (first
row \emph{and} first column the identity), so its second row satisfies the
extra constraint $\sigma(1)=2$, whereas $\PP_n$ is defined from squares
normalised in the first row only.  This costs nothing: sorting the rows of
a first-row-normalised square by their first-column entries is an
$(n-1)!$-to-one map onto reduced squares which is exchangeable in the rows
$2,\dots,n$, so the second row of a uniform loop is distributed as
$\sigma$ \emph{conditioned on} $\{\sigma(1)=2\}$.  By
Lemma~\ref{lem:type} both $\PP_n$ and $\QQ_n$ are uniform on each
derangement class, and on any fixed class $\sigma(1)$ is uniform on
$\{2,\dots,n\}$ (conjugate by permutations fixing $1$); hence
$\Pr(\sigma(1)=2\mid\lambda)=\tfrac1{n-1}$ independently of the type,
the cycle-type law of the loop's second row is \emph{exactly} $\PP_n$,
and the total-variation distance between the conditioned laws equals
$\dtv(\PP_n,\QQ_n)$ exactly.  (Verified by complete enumeration at
$n=5,6$: the three laws --- second row, loop second row, conditioned
second row --- coincide as exact rationals.)  All statements in these
notes therefore apply verbatim to loops.
\end{remark}

%======================================================================
\section{The refined picture: exact data analysis for \texorpdfstring{$n\le11$}{n<=11}}\label{sec:refined}

All statements in this section are \emph{exact} computations from the CGW
tables of $\CC_n(\lambda)$, $4\le n\le 11$, re-verified in two independent
ways: the published fractions $\PP_n(\lambda)$ are reproduced exactly, and
$\sum_\lambda \gamma(\lambda)\CC_n(\lambda)\,(n-2)!\,n! = L_n$ reproduces the
known numbers of Latin squares up to $n=11$ \cite{McKW05}.

\subsection{Additivity over cycles}\label{sec:additivity}

The completion counts are astonishingly rigid.  Along the chain of types
$(n)\to(2,n-2)\to(2,2,n-4)\to\cdots$, each added $2$-cycle multiplies
$\CC_n$ by almost exactly the same factor:

\begin{center}
\small
\begin{tabular}{llll}
\toprule
$n$ & step & $\CC_n(\mu')/\CC_n(\mu)-1$ & $\times\,n^3$ \\
\midrule
10 & $(10)\to(2,8)$          & $+1.320265\times10^{-3}$ & $1.3203$\\
10 & $(2,8)\to(2^2,6)$       & $+1.321183\times10^{-3}$ & $1.3212$\\
10 & $(2^2,6)\to(2^3,4)$     & $+1.323079\times10^{-3}$ & $1.3231$\\
10 & $(2^3,4)\to(2^5)$       & $+2.647253\times10^{-3}$ & $2.6473$\\
11 & $(11)\to(2,9)$          & $+9.609945\times10^{-4}$ & $1.2791$\\
11 & $(2,9)\to(2^2,7)$       & $+9.620287\times10^{-4}$ & $1.2805$\\
11 & $(2^2,7)\to(2^3,5)$     & $+9.610360\times10^{-4}$ & $1.2791$\\
11 & $(2^3,5)\to(2^4,3)$     & $+9.750180\times10^{-4}$ & $1.2977$\\
\bottomrule
\end{tabular}
\end{center}

\noindent
The one apparent exception, the last step at $n=10$, changes the type by
\emph{two} $2$-cycles ($\lambda_2\colon 3\to5$, removing a $4$-cycle), and its
increment is exactly twice the others; the last step at $n=11$ adds a
$3$-cycle as well, and the excess $+1.4\times10^{-5}$ matches the fitted
$3$-cycle coefficient below.  This motivates:

\begin{observation}[additive fit]\label{obs:fit}
Fit $\log \CC_n(\lambda) = c_0(n)+\sum_{\ell\ge2} h_n(\ell)\lambda_\ell$ by
least squares (the gauge freedom $h(\ell)\mapsto h(\ell)+t\ell$ is fixed by
$h_n(n)=0$).  Then, using all types at each order:
\[
\begin{array}{l|lll|l}
n & h_n(2) & h_n(3) & h_n(4) & \max_\lambda|\mathrm{resid}| \\\hline
10 & 1.323\times10^{-3} & 5.12\times10^{-5} & -6.9\times10^{-6} & 3.2\times10^{-5}\\
11 & 9.605\times10^{-4} & 1.544\times10^{-5} & 1.2\times10^{-6} & 1.8\times10^{-6}
\end{array}
\]
At $n=11$ the additive model reproduces all fourteen $\log\CC_{11}(\lambda)$
values to within $1.8\times10^{-6}$, while the leading coefficient is
$9.6\times10^{-4}$: additivity holds to relative precision
$\approx 0.2\%$, and improves rapidly with $n$.
\end{observation}

There is \emph{no} detectable dependence on the sign of $\sigma$: adding a
$\pm\sgn(\sigma)$ term to the fit yields a coefficient
($4.5\times10^{-7}$ at $n=11$) below the residual noise.  Consistently, the
parity of $\sigma$ is explained entirely by the derangement count
disparity (see \S\ref{sec:consequences}).

\subsection{Scaling of the coefficients}

The fitted $h_n(2)$, and two independent estimators of it --- the chain
increments above, and the tilt of the expected number of rows-$(1,2)$
intercalates, $2n^3(\E_{\PP_n}[\lambda_2]-\tfrac12)$ --- agree and scale as
$n^{-3}$:
\[
\begin{array}{l|llllll}
n & 6 & 7 & 8 & 9 & 10 & 11\\\hline
n^3\,h_n(2)\text{-proxy (chain)} & 10.3 & 2.01 & 1.71 & 1.16 & 1.32 & 1.28\\
2n^3(\E_{\PP_n}\lambda_2-\tfrac12) & 22.0 & 0.93 & 2.55 & 1.06 & 1.33 & 1.28
\end{array}
\]
Odd and even $n$ approach the limit from different sides (as CGW already
noticed for their coefficient of variation), and $n\le11$ cannot pin the
limit precisely; the two largest orders give
\begin{equation}\label{eq:c2}
c_2 \;:=\; \lim_{n\to\infty} n^{3}h_n(2)\;\approx\;1.30\pm 0.05
\qquad(\text{numerically; existence of the limit is conjectural}).
\end{equation}
The $3$-cycle coefficient is much smaller
($n^4 h_n(3)=0.69,\,0.51,\,0.23$ at $n=9,10,11$), consistent with
$h_n(3)=O(n^{-4})$ and, in view of the layer analysis of
\S\ref{sec:layers}, conjecturally $\asymp n^{-5}$; coefficients for
$\ell\ge4$ are below the resolution of the $n\le 11$ data.

%======================================================================
\section{Precise forms of the conjecture}\label{sec:conjectures}

We can now state the proposed quantitative hierarchy.  Recall
$c_2\approx1.30$ from \eqref{eq:c2}.

\begin{conjecture}[refined law]\label{conj:refined}
There are quantities $h_n(\ell)$, $2\le\ell\le n$, with
\[
\log \CC_n(\lambda)\;=\;c_0(n)\;+\;\sum_{\ell\ge2}h_n(\ell)\,\lambda_\ell
\;+\;O\!\bigl(n^{-4}\bigr)
\quad\text{uniformly in }\lambda,
\]
such that $h_n(2)=c_2n^{-3}(1+o(1))$ for a constant $c_2>0$ (numerically
$c_2\approx1.30$), and, for each fixed $\ell\ge3$, $h_n(\ell)=O(n^{-4})$
(conjecturally $h_n(\ell)= c_\ell\,n^{1-2\ell}(1+o(1))$).
\end{conjecture}

\begin{conjecture}[rates]\label{conj:rates}
As $n\to\infty$:
\begin{enumerate}[label=(\alph*),leftmargin=2em]
\item $\displaystyle
\max_{\sigma,\tau\in\DD_n}\frac{N_n(\sigma)}{N_n(\tau)}
=1+\frac{c_2}{2}\,n^{-2}(1+o(1))$, the extremes being attained at
$\lambda=(2^{\lfloor n/2\rfloor}\cdots)$ and $\lambda=(n)$;
\item $\displaystyle
\dtv(\PP_n,\QQ_n)\;=\;\frac{c_2\,e^{-1/2}}{2}\,n^{-3}(1+o(1))
\;\approx\;0.394\,n^{-3}$,
and the same for the total variation distance between the law of the second
row and the uniform distribution on $\DD_n$.
\end{enumerate}
\end{conjecture}

Part (b) follows from Conjecture~\ref{conj:refined} by a direct computation:
to leading order $\PP_n(\lambda)\propto \QQ_n(\lambda)e^{h_n(2)\lambda_2}$,
$\lambda_2$ is asymptotically Poisson$(1/2)$ under $\QQ_n$, and
$\E\,|X-\E X| = e^{-1/2}$ for $X\sim\mathrm{Poisson}(1/2)$.
The exact $n\le11$ data match both predictions strikingly well:
\[
\begin{array}{l|llll}
n & 8 & 9 & 10 & 11\\\hline
n^3\,\dtv(\PP_n,\QQ_n) & 0.62 & 0.34 & 0.40 & 0.39\\
n^2\,(\max/\min-1) & 1.65 & 0.37 & 0.66 & 0.47
\end{array}
\qquad
\begin{array}{l}
\text{predicted limits:}\\
0.394\ \text{and}\ c_2/2\approx0.65\ (n\ \text{even}),\\
\text{smaller along odd $n$ (no type } (2^{n/2})).
\end{array}
\]

\begin{remark}[``very rapidly'']
The convergence is thus \emph{polynomial}, with small constants --- not
exponential.  At $n=7$ the extreme completion counts already differ by only
$1.4\%$ (Cameron's observation), yet the ratio $1+\Theta(n^{-2})$ decays
slowly thereafter: uniformity to $0.1\%$ in sup-ratio requires $n\gtrsim36$.
For every \emph{fixed} set $S\subseteq\DD_n$ of derangements, however,
$\Pr(\sigma\in S)=|S|/D_n+O(n^{-3})$ --- rapid indeed.
\end{remark}

\subsection{Consequences and cross-checks}\label{sec:consequences}

\paragraph{Parity.}  The number of even minus odd derangements of $[n]$ is
$(-1)^{n-1}(n-1)$, so under $\QQ_n$,
$\Pr(\text{even})-\tfrac12=(-1)^{n-1}(n-1)/(2D_n)$, of order $n^{-1}e\,/(n-1)!$.
If Conjecture~\ref{conj:rates}(b) holds \emph{at the level of the sign
functional} with its natural precision, then the same asymptotics should hold
under $\PP_n$.  The exact data agree:
\[
\begin{array}{l|lllll}
n & 7 & 8 & 9 & 10 & 11\\\hline
\Pr_{\PP_n}(\mathrm{even})-\tfrac12 & +1.813\cdot10^{-3} & -1.998\cdot10^{-4}
& +2.952\cdot10^{-5} & -3.168\cdot10^{-6} & +3.250\cdot10^{-7}\\
(-1)^{n-1}\frac{n-1}{2D_n} & +1.618\cdot10^{-3} & -2.360\cdot10^{-4}
& +2.996\cdot10^{-5} & -3.371\cdot10^{-6} & +3.405\cdot10^{-7}\\
\text{ratio} & 1.12 & 0.85 & 0.985 & 0.94 & 0.95
\end{array}
\]
This suggests the sharp parity conjecture
$\Pr(\sigma\ \text{even})-\tfrac12\sim(-1)^{n-1}(n-1)/(2D_n)$ --- a
statement finer, for this particular functional, than
Conjecture~\ref{conj:rates}(b), since the parity imbalance is
super-exponentially small while $\dtv\asymp n^{-3}$.  (It is the
$\sigma_{1,2}$-analogue of the row-parity phenomena studied in \cite{KPS25}.)

\paragraph{Intercalates.}  Since the intercalates inside rows $(1,2)$ number
exactly $\lambda_2$, symmetry over row pairs gives
$\E[\mathbf N]=\binom n2\,\E_{\PP_n}[\lambda_2]$ for the total intercalate
count $\mathbf N$.  Under $\QQ_n$, $\E[\lambda_2]=\binom
n2D_{n-2}/D_n=\tfrac12+O(1/(n-2)!)$.  Conjecture~\ref{conj:refined} thus
predicts a third-order term in the intercalate expectation:
\[
\E[\mathbf N]\;=\;\frac{n^2}{4}-\frac n4+\frac{c_2+o(1)}{4n}.
\]
The known asymptotics $\E[\mathbf N]\sim n^2/4$ \cite{KSSS22} match the
leading term; the finer terms appear to be new and give an independent
interpretation of $c_2$.

\paragraph{Loops.}  By \eqref{eq:setTV} and \S\ref{sec:known}, even the weak
Conjecture~\ref{conj:cgw} implies that almost all loops have multiplication
group $S_n$ and trivial character theory; Conjecture~\ref{conj:rates}(b)
would give an $O(n^{-3})$ (plus {\L}uczak--Pyber) error term.

%======================================================================
\section{The third-row layer: exact results}\label{sec:thirdrow}

Why $n^{-3}$?  The mechanism becomes visible if one builds the Latin square
row by row.  Write
\[
N_n(\sigma)\;=\;\prod_{k=2}^{n-1}\frac{N^{(k+1)}_n(\sigma)}{N^{(k)}_n(\sigma)},
\qquad
N^{(k)}_n(\sigma)=\#\{k\times n \text{ Latin rectangles with rows
}1,2=(\mathrm{id},\sigma)\},
\]
so that $\log N_n(\sigma)$ is a sum of $n-2$ \emph{layer contributions}.  The
first layer is exactly solvable.

For $k=2$, $N^{(3)}_n(\sigma)=\per(J-I-P_\sigma)$ counts the permutations
\emph{discordant} with both $\mathrm{id}$ and $\sigma$.  The forbidden
bipartite graph $I\cup P_\sigma$ is a disjoint union of even cycles, one
cycle $C_{2\ell}$ per $\ell$-cycle of $\sigma$; by inclusion--exclusion
\begin{equation}\label{eq:IE}
\per(J-B)\;=\;\sum_{k\ge0}(-1)^k m_k(B)\,(n-k)!,
\end{equation}
where $m_k(B)$ is the number of $k$-matchings of $B$, and the matching
generating polynomial $M_B(x)=\sum_k m_k(B)x^k$ factorises over the cycles,
with
\[
M_{C_m}(x)=\alpha_+^m+\alpha_-^m,\qquad
\alpha_\pm=\tfrac12\bigl(1\pm\sqrt{1+4x}\bigr),\qquad
\alpha_+\alpha_-=-x .
\]

\begin{lemma}[fusion; Touchard--Riordan]\label{lem:fusion}
For $a\ge b\ge1$,
\(
M_{C_{2a}}(x)\,M_{C_{2b}}(x)=M_{C_{2(a+b)}}(x)+x^{2b}M_{C_{2(a-b)}}(x),
\)
with the convention $M_{C_0}\equiv2$.
\end{lemma}

\begin{proof}
Expand both products in $\alpha_\pm$:
$(\alpha_+^{2a}+\alpha_-^{2a})(\alpha_+^{2b}+\alpha_-^{2b})
=\alpha_+^{2(a+b)}+\alpha_-^{2(a+b)}
+(\alpha_+\alpha_-)^{2b}(\alpha_+^{2(a-b)}+\alpha_-^{2(a-b)})$,
and $(\alpha_+\alpha_-)^{2b}=x^{2b}$.
\end{proof}

This is classical: in Riordan's normalisation of \emph{associated} m\'enage
polynomials $m_k(x)=x^kM_{C_{2k}}(-x^{-1})$ one has the Chebyshev form
$m_k(x)=2\,T_{2k}(\sqrt x/2)$ and the factor-free product rule
$m_j\,m_k=m_{j+k}+m_{j-k}$
\cite[Ch.~8, eqs.~(24)--(25)]{Riordan58}.

Let $M_n$ denote the m\'enage numbers
$0,1,2,13,80,579,4738,43387,439792,4890741,\dots$ ($n=2,3,\dots$), i.e.\
$M_n=\per(J-I-P_{c})$ for an $n$-cycle $c$; classical asymptotics
(e.g.\ via \eqref{eq:IE}) give $M_n=n!\,e^{-2}\bigl(1-\tfrac1n+O(n^{-2})\bigr)$.

\begin{theorem}[third row, exact; essentially Touchard--Riordan]\label{thm:thirdrow}
Let $2\le\ell\le n/2$ and let $\sigma$ have cycle type $(\ell,\,n-\ell)$.
Then
\[
N^{(3)}_n(\sigma)\;=\;M_n+M_{n-2\ell}
\qquad(\text{with }M_0:=2\text{ when }\ell=n/2).
\]
More generally, for any type $\lambda\cup\{a,b\}$ with $a\ge b$, $a-b\ge2$,
\[
N^{(3)}_n(\lambda\cup\{a,b\})
=N^{(3)}_n(\lambda\cup\{a+b\})
+N^{(3)}_{n-2b}(\lambda\cup\{a-b\}),
\]
and when $a=b$ the last term is replaced by
$2\,N^{(3)}_{n-2b}(\lambda)$.
\end{theorem}

\begin{proof}
Insert Lemma~\ref{lem:fusion} into \eqref{eq:IE}: the factor $x^{2b}$ shifts
the matching index by $2b$, turning $(n-k)!$ into $((n-2b)-k')!$, which
reproduces the inclusion--exclusion sum \eqref{eq:IE} on the ground set of
size $n-2b$ for the reduced type.  (Verified by exhaustive computation for
all cycle types for $n\le25$.)
\end{proof}

\begin{corollary}[per-cycle effect at the third row]\label{cor:thirdrow}
For fixed $\ell\ge2$,
\[
\rho_\ell(n):=\frac{N^{(3)}_n((\ell,n-\ell))}{N^{(3)}_n((n))}-1
=\frac{M_{n-2\ell}}{M_n}
=\frac{1}{(n)_{2\ell}}\Bigl(1-\frac{2\ell}{n^2}+O(n^{-3})\Bigr)
=n^{-2\ell}\Bigl(1+\binom{2\ell}{2}\frac1n+O(n^{-2})\Bigr).
\]
\end{corollary}

Numerically (exact rational computations up to $n=400$, Richardson
extrapolation): $n^4\rho_2 = 1+6/n+21/n^2+42/n^3+\cdots$ and
$n^6\rho_3=1+15/n+134/n^2+\cdots$, matching the corollary
($\binom42=6$, $\binom62=15$; the third coefficients pick up the
$-2\ell/n^2$ correction from the m\'enage asymptotics).

Moreover the \emph{additivity} of \S\ref{sec:additivity} is already present,
and provable, at this layer: iterating Theorem~\ref{thm:thirdrow} shows that
\[
\log N^{(3)}_n(\lambda)=\log M_n+\sum_{\ell\ge2}\lambda_\ell\,
\frac{M_{n-2\ell}}{M_n}+(\text{pair interactions of smaller order}),
\]
and an exact least-squares check at $n=30$ (all $1039$ cycle types) gives a
maximal additivity residual $1.1\times10^{-10}$ against a leading coefficient
$h^{(3)}(2)=M_{26}/M_{30}=1.5\times10^{-6}$.

\begin{remark}[attribution]\label{rem:attribution}
Permutations discordant with two given permutations are a classical topic
going back to Touchard \cite{Touchard53} and Riordan \cite{Riordan54}, and
Chapter~8, \S3 of Riordan's book \cite{Riordan58} contains everything needed
for Theorem~\ref{thm:thirdrow}: the rook polynomial of the cycle class
$(k)=(1^{k_1}2^{k_2}\cdots)$ factorises over cycles (his eq.~(22)), the
associated polynomials satisfy the Chebyshev product rule (eqs.~(24)--(25)),
and his Theorem~2 reduces the hit polynomial of any class to a linear
combination $\sum_jA_j(1-t)^jU_{n-j}(t)$ of m\'enage hit polynomials.
Our identity $N^{(3)}_n((\ell,n-\ell))=M_n+M_{n-2\ell}$ is the evaluation at
$t=0$ of the case $m_\ell m_{n-\ell}=m_n+m_{n-2\ell}$, and the general
recursion of Theorem~\ref{thm:thirdrow} follows from eq.~(26) there; Riordan
attributes the substance to Touchard \cite{Touchard53}.  What we add in this
section is only the quantitative reading relevant to Cameron's conjecture:
the per-$\ell$-cycle effect $\rho_\ell(n)=M_{n-2\ell}/M_n\sim n^{-2\ell}$
with its full expansion, and the additivity of $\log N^{(3)}$ over cycles.
Section 8.4 of \cite{Riordan58} carries the classical theory one layer
further in aggregate form --- the three-line Latin rectangle count
$K(3,n)=\sum_k\binom nkD_{n-k}D_kU_{n-2k}$ --- and the ``$k$ discordant
permutations'' literature (Moser's very reduced $4\times n$ rectangles
\cite{Moser67}; Whitehead \cite{Whitehead79}) is the natural classical
toolkit for the layer-$4$ analogue raised in \S\ref{sec:open}.
\end{remark}

%======================================================================
\section{Layer accumulation: the \texorpdfstring{$n^{-3}$}{n\^{}-3} mechanism}\label{sec:layers}

Corollary~\ref{cor:thirdrow} gives a per-$2$-cycle effect $\asymp n^{-4}$ at
the first layer.  The full effect is $h_n(2)\asymp n^{-3}$ --- precisely $n$
times larger.  This suggests that \emph{every} layer contributes
$\asymp n^{-4}$, the influence of the fixed first two rows persisting
(neither amplified nor washed out) as the rectangle grows.  The fourth layer
can still be computed exactly (summing $\per(J-I-P_\sigma-P_\tau)$ over all
third rows $\tau$; Ryser--Gray-code permanents):

\begin{center}
\small
\begin{tabular}{c|cc|cc}
\toprule
 & \multicolumn{2}{c|}{$2$-cycle effect} & \multicolumn{2}{c}{$3$-cycle effect}\\
$n$ & layer 3 ($=M_{n-4}/M_n$) & layer 4 & layer 3 ($=M_{n-6}/M_n$) & layer 4\\
\midrule
8  & $4.221\times10^{-4}$ & $5.015\times10^{-4}$ & $0$ & $-7.0\times10^{-5}$\\
9  & $2.996\times10^{-4}$ & $3.134\times10^{-4}$ & $2.305\times10^{-5}$ & $1.974\times10^{-5}$\\
10 & $1.819\times10^{-4}$ & $1.872\times10^{-4}$ & $4.548\times10^{-6}$ & $4.093\times10^{-6}$\\
11 & $1.184\times10^{-4}$ & $1.205\times10^{-4}$ & $2.658\times10^{-6}$ & $2.469\times10^{-6}$\\
\bottomrule
\end{tabular}
\end{center}

\noindent
(The layer-$k$ effect is
$\log\bigl[L_k((\ell,n-\ell))/L_k((n))\bigr]$ up to lower order, with
$L_k(\lambda)=N^{(k+1)}/N^{(k)}$; at $n=9$ the exact $3$-cycle layer-3 value
is $1/M_9=1/43387$, a pleasing instance of Theorem~\ref{thm:thirdrow}.)

The layer-4 effect equals the layer-3 effect within $1\text{--}5\%$ at
$n=10,11$; the layer-5 effect (now exactly computed for $n\le9$,
\S\ref{sec:fifth}) equals the layer-4 effect within $3\%$ at $n=9$.
Summing the exactly known layers 3, 4 and comparing with the
exactly known total at $n=11$: layers $3+4$ contribute
$2.39\times10^{-4}$ of $h_{11}(2)=9.61\times10^{-4}$; the remaining seven
layers average $1.03\times10^{-4}\approx(1.51)\,n^{-4}$ --- the same order,
declining towards the final (forced) layers.  Writing the conjectural layer
profile as $t_k(n)= \gamma(k/n)\,n^{-4}(1+o(1))$ with $\gamma(0)=1$ (by
Corollary~\ref{cor:thirdrow}), the constant in \eqref{eq:c2} is
\[
c_2=\int_0^1 \gamma , \qquad \gamma(0)=1 .
\]
Determining $\gamma$ --- e.g.\ via permanent asymptotics for complements of
random $k\times n$ rectangles containing a planted intercalate --- is, we
believe, the most promising route to the sharp constant.  The same mechanism
predicts $h_n(\ell)\asymp n\cdot n^{-2\ell}=n^{1-2\ell}$ for each fixed
$\ell\ge3$, which is consistent with (though not yet demanded by) the
$n\le11$ data.

\begin{remark}[a finite-size caveat on $c_2$]\label{rem:c2caveat}
The estimate $c_2\approx1.30$ in \eqref{eq:c2} is the value of $n^3h_n(2)$
at $n=10,11$.  The layer analysis shows that this quantity carries a
substantial finite-size inflation: the early-layer weights track
$1/(n)_4=n^{-4}(1+6/n+\cdots)$, whose $n\le11$ values exceed $n^{-4}$ by
$50\text{--}70\%$.  If, say, $\gamma\equiv1$, then
$n^3h_n(2)\approx (n-2)\,n^3/(n)_4$ would still be $\approx1.5$ at $n=11$
while converging to $c_2=1$.  The exact totals at $n=11$ (layer sum $14.1\,
n^{-4}$ against $9\times1.73\,n^{-4}=15.6\,n^{-4}$ for nine constant-weight
layers) show the late layers contributing $\approx0.9$ of the early-layer
weight on average at that size.  We therefore regard the \emph{limit} $c_2$
as substantially uncertain --- plausibly anywhere in $[0.8,1.5]$ --- while
$n^3h_n(2)\approx1.3$ remains an accurate description of the exact data at
$n=10,11$ (and, through Conjecture~\ref{conj:rates}(b), of
$\dtv$ at those orders).  Determining $\gamma$ in the bulk is the key open
computation; \S\ref{sec:pair} settles its behaviour at the fourth layer.
\end{remark}

%======================================================================
\section{A pair calculus for the fourth row}\label{sec:pair}

The fourth layer is governed by
$N^{(4)}_n(\sigma)=\sum_\tau \per(J-B_2-P_\tau)$, a sum over third rows of
permanents of complements of \emph{three}-regular boards, for which no
product formula is available.  The following identity nevertheless reduces
everything to the Touchard--Riordan calculus, extended to punctured
(path-and-cycle) boards.

For a partial matching $M$ of the complement of $B_2=I\cup P_\sigma$ (a set
of cells, no two in a row or column, all avoiding $B_2$), let
\[
R(M;\sigma)\;=\;\#\{\tau\in S_n:\ \tau\ \text{avoids}\ B_2,\ \tau\supseteq M\}.
\]
Deleting the rows and columns of $M$ cuts the cycles of $B_2$ into disjoint
paths and cycles; with $m_k$ the matching numbers of the resulting board and
$N=n-|M|$,
\begin{equation}\label{eq:punctured}
R(M;\sigma)\;=\;\sum_k(-1)^k m_k\,(N-k)! ,
\end{equation}
computable from the matching polynomials of paths
($m_k(P_e)=\binom{e-k+1}{k}$, the Chebyshev partner of \eqref{eq:IE}) and
cycles.  We verified \eqref{eq:punctured} against exhaustive enumeration for
all tested puncture sets at $n=6,7$.

\begin{proposition}[pair inclusion--exclusion]\label{prop:pairIE}
\[
N^{(4)}_n(\sigma)\;=\;\sum_{j\ge0}(-1)^j\!\!\sum_{M\in\mathcal M_j}
R(M;\sigma)^2 ,
\]
where $\mathcal M_j$ is the set of $j$-cell partial matchings of the
complement of $B_2$.
\end{proposition}

\begin{proof}
$N^{(4)}$ counts ordered pairs $(\tau,\eta)$ of $B_2$-avoiding permutations
with no common cell.  Inclusion--exclusion over the common-cell set (which
is automatically a partial matching) gives the sum; for fixed $M$, the pairs
with $\tau,\eta\supseteq M$ number $R(M;\sigma)^2$ by independence.
(Verified exactly at $n=6,7$: e.g.\ $N^{(4)}_7((7))=83600$ both ways.)
\end{proof}

Probabilistically: if $\tau,\eta$ are \emph{independent} uniform
$B_2$-avoiding rows and $X=|\tau\cap\eta|$ their number of common cells,
then with $t_j=T_j/T_0$, $T_j=\sum_{M\in\mathcal M_j}R(M)^2$:
\[
\Phi(\sigma):=\frac{N^{(4)}_n(\sigma)}{N^{(3)}_n(\sigma)^2}
=\Pr(X=0),\qquad t_j=\E\binom Xj,
\qquad
\underbrace{r_4-r_3}_{\text{layer-4 minus layer-3 effect}}=\Delta_\lambda\log\Pr(X=0),
\]
where $\Delta_\lambda$ denotes the difference between $\sigma$ of types
$(2,n-2)$ and $(n)$.  The layer-4 effect thus equals the layer-3 effect
\emph{plus the cycle-type sensitivity of a pairwise-clash probability} ---
and the latter is very small:

\begin{observation}[near-uniformity of cell probabilities]\label{obs:cells}
Writing $p_c=R(\{c\};\sigma)/N^{(3)}(\sigma)$ for the probability that a
random $B_2$-avoiding row passes through the admissible cell $c$
(so $\sum_cp_c=n$ over the $n(n-2)$ admissible cells),
\[
\mu:=\E X=t_1=\sum_cp_c^2\;=\;\frac n{n-2}\;+\;n(n-2)\,
\operatorname{Var}_c(p_c)\;\ge\;\frac n{n-2},
\]
and exact computation shows the variance term to be astonishingly small:
$t_1-\tfrac n{n-2}\approx 2.1\times10^{-7}$ at $n=12$ (type $(12)$),
decaying like $n^{-6\pm}$ and depending on the cycle type only in its last
digits.  Equivalently, the cell probabilities are uniform to relative
$O(n^{-3})$: at $n=20$ (type $(20)$) the exact profile is flat to
$10^{-8}$ across all cells except at $\omega$-distance $2$ (and its mirror
$n-1$), where a single-cell deficit $\approx-1.4\,n^{-3}$ appears: the value
$\tau(i)=\sigma^2(i)$ is slightly disfavoured.  (Both facts are provable by
path fusion; we record them as computed.)  This is another appearance of the
$n^{-3}$ scale of the master conjecture.
\end{observation}

\subsection{The fourth layer, extended}

Truncations $S_J=\sum_{j\le J}(-1)^jT_j$ of Proposition~\ref{prop:pairIE}
give layer-4 effects with rapidly cancelling errors: against the exact
values (available for $n\le11$ from \S\ref{sec:layers}), the $J=2$ and
$J=3$ estimates bracket the truth within a few percent and their average is
accurate to $\approx0.4\%$.  This extends the layer-4 table:

\begin{center}
\small
\begin{tabular}{c c c c c}
\toprule
$n$ & $r_3=M_{n-4}/M_n$ & $r_4$ & $r_4\cdot(n)_4$ & $r_4/r_3$\\
\midrule
8  & $4.2212\times10^{-4}$ & $5.0153\times10^{-4}$ & $0.843$ & $1.188$\\
9  & $2.9963\times10^{-4}$ & $3.1340\times10^{-4}$ & $0.948$ & $1.046$\\
10 & $1.8190\times10^{-4}$ & $1.8715\times10^{-4}$ & $0.943$ & $1.029$\\
11 & $1.1839\times10^{-4}$ & $1.2049\times10^{-4}$ & $0.954$ & $1.018$\\
12 & $8.0011\times10^{-5}$ & $8.074\times10^{-5}\,^*$ & $0.959$ & $1.009$\\
13 & $5.5940\times10^{-5}$ & $5.629\times10^{-5}\,^*$ & $0.966$ & $1.006$\\
14 & $4.0247\times10^{-5}$ & $4.043\times10^{-5}\,^*$ & $0.971$ & $1.005$\\
15 & $2.9676\times10^{-5}$ & $2.978\times10^{-5}\,^*$ & $0.976$ & $1.003$\\
16 & $2.2351\times10^{-5}$ & $2.241\times10^{-5}\,^*$ & $0.979$ & $1.003$\\
17 & $1.7149\times10^{-5}$ & $1.718\times10^{-5}\,^*$ & $0.981$ & $1.002$\\
18 & $1.3374\times10^{-5}$ & $1.340\times10^{-5}\,^*$ & $0.984$ & $1.002$\\
\bottomrule
\end{tabular}
\end{center}

\noindent
($^*$: series estimate $\tfrac12(\mathrm{eff}(2)+\mathrm{eff}(3))$, with a
validated systematic of $\approx-0.4\%$; $n\le11$ rows are exact.)  The
evidence for
\begin{equation}\label{eq:r4}
r_4(n)\;=\;\frac{1+o(1)}{(n)_4}\;=\;r_3(n)\,\bigl(1+O(n^{-1})\bigr)
\end{equation}
is now overwhelming; the exact $n\le11$ values show the relative excess
$r_4/r_3-1$ decaying from $4.6\%$ ($n=9$) to $1.8\%$ ($n=11$), already at
the boundary of the series systematics beyond that.  The decomposition of
$r_4-r_3=\Delta_\lambda\log\Pr(X=0)$ through the truncations shows the
difference is \emph{not} dominated by $\Delta\mu$ (which contributes only
$\approx5\%$ of it at $n=11$) but spread across the higher coincidence
statistics --- consistent with its smallness and rapid decay.
Section~\ref{sec:switch} turns this evidence into an exact ``switch
calculus'' and a proof of \eqref{eq:r4} (Theorem~\ref{thm:layer4}).

%======================================================================
\section{An exact switch calculus for the fourth layer}\label{sec:switch}

This section develops the machinery for TODO~1 of the programme: a proof
that the fourth-layer effect equals the third-layer effect to relative
$O(1/n)$,
\begin{equation}\label{eq:r4thm}
r_4(n)\;=\;r_3(n)\,\bigl(1+O(n^{-1})\bigr),
\qquad r_3(n)=\frac{M_{n-4}}{M_n}.
\end{equation}
The idea is to compare the two cycle types through \emph{adjacent
representatives}: derangements $\sigma$ (type $(n)$) and $\sigma'$ (type
$(2,n-2)$) whose boards differ by a single $2\times2$ switch.  All
$\sigma$-dependent quantities then differ by four-term brackets of
punctured counts, and these brackets are governed \emph{exactly} by
Wronskian-type identities for path matching polynomials --- the punctured
analogues of the fusion identity (Lemma~\ref{lem:fusion}).  The programme is carried out in full: every step is an exact identity,
an exhaustively verified finite case analysis, or a uniform Bonferroni
estimate, and each is cross-checked against exact computation.

\subsection{Adjacent representatives and the switch decomposition}

Since $N^{(k)}_n$ depends only on the cycle type
(Lemma~\ref{lem:type}), we may choose
\[
\sigma=(1\,2\,3\cdots n),\qquad \sigma'=(1\,2)(3\,4\cdots n).
\]
Then $\sigma(i)=\sigma'(i)$ except at $i=2,n$:
\[
B_2(\sigma)\setminus B_2(\sigma')=\{d_1,d_2\},\quad
B_2(\sigma')\setminus B_2(\sigma)=\{e_1,e_2\},
\]
\[
d_1=(2,3),\ d_2=(n,1),\qquad e_1=(2,1),\ e_2=(n,3).
\]
Let $B_0=B_2(\sigma)\cap B_2(\sigma')$, a board of $2n-2$ cells whose
bipartite graph is a disjoint union of two paths,
\[
P^{(1)}:\ c_1-r_1-c_2-r_2\quad(3\ \text{edges}),\qquad
P^{(2)}:\ c_3-r_3-c_4-\cdots-c_n-r_n\quad(2n-5\ \text{edges}),
\]
where $r_i$, $c_j$ denote row and column vertices.  The four cells
$d_1,d_2,e_1,e_2$ avoid $B_0$ and connect the four path-ends pairwise:
$d_1=r_2c_3$ and $d_2=r_nc_1$ join the two paths (\emph{crosswise}),
while $e_1=r_2c_1$ and $e_2=r_nc_3$ close each path onto itself.
Adding $\{d_1,d_2\}$ to $B_0$ yields the single cycle $C_{2n}$ of
$B_2(\sigma)$; adding $\{e_1,e_2\}$ yields the $C_4\cup C_{2n-4}$ of
$B_2(\sigma')$.

Let $A_0$ be the set of permutations avoiding $B_0$, $N_0=|A_0|$, and for
cells $z_1,\dots$ let $R_0(M)$ denote the number of $\tau\in A_0$ with
$\tau\supseteq M$ (a punctured count, computable by
\eqref{eq:punctured}; all components of the punctured board are
\emph{paths}).  Write
\[
N^{(4)}_0=\#\{(\tau,\eta)\in A_0^2:\ \tau\cap\eta=\emptyset\},\qquad
G(z)=\#\{(\tau,\eta)\in A_0^2:\ \tau\cap\eta=\emptyset,\ z\in\tau\},
\]
and analogously $G_2(z,z')$ ($z,z'\in\tau$) and $H(z,z')$ ($z\in\tau$,
$z'\in\eta$).

\begin{lemma}[switch decomposition; exact]\label{lem:switchIE}
With $\Delta N^{(4)}=N^{(4)}_n(\sigma')-N^{(4)}_n(\sigma)$,
\begin{align}
\Delta N^{(4)}
&=2\,\mathrm{Br}_G
+2\bigl[G_2(e_1,e_2)-G_2(d_1,d_2)\bigr]
+2\bigl[H(e_1,e_2)-H(d_1,d_2)\bigr],
\label{eq:switchIE}\\
\mathrm{Br}_G&:=G(d_1)+G(d_2)-G(e_1)-G(e_2).\notag
\end{align}
\end{lemma}

\begin{proof}
A pair of disjoint rows avoiding $B_2(\sigma)$ is a pair of disjoint
members of $A_0$ avoiding $d_1,d_2$; inclusion--exclusion over the events
$d_i\in\tau\cup\eta$ gives
$N^{(4)}(\sigma)=N^{(4)}_0-U(d_1)-U(d_2)+U(d_1,d_2)$ with
$U(z)=2G(z)$ (a cell lies in at most one row of a disjoint pair) and
$U(z,z')=2G_2(z,z')+2H(z,z')$ (using $H(z,z')=H(z',z)$, by swapping
$\tau,\eta$).  Similarly for $\sigma'$ with $e_1,e_2$; subtract.
(Verified by complete enumeration at $n=6$: $\Delta N^{(4)}=10=
2\cdot6+2\cdot(-1)+2\cdot0$; and $n=7$: $400=2\cdot224+2\cdot4-2\cdot28$;
\texttt{switch.py}.)
\end{proof}

\begin{remark}[a hidden symmetry]\label{rem:psi}
The map $\Psi(i,c)=(3-c,\,3-i)\bmod n$ is an involution of the cell grid
that preserves $B_0$, swaps $d_1\leftrightarrow d_2$, and fixes $e_1$ and
$e_2$.  It induces the involution
$\tau\mapsto\Psi(\tau)$, $\Psi(\tau)(i)=3-\tau^{-1}(3-i)$, of $A_0$,
which preserves disjointness of pairs.  Hence $G(d_1)=G(d_2)$
\emph{exactly} (confirmed by exact computation at $n=8,\dots,11$), while
$G(e_1)$ and $G(e_2)$ are merely each $\Psi$-invariant.
\end{remark}

\subsection{The path toolkit}

Let $p_e(x)=\sum_k m_k(P_e)x^k$ be the matching polynomial of the path
with $e$ edges ($p_0=1$, $p_1=1+x$, $p_e=p_{e-1}+xp_{e-2}$), so that with
$\alpha_\pm$ as in \S\ref{sec:thirdrow},
$p_e=(\alpha_+^{\,e+2}-\alpha_-^{\,e+2})/(\alpha_+-\alpha_-)$.
Extend to negative indices by this formula:
\[
p_{-1}=1,\qquad p_{-2}=0,\qquad
p_{-e}=-(-x)^{\,2-e}\,p_{e-4}\ (e\ge3).
\]

\begin{lemma}[path Wronskian]\label{lem:wronskian}
For all integers $a,b$,
\begin{equation}\label{eq:wronskian}
p_a\,p_b-p_{a-1}\,p_{b+1}\;=\;-(-x)^{\,b+2}\,p_{a-b-3}.
\end{equation}
In particular $p_2p_{C-1}-p_1p_C=-x^3p_{C-4}$ and
$p_2p_C-p_3p_{C-1}=-x^4p_{C-5}$.
\end{lemma}

\begin{proof}
Put $d_m=(\alpha_+^m-\alpha_-^m)/(\alpha_+-\alpha_-)$, so $p_e=d_{e+2}$.
In the product $(\alpha_+-\alpha_-)^2(p_ap_b-p_{a-1}p_{b+1})$ the
``diagonal'' terms $\alpha_\pm^{a+b+4}$ cancel and the cross terms give
\[
-(\alpha_+\alpha_-)^{b+2}\Bigl[\alpha_+^{a-b}+\alpha_-^{a-b}
-\alpha_+^{a-b-1}\alpha_--\alpha_-^{a-b-1}\alpha_+\Bigr]
=-(-x)^{b+2}(\alpha_+-\alpha_-)\bigl(\alpha_+^{a-b-1}-\alpha_-^{a-b-1}\bigr),
\]
using $\alpha_+\alpha_-=-x$; divide by $(\alpha_+-\alpha_-)^2$.
(Verified symbolically for $a,b\in[0,11]$; \texttt{fusion\_lemmas.py}.)
\end{proof}

The cycle polynomial decomposes over paths:
$c_m:=\alpha_+^m+\alpha_-^m=d_{m+1}+x\,d_{m-1}$, i.e.
\begin{equation}\label{eq:cycletopath}
M_{C_{2m}}(x)=p_{2m-1}(x)+x\,p_{2m-3}(x).
\end{equation}

For a polynomial $f$ and integer $m\ge0$ define the \emph{IE evaluation}
$\langle f\rangle_m=\sum_k(-1)^k[x^k]f\cdot(m-k)!$, so that the count of
permutations of an $m$-set avoiding a board with matching polynomial $f$
is $\langle f\rangle_m$, and
\begin{equation}\label{eq:shift}
\langle x^s f\rangle_m=(-1)^s\langle f\rangle_{m-s}.
\end{equation}

\subsection{The empty bracket: an exact identity}

Define, for a partial matching $M$ avoiding $B_0$ and compatible with a
pin $z$ (i.e.\ $M\cup\{z\}$ is again a matching avoiding $B_0$), the
\emph{bracket}
\[
\mathrm{Br}(M)\;=\;R_0(M\cup d_1)+R_0(M\cup d_2)-R_0(M\cup e_1)-R_0(M\cup e_2).
\]

\begin{proposition}[empty bracket]\label{prop:S0}
$\mathrm{Br}(\emptyset)=M_{n-4}$.
\end{proposition}

\begin{proof}
Pinning a cell $(i,c)$ deletes the vertices $r_i,c_c$ from $B_0$.  Each
of the four special vertices $r_2,c_1,c_3,r_n$ has degree $1$ in $B_0$,
so each pin removes exactly two edges; the resulting boards are unions of
two or three paths.  Grouping the bracket as
$(R_0(d_1)-R_0(e_1))+(R_0(d_2)-R_0(e_2))$, both differences compare
end-deletions of two paths within a common board, and
Lemma~\ref{lem:wronskian} gives, as polynomial identities on the boards,
\[
f_{d_1}-f_{e_1}=-x^3p_{2n-9},\qquad
f_{d_2}-f_{e_2}=-x^4p_{2n-11}
\]
(the computation is the generic case of Lemma~\ref{lem:master} below with
$A=2$, $C=2n-5$, resp.\ $A''=3$, $C''=2n-6$).  Hence by \eqref{eq:shift},
\[
\mathrm{Br}(\emptyset)
=\langle p_{2n-9}\rangle_{n-4}-\langle p_{2n-11}\rangle_{n-5}
=\bigl\langle p_{2(n-4)-1}+x\,p_{2(n-4)-3}\bigr\rangle_{n-4}
=\langle M_{C_{2(n-4)}}\rangle_{n-4}=M_{n-4},
\]
using \eqref{eq:cycletopath}.  (Verified exactly for $6\le n\le14$;
\texttt{switch.py}, check~1.)
\end{proof}

Applying the same identity to single rows instead of pairs gives
$\Delta N^{(3)}=g(d_1)+g(d_2)-g(e_1)-g(e_2)=M_{n-4}$ with
$g(z)=R_0(\{z\})$ --- an alternative, ``local'' proof of the $\ell=2$
case of Theorem~\ref{thm:thirdrow}: the two double-pin terms
$R_0(\{d_1,d_2\})$ and $R_0(\{e_1,e_2\})$ coincide because both pin pairs
delete the \emph{same} four vertices $r_2,r_n,c_1,c_3$.

\subsection{Structure of the bracket for general punctures}

Fix a matching $M$ avoiding $B_0$ with $\mathrm{rows}(M)\cap\{2,n\}=
\mathrm{cols}(M)\cap\{1,3\}=\emptyset$ (so that all four pins are
compatible; the remaining cases are the \emph{incompatible classes}
treated in \S\ref{sec:boundary}).  Let $G_M$ denote $B_0$ punctured by
$M$ and let $f_{M\cup z}$ be the matching polynomial of $G_M$ minus the
two vertices of $z$, so $R_0(M\cup z)=\langle f_{M\cup z}\rangle_{n-|M|-1}$
and the bracket polynomial is
\[
B_M\;=\;f_{M\cup d_1}+f_{M\cup d_2}-f_{M\cup e_1}-f_{M\cup e_2}.
\]

\begin{lemma}[valuation]\label{lem:valuation}
$[x^0]B_M=[x^1]B_M=0$, and
\[
[x^2]B_M=\bigl(\mathbf 1_{c_2\text{-column}\notin M}-\mathbf 1_{c_n\text{-column}\notin M}\bigr)
\bigl(\mathbf 1_{r_3\text{-row}\notin M}-\mathbf 1_{r_1\text{-row}\notin M}\bigr).
\]
In particular the valuation of $B_M$ is at least $3$ unless $M$ meets
exactly one of the columns $\{2,n\}$ \emph{and} exactly one of the rows
$\{1,3\}$, and it is always at least $2$.
\end{lemma}

\begin{proof}
$[x^0]$: each $f$ has constant term $1$.  $[x^1]$: $m_1$ of a vertex-deleted
graph is $E-\deg u-\deg v$ for nonadjacent $u,v$; the four deletions use
each of $r_2,r_n,c_1,c_3$ twice with opposite signs, so both the $E$- and
the degree-terms cancel.  $[x^2]$: $m_2(G\smallsetminus\{u,v\})=
\binom{E-\deg u-\deg v}{2}-A(G\smallsetminus\{u,v\})$ where $A$ counts
adjacent edge-pairs.  Writing $a=\deg r_2$, $b=\deg c_3$, $c=\deg r_n$,
$d=\deg c_1$ (degrees in $G_M$), the binomial terms contribute
$(a-c)(b-d)$ by the quadratic identity
$f(a{+}b)+f(c{+}d)-f(a{+}d)-f(c{+}b)=(a-c)(b-d)$ for
$f(t)=\binom{E-t}{2}$.  Since the graph is bipartite, a row-vertex and a
column-vertex have no common neighbour, so
$A(G\smallsetminus\{u,v\})=A(G)-\binom{\deg u}{2}-\binom{\deg v}{2}
-\sum_{w\sim u}(\deg w-1)-\sum_{w\sim v}(\deg w-1)$; every one of these
terms appears twice with opposite signs in the bracket, so the $A$-terms
cancel identically.  Finally $\deg r_2=\mathbf 1[\text{column
}2\notin\mathrm{cols}(M)]$ (its only $B_0$-edge is the cell $(2,2)$),
and similarly $\deg r_n=\mathbf 1[c_n\ \mathrm{free}]$,
$\deg c_3=\mathbf 1[r_3\ \mathrm{free}]$, $\deg c_1=\mathbf 1[r_1\ \mathrm{free}]$.
(Exhaustive verification for all $|M|\le2$ at $n=9,10$:
minimal valuation $2$, and valuation $2$ occurs exactly in the predicted
classes; \texttt{valuation\_scan.py}.)
\end{proof}

\begin{lemma}[master formula]\label{lem:master}
Let $M$ be as above.  In $G_M\smallsetminus r_2$, let $A$ be the edge-length
of the path-piece containing $c_1$ and $C$ the edge-length of the piece
containing $c_3$; in $G_M\smallsetminus r_n$, let $A''$ and $C''$ be the
edge-lengths of the pieces containing $c_1$ and $c_3$.  Then, as an
identity of polynomials,
\begin{equation}\label{eq:master}
B_M\;=\;(-x)^{A+1}\,p_{C-A-2}\cdot f^{\mathrm{rest}}_1
\;-\;(-x)^{A''+1}\,p_{C''-A''-2}\cdot f^{\mathrm{rest}}_2 ,
\end{equation}
where $f^{\mathrm{rest}}_i$ are the matching polynomials of the remaining
pieces.  Moreover $A\le 2$ and $A''\le3$ always, and in the
\emph{window-avoiding} case --- $M$ meets none of rows $1,3$, columns
$2,n$, and no piece-end lies within distance $8$ of $c_3$ or $r_n$ ---
one has $A=2$, $A''=3$, $C''\in\{C,C-1\}$, and
\begin{equation}\label{eq:generic}
B_M=-x^3\,p_{C-4}\,p_D\,\textstyle\prod_{\mathrm{mid}}
-x^4\,p_{C-5}\,p_{D-1}\,\prod_{\mathrm{mid}},
\end{equation}
with $D$ the edge-length of the piece containing $r_n$ and
$\prod_{\mathrm{mid}}$ the product over the interior pieces.
\end{lemma}

\begin{proof}
Group $B_M=(f_{M\cup d_1}-f_{M\cup e_1})+(f_{M\cup d_2}-f_{M\cup e_2})$.
The first difference deletes $r_2$ from $G_M$ and then compares deletion
of the end-vertices $c_3$ (in its piece of length $C$) and $c_1$ (in its
piece of length $A$); these lie in distinct components (they descend from
the distinct components $P^{(2)},P^{(1)}$ of $B_0$), so the difference is
$\bigl[p_{C-1}p_A-p_Cp_{A-1}\bigr]\cdot f^{\mathrm{rest}}_1$, which is
$(-x)^{A+1}p_{C-A-2}f^{\mathrm{rest}}_1$ by \eqref{eq:wronskian}.  The
second difference is handled identically after deleting $r_n$ (with the
roles of $c_1,c_3$ reversed, whence the minus sign).  $A\le2$, $A''\le3$
because every piece of $P^{(1)}$ has at most $3$ edges.  In the
window-avoiding case $P^{(1)}$ is intact, so $A=2$ ($P^{(1)}\smallsetminus
r_2=c_1r_1c_2$) and $A''=3$; deleting $r_n$ shortens only the
$r_n$-piece, so $C''=C$ if the $c_3$- and $r_n$-pieces are distinct
($C''=C-1$ if $M$ leaves $P^{(2)}$ uncut).  Substituting gives
\eqref{eq:generic}.  (Verified as exact integer identities on random
generic $M$ with $1\le|M|\le4$ at $n=9,11,13$;
\texttt{fusion\_lemmas.py}.)
\end{proof}

\subsection{A uniform estimate for punctured path boards}

\begin{lemma}[uniform punctured estimate]\label{lem:uniform}
There are absolute constants $m_0, C_0$ with the following property.
Let $F$ be a disjoint union of paths and cycles with $E$ edges,
considered as a board in an $m\times m$ grid, $E\le 2m$, $m\ge m_0$.  Then
\[
\langle M_F\rangle_m\;=\;m!\,e^{-E/m}\,(1+\theta),
\qquad |\theta|\le C_0/m .
\]
\end{lemma}

\begin{proof}
Write $S=\langle M_F\rangle_m/m!=\sum_k(-1)^k m_k/(m)_k$.  By Bonferroni
(the summands are the inclusion--exclusion of the events ``the
permutation hits a fixed cell of $F$''), the partial sums alternate
around $S$, so for every $K$,
$\bigl|S-\sum_{k<K}(-1)^km_k/(m)_k\bigr|\le m_K/(m)_K$.
For $k\le K:=\lceil2\log m\rceil$ we use the two-sided bounds
$\frac1{k!}\prod_{i<k}(E-3i)\le m_k\le E^k/k!$ (each edge of a
max-degree-$2$ graph is adjacent to at most two others) and
$(m)_k=m^ke^{-k^2/2m+O(k^3/m^2)}$, giving
$m_k/(m)_k=\frac{(E/m)^k}{k!}(1+\delta_k)$ with
$|\delta_k|\le Ck^2(1/E+1/m)$ when $3K^2\le E$, and
$m_k/(m)_k\le\frac{(E/m)^k}{k!}e^{k^2/m}$ always.  Summing,
$\sum_{k<K}(-1)^k\frac{(E/m)^k}{k!}(1+\delta_k)=e^{-E/m}+O(1/m)$
(the weighted error $\sum_k\frac{(E/m)^k}{k!}k^2(1/E+1/m)=O(1/m)$
uniformly in $E\le2m$, including small $E$), the truncation tails of the
exponential series and the Bonferroni tail being
$O((2e/K)^K)=O(m^{-2})$.  If $E<3K^2$ the same computation with the
crude bound $|\delta_k|\le Ck^2/E$ for $k\le\sqrt{E/3}$ and Bonferroni
truncation there gives the claim directly.
\end{proof}

\begin{corollary}[weights and moments]\label{cor:moments}
For $n\ge n_0$ and any matching $M$ avoiding $B_0$ with $|M|=j\le n/2$:
$q(M):=R_0(M)/N_0\le 2e^{6j/n}/(n)_j$, and consequently
\[
\sum_{j\ge0}(j+1)^4\!\!\sum_{M\in\mathcal M_j}\!q(M)^2\;\le\;C_1
\qquad\text{(absolute constant)} .
\]
Moreover $N^{(4)}_0=\sum_M(-1)^{|M|}R_0(M)^2$ (pair inclusion--exclusion,
Proposition~\ref{prop:pairIE} on $B_0$), and $c_2^{-1}N_0^2\le
N^{(4)}_0\le N_0^2$ for an absolute $c_2$ and $n\ge n_0$.
\end{corollary}

\begin{proof}
The weight bound follows from Lemma~\ref{lem:uniform} applied to numerator
and denominator ($E$ changes by at most $4j$ under puncturing).  The
moment bound follows from $|\mathcal M_j|\le (n(n-2))^j/j!$.  The lower
bound on $N^{(4)}_0$ follows from Bonferroni in the pair
inclusion--exclusion together with
$t_j:=\sum_{\mathcal M_j}q(M)^2=\mu^j/j!\,(1+O((j{+}1)/n))$ for
$j\le5$, $\mu=t_1\in[1,1.5]$, which is a direct consequence of
Lemma~\ref{lem:uniform}.
\end{proof}

\subsection{The generic stratum}

\begin{proposition}[generic per-puncture estimate]\label{prop:generic}
There are $n_0,C$ such that for $n\ge n_0$, every window-avoiding $M$
with $|M|=j\le n/2$ satisfies
\[
\langle B_M\rangle_{n-1-j}
=\frac{M_{n-4}}{M_n}\,R_0(M)\,\bigl(1+\delta_M\bigr),
\qquad|\delta_M|\le C\,\frac{j+1}{n}.
\]
\end{proposition}

\begin{proof}
By \eqref{eq:generic} and \eqref{eq:shift},
$\langle B_M\rangle_{n-1-j}
=\langle p_{C-4}p_D\prod_{\mathrm{mid}}\rangle_{n-4-j}
-\langle p_{C-5}p_{D-1}\prod_{\mathrm{mid}}\rangle_{n-5-j}$.
Apply Lemma~\ref{lem:uniform} to both terms and to
$R_0(M)=\langle p_3\,p_C\,p_D\prod_{\mathrm{mid}}\rangle_{n-j}$; the
edge counts of the two bracket boards are $E(G_M)-7$ and $E(G_M)-9$, so
the exponential factors agree with that of $R_0(M)$ up to
$e^{O((j+1)/n)}$, while the factorial prefactors give
\[
\frac{(n-4-j)!-(n-5-j)!\,e^{O(1/n)}}{(n-j)!}
=\frac{1}{(n-j)_4}\Bigl(1+O\bigl(\tfrac1{n-j}\bigr)\Bigr).
\]  Since $M_{n-4}/M_n=(n)_4^{-1}(1+O(n^{-2}))$ and
$(n)_4/(n-j)_4=1+O(j/n)$ for $j\le n/2$, the claim follows.
\end{proof}

\subsection{Boundary and incompatible classes; the remaining lemma}
\label{sec:boundary}

Expanding each $G$-term of $\mathrm{Br}_G$ by pair inclusion--exclusion
over common cells $M\subseteq\tau\cap\eta$ (which automatically avoid the
pinned row and column) gives the exact stratified form
\begin{equation}\label{eq:strata}
\mathrm{Br}_G=\sum_{j\ge0}(-1)^j S_j,\qquad
S_j=\sum_{M\in\mathcal M_j}\Bigl[\textstyle\sum_{z\in\{d_1,d_2\}}
R_0^\dagger(M\cup z)-\sum_{z\in\{e_1,e_2\}}R_0^\dagger(M\cup z)\Bigr]R_0(M),
\end{equation}
where $R_0^\dagger(M\cup z)=R_0(M\cup z)$ if $M\cup\{z\}$ is a matching,
$=R_0(M)$ if $z\in M$, and $=0$ otherwise.  ($S_0=M_{n-4}N_0$ by
Proposition~\ref{prop:S0}; verified together with $S_1,S_2,S_3$ against
direct enumeration at $n=8,9$; \texttt{switch\_j1.py}.)
The matchings $M$ fall into three classes:
\begin{itemize}[leftmargin=2em]
\item[\textup{(a)}] \emph{window-avoiding} $M$: controlled by
Proposition~\ref{prop:generic};
\item[\textup{(b)}] \emph{compatible boundary} $M$ (all pins compatible
but $M$ meets the window): write $B_M=(f_{M\cup d_1}-f_{M\cup e_1})
+(f_{M\cup d_2}-f_{M\cup e_2})$; by Lemma~\ref{lem:master} each group is
a single $x^{s_i}$-shifted path product with $s_1,s_2\ge2$.  If both
$s_i\ge3$, Lemma~\ref{lem:uniform} gives
$|\langle B_M\rangle_{n-1-j}|\le C(n-j)^{-4}R_0(M)$.  If some $s_i=2$
then, unless $M$ makes the two-hit configuration of
Lemma~\ref{lem:valuation} (extra measure $O((j{+}1)^2/n^2)$, crude bound
$C(n-j)^{-3}R_0(M)$), the $x^2$ coefficients of the two groups cancel
($[x^2]B_M=0$), the two shifted products have edge counts differing by
$O(1)$, and the difference of their Lemma~\ref{lem:uniform} evaluations
is again $O((n-j)^{-4})R_0(M)$.  With Lemma~\ref{lem:linemeasure}, the
class-(b) total is $O((j{+}1)^2/n)\cdot
(M_{n-4}/M_n)\sum_{\mathcal M_j}R_0(M)^2$ per stratum;
\item[\textup{(c)}] \emph{incompatible classes} ($M$ meets row $2$, row
$n$, column $1$, or column $3$, or contains a pin cell): here the bracket
degenerates to two-term or one-term combinations whose individual sizes
are as large as $(n-j)^{-2}R_0(M)$ (two-term, $x^1$-level) or
$(n-j)^{-1}R_0(M)$ (one-term, $x^0$-level).
\end{itemize}

The class-(c) terms are \emph{not} small individually, but they cancel
in structured pairs.  The key device is a family of \emph{line-swap
pairings}: bijections $M\mapsto\varphi(M)$ between two classes which
leave the deleted row- and column-sets of the \emph{pinned} counts
unchanged --- so that the pinned factors agree exactly --- while the
free factors $R_0(M)$, $R_0(\varphi M)$ differ by a Wronskian
(Lemma~\ref{lem:wronskian}), hence by a relative $O(n^{-3})$.  We first
record the measure bound used throughout.

\begin{lemma}[line measure]\label{lem:linemeasure}
For any fixed line (row or column) $\ell$ and $j\le n/2$,
$\sum_{M\in\mathcal M_j,\,M\text{ meets }\ell}R_0(M)^2
\le C\,\tfrac{j}{n}\,\tilde C^{\,j}N_0^2/j!$ with absolute $C,\tilde C$.
\end{lemma}

\begin{proof}
By Lemma~\ref{lem:uniform}, $R_0(M)\le C'R_0(M\smallsetminus z)/(n-j)$
for $z\in M$; sum over the $\le n$ cells $z$ of $\ell$ and over
$M'=M\smallsetminus z\in\mathcal M_{j-1}$, and use
Corollary~\ref{cor:moments}.
\end{proof}

\begin{lemma}[pairing]\label{lem:pairing}
The total contribution to $\sum_j(-1)^jS_j$ of the classes
\textup{(c-2)} ($M$ meets one pinned row and one pinned column) and
\textup{(c-3)} ($M$ contains a pin cell), and the difference
$G_2(e_1,e_2)-G_2(d_1,d_2)$, are all
$O\bigl(n^{-1}(M_{n-4}/M_n)N_0^2\bigr)$.
\end{lemma}

\begin{proof}
\emph{(c-2).}  Consider the class $\{$row $2$, column $3\}$: all pins
except $d_2$ are incompatible, so the stratum term is
$+R_0(M\cup d_2)R_0(M)$.  Map $\varphi$: replace the column-$3$ cell
$(\rho,3)\in M$ by $(\rho,1)$; then $\varphi M$ lies in the class
$\{$row $2$, column $1\}$, whose only compatible pin is $e_2$, with
term $-R_0(\varphi M\cup e_2)R_0(\varphi M)$.  The deletions of
$M\cup d_2$ and $\varphi M\cup e_2$ are the \emph{same} rows and the
same columns ($\{3\}\cup\{1\}$ in both), so
$R_0(M\cup d_2)=R_0(\varphi M\cup e_2)$ \emph{exactly} (verified on
random instances; \texttt{switch.py} suite).  The free factors differ by
the deletion of $c_3$ versus $c_1$ from the same board with $r_2$
already removed, which is the master-formula comparison:
$R_0(M)-R_0(\varphi M)=\langle-x^3p_{C-4}\cdot
f^{\mathrm{rest}}\rangle$, a relative $O(n^{-3})$.  Since the class
carries measure $O((j{+}1)^2/n^2)$ (Lemma~\ref{lem:linemeasure} twice)
and the pinned factor a further $(n-j)^{-1}$, the paired total is
$O(n^{-2}\cdot n^{-1}\cdot n^{-3}\cdot n\,N_0^2)$ summed over strata,
i.e.\ $O(n^{-5})N_0^2=O(n^{-1}(M_{n-4}/M_n))N_0^2$.  The pairs
$\{$row $n$, col $1\}$ / $\{$row $n$, col $3\}$ (surviving pins
$d_1$/$e_1$) are handled by the same map; $\varphi$-exceptional cases
($\rho=1$, where $(\rho,1)\in B_0$) are paired instead by the row swap
$(2,\gamma)\mapsto(n,\gamma)$ against the class $\{$row $n$, col $3\}$,
with the same exact pinned-factor equality and an $x^2$ Wronskian, and
measure smaller by a further $n^{-1}$; deeper exceptional cells carry
measure $O(n^{-4})$ and the crude one-term bound suffices.
Exact totals confirm the mechanism: the two class totals (all strata
$j\le4$) are $+4.277$ vs $-4.276$ in units of $M_{n-4}N_0$ at $n=8$, and
$+3.9012$ vs $-3.9025$ at $n=9$ --- cancellation to relative
$O(n^{-2})$.

\emph{(c-3).}  For $M\ni d_1$ the surviving terms are
$+R_0(M)^2+R_0(M\cup d_2)R_0(M)$ (pins $e_1,e_2$ are incompatible with
$d_1$'s row and column); for $M\ni e_1$ they are
$-R_0(M)^2-R_0(M\cup e_2)R_0(M)$ ($d_1,d_2$ incompatible).  Pair
$M\leftrightarrow M'=M\smallsetminus d_1\cup e_1$ (same rows; columns
$3\leftrightarrow1$).  Again the double-deletion sets of
$M\cup d_2$ and $M'\cup e_2$ coincide, so the cross terms cancel up to
the free-factor Wronskian; and
$R_0(M)^2-R_0(M')^2=(R_0(M)-R_0(M'))(R_0(M)+R_0(M'))$ with
$R_0(M)-R_0(M')$ the \emph{same} $c_3$-vs-$c_1$ comparison with $r_2$
deleted, again $-x^3p_{C-4}$-suppressed.  With class measure
$O(n^{-2})$ this gives $O(n^{-2}\cdot n^{-3})N_0^2$ per stratum,
$O(n^{-5})N_0^2$ in total; same for the $d_2/e_2$ pairing.  (Exact
check at $n=9$, strata $j\le3$: the four pin-cell class totals are
$-33.15,\,+33.00,\,-33.15,\,+33.18$ in units of $M_5N_0$, pairwise sums
$-0.15$ and $+0.03$.)

\emph{$G_2$-difference.}  In the pair inclusion--exclusion for
$G_2(d_1,d_2)$ and $G_2(e_1,e_2)$ the compatible-$M$ ranges coincide
(both pin pairs delete rows $2,n$ and columns $1,3$), and the pinned
factors $R_0(M\cup\{d_1,d_2\})=R_0(M\cup\{e_1,e_2\})$ agree exactly; so
all terms cancel except those with $M$ containing a pin cell, which are
handled exactly as in (c-3).
\end{proof}

\begin{lemma}[$H$-difference]\label{lem:Hdiff}
$H(e_1,e_2)-H(d_1,d_2)=O\bigl(n^{-1}(M_{n-4}/M_n)N_0^2\bigr)$.
\end{lemma}

\begin{proof}
Both $H$'s expand over the \emph{same} $M$-range ($M$ must avoid rows
$2,n$ and columns $1,3$, and cannot contain a pin cell since the pins
lie in $\tau\smallsetminus\eta$ resp.\ $\eta\smallsetminus\tau$), with
terms $R_0(M\cup e_1)R_0(M\cup e_2)-R_0(M\cup d_1)R_0(M\cup d_2)$.
Writing the difference as
$R_0(M\cup e_1)[R_0(M\cup e_2)-R_0(M\cup d_2)]
+R_0(M\cup d_2)[R_0(M\cup e_1)-R_0(M\cup d_1)]$, each bracket is a
single pin-difference, i.e.\ an $x^{s}$-shifted single path product by
Lemma~\ref{lem:master} with $s\ge3$ off the window and $s\ge2$ on it
(Lemma~\ref{lem:valuation} logic); with the two pinned factors each
$O(R_0(M)/(n-j))$ this gives $O(R_0(M)^2 n^{-5})$ per generic $M$ and
$O(R_0(M)^2n^{-4})$ on the window (measure $O((j{+}1)/n)$); summing with
Corollary~\ref{cor:moments} gives the claim.
\end{proof}

It remains to control class \textup{(c-1)}: $M$ meets exactly one
pinned line, say row $2$ (the cases row $n$, column $1$, column $3$
being its images under $\Psi$ and under transposition of the roles of
$d$- and $e$-cells), so that the bracket is the two-term
$R_0(M\cup d_2)-R_0(M\cup e_2)$.  The dangerous sub-case is a
\emph{doubly degenerate} window hit, where this bracket is only
$x^1$-suppressed; it is handled by an exact \emph{column-exchange
identity} and one more line-swap pairing.

\begin{lemma}[column exchange]\label{lem:colexchange}
Let $M$ avoid rows $n,3$ and columns $1,3$ and meet both rows $1$ and
$2$.  Then, exactly,
\[
R_0(M\cup d_2)-R_0(M\cup e_2)\;=\;-\,R_0\bigl(M\cup\{(3,1),(n,3)\}\bigr).
\]
Dually, if $M$ avoids rows $n,1$ and columns $1,3$ and meets rows $2$
and $3$, then
$R_0(M\cup d_2)-R_0(M\cup e_2)=+\,R_0(M\cup\{(1,3),(n,1)\})$.
And if $M$ meets all three rows $1,2,3$, the bracket vanishes
identically.
\end{lemma}

\begin{proof}
Since rows $1,2$ meet $M$, the cells $(1,1),(1,2),(2,2)$ of $B_0$ are
punctured and column $1$ carries no constraint in the reduced system;
both counts live on the rows $\ne n$, and $R_0(M\cup d_2)$
(resp.\ $R_0(M\cup e_2)$) counts the assignments leaving column $1$
(resp.\ column $3$) unused.  Given $\tau$ with column $1$ unused, move
the row $\tau^{-1}(3)$ from column $3$ to column $1$: the new cell is
always admissible (column $1$ is free of constraints), so this is an
injection into the column-$3$-unused family; its image misses exactly
the $\tau'$ with $\tau'(3)=1$, whose count is
$R_0(M\cup\{(3,1),(n,3)\})$ --- note $(3,1)$ avoids $B_0$ and the
deleted lines, and leaving column $3$ unused while deleting row $n$ is
the deletion set of $\{(n,3)\}$.  The dual statement is the same
exchange with the roles of columns $1$ and $3$ reversed (column $3$ is
constraint-free exactly when its unique $B_0$-cell $(3,3)$ is punctured,
i.e.\ row $3$ meets $M$).  If rows $1,2,3$ all meet $M$, both columns
are constraint-free and the two counts coincide.  (All three statements
verified as exact integer identities on random class-$M$ at $n=9,11,13$;
\texttt{switch.py} suite.)
\end{proof}

\begin{lemma}[boundary]\label{lem:boundary}
The total contribution of class \textup{(c-1)} to $\sum_j(-1)^jS_j$ is
$O\bigl(n^{-1}(M_{n-4}/M_n)N_0^2\bigr)$.
\end{lemma}

\begin{proof}
Sub-classify by the window hits accompanying the row-$2$ hit.
\emph{(i) No other special hit:} $A''=2$ in Lemma~\ref{lem:master}, an
$x^3$-shifted single product: $|\mathrm{bracket}|\le C(n-j)^{-4}
R_0(M)\cdot(n-j)$, and Lemma~\ref{lem:linemeasure} gives a class total
$O((j{+}1)/n)(M_{n-4}/M_n)$ per stratum, as in class (b).
\emph{(ii) One additional hit making $A''=1$ (column $2$) or shortening
the $c_3$-side pieces:} valuation $\ge2$ with measure
$O((j{+}1)^2/n^2)$; crude bounds suffice.
\emph{(iii) Row $1$ also hit ($A''=0$, the $x^1$ case):} by
Lemma~\ref{lem:colexchange} the stratum term is exactly
$-R_0(M)R_0(M\cup\{(3,1),(n,3)\})$.  Pair this class with the class
``rows $2,3$ hit'' via the row swap $\varphi:(1,\gamma)\mapsto
(3,\gamma)$ on the row-$1$ cell: $\varphi$ maps one class onto the
other (for the $O(1)$ exceptional columns $\gamma$ where $(3,\gamma)\in
B_0$ or collisions occur, see below), and by the dual identity the
partner's term is $+R_0(\varphi M)R_0(\varphi M\cup\{(1,3),(n,1)\})$.
The two pinned factors agree \emph{exactly}: the deletion sets of
$M\cup\{(3,1),(n,3)\}$ and $\varphi M\cup\{(1,3),(n,1)\}$ are the same
rows and the same columns.  Hence the paired contribution is
$[R_0(M)-R_0(\varphi M)]R_0(M\cup\{(3,1),(n,3)\})$, and
$R_0(M)-R_0(\varphi M)$ compares the deletion of $r_1$ versus $r_3$
from a common board --- a Wronskian difference, of relative size
$O(1/n)$ by Lemmas~\ref{lem:wronskian} and~\ref{lem:uniform} (measured:
$\approx-2/n^2$).  Since the paired class carries measure
$O((j{+}1)^2/n^2)$ and the pinned factor two powers, the total is
$O(n^{-2}\cdot n^{-2}\cdot n^{-1})N_0^2=O(n^{-5})N_0^2$ per stratum
family, summable by Corollary~\ref{cor:moments}.  The
$\varphi$-exceptional column slices ($\gamma$ with $(3,\gamma)\in B_0$,
i.e.\ $\gamma\in\{3,4\}$, or $\gamma$ colliding with $M$'s columns) are
single-column sub-families of a class of size $O(n^{-4})N_0^2$, hence
$O(n^{-5})N_0^2$ by the crude bounds, as are the deeper multi-hit
sub-classes; and the three-row class ``rows $1,2,3$'' contributes $0$
identically by Lemma~\ref{lem:colexchange}.  The mirror families (row
$n$; columns $1$, $3$) are identical under $\Psi$ and under exchanging
the roles of the $d$- and $e$-pins.
\end{proof}

\noindent
Exact class computations confirm the mechanism: the rows-$\{1,2\}$
family totals $-0.0611\,M_4N_0$ at $n=8$ and $-0.0528\,M_5N_0$ at $n=9$
($\approx-0.5/n$ in main-term units) while its individual strata are
$\Theta(n)$ larger; the pinned-factor equalities and both column-exchange
identities are integer-exact on all tested instances
(\texttt{switch.py} suite).

\medskip
\noindent
The full quantity is measured directly by exact permanent sums: with
$E^\pm[A]$ the mean number of disjoint partners over $\tau\ni d$-cells
resp.\ $e$-cells,
\[
\frac{\mathrm{Br}_G\,N_0}{M_{n-4}\,N^{(4)}_0}
=0.8919,\ 0.8685,\ 0.8818,\ 0.8929
\qquad(n=8,9,10,11),
\]
and the deviation decomposes as
$m^+(E^+[A]-E^-[A])/(M_{n-4}E[A])=-1.17/n\pm2\%$ ($n=9,10,11$) plus
$(E^-[A]-E[A])/E[A]=O(n^{-2})$ (\texttt{anatomy.c}), exactly the
$1-O(1/n)$ shape required by Theorem~\ref{thm:layer4}.

\begin{theorem}[fourth layer]\label{thm:layer4}
For $n\ge n_0$,
\[
\log\frac{N^{(4)}_n((2,n{-}2))}{N^{(4)}_n((n))}
=2\,\frac{M_{n-4}}{M_n}+O(n^{-5}),
\qquad\text{hence}\qquad
r_4(n)=r_3(n)\bigl(1+O(n^{-1})\bigr).
\]
\end{theorem}

\begin{proof}
Combine \eqref{eq:switchIE} and \eqref{eq:strata}.  Strata with
$j>n/2$ are bounded crudely by
$|\mathcal M_j|\,(n-j)!\,(n-j-1)!$, which is superexponentially small
against $N_0^2$; assume $j\le n/2$ throughout.  Classes (a) and (b):
by Proposition~\ref{prop:generic} and the class-(b) bound, using
Corollary~\ref{cor:moments} to sum unsigned errors,
\[
\sum_j(-1)^jS_j^{(a)+(b)}
=\frac{M_{n-4}}{M_n}\Bigl[\sum_M(-1)^{|M|}R_0(M)^2\Bigr]\bigl(1+O(n^{-1})\bigr)
=\frac{M_{n-4}}{M_n}N^{(4)}_0\bigl(1+O(n^{-1})\bigr),
\]
where re-completing the class-(c) matchings into the full signed sum
$\sum_M(-1)^{|M|}R_0(M)^2=N^{(4)}_0$ costs another
$O(n^{-1})$-relative term ($O(1/n)$ of the pair measure lies in class
(c)).  Class (c) and the $G_2$-, $H$-differences are
$O(n^{-1}(M_{n-4}/M_n)N_0^2)$ by Lemmas~\ref{lem:pairing},
\ref{lem:Hdiff} and~\ref{lem:boundary}, and $N_0^2\le e^{\mu+o(1)}
N^{(4)}_0$ by Corollary~\ref{cor:moments}.  Hence
$\Delta N^{(4)}=2(M_{n-4}/M_n)N^{(4)}_0(1+O(1/n))$.  Finally
$N^{(4)}_0=N^{(4)}_n((n))(1+O(1/n))$ (the switch decomposition for
$N^{(4)}(\sigma)$ against $N^{(4)}_0$, with $G(z)\le e^{\mu}N^{(4)}_0/
(n-2)$), so $\Delta\log N^{(4)}=2M_{n-4}/M_n+O(n^{-5})$, and
$r_4-r_3=\Delta\log N^{(4)}-2\,\Delta\log N^{(3)}
=O(n^{-5})=r_3\cdot O(n^{-1})$.
\end{proof}

\begin{remark}[general $\ell$; the layer-$k$ programme]
Choosing $\sigma=(1\cdots n)$ and $\sigma'=(1\cdots \ell)(\ell{+}1\cdots
n)$ gives boards differing by one switch at distance $\ell$; the same
calculus applies with $P^{(1)}$ of length $2\ell-1$, and the empty
bracket is $M_{n-2\ell}$ by the same Wronskian computation, predicting
$r_4^{(\ell)}=\rho_\ell(1+O(1/n))$ for the $\ell$-cycle layer-4 effect.
More significantly, nothing in the preceding subsections is specific to
\emph{two} added rows: the switch decomposition and the bracket calculus apply
verbatim to $N^{(k)}$ for any fixed $k$, with pair counts replaced by
$(k{-}2)$-tuple counts and Lemma~\ref{lem:uniform} applied to the
correspondingly punctured boards.  Section~\ref{sec:fifth} carries this
out in full for $k=5$.  The obstruction to
$r_k=r_3(1+o(1))$ for growing $k$ is that the boards seen by the last
rows are no longer unions of two paths --- this is where the cluster
expansion of TODO~4 takes over.
\end{remark}



%======================================================================
\section{The fifth layer: the switch calculus transfers}\label{sec:fifth}

The remark closing Section~\ref{sec:switch} promised that nothing in the
switch calculus is specific to \emph{two} added rows.  This section makes
that promise concrete for the fifth layer: we derive the triple
inclusion--exclusion identity, transfer the switch decomposition and the
bracket calculus, and obtain
\begin{equation}\label{eq:r5thm}
r_5(n)=r_3(n)\bigl(1+O(n^{-1})\bigr),
\end{equation}
where $r_5$ is the fifth-layer effect
$\Delta_\lambda\log\bigl(N^{(5)}/N^{(4)}\bigr)$ between types $(2,n{-}2)$
and $(n)$.  Exact computations (independent of the calculus) confirm the
result at $n=8,9$ with striking precision.  Everything is written at the
same working-note rigor as Section~\ref{sec:switch}: exact identities and
finite verifications are labelled as such, and the two genuinely new
bookkeeping steps (the spectator factors in the strata analysis) are
flagged explicitly.

\subsection{A triple inclusion--exclusion identity}

$N^{(5)}(\sigma)$ is the number of ordered triples $(\tau,\eta,\theta)$
of $B_2$-avoiding permutations that are pairwise cell-disjoint.
Expanding each of the three pairwise-disjointness indicators by
inclusion--exclusion over subsets of that pair's coincidence set --- the
three expansions are independent --- gives the analogue of
Proposition~\ref{prop:pairIE}.

\begin{proposition}[triple inclusion--exclusion]\label{prop:tripleIE}
Let $M_1,M_2,M_3$ range over matchings of the complement of $B_2$
(the coincidence sets of $(\tau,\eta)$, $(\tau,\theta)$,
$(\eta,\theta)$ respectively), subject to \emph{compatibility}: each
pairwise union $M_i\cup M_j$ is again a matching (shared cells are
allowed; two distinct cells of $M_i$, $M_j$ may not share a line).
Then
\begin{equation}\label{eq:tripleIE}
N^{(5)}(\sigma)=\sum_{(M_1,M_2,M_3)}(-1)^{|M_1|+|M_2|+|M_3|}
R(M_1{\cup}M_2)\,R(M_1{\cup}M_3)\,R(M_2{\cup}M_3),
\end{equation}
with $R(\cdot)$ the punctured counts of \eqref{eq:punctured}.
Equivalently, in resummed form: compatibility of the triple is
equivalent to $U=M_1\cup M_2\cup M_3$ being a matching, and collecting
the label multiplicities cell by cell,
\begin{equation}\label{eq:tripleIEU}
N^{(5)}(\sigma)=\sum_{U}\ \sum_{U=U_{ab}\sqcup U_{ac}\sqcup U_{bc}\sqcup U_{abc}}
(-1)^{|U_{ab}|+|U_{ac}|+|U_{bc}|}\,2^{|U_{abc}|}\,
R(A)\,R(B)\,R(C),
\end{equation}
where $A=U_{ab}\cup U_{ac}\cup U_{abc}$,
$B=U_{ab}\cup U_{bc}\cup U_{abc}$, $C=U_{ac}\cup U_{bc}\cup U_{abc}$
are the constraint sets of $\tau,\eta,\theta$: \emph{every coincidence
cell lies in at least two of the three constraint sets}.
\end{proposition}

\begin{proof}
For fixed $\tau,\eta,\theta$ avoiding $B_2$, the product of the three
indicators $\prod_{\text{pairs}}\mathbf 1[\text{disjoint}]$ expands as
$\sum_{M_1\subseteq\tau\cap\eta}(-1)^{|M_1|}\cdots$, independently in
the three pairs; exchanging sums gives \eqref{eq:tripleIE}, the count of
triples with $\tau\supseteq M_1\cup M_2$ etc.\ factorizing since the
three rows are constrained separately (and vanishing unless each union
is a matching).  For \eqref{eq:tripleIEU}: if two cells of the $M_i$
share a line, compatibility of the relevant pair forces them equal, so
pairwise compatibility is the same as $U$ being a matching; each
$u\in U$ carries the nonempty label set $T(u)\subseteq\{1,2,3\}$ of the
$M_i$ containing it, the constraint-set membership of $u$ depends only
on $T(u)$ ($|T|=1$ gives exactly two memberships, $|T|\ge2$ gives all
three), and summing $(-1)^{|T|}$ over the label sets with a given
membership pattern gives weight $-1$ for each two-membership pattern
and $+3-1=+2$ for the full pattern.
(Both forms verified against brute-force enumeration of all triples at
$n=5$, and \eqref{eq:tripleIEU} additionally at $n=6,7$, both cycle
types: $N^{(5)}$ values $36$; $4032$, $4224$; $2185152$, $2198160$.
\texttt{triple\_ie.py}.)
\end{proof}

\begin{remark}[the truncated series is slowly convergent]
Truncating \eqref{eq:tripleIEU} at total order
$q=|M_1|+|M_2|+|M_3|\le Q$ gives an alternating series whose natural
parameter is the \emph{total} coincidence intensity of the three pairs,
$3\mu\approx6$ --- three times that of the pair series of
Section~\ref{sec:pair}.  Convergence in $Q$ is correspondingly slower:
at $n=8$ the $Q\le3$ truncation underestimates the fifth-layer effect
by $24\%$ (exact anchor below), at $n=9$ by $6\%$; the truncation error
shrinks rapidly with $n$ ($Q=3$ values $0.922$, $0.942$, $0.964$,
$0.970$, $0.974$ at $n=10,\dots,14$ in
the $r_5\cdot(n)_4$ normalisation of the table below, each presumably a
few percent low, approaching $1$ monotonically as the layer-4 sequence
did).  We therefore use the truncated series only as a
consistency check at larger $n$ and rely on exact enumeration
(\texttt{n5\_direct2.c}) at $n\le9$; the switch calculus below is what
replaces series estimation as the source of the theorem.
(\texttt{layer5\_production.py}.)
\end{remark}

\subsection{The switch decomposition for triples}

Adopt the setting of Section~\ref{sec:switch}: adjacent representatives
$\sigma,\sigma'$, common board $B_0$ (two paths), pins
$d_1=(2,3)$, $d_2=(n,1)$, $e_1=(2,1)$, $e_2=(n,3)$, and $A_0$ the set of
$B_0$-avoiding permutations.  Write
\[
N^{(5)}_0=\#\{(\tau,\eta,\theta)\in A_0^3\ \text{pairwise disjoint}\},
\qquad
G^{(5)}(z)=\#\{\text{such triples with }z\in\tau\},
\]
and analogously $G_2^{(5)}(z,z')$ ($z,z'\in\tau$) and $H^{(5)}(z,z')$
($z\in\tau$, $z'\in\eta$).

\begin{lemma}[switch decomposition, fifth layer; exact]\label{lem:switchIE5}
With $\Delta N^{(5)}=N^{(5)}(\sigma')-N^{(5)}(\sigma)$,
\begin{align}
\Delta N^{(5)}
&=3\,\mathrm{Br}^{(5)}_G
+3\bigl[G_2^{(5)}(e_1,e_2)-G_2^{(5)}(d_1,d_2)\bigr]
+6\bigl[H^{(5)}(e_1,e_2)-H^{(5)}(d_1,d_2)\bigr],
\label{eq:switchIE5}\\
\mathrm{Br}^{(5)}_G&:=G^{(5)}(d_1)+G^{(5)}(d_2)-G^{(5)}(e_1)-G^{(5)}(e_2).
\notag
\end{align}
\end{lemma}

\begin{proof}
As for Lemma~\ref{lem:switchIE}: inclusion--exclusion over the events
$d_i\in\tau\cup\eta\cup\theta$, with $U(z)=3G^{(5)}(z)$ (choice of the
row through $z$) and $U(z,z')=3G_2^{(5)}(z,z')+6H^{(5)}(z,z')$ (one row
through both, or an ordered pair of distinct rows).  Verified by
complete enumeration at $n=6$
($\Delta N^{(5)}=192=288-24-72$) and $n=7$
($13008=15384+2496-4872$); and at $n=8$ the pieces computed on $B_0$
(\texttt{anatomy5.c}: $3\mathrm{Br}^{(5)}_G=2649312$,
$3\,\Delta G_2^{(5)}=159840$, $6\,\Delta H^{(5)}=-596544$) sum to
$2212608$, exactly the value of $\Delta N^{(5)}$ from the independent
direct enumeration of $N^{(5)}$ (\texttt{n5\_direct2.c}); likewise at
$n=9$: $1632687552+69975336-375847080=1326815808=\Delta N^{(5)}$.
\end{proof}

\begin{remark}[$\Psi$ acts diagonally]
The involution $\Psi$ of Remark~\ref{rem:psi} preserves $A_0$ and
disjointness, hence acts on triples coordinatewise; consequently
$G^{(5)}(d_1)=G^{(5)}(d_2)$ \emph{exactly} (confirmed at $n=7,8,9$:
e.g.\ $G^{(5)}(d_1)=G^{(5)}(d_2)=357845442448$ at $n=9$).
\end{remark}

\subsection{The stratified bracket; spectator factors}

Expand $G^{(5)}(z)$ by Proposition~\ref{prop:tripleIE} on the board
$B_0$ with the pin $z$ carried by $\tau$'s factor.  In the notation of
\eqref{eq:tripleIEU} ($A$ = $\tau$'s constraint set, $B,C$ the other
two),
\begin{equation}\label{eq:BrG5strata}
\mathrm{Br}^{(5)}_G=\sum_{A}\,B_A\,w(A),
\qquad
w(A):=\!\!\sum_{\substack{(M_1,M_2,M_3)\\ M_1\cup M_2=A}}\!\!
(-1)^{|M_1|+|M_2|+|M_3|}R_0(B)\,R_0(C),
\end{equation}
where $B_A=\sum_{z\in\{d_1,d_2\}}R_0^\dagger(A\cup z)
-\sum_{z\in\{e_1,e_2\}}R_0^\dagger(A\cup z)$ is \emph{the same}
four-term bracket as in \eqref{eq:strata}, now indexed by $A$.  Two
exact identities anchor \eqref{eq:BrG5strata}:
\begin{itemize}[leftmargin=2em]
\item[(i)] $\sum_A R_0(A)\,w(A)=N^{(5)}_0$ (the regrouped
\eqref{eq:tripleIEU} on $B_0$);
\item[(ii)] $w(\emptyset)=\sum_{M_3}(-1)^{|M_3|}R_0(M_3)^2=N^{(4)}_0$
(the pair inclusion--exclusion on $B_0$), so the $A=\emptyset$ stratum
equals $M_{n-4}\,N^{(4)}_0$ \emph{exactly} by
Proposition~\ref{prop:S0}.
\end{itemize}
(Both identities verified exactly on $B_0$ at $n=6,7$;
\texttt{switch5.py}, \texttt{check3}.)
Since $B_A$ is unchanged, the entire bracket toolkit of
Section~\ref{sec:switch} --- Lemma~\ref{lem:wronskian},
Lemma~\ref{lem:valuation}, Lemma~\ref{lem:master},
Proposition~\ref{prop:generic} --- applies verbatim in the $A$-slot.
What is new is the weight: $R_0(M)$ is replaced by the spectator weight
$w(A)$.  The transfer rests on two observations.

\emph{First observation.}  Every cell of $A$ lies in $B$ or $C$
(cells of $M_1$ lie in $B$, cells of $M_2$ in $C$).  Hence a line-swap
pairing $\varphi$ moving a cell of $A$ (as in Lemma~\ref{lem:pairing}
or Lemma~\ref{lem:colexchange}) moves the \emph{same} cell in exactly
one spectator factor when the cell has a single label, and in both when
it is shared; the third factor is untouched in the first case.  The
pinned-factor equalities of the pairing lemmas --- which involve only
$B_A$, through the deleted row- and column-\emph{sets} of the pinned
counts --- are therefore untouched, while the free-factor difference
becomes, by the product rule,
\[
R_0(B)R_0(C)-R_0(\varphi B)R_0(\varphi C)
=\bigl[R_0(B)-R_0(\varphi B)\bigr]R_0(C)
+R_0(\varphi B)\bigl[R_0(C)-R_0(\varphi C)\bigr],
\]
each bracket a single line-swap difference controlled by
Lemma~\ref{lem:wronskian} exactly as at the fourth layer, i.e.\ a
relative $O(n^{-3})$ per factor.

\emph{Second observation (spectator stability).}  The pairing must also
re-index the spectator sum over $M_3$: an $M_3$ compatible with
$(M_1,M_2)$ maps to one compatible with $\varphi(M_1,M_2)$ by the
identity map except on the $O(1/n)$-measure set of $M_3$ meeting the
two swapped lines, which is absorbed by (the $B_0$-analogue of)
Lemma~\ref{lem:linemeasure} applied inside $w$.  Thus
$|w(A)-w(\varphi A)|=O(n^{-1})\,\overline w(A)$ with
$\overline w$ the unsigned weight, and the paired-class cancellations
of Lemmas~\ref{lem:pairing}, \ref{lem:Hdiff}, \ref{lem:colexchange}
and~\ref{lem:boundary} survive with the same $O(n^{-1})$ headroom.
Finally, Lemma~\ref{lem:uniform} and Corollary~\ref{cor:moments}
apply to $R_0(B)$, $R_0(C)$ individually (any max-degree-2 board), so
the unsigned sums $\sum_A|B_A|\,\overline w(A)$ obey the same stratum
bounds as at the fourth layer with $N_0^2$ replaced by an
$N^{(5)}_0$-scale; the $G_2$- and $H$-differences of
\eqref{eq:switchIE5} are handled by the same pairings with one or two
pins carried by spectator slots.

These two observations are the only new content; we record the outcome.

\begin{theorem}[fifth layer]\label{thm:layer5}
For $n\ge n_0$,
\[
\log\frac{N^{(5)}_n((2,n{-}2))}{N^{(5)}_n((n))}
=3\,\frac{M_{n-4}}{M_n}+O(n^{-5}),
\qquad\text{hence}\qquad
r_5(n)=r_4(n)\bigl(1+O(n^{-1})\bigr)=r_3(n)\bigl(1+O(n^{-1})\bigr).
\]
\end{theorem}

\begin{proof}
Run the proof of Theorem~\ref{thm:layer4} on \eqref{eq:switchIE5} and
\eqref{eq:BrG5strata}, with Proposition~\ref{prop:generic} in the
$A$-slot, weights $w(A)$ in place of $R_0(M)$, and the two observations
above supplying the paired-class and unsigned-sum bounds; the generic
strata resum, via identity (i), to
$\mathrm{Br}^{(5)}_G=(M_{n-4}/M_n)\,N^{(5)}_0\bigl(1+O(1/n)\bigr)$, and
the pin-pin differences of \eqref{eq:switchIE5} are
$O\bigl(n^{-1}(M_{n-4}/M_n)N^{(5)}_0\bigr)$.  Then
$N^{(5)}_0=N^{(5)}((n))(1+O(1/n))$ as before, giving
$\Delta\log N^{(5)}=3M_{n-4}/M_n+O(n^{-5})$, and
$r_5-r_3=\Delta\log N^{(5)}-\Delta\log N^{(4)}-\Delta\log N^{(3)}
=O(n^{-5})$.
\end{proof}

\subsection{Exact fifth-layer data}

Direct enumeration (64/128-bit cell masks; \texttt{n5\_direct.c},
\texttt{n5\_direct2.c}) gives $N^{(5)}$ exactly at $n\le9$; the $N^{(3)}$
and $N^{(4)}$ values reproduce the known exact tables, certifying the
enumeration.

\begin{center}
\begin{tabular}{r|rr|rrr|rr}
$n$ & $N^{(5)}((n))$ & $N^{(5)}((2,n{-}2))$ &
$r_3(n)_4$ & $r_4(n)_4$ & $r_5(n)_4$ &
$\frac{\Delta\log N^{(4)}}{2M_{n-4}/M_n}$ &
$\frac{\Delta\log N^{(5)}}{3M_{n-4}/M_n}$\\\hline
$6$ & $4032$ & $4224$ & $0$ & $2.454$ & $14.29$ & --- & ---\\
$7$ & $2185152$ & $2198160$ & $1.450$ & $2.560$ & $0.976$ & $1.382$ & $1.146$\\
$8$ & $1539331584$ & $1541544192$ & $0.709$ & $0.842$ & $0.862$ & $1.094$ & $1.134$\\
$9$ & $1417404643008$ & $1418731458816$ & $0.906$ & $0.948$ & $0.976$ & $1.023$ & $1.041$\\
\end{tabular}
\end{center}

\noindent Here $r_k(n)_4$ abbreviates $r_k\cdot n(n{-}1)(n{-}2)(n{-}3)$,
and $r_k=\Delta_\lambda\log(N^{(k)}/N^{(k-1)})$.  At $n=8,9$:
$r_5/r_4=1.023$, $1.030$ --- the fifth-layer effect tracks the fourth
to $3\%$ already at $n=9$ (small $n=6,7$ are off-trend, as for every
other layer statistic).  On the anatomy side,
$\mathrm{Br}^{(5)}_G=5128$, $883104$, $544229184$ at $n=7,8,9$ with
$\mathrm{Br}^{(5)}_G N_0/(M_{n-4}N^{(5)}_0)=0.612$, $0.696$, $0.722$,
consistent with a $1-c/n$ approach with $c\approx2.5$ ---
about twice the fourth-layer constant $c\approx1.17$, as the two
spectator rows suggest.

\begin{remark}[all types at once]\label{rem:alltypes}
The direct enumeration is cheap enough at $n\le8$ to cover \emph{every}
derangement type (\texttt{n5\_types.c}).  Writing
$\Delta_\lambda=\log N^{(k)}(\lambda)-\log N^{(k)}((n))$, the exact
$n=8$ data are:
\begin{center}
\begin{tabular}{l|rrr|cc}
$\lambda$ & $\Delta_\lambda\log N^{(3)}$ & $\Delta_\lambda\log N^{(4)}$
& $\Delta_\lambda\log N^{(5)}$ &
$\frac{\Delta\log N^{(4)}}{2\Delta\log N^{(3)}}$ &
$\frac{\Delta\log N^{(5)}}{3\Delta\log N^{(3)}}$\\\hline
$(2,6)$ & $+4.220\cdot10^{-4}$ & $+9.234\cdot10^{-4}$ & $+1.436\cdot10^{-3}$ & $1.094$ & $1.134$\\
$(3,5)$ & $0$ & $-6.95\cdot10^{-5}$ & $-2.14\cdot10^{-4}$ & --- & ---\\
$(4,4)$ & $+4.220\cdot10^{-4}$ & $+8.570\cdot10^{-4}$ & $+1.314\cdot10^{-3}$ & $1.015$ & $1.038$\\
$(2,2,4)$ & $+1.266\cdot10^{-3}$ & $+2.714\cdot10^{-3}$ & $+4.213\cdot10^{-3}$ & $1.072$ & $1.110$\\
$(2,3,3)$ & $+4.220\cdot10^{-4}$ & $+7.865\cdot10^{-4}$ & $+1.022\cdot10^{-3}$ & $0.932$ & $0.807$\\
$(2,2,2,2)$ & $+2.951\cdot10^{-3}$ & $+6.384\cdot10^{-3}$ & $+1.006\cdot10^{-2}$ & $1.082$ & $1.136$\\
\end{tabular}
\end{center}
Three reads.  (i) \emph{Uniformity over types}: every $2$-cycle-bearing
type shows the same finite-$n$ factor $\approx1.09$ (layer 4) and
$\approx1.11\text{--}1.14$ (layer 5) as the adjacent pair
$(2,6)$ vs.\ $(8)$, with an $O(1/n)$-size spread --- the per-type version
of Theorems~\ref{thm:layer4}--\ref{thm:layer5}.  (ii) \emph{The
$M_{n-2\ell}$ pattern}: a $3$-cycle has $M_{n-6}=M_2=0$, and indeed its
layer-4,5 effects are not zero but an order smaller (and negative) than
the $2$-cycle effects.  (iii) \emph{The $M_0=2$ convention}: for
$\lambda=(4,4)$ the leading bracket is $M_{n-2\ell}=M_0=2$, the same
value as the $2$-cycle's $M_4=2$ at $n=8$ --- and indeed $(4,4)$ tracks
$(2,6)$ closely at layers 3, 4 \emph{and} 5, a three-layer confirmation
of the convention entering Lemma~\ref{lem:fusionexp}.  The deviation of
$(2,3,3)$ is itself structured: its layer-4 effect equals
$\Delta\log N^{(4)}(2,6)+2\,\Delta\log N^{(4)}(3,5)$ to $2\cdot10^{-6}$
($0.3\%$ of its value; $1.3\cdot10^{-5}$, i.e.\ $1.3\%$, at layer 5) ---
additivity over cycles at the fourth and fifth layers, within the
accuracy that the neglected $\ell\ge5$ cycle terms allow one to test at
a single $n$.
\end{remark}

\begin{remark}[general $k$; the profile at the origin]\label{rem:layerk}
By induction on $k$, the identical transfer (a $(k{-}2)$-tuple
inclusion--exclusion; one more spectator factor per layer) gives, for
each \emph{fixed} $k$,
\[
\Delta_\lambda\log N^{(k)}=(k-2)\,\frac{M_{n-4}}{M_n}+O_k(n^{-5}),
\qquad
r_k(n)=r_3(n)\bigl(1+O_k(1/n)\bigr),
\]
with constants growing with $k$ (each spectator contributes its own
$O(1/n)$ bookkeeping).  In the language of the layer profile
$\gamma$ of Section~\ref{sec:layers}, $\gamma$ is flat at the origin:
every fixed layer contributes the same $(1+o(1))\,M_{n-4}/M_n$ per
$2$-cycle.  The uniformity in $k$ needed for the bulk of the profile
(layers $k\asymp tn$) remains the province of the cluster expansion
(TODO~4): there the boards are no longer two paths and the calculus
here does not apply.
\end{remark}


%======================================================================
\section{Three-row rectangles: exact moments and a joint Poisson limit}\label{sec:threerow}

The same calculus resolves the ``cycle-type-refined $K(3,n)$'' question at
the level of $2$-cycles.  Let $R$ be a uniform $3\times n$ Latin rectangle
and let $X_{12},X_{13},X_{23}$ be the numbers of intercalates between the
three pairs of rows (equivalently, the $2$-cycle counts of the three row
permutations $\sigma_{ab}$).

\begin{lemma}[exchangeability]\label{lem:exch}
$(X_{12},X_{13},X_{23})$ is exchangeable under the full symmetric group on
the three pair labels.
\end{lemma}

\begin{proof}
Permuting the rows of a rectangle is a bijection of rectangles and induces
the corresponding permutation of the three unordered pairs.
\end{proof}

\begin{theorem}[joint Poisson limit]\label{thm:poisson}
For each fixed $(a,b,c)$,
\[
\E\bigl[(X_{12})_a\,(X_{13})_b\,(X_{23})_c\bigr]
=\Bigl(\frac12\Bigr)^{a+b+c}\bigl(1+O_{a,b,c}(n^{-1})\bigr),
\]
so $(X_{12},X_{13},X_{23})$ converges jointly to three independent
$\mathrm{Poisson}(1/2)$ variables.
\end{theorem}

The proof rests on two uniform estimates, both immediate from
Lemma~\ref{lem:uniform}, whose Bonferroni proof applies verbatim to any
board of maximum degree $2$ (unions of paths \emph{and cycles}).

\begin{lemma}[uniform third-row estimates]\label{lem:thirdunif}
Uniformly over derangement types, $N^{(3)}(\sigma)=n!\,e^{-2}(1+O(1/n))$;
and for every partial matching $M$ avoiding $B_2(\sigma)$ with
$|M|=j\le n/2$,
\[
R(M;\sigma)=(n-j)!\,e^{-E_M/(n-j)}\bigl(1+O(\tfrac1{n-j})\bigr)
=(n-j)!\,e^{-2}\bigl(1+O(\tfrac{j+1}{n})\bigr),
\]
where $E_M\in[2n-4j,2n-2j]$ is the number of cells of the punctured
board.  In particular
$q(M):=R(M;\sigma)/N^{(3)}(\sigma)\le C\,e^{6j/n}/(n)_j$.
\end{lemma}

\begin{proof}[Proof of Theorem~\ref{thm:poisson}]
Normalise the first row to $\mathrm{id}$; the uniform measure is then
uniform over pairs $(\sigma,\tau)$ with $\sigma\in\DD_n$ and $\tau$
avoiding $B_2(\sigma)$, and $X_{12}=\lambda_2(\sigma)$,
$X_{13}=\lambda_2(\tau)$, $X_{23}=\lambda_2(\tau\sigma^{-1})$.  Since
distinct $2$-cycles of a permutation are disjoint,
\[
\E\bigl[(X_{12})_a(X_{13})_b(X_{23})_c\bigr]
=\sum_{(A,B,C)}\Pr\bigl(A\subseteq\mathrm{cyc}_2(\sigma),\
B\subseteq\mathrm{cyc}_2(\tau),\ C\subseteq\mathrm{cyc}_2(\tau\sigma^{-1})\bigr),
\]
the sum over ordered tuples $A$ of $a$ pairwise disjoint $2$-sets, and
similarly $B$ ($b$ sets) and $C$ ($c$ sets).  A $2$-cycle $\{u,v\}$ of
$\tau\sigma^{-1}$ prescribes the $\tau$-cells
$(\sigma^{-1}(u),v),(\sigma^{-1}(v),u)$; a $2$-cycle $\{i,j\}$ of $\tau$
prescribes $(i,j),(j,i)$.

\emph{Generic configurations.}  Call $(A,B,C)$ generic if the
$2(a+b+c)$ indices appearing in $A\cup B\cup C$ are pairwise distinct.
There are $\binom n2^{a+b+c}(1+O_{a,b,c}(1/n))$ generic tuples.  Fix one.
First, $\Pr(A\subseteq\mathrm{cyc}_2(\sigma))=D_{n-2a}/D_n
=(n)_{2a}^{-1}(1+O(1/(n-2a)!))$.  Conditionally on
$\{A\subseteq\mathrm{cyc}_2(\sigma)\}$, the derangement $\sigma$
restricted to the complement of the support of $A$ is uniform; the
$B$- and $C$-prescriptions form a matching $M(\sigma)$ of exactly
$2(b+c)$ $\tau$-cells avoiding $B_2(\sigma)$ \emph{unless} $\sigma$
satisfies one of $O_{b,c}(1)$ degeneracy events, each of conditional
probability $O(1/n)$: a prescribed cell may land on $B_2(\sigma)$
(e.g.\ $\sigma(i)=j$ for a $B$-pair $\{i,j\}$, or $v=\sigma^{-1}(u)$ for
a $C$-pair), or two prescriptions may conflict in a row or column
(e.g.\ $\sigma^{-1}(u)$ in the support of $B$); for a generic tuple all
index slots are distinct, so no two prescribed cells can \emph{coincide},
and on every degeneracy event the prescriptions are inconsistent,
$q=0$.  Off these events, Lemma~\ref{lem:thirdunif} gives
$\Pr(\tau\supseteq M(\sigma)\mid\sigma)=q(M(\sigma))
=(n)_{2(b+c)}^{-1}(1+O(1/n))$.  Hence each generic tuple contributes
$(n)_{2a}^{-1}(n)_{2(b+c)}^{-1}(1+O_{a,b,c}(1/n))$, and the generic sum
is $\binom n2^{a+b+c}(n)_{2a}^{-1}(n)_{2(b+c)}^{-1}(1+O(1/n))
=2^{-(a+b+c)}(1+O_{a,b,c}(1/n))$.

\emph{Collision configurations.}  The remaining tuples are classified by
their collision pattern: which equalities hold among the $2(a+b+c)$
index slots --- finitely many patterns for fixed $(a,b,c)$.  A pattern
with $s\ge1$ independent equalities has $O(n^{2(a+b+c)-s})$ tuples, so it
suffices to show that the per-tuple contribution is
$O(n^{-2(a+b+c)}\cdot n^{\,s-1})$.  Bound it as follows:
$\Pr(A\subseteq\mathrm{cyc}_2(\sigma))\le C n^{-2a}$; conditionally on
$\sigma$, if the prescriptions are consistent they determine a matching
of $2(b+c)-t$ cells, where $t$ counts coincidences among prescribed
cells, and Lemma~\ref{lem:thirdunif} gives
$q\le Cn^{-2(b+c)+t}$.  Two $B$-cells, or two $C$-cells, never coincide
(the pairs within $B$, within $C$, are disjoint); a coincidence of the
$B$-cell $(i,j)$ with the $C$-cell $(\sigma^{-1}(u),v)$ requires the
index equality $v=j$ (counted in $s$) \emph{and} the $\sigma$-equation
$u=\sigma(i)$, of conditional probability $O(1/n)$; and $t$ such
coincidences involve $t$ distinct such $\sigma$-equations, of joint
probability $O(n^{-t})$.  (A full coincidence of a $B$-pair with a
$C$-pair is impossible: it would force $\sigma(i)=i$.)  Hence the
per-tuple bound is $Cn^{-2a}\cdot n^{-t}\cdot n^{-2(b+c)+t}
=Cn^{-2(a+b+c)}$, which beats the required bound by $n^{s-1}\ge1$
--- with room to spare unless $s=1$, where it is exactly what is needed.

Summing, $\E[(X_{12})_a(X_{13})_b(X_{23})_c]
=2^{-(a+b+c)}(1+O_{a,b,c}(1/n))$, and the joint Poisson limit follows by
the method of moments.
\end{proof}

The toolkit computes the low joint moments \emph{exactly} (certified against
brute-force enumeration at $n=7$):

\begin{center}
\small
\begin{tabular}{c c c c c c}
\toprule
$n$ & $\E X_{12}-\tfrac12$ & ratio to $\dfrac{M_{n-4}}{2M_n}$ &
$\E[X_{12}X_{13}]$ & $n^2\operatorname{Cov}(X_{12},X_{13})$ &
$\E\binom{X_{13}}2$\\
\midrule
8  & $+5.162\times10^{-4}$ & $2.45$ & $0.26230$ & $0.754$ & $0.12761$\\
9  & $+1.143\times10^{-4}$ & $0.76$ & $0.25873$ & $0.698$ & $0.12466$\\
10 & $+9.467\times10^{-5}$ & $1.041$ & $0.25689$ & $0.679$ & $0.12511$\\
11 & $+5.886\times10^{-5}$ & $0.9944$ & $0.25552$ & $0.661$ & $0.12502$\\
12 & $+4.005\times10^{-5}$ & $1.0010$ & $0.25453$ & $0.646$ & $0.12502$\\
13 & $+2.797\times10^{-5}$ & $1.0000$ & $0.25378$ & $0.633$ & $0.12501$\\
\bottomrule
\end{tabular}
\end{center}

\noindent
(All three $\E X_{ab}$ agree exactly, as do the three mixed products ---
Lemma~\ref{lem:exch} at the level of exact rationals.)  Two structural
features deserve note.

The next two statements need a quantitative form of the fusion
expansion of $N^{(3)}$, of independent interest: it proves the
additivity of $\log N^{(3)}$ over cycles (\S\ref{sec:thirdrow}) with an
explicit second-order remainder.

\begin{lemma}[fusion expansion with remainder]\label{lem:fusionexp}
There are $n_0,C$ such that for all $n\ge n_0$ and every derangement
type $\lambda\vdash n$ with $\kappa=\kappa(\lambda)$ parts,
\[
N^{(3)}(\lambda)\;=\;M_n+\sum_{\ell\ge2}\lambda_\ell\,M_{n-2\ell}
+E(\lambda),\qquad |E(\lambda)|\;\le\;C\,\kappa^2\,(n-8)!\,,
\]
with the conventions $M_0=2$, $M_1=-1$ (the IE values
$\langle M_{C_{2m}}\rangle_m$) and $M_m=0$ for $m<0$.  Exhaustive
computation over all types gives
$\max_\lambda|E(\lambda)|/(\kappa^2(n-8)!)\approx0.06$ at $n=12,16,20$,
and $E(\lambda)\ge0$ except for the single type $(\tfrac n2,\tfrac n2)$,
where $E=-2$.
\end{lemma}

\begin{proof}
Induction on $\kappa$; the case $\kappa=1$ is trivial ($E=0$).  For
$\kappa\ge2$ pick two parts as follows: two equal parts $b=a$ if any
exist, otherwise the smallest and the largest, $b<a$ (then $a-b\ge2$
unless $\kappa=2$ and $\lambda=(m,m+1)$, a single boundary case computed
directly from $M_{C_2}=1+2x$: there
$N^{(3)}=M_n+\langle x^{2m}(1+2x)\rangle_n=M_n-1$ and $E=0$ under the
stated conventions).  Theorem~\ref{thm:thirdrow} gives
$N^{(3)}(\lambda)=N^{(3)}(\mu)+N^{(3)}_{n-2b}(\nu)$, where $\mu$
replaces $\{a,b\}$ by $a+b$ and $\nu$ (on $n-2b$ points) replaces them
by $a-b$ (for $a=b$: $2N^{(3)}_{n-2b}(\lambda\smallsetminus\{a,b\})$,
same bookkeeping).  Apply the induction hypothesis to both terms.  In
the sum, the first-order terms reassemble as
$\sum_\ell\lambda_\ell M_{n-2\ell}$ \emph{plus} the correction terms
\[
M_{n-2(a+b)}\;+\;\sum_{\ell\text{-parts of }\nu\smallsetminus\{a-b\}}
M_{(n-2b)-2\ell}\;\le\;M_{n-8}+\kappa\,M_{n-8}
\]
(each index is at most $n-8$ because $a,b,\ell\ge2$), the term
$M_{(n-2b)-2(a-b)}=M_{n-2a}$ cancelling the deficit of $\mu$'s
first-order sum.  Hence
$|E(\lambda)|\le|E(\mu)|+|E(\nu)|+C'(\kappa+1)(n-8)!
\le C(\kappa-1)^2(n-8)!+C\kappa^2(n-2b-8)!+C'(\kappa+1)(n-8)!$,
and since $(n-2b-8)!\le(n-12)!\le(n-8)!/n^3$, the right side is at most
$C\kappa^2(n-8)!$ once $C\ge2C'$ and $n\ge n_0$.
\end{proof}

We also record the standard moment computation for derangement cycle
counts: for $s_\ell\ge0$ with $m=\sum_\ell\ell s_\ell\le n/2$,
\begin{equation}\label{eq:derangemoments}
\E_{Q_n}\Bigl[\prod_\ell(\lambda_\ell)_{s_\ell}\Bigr]
=\prod_\ell\ell^{-s_\ell}\cdot\frac{(n)_m D_{n-m}}{D_n}
=\prod_\ell\ell^{-s_\ell}\,\Bigl(1+O\bigl(\tfrac1{(n-m)!}\bigr)\Bigr),
\end{equation}
by counting derangements with an ordered tuple of marked cycles.

\begin{proposition}[the third row tilts the second]\label{prop:tilt}
$\displaystyle
\E[X_{12}]-\tfrac12=\frac{M_{n-4}}{2M_n}\,\bigl(1+O(n^{-2})\bigr).$
\end{proposition}

\begin{proof}
$X_{12}=\lambda_2(\sigma)$ where $\Pr(\sigma)\propto N^{(3)}(\sigma)$ on
$\DD_n$, so with $w(\lambda)=N^{(3)}(\lambda)/M_n$ and $\E_Q$ the
uniform-derangement mean,
\[
\E[X_{12}]-\E_Q[\lambda_2]
=\frac{\operatorname{Cov}_Q(\lambda_2,w)}{\E_Q[w]} .
\]
By Lemma~\ref{lem:fusionexp}, $w=1+\sum_\ell\rho_\ell\lambda_\ell+
E(\lambda)/M_n$ with $\rho_\ell=M_{n-2\ell}/M_n$.  By
\eqref{eq:derangemoments}, $\operatorname{Var}_Q(\lambda_2)=\tfrac12+
O(1/(n-4)!)$ and $\operatorname{Cov}_Q(\lambda_2,\lambda_\ell)=
O(1/(n-2-\ell)!)$ for $\ell\ge3$ (both are differences of products of
factorial moments, and the Poisson values cancel exactly).  For the
remainder, $\kappa\le n/2$ gives the crude bound
$|\operatorname{Cov}_Q(\lambda_2,E/M_n)|\le C n^2(n-8)!/M_n=O(n^{-6})$.
Hence
\[
\operatorname{Cov}_Q(\lambda_2,w)=\rho_2\cdot\tfrac12+O(n^{-6}),
\qquad
\E_Q[w]=1+O(n^{-4}),
\]
and $\E_Q[\lambda_2]=\tfrac12+O(1/(n-2)!)$ by
\eqref{eq:derangemoments}; assembling,
$\E[X_{12}]-\tfrac12=\tfrac12\rho_2(1+O(n^{-2}))$, as claimed (the
$O(n^{-2})$ is dominated by $O(n^{-6})/\rho_2$).
The exact table above confirms the formula to $10^{-4}$ relative by
$n=13$ --- a sharp quantitative cross-check of the layer framework:
\emph{conditioning on one extra row tilts the second row's $2$-cycle
count by exactly the layer-3 coefficient times the variance}.
\end{proof}

\begin{proposition}[pair--pair covariance]\label{prop:cov}
$\displaystyle \operatorname{Cov}(X_{12},X_{13})
=\frac{1/2+O(n^{-1}\log n)}{n^{2}}\;>0 .$
\end{proposition}

\begin{proof}
Write $h(\sigma)=\E[X_{13}\mid\sigma]=\sum_{\{i,j\}\,\mathrm{adm}}
q_{ij}(\sigma)$, where $q_{ij}=R(M_{ij};\sigma)/N^{(3)}(\sigma)$,
$M_{ij}=\{(i,j),(j,i)\}$, and the sum runs over the \emph{admissible}
pairs (both cells off $B_2(\sigma)$); their number is
$\binom n2-n+\lambda_2(\sigma)$ exactly (the excluded pairs are the
$\{i,\sigma(i)\}$, of which $\lambda_2$ coincide in pairs).  A $2$-cycle
$\{i,j\}$ of $\sigma$ thus makes $\{i,j\}$ doubly forbidden but
\emph{frees} the two pairs that $i$ and $j$ would otherwise each
exclude.  Split
\[
h(\sigma)=\bigl[\tbinom n2-n+\lambda_2(\sigma)\bigr]\,(n)_2^{-1}
+\Delta(\sigma),\qquad
\Delta(\sigma)=\sum_{\mathrm{adm}}\bigl[q_{ij}(\sigma)-(n)_2^{-1}\bigr],
\]
so that, under the $N^{(3)}$-tilted measure of the three-row ensemble,
\[
\operatorname{Cov}(X_{12},X_{13})
=(n)_2^{-1}\operatorname{Var}(\lambda_2)
+\operatorname{Cov}(\lambda_2,\Delta).
\]
As in Proposition~\ref{prop:tilt} (second factorial moments via
Lemma~\ref{lem:fusionexp} and \eqref{eq:derangemoments}),
$\operatorname{Var}(\lambda_2)=\tfrac12+O(n^{-2})$.

$\Delta(\sigma)$ depends only on the cycle type, and we claim the
\emph{join-stability} bound: for types $\lambda,\mu$ one join apart,
$|\Delta(\lambda)-\Delta(\mu)|\le Cn^{-3}$.  Choose adjacent
representatives as in \S\ref{sec:switch} (boards differing in one
$2\times2$ switch).  For each admissible pair,
$R(M_{ij};\sigma')-R(M_{ij};\sigma)$ equals the four-term bracket
$\mathrm{Br}(M_{ij})$ of \S\ref{sec:switch} (the double-pin terms cancel
as identical punctured boards, and if $M_{ij}$ meets a pinned line the
bracket degenerates to the compatible terms).  By
Lemmas~\ref{lem:valuation} and~\ref{lem:master}: for the
$O(n^2)$ pairs avoiding the switch window the bracket is $x^3$-shifted
and pinned, giving relative differences $O(n^{-4})$; for the $O(n)$
pairs meeting one window line, valuation $\ge2$ gives $O(n^{-3})$; for
the $O(1)$ pairs meeting two pinned lines the crude one-term bound gives
$O(n^{-1})$.  Together with the $N^{(3)}$-normalisation
($\Delta_\lambda\log N^{(3)}=O(n^{-4})$) and the $O(1)$ pairs entering
or leaving the admissible set (each contributing
$q-(n)_2^{-1}=O(n^{-3})$), this sums to
$|\Delta(\lambda)-\Delta(\mu)|\le C(n^2\cdot n^{-4}+n\cdot n^{-3}+1\cdot
n^{-1})\cdot(n)_2^{-1}+O(n^{-4})\le Cn^{-3}$.

Any type with $\kappa$ parts is $\kappa-1$ joins from $(n)$, so
$|\Delta(\lambda)-\Delta((n))|\le C\kappa n^{-3}$, and since the tilted
density is bounded, $\E[\kappa^2]\le C\log^2 n$; Cauchy--Schwarz gives
$|\operatorname{Cov}(\lambda_2,\Delta)|\le Cn^{-3}\log n$.  Assembling,
$\operatorname{Cov}(X_{12},X_{13})=\tfrac12(n)_2^{-1}(1+O(n^{-1}\log n))$.
The exact values $n^2\mathrm{Cov}=0.661,\,0.646,\,0.633$ at $n=11,12,13$
extrapolate to $\approx0.48+2.0/n$, consistent with the constant
$\tfrac12$ at these small orders.
\end{proof}

\begin{remark}
Theorem~\ref{thm:poisson} is the $k=3$ instance of the expectation that
intercalate structure across any $O(1)$ rows is asymptotically that of an
independent model; it may be seen as the (easy) rectangle-regime companion
of the full-square intercalate results of \cite{KSSS22}.  We have not found
it stated in the literature; the classical aggregate $K(3,n)$ theory
(\cite[Ch.~8 \S4]{Riordan58}) does not resolve joint pair statistics.
\end{remark}


%======================================================================
\section{An exact switching identity and the B-pair reformulation}\label{sec:identity}

The switchings of \cite{CGW08} do more than bound ratios: their multigraph
is exactly balanced on one side, and this yields an \emph{identity}.  We
briefly recall the notions (see \cite[\S2--3]{CGW08} for details).

Fix a Latin square $L$ (of order $n$; we take the CGW parameters $m=n$,
$F=\emptyset$).  The first two rows decompose into \emph{row cycles}
(corresponding to the cycles of $\sigma_{1,2}$).  For two columns $j\ne j'$
with $j'\notin\{\omega(j),\omega^{-1}(j)\}$ --- where $\omega$ is the column
map defined by $L(1,\omega(j))=L(2,j)$ --- the pair $\{j,j'\}$ is an
\emph{$A$-pair} if rows $1,2$ lie in different column cycles of the column
pair, and a \emph{$B$-pair} otherwise.  The \emph{switch} of $L$ at
$\{j,j'\}$ trades the row cycle (among the two through $j,j'$) containing
the minimal symbol; the \emph{flip} at an $A$-pair $\{j,j'\}$ whose columns
lie in different row cycles of lengths $\alpha,\beta$ trades the column
cycle of $\{j,j'\}$ through row 2, and joins the two row cycles into one of
length $\alpha+\beta$.

Let $\mu\vdash n$ contain parts $\alpha\neq\beta$ ($\alpha,\beta\ge2$) and
let $\lambda$ be $\mu$ with $\alpha,\beta$ merged into $\alpha+\beta$.  Let
$L$ be uniform on $S(\mu)=\{L:\ \sigma_{1,2}(L)\text{ has type }\mu\}$,
pick an $\alpha$-cycle and a $\beta$-cycle of $L$ uniformly, and columns
$j,j'$ uniformly in them.  Perform the CGW joining step:
if $\{j,j'\}$ is a $B$-pair, switch at $\{j,j'\}$ (making it an $A$-pair);
then flip at $\{j,j'\}$, obtaining $L'\in S(\lambda)$ in which $\{j,j'\}$
lies in a single $(\alpha+\beta)$-cycle with
$(\omega')^{\beta}(j)=j'$, where $\omega'=\omega_{L'}$.  Define the two
\emph{post-flip events}
\[
B_1=\bigl\{\{\omega'^{-1}(j),\omega'^{-1}(j')\}\text{ is a $B$-pair of
}L'\bigr\},
\qquad
B_2=\bigl\{\{\omega'(j),\omega'(j')\}\text{ is a $B$-pair of }L'\bigr\},
\]
(the shifted pairs are automatically admissible when $\alpha\ne\beta$), and
put $\bar q=\bar q(\mu\to\lambda)=\Pr(B_1)+\Pr(B_2)\in[0,2]$.

\begin{theorem}[exact identity; $\alpha\ne\beta$]\label{thm:identity}
With the notation above,
\[
\frac{|S(\lambda)|}{|S(\mu)|}
=\frac{\gamma(\lambda)}{\gamma(\mu)}\cdot\frac{1+\bar q}{2},
\qquad\text{equivalently}\qquad
\boxed{\;\bar q \;=\; 2\,\frac{\CC_n(\lambda)}{\CC_n(\mu)}\;-\;1.\;}
\]
\end{theorem}

\begin{proof}[Proof (source and verification)]
This is the exact content of the CGW splitting/joining multigraph
$G_S=G_J$ \cite[\S3]{CGW08}: the splitting side is \emph{regular} --- every
$L'\in S(\lambda)$ produces exactly two edges for each of its
$\lambda_{\alpha+\beta}$ cycles of length $\alpha+\beta$ and each of the
$\alpha+\beta$ column pairs at $\omega$-distance $\beta$ within the cycle
(\cite[Eq.~(2)]{CGW08}) --- while the joining side produces, for each of the
$\mu_\alpha\mu_\beta\alpha\beta$ instances, $1+\mathbf 1_{B_1}+\mathbf
1_{B_2}$ edges.  Equating the edge counts and using
$\gamma(\lambda)/\gamma(\mu)=
\mu_\alpha\mu_\beta\alpha\beta/((\mu_{\alpha+\beta}+1)(\alpha+\beta))$ gives
the identity.  Because the bookkeeping is delicate, we certified the exact
statement independently: by \emph{complete enumeration} of all Latin squares
of orders $n=5$ ($\mu=(2,3)$: both sides equal $2$) and $n=6$
($\mu=(2,4)$: both sides equal $10/11$, from $3{,}041{,}280$ instances), and
statistically at $n=8$ (\S\ref{sec:numerics}).  The min-symbol tie-break in
the switch is averaged over symbol relabellings, which is the invariant
formulation of the CGW rule.
\end{proof}

\begin{remark}[the case $\alpha=\beta$]\label{rem:albe}
For $\alpha=\beta=2$ the two shifted pairs coincide and, with the event
counted twice, the identity again holds exactly (verified at $n=6$,
$\mu=(2,2,2)$: both sides $4/7$).  For $\alpha=\beta\ge3$ our reading of the
joining bookkeeping does \emph{not} reproduce the splitting count (at $n=6$,
$\mu=(3,3)$: joining average $5/6$ against the required $3/4$), so the exact
identity in that case needs a corrected accounting; we do not use it ---
see the routing argument in \S\ref{sec:reduction}.  (The CGW
\emph{inequalities} are unaffected.)
\end{remark}

\begin{lemma}[symmetry of the two events]\label{lem:involution}
$\Pr(B_1)=\Pr(B_2)$ exactly, for every $n,\mu,\alpha,\beta$.  Hence
$\Pr(B_i)=\bar q/2 = \CC_n(\lambda)/\CC_n(\mu)-\tfrac12$.
\end{lemma}

\begin{proof}[Proof sketch]
The involution $\iota$ that transposes rows $1$ and $2$ preserves $S(\mu)$
and the instance measure (it fixes the row-cycle column sets, the $A/B$
status of every pair, and the min-symbol rule, and conjugates
$\omega\mapsto\omega^{-1}$), and it exchanges the two shifted pairs, mapping
the event $B_1$ for the instance of $L$ onto the event $B_2$ for the
corresponding instance of $\iota(L)$.  We verified the exchange
\emph{instance-by-instance} over $7{,}000$ instances at $n=6$ (both $A$- and
$B$-instances, both switch variants); the exact equalities
$\Pr(B_1)=\Pr(B_2)=5/11,\,5/12,\,2/7$ in the three $n=6$ enumerations confirm
it in aggregate.  A self-contained algebraic proof (chasing $\iota$ through
the flip, whose defining column cycle is the one through row 2) is routine
but fiddly; we record it as such.
\end{proof}

Theorem~\ref{thm:identity} converts the CGW exact tables into exact values
of a structural probability of random Latin squares.  For the join
$(2,n-2)\to(n)$:
\[
\begin{array}{l|lllll}
n & 7 & 8 & 9 & 10 & 11\\\hline
\bar q = 2\CC_n((n))/\CC_n((2,n-2))-1
& 0.988372 & 0.993345 & 0.996822 & 0.997363 & 0.998080\\
\tfrac12(1-\bar q)\,n^{3} & 1.99 & 1.70 & 1.16 & 1.32 & 1.28
\end{array}
\]
Thus each post-flip event has probability
$\tfrac12-\Theta(n^{-3})$: the deficit is precisely the refined-law
coefficient.  By contrast a \emph{generic} admissible column pair is a
$B$-pair with probability $\tfrac12+\Theta(n^{-2})$
(\S\ref{sec:numerics}), and two marked points of a uniform random
derangement lie in a common cycle with probability
$\tfrac12+(1+o(1))/n$.  The post-flip pairs are extraordinarily well
balanced, and this balance \emph{is} the conjecture:

\begin{hypothesis}[B-pair equidistribution, $\mathrm{EQ}(\delta)$]\label{hyp:eq}
For every pair $\mu\to\lambda$ as above with $\alpha\neq\beta$,
$\bigl|\bar q(\mu\to\lambda)-1\bigr|\le\delta(n)$.
\end{hypothesis}

By Theorem~\ref{thm:identity}, $\mathrm{EQ}(\delta)$ with
$\delta(n)=O(n^{-3+\varepsilon})$ is \emph{equivalent} to the corresponding
precision in Conjecture~\ref{conj:refined} across $\alpha\neq\beta$ joins;
$\mathrm{EQ}(o(1/\log n))$ already implies the weak conjecture
(\S\ref{sec:reduction}).  In \S\ref{sec:marked} we prove an identity
that removes the joining measure and the events $B_1,B_2$ from this
picture entirely.

%======================================================================
\section{The marked-pair identity: eliminating the post-flip events}
\label{sec:marked}

The identity of Theorem~\ref{thm:identity} expresses the completion
ratio through the post-flip events $B_1,B_2$ of the \emph{joining}
measure.  In this section we prove a cleaner exact identity that
eliminates the joining measure, the switch tie-break, and the events
$B_1,B_2$ altogether: the ratio $\CC_n(\lambda)/\CC_n(\mu)$ is the
$A$-versus-$B$ \emph{odds of a single marked column pair in a single
uniform square of} $S(\lambda)$.  The proof is self-contained (it does
not use the multigraph of \cite{CGW08}), covers $\alpha=\beta$ without
any modification --- resolving the bookkeeping discrepancy of
Remark~\ref{rem:albe} --- and yields the lower bound
$\CC_n(\lambda)/\CC_n(\mu)>\tfrac12$ for free.

Throughout, squares are normalised to have first row $\mathrm{id}$, so
the column map is $\omega=\sigma$ ($1\circ\omega(j)=2\circ j$ with the
first row the identity gives $\omega(j)=\sigma(j)$).  For the proofs it
is convenient to work with \emph{unnormalised} squares (equivalently,
symbol-relabelling classes): every instance count below then acquires a
common factor $n!$, which cancels from every ratio, and the maps we use
(flip, switch) act on unnormalised squares without any renormalisation
step.  Statuses, cycle types, and column positions are
relabelling-invariant, so we freely state results in the normalised
picture.

\begin{definition}[marked instances]\label{def:marked}
Let $\lambda\vdash n$ contain a part $m=\alpha+\beta$
($\alpha,\beta\ge2$), and let $\mu$ be $\lambda$ with one part $m$
split into $\alpha,\beta$.  The \emph{splitting instance space} is
\[
X \;=\; X(n,\lambda;\alpha,\beta)
\;=\;\bigl\{(L',P)\;:\;L'\in S(\lambda),\;
P=\{j,\sigma^\alpha(j)\}\text{ with } j\text{ in an $m$-cycle of }
\sigma(L')\bigr\},
\]
unordered pairs counted once.  Thus every $L'$ carries exactly
$\lambda_m\, m$ marked pairs when $\alpha\neq\beta$ and
$\lambda_m\,\alpha$ when $\alpha=\beta$.  Since $\alpha,\beta\ge 2$
every marked pair is admissible, and $X=X_A\sqcup X_B$ according to the
$A$/$B$ status of $P$ in $L'$.  The \emph{joining instance space} is
\[
J \;=\; J(n,\mu;\alpha,\beta)
\;=\;\bigl\{(L,P)\;:\;L\in S(\mu),\;P=\{j,j'\},\;
j\text{ in an $\alpha$-cycle},\; j'\text{ in a different
$\beta$-cycle}\bigr\},
\]
again with $J=J_A\sqcup J_B$ by the status of $P$ in $L$.
\end{definition}

The counting arithmetic is immediate: with
$D_\nu$ the number of derangements of cycle type $\nu$ and
$N(\nu)=\CC_n(\nu)(n-2)!$,
\begin{equation}\label{eq:Xarith}
\frac{2|X|}{|J|}
=\frac{2\,D_\lambda N(\lambda)\,\lambda_m m}
       {D_\mu N(\mu)\,\mu_\alpha\mu_\beta\,\alpha\beta}
=\frac{2\CC_n(\lambda)}{\CC_n(\mu)}
\qquad(\alpha\neq\beta),
\end{equation}
using $D_\lambda/D_\mu=\alpha\mu_\alpha\,\beta\mu_\beta/(m\lambda_m)$;
for $\alpha=\beta$ one has
$D_\lambda/D_\mu=\alpha^2\mu_\alpha(\mu_\alpha-1)/(2m\lambda_m)$,
$|J|=|S(\mu)|\binom{\mu_\alpha}{2}\alpha^2$ and
$|X|=|S(\lambda)|\lambda_m\alpha$, and \eqref{eq:Xarith} holds
verbatim.  The content of the identity is therefore the following
count.

\begin{lemma}[cycle trades preserve pair status]\label{lem:cycletrade}
Let $P=\{j,j'\}$ be an admissible column pair of a Latin square $L$ and
let $C$ be one column cycle of $P$.  Trading $C$ (swapping the entries
of columns $j,j'$ in the rows of $C$) inverts the pair permutation
$\pi_P$ on $C$ and fixes it elsewhere; the partition of rows into
column cycles of $P$ is unchanged.  In particular the $A$/$B$ status of
$P$ itself is invariant under the CGW flip at $P$, which trades the
full column cycle through row~$2$.
\end{lemma}

\begin{proof}
Write $\pi=\pi_P$, so $\pi(r)$ is the row of column $j'$ holding the
symbol $L(r,j)$; $C$ is $\pi$-closed and the trade swaps
$L(r,j)\leftrightarrow L(r,j')$ for $r\in C$.  Fix $r\in C$ and let
$\pi'$ be the new pair permutation.  The new $L(r,j)$ is the old
$L(r,j')$; the row of the new column $j'$ holding that symbol is found
by solving (new $L(s,j')$) $=$ (old $L(r,j')$): for $s\in C$ the new
$L(s,j')$ is the old $L(s,j)$, and old $L(s,j)=$ old $L(r,j')$ exactly
when $\pi(s)=r$.  Hence $\pi'=\pi^{-1}$ on $C$; off $C$ nothing
changes.  A cycle inverted in place has the same support, so the
partition into cycles is unchanged.
\end{proof}

\begin{lemma}[switch toggle]\label{lem:switchtoggle}
Let $(L,P)\in J$ and let $Q$ be the row cycle through one column of
$P$ (the two row cycles are distinct).  Trading $Q$ (the CGW switch)
changes $\pi_P$ by precomposition with the transposition of rows
$1,2$: the cycles of $\pi_P$ through rows $1,2$ are split if rows
$1,2$ are in one cycle, and joined if in two.  Hence the switch
toggles the status of $P$, preserves $S(\mu)$ and the instance data
(it inverts one $\sigma$-cycle in place), and is an involution:
the trade preserves the traded cycle's symbol set (the two rows carry
the same symbols on its columns), so the min-symbol rule re-selects the
same cycle.  Consequently $|J_A|=|J_B|$, exactly.
\end{lemma}

\begin{proof}
Say $Q$ is the row cycle through $j$; its columns include $j$ but not
$j'$.  The trade swaps the row-1 and row-2 entries in the columns of
$Q$; for $\pi_P$ only column $j$ matters, where the entries of rows
$1,2$ are exchanged.  Thus $\pi_P'=\pi_P\circ\tau_{12}$, and
multiplying by a transposition splits or joins the cycles through its
two points.  Rows $1,2$ lie in one cycle of $\pi_P$ exactly when $P$
is a $B$-pair.  The trade swaps the two rows' entries along the
columns of $Q$, which inverts the corresponding $\sigma$-cycle in
place: the cycle type and the column sets of all row cycles are
preserved, so $(L,P)\mapsto(\mathrm{sw}(L),P)$ maps $J$ to $J$.
\end{proof}

\begin{theorem}[marked-pair identity]\label{thm:marked}
For every $n$, $\lambda$, and split $m=\alpha+\beta$ with
$\alpha,\beta\ge2$ (including $\alpha=\beta$),
\[
|J_A|\;=\;|J_B|\;=\;|X_A|,
\qquad\text{hence}\qquad
\boxed{\;\frac{\CC_n(\lambda)}{\CC_n(\mu)}\;=\;\frac{|X|}{2|X_A|},
\qquad
2\,\frac{\CC_n(\lambda)}{\CC_n(\mu)}-1\;=\;\frac{|X_B|}{|X_A|}.\;}
\]
In particular $\CC_n(\lambda)>\tfrac12\,\CC_n(\mu)$ unconditionally,
and Hypothesis~\ref{hyp:eq} ($\mathrm{EQ}(\delta)$) is equivalent to
\[
\Bigl|\;\Pr\bigl(P\text{ is a $B$-pair}\bigr)-\tfrac12\;\Bigr|
\;\le\;\tfrac{\delta}{4}(1+o(1)),
\]
for $(L',P)$ uniform on $X$: the law of a single bit --- whether rows
$1,2$ lie in one column cycle of the marked pair
$\{j,\sigma^\alpha(j)\}$ --- of a single uniform square of
$S(\lambda)$.
\end{theorem}

\begin{proof}
By \eqref{eq:Xarith} it suffices to prove $|J|=2|X_A|$.
The flip at an $A$-pair $P$ whose columns lie in one row cycle at
$\sigma$-distances $(\alpha,\beta)$ splits that cycle into an
$\alpha$- and a $\beta$-cycle and maps $S(\lambda)\to S(\mu)$; by
Lemma~\ref{lem:cycletrade} the pair $P$ is still an $A$-pair
afterwards, so $x\mapsto\mathrm{unflip}(x)$ maps $X_A$ into $J_A$.
The flip at a fixed pair is an involution (the traded cycle again
contains row 2, with the same row support), so
$\mathrm{join}=\mathrm{flip}$ inverts it: $X_A\leftrightarrow J_A$ is
a bijection.  Composing with the switch of
Lemma~\ref{lem:switchtoggle} gives a bijection
$X_A\leftrightarrow J_B$, namely
$x\mapsto \mathrm{sw}(\mathrm{unflip}(x))$ with inverse
$(L,P)\mapsto \mathrm{flip}(\mathrm{sw}(L))$ --- the map
``switch-then-flip'', which is exactly the CGW joining step at a
$B$-pair.  Hence $|J|=|J_A|+|J_B|=2|X_A|$.
\end{proof}

\begin{remark}[what became of the events $B_1,B_2$]
Theorem~\ref{thm:marked} recovers Theorem~\ref{thm:identity} and
sharpens it: the joining measure, the min-symbol averaging, and the
post-flip events have disappeared.  The mechanism is that
\emph{every} joining edge lands in $X_A$ (Lemma~\ref{lem:cycletrade}),
each marked-$A$ instance receiving exactly two (one from an $A$-join,
one from its switched partner), while $X_B$ is never reached: the
$B_1,B_2$ bookkeeping of \cite{CGW08} was the joining-side accounting
of $|X_B|$.  On $X_B$ itself the only known measure-preserving
involution is the cross-switch of \cite[Def.~3.9]{CGW08}.  By
\cite[Lem.~3.10]{CGW08} it toggles the status of every pair
$\{\omega^{k}(j),\omega^{k'}(j')\}$ with $1\le k<\alpha$,
$1\le k'<\beta$ --- a family containing \emph{both} shifted pairs, at
$(k,k')=(1,1)$ and $(\alpha-1,\beta-1)$ --- while (empirically: all
$2880$ marked-$B$ instances at $n=5$) fixing the marked status
itself.  This yields
$\Pr(B_1\mid\text{marked }B)=\Pr(B_2\mid\text{marked }B)=\tfrac12$
exactly, explaining the extraordinary balance of the post-flip events
observed in \S\ref{sec:identity}.
\end{remark}

\begin{remark}[the case $\alpha=\beta$, and Remark~\ref{rem:albe}]
Nothing in the proof distinguishes $\alpha=\beta$: marked pairs
$\{j,\sigma^\alpha(j)\}$ in a $2\alpha$-cycle are counted once
($\lambda_m\alpha$ per square), and \eqref{eq:Xarith} holds with the
$\alpha=\beta$ conventions.  Exhaustive verification at $n=6$:
$\mu=(3,3)$ gives $|X_B|/|X_A|=622080/829440=3/4$, exactly the value
$2\CC_6((6))/\CC_6((3,3))-1$ that the joining-side bookkeeping of
Remark~\ref{rem:albe} failed to reproduce.  The $\alpha=\beta\ge3$
anomaly is thus resolved: the identity was never about the joining
events.
\end{remark}

\paragraph{Verification.}  By complete enumeration
(\texttt{pf\_bridge.py}, \texttt{xside.c}): at $n=5$,
$\mu=(2,3)$: $|J_A|=|J_B|=|X_A|=1440$, $|X_B|=2880$, ratio $2$
$=\bar q$ exactly.  At $n=6$: $(2,4)\to(6)$:
$1520640/1520640/1520640$, $|X_B|=1382400$, ratio $10/11$;
$(2,2,2)\to(2,4)$: ratio $4/7$; $(3,3)\to(6)$: ratio $3/4$ --- all
equal to $2\CC_6(\lambda)/\CC_6(\mu)-1$ exactly.  At $n=7$, fixing one
representative $\sigma_0$ per type (conjugation symmetry) and
enumerating all $\approx6.6\times10^6$ completions per type
(\texttt{xside2.c}), the $X$-side odds reproduce
$2\CC_7(\lambda)/\CC_7(\mu)-1$ exactly for all four admissible merges:
$(2,5)\to(7)$: $22848000/23116800=85/86$;
$(3,4)\to(7)$: $23063040/22901760=143/142$ (note the $B$-excess);
$(2,2,3)\to(2,5)$: $107/108$; $(2,2,3)\to(3,4)$: $35/36$.
The bijections themselves
(unflip status, roundtrips, hit multiplicities, orphan structure) were
verified instance-by-instance at $n=5,6$.


%======================================================================
\section{A reduction theorem}\label{sec:reduction}

\begin{theorem}\label{thm:reduction}
Assume $\mathrm{EQ}(\delta)$ with $\delta=\delta(n)\to0$.
Then
\[
\dtv(\PP_n,\QQ_n)\;=\;O\bigl(\delta\,\log n\;+\;n^{-2/3}\bigr).
\]
In particular $\mathrm{EQ}(o(1/\log n))$ implies
Conjecture~\ref{conj:cgw}, hence also that almost all loops have
multiplication group $S_n$ and trivial character theory.
\end{theorem}

\begin{proof}
Call $\lambda$ \emph{plain} if not all its parts are equal (so the only
non-plain types are $(\ell^{\kappa})$, $\kappa\ge2$).

\emph{Routing.}  Every plain $\lambda$ with $\kappa=\kappa(\lambda)$ parts is
reachable from $(n)$ by a chain of $\kappa-1$ splits, each with
$\alpha\ne\beta$.  Build the chain in reverse, splitting parts of $\lambda$
off a shrinking residual, in a suitable order; a collision means splitting a
part $p$ off a residual $R$ with $p=R-p$.  If the two largest parts of
$\lambda$ are unequal, split smallest-first and let them be the final pair:
while at least three parts remain, the split-off minimum $p$ satisfies
$R-p\ge2p>p$ (all parts are $\ge2$ and at least two other parts remain), so
no collision occurs, and the final split is unequal by hypothesis.  If
instead the two largest parts are equal, say to $\ell_{\max}$, then
$\ell_{\max}\neq n/2$ (else $\lambda=(n/2,n/2)$ is not plain), so we may
split one copy of $\ell_{\max}$ off first without collision, then proceed
smallest-first among the remaining parts, reserving for the final split the
pair $(b,\ell_{\max})$ where $b<\ell_{\max}$ is the smallest part: at every
intermediate step the residual minus the split-off minimum still contains
the reserved $\ell_{\max}$ and the reserved $b$, hence exceeds the
minimum (also when a further copy of $\ell_{\max}$ is the split-off
minimum: $R-p\ge\ell_{\max}+b>p$), and the final split is unequal.

For each step of such a chain, Theorem~\ref{thm:identity} and
$\mathrm{EQ}(\delta)$ give
$|S(p_{i+1})|/|S(p_i)|=(\gamma(p_{i+1})/\gamma(p_i))(1+\varepsilon_i)$ with
$|\varepsilon_i|\le\delta$.  Multiplying along the chain,
\[
\frac{\PP_n(\lambda)}{\PP_n((n))}
=\frac{\gamma(\lambda)}{\gamma((n))}\,(1+\varepsilon)^{\pm(\kappa-1)},
\qquad |\varepsilon|\le\delta,
\]
so for all plain $\lambda$ with $\kappa(\lambda)\le K$,
$\PP_n(\lambda)=\QQ_n(\lambda)\,Z^{-1}(1+O(K\delta))$ where $Z>0$ is a
normalising constant independent of $\lambda$ (namely
$D_n\CC_n((n))/\sum\gamma\CC$).

Let $K=\lceil A\log n\rceil$ with $A$ a large absolute constant, and let
$\mathcal B$ be the set of types that are either non-plain or have
$\kappa>K$.  Under $\QQ_n$, Markov's inequality with the generating function
$\sum_{\pi\in\DD_n}2^{\kappa(\pi)}=n!\,[x^n]e^{-2x}(1-x)^{-2}\le n\cdot n!$
(cf.\ \cite[proof of Lemma~4.4]{CGW08}) gives
$\QQ_n(\kappa>K)\le e\,n\,2^{-K}=O(n^{-2})$ for $A\log2\ge3$; also
$\QQ_n(\text{non-plain})\le\sum_{\ell\mid n,\ \ell<n}
\gamma((\ell^{n/\ell}))/D_n=O(n^{-1})$.  Under $\PP_n$, the unconditional
CGW bound \cite[Cor.~4.5]{CGW08},
$\PP_n(\lambda)\le n^{1/3}(\gamma(\lambda)/n!)\,2^{\kappa(\lambda)}$,
combined with
$\sum_{\pi\in\DD_n}4^{\kappa(\pi)}=n!\,[x^n]e^{-4x}(1-x)^{-4}
\le \binom{n+3}{3}\,n!$, gives
\[
\PP_n(\kappa>K)\;\le\;n^{1/3}\,2^{-K}\sum_\lambda
\frac{\gamma(\lambda)}{n!}\,4^{\kappa(\lambda)}
\;=\;O\bigl(n^{10/3}2^{-K}\bigr)\;=\;O(n^{-2/3})
\]
for $A\log2\ge4$, and similarly
$\PP_n(\text{non-plain})\le n^{1/3}\sum_{k\ge2,\ k\mid n}
4^{k}/(k!\,(n/k)^{k})=O(n^{-5/3})$.

On the complement of $\mathcal B$, summing
$\PP_n=\QQ_n Z^{-1}(1+O(K\delta))$ over $\lambda\notin\mathcal B$ and using
the tail bounds shows $Z^{-1}=1+O(K\delta)+O(n^{-2/3})$, whence
$\sum_{\lambda\notin\mathcal B}|\PP_n-\QQ_n|=O(K\delta)+O(n^{-2/3})$.
Combining with the tails,
$\dtv(\PP_n,\QQ_n)=O(\delta\log n+n^{-2/3})$.  The loops statement follows
from \eqref{eq:setTV}, Remark~\ref{rem:reduced} (the passage to reduced
squares is exact), and \S\ref{sec:known} as in \cite[Lemma 6.2]{CGW08}.
\end{proof}

\begin{remark}
The proof only uses $\mathrm{EQ}$ for splits along the specific chains, a
strictly weaker hypothesis; and the numerically observed truth is
$\mathrm{EQ}(Cn^{-3})$ on those chains.  The hypothesis is a statement about
one bit of structure --- \emph{rows 1 and 2 lie in the same column cycle of a
shifted post-flip column pair} --- for a uniformly random element of
$S(\lambda)$; it seems genuinely more tractable than the conjecture itself,
being an equidistribution assertion for a single event rather than a global
comparison of measures.  A natural route is a second-moment/switching
argument within $S(\lambda)$ itself; another is the transpose duality: the
$B$-status of column pairs of $L$ is the same-cycle statistic of row pairs
of $L^{\mathsf T}$, whose cycle-type law is again $\PP_n$ --- a
self-referential structure that may admit a bootstrap.
\end{remark}

%======================================================================
\section{Monte Carlo verification at larger \texorpdfstring{$n$}{n}}\label{sec:numerics}

We implemented the Jacobson--Matthews chain \cite{JM96} (with the
$\pm1$-tensor moves; uniform stationary distribution on proper squares) and
measured the quantities of \S\ref{sec:identity} at orders beyond the exact
tables.  Two methodological points deserve record.

\begin{enumerate}[leftmargin=2em]
\item \emph{Sampling bias of ``first proper state after $T$ moves''.}
Emitting the first proper state reached after a fixed number of moves is
\emph{not} uniform: the entry distribution into the proper set is biased by
the improper-excursion flux, which varies across squares (squares rich in
intercalates admit more proper-to-proper moves).  At $n=4$ this distorts the
type-$(2,2)$ probability from $1/2$ to $0.3887$; at $n=8$ it produced a
cycle-type $\chi^2\approx700$ on $6$ d.o.f.\ against the exact distribution.
Correct procedure: sample on a fixed grid of times and \emph{discard}
improper checkpoints.  With this correction the $n=4$ marginal is uniform to
$5\times10^{-4}$ over $2\times10^6$ samples.
\item \emph{Validation against exact data.}  With the corrected sampler:
at $n=8$ ($4\times10^6$ samples) the seven cycle-type frequencies match the
exact distribution with $\chi^2=2.7$ on $6$ d.o.f.; at $n=11$
($2\times10^6$ samples) all fourteen types give $\chi^2=10.4$ on $13$
d.o.f.
\end{enumerate}

\begin{center}
\small
\begin{tabular}{c c c c c c}
\toprule
$n$ & samples & $\Pr(B_1)$ & $\Pr(B_2)$ & $\bar q$ (measured) & $\bar q$ (exact, Thm~\ref{thm:identity})\\
\midrule
8  & $4.0\times10^6$ & $0.4971(5)$ & $0.4961(5)$ & $0.9932(8)$ & $0.993345$\\
11 & $2.0\times10^6$ & $0.4987(9)$ & $0.4992(9)$ & $0.9979(15)$ & $0.998080$\\
12 & $2.0\times10^6$ & $0.5009(10)$ & $0.4986(10)$ & $0.9995(15)$ & [$0.99850$ predicted]\\
16 & $1.5\times10^6$ & $0.5001(13)$ & $0.4983(13)$ & $0.9984(21)$ & [$0.99937$ predicted]\\
\bottomrule
\end{tabular}
\end{center}

\noindent
(Joins $\mu=(2,n-2)\to(n)$ throughout; ``predicted'' extrapolates
$\bar q=1-2c_2n^{-3}$.)  At $n=8$ the measurement is an $8\sigma$ detection
of $\bar q<1$, in exact agreement with the identity; at $n=11$ it again
matches the exact value; from $n=12$ on, the conjectured deviation
$1-\bar q\asymp2.6\,n^{-3}$ sits below Monte Carlo resolution, and the
measurements are consistent both with it and with
Hypothesis~\ref{hyp:eq} at the precision $|\bar q-1|\lesssim2\times10^{-3}$.

For a \emph{generic} admissible pair the measured $B$-probabilities are
$0.4993(3)$, $0.4997(4)$, $0.5002(4)$, $0.5007(4)$ at $n=8,11,12,16$; exact
enumeration
gives $3/5$ at $n=5$ and $37/78\approx0.474$ at $n=6$.  Two marked points of
a uniform derangement share a cycle with probability
$s_n=\tfrac12+(1+o(1))/n$ ($\approx0.55$ at $n=20$), so a naive analogy
would predict an excess $\approx+1/n$; but a uniform admissible column pair
excludes the $\approx n$ pairs $\{j,\omega(j)\}$, which are \emph{always}
same-cycle, and this exclusion cancels the $1/n$ term exactly, leaving
$\Pr(B\mid\text{admissible})=\tfrac12+O(n^{-2})$ --- as observed.  The
post-flip pairs of \S\ref{sec:identity} are then closer still to balance,
with deviation $\Theta(n^{-3})$: a three-level hierarchy
$\tfrac12+\Theta(n^{-1})\ \to\ \tfrac12+O(n^{-2})\ \to\ \tfrac12-\Theta(n^{-3})$
of same-cycle statistics.

%======================================================================
\section{Can the modern machinery reach the conjecture?}\label{sec:kps}

The strongest current toolkit for random Latin squares is that of
\cite{KSSS22,KPS25}: comparison with the triangle-removal process (TRP), and
re-randomisation via intercalate switchings.

\paragraph{Transfer costs.}  The TRP comparison theorems transfer events at
multiplicative cost $e^{O(n\log^2n)}$.  This is decisive for
large-deviation-scale statements (probabilities $e^{-\Theta(n^{4/3}\log n)}$
or smaller) but is far too lossy for Conjecture~\ref{conj:rates}, whose
content is at scale $1\pm n^{-3}$; even Conjecture~\ref{conj:cgw}
(constant-scale total variation) is out of its direct reach.

\paragraph{Re-randomisation.}  More promising is the \cite{KPS25}
re-randomisation philosophy.  An intercalate switch involving row $2$ and
one other row composes $\sigma_{1,2}$ with a transposition, i.e.\ acts on
the cycle type by an elementary split or merge --- these are exactly
degenerate (length-4 column cycle) cases of the CGW flip.  The row-parity
theorem of \cite{KPS25} shows that enough disjoint, well-separated
(``stable'') intercalates exist, and mix well enough, to drive a
$\pm1$ functional (the parity) to near-perfect equidistribution.  The cycle
type of $\sigma_{1,2}$ is a richer functional, but the induced dynamics on
types is a split/merge chain whose stationary law on cycle types, were the
switch statistics exactly balanced, would be $\QQ_n$-like.

\paragraph{A concrete programme via the marked-pair identity.}
Theorem~\ref{thm:marked} makes the re-randomisation target precise and,
we believe, attackable with the machinery of \cite{KPS25}.  The target
bit is now: for uniform $L\in S(\lambda)$ and a marked pair
$P=\{j,\sigma^\alpha(j)\}$, whether rows $1,2$ lie in one column cycle
of $P$.  Three structural facts line up favourably.

(i) \emph{The conditioning is free.}  Any intercalate avoiding rows
$1,2$ can be switched without touching $\sigma$: subset-switches of a
stable collection $\mathcal I_0(L)$ of disjoint intercalates in rows
$3..n$ (in the sense of \cite[\S7]{KPS25}: the collection is invariant
under such switches) act \emph{within} $S(\lambda)$, preserve the
marking, and re-randomise exactly ($\Pr(L\to L')=2^{-|\mathcal I_0|}$
both ways): conditionally on the orbit, the configuration is uniform
on $\{0,1\}^{\mathcal I_0}$.  No comparison between $S(\lambda)$ and
the uniform measure is needed at this step.

(ii) \emph{Switch action on the bit.}  An intercalate on rows
$\{r,r'\}\subseteq\{3..n\}$ meeting exactly one of the columns $j,j'$
acts on the pair permutation $\pi_P$ by precomposition with
$\tau_{rr'}$: it splits or joins the $\pi_P$-cycles through $r,r'$,
and flips the target bit precisely when $r,r'$ ``interleave'' rows
$1,2$ (one lies on the segment between them).  Intercalates meeting
\emph{both} or \emph{neither} of the columns act by conjugation and
fix the bit.  So the bit is non-linear in the switch configuration but
locally computable, and each useful intercalate flips it with
probability of constant order in a random-like configuration.

(iii) \emph{The quantifiers fit the budget.}  The intercalate-expander
lemma of \cite[Lem.~8.3, Def.~8.1]{KPS25} (with $\ell=\log^{11}n$)
implies that all but at most $\ell$ columns carry at least $6\ell$
disjoint stable intercalates (their construction of the row set $R$,
transposed), with rows and partner columns avoidable in prescribed
$\beta n$-sets --- in particular avoiding rows $1,2$ and column $j'$.
The exceptional fraction of marks is $O(\log^{11}n/n)=o(1/\log n)$:
exactly within the $\mathrm{EQ}(o(1/\log n))$ budget that the weak
conjecture requires.  Moreover their approximation costs are
affordable here: $\Pr(S(\lambda))\ge e^{-O(n\log n)}$ for every type
$\lambda$ (worst case all-$2$-cycles), while the expander lemma fails
with probability $e^{-\omega(n\log^2 n)}$, so the expander property
holds for $L$ conditioned on $S(\lambda)$, for every $\lambda$, by a
union-of-costs argument.

What remains --- the genuinely open core --- is the \emph{non-linear
orbit-mixing lemma}: for a $2^{\mathcal I_0}$-uniform configuration
and $k\ge 6\log^{11}n$ disjoint stable intercalates through column $j$
avoiding rows $1,2$, show that the conditional law of the bit given
the orbit is $\tfrac12+o(1/\log n)$ for all but an $o(1/\log n)$
fraction of orbits.  Unlike the parity bits of \cite{KPS25} the bit is
not $\mathbb F_2$-linear in the configuration, so their kernel
computation must be replaced by a mixing argument for the
cycle-structure bit under random subsets of transpositions --- the
orbit dynamics on $\pi_P$ is a coagulation--fragmentation process
driven by (a random subset of) disjoint transpositions, for which the
classical toolbox is that of random transposition walks
\cite{Sch05} --- but the
per-switch flip probability is of constant order and the budget
$o(1/\log n)$ is far weaker than their polynomial error terms.  We
record this as the sharpest currently-visible route to
Conjecture~\ref{conj:cgw}.

\medskip\noindent
\emph{[Session-6 update: the non-linear orbit-mixing lemma just
described has been reformulated exactly.  \S\ref{sec:orbit} proves
that the conditional bit given an orbit is an explicit
$\mathbb{F}_2$-rank event in a fixed interlacement matrix
(Theorem~\ref{thm:master}), derives the supply and conditional
expander lemmas, shows the orbit-by-orbit form of the mixing lemma
stated above is too strong, and replaces it by the annealed
Problem~\ref{prob:annealed}.]}

%======================================================================
\section{The orbit bit is a rank condition: an exact interlacement
formula}\label{sec:orbit}

The open core of the transfer programme of \S\ref{sec:kps} is the
``non-linear orbit-mixing lemma'': conditionally on a switching orbit,
the marked bit --- rows $1,2$ in one column cycle of the marked pair
--- is a cycle-structure functional of a product of a fixed
permutation with a uniformly random subset of disjoint transpositions,
and one must show its conditional law is close to $\mathrm{Ber}(1/2)$.
In this section we prove that this functional is, exactly, a rank
condition over $\mathbb{F}_2$ in a matrix determined by the orbit ---
so the ``non-linear'' bit is a (non-linear but completely explicit)
linear-algebraic event after all, directly comparable to the kernel
computation that powers the parity theorem of \cite{KPS25} (their
Lemma~8.5).  The formula turns the orbit-mixing problem into a
question about ranks of random principal submatrices of a fixed
interlacement matrix; we then (i) derive an exact identity for the
conditional bias, (ii) exhibit the extremal configurations, (iii)
show computationally that in generic position the bias decays like
$k^{-1/2}$ in the number $k$ of available switches --- far inside the
$o(1/\log n)$ budget when $k\ge6\log^{11}n$ --- and (iv) reformulate
the open core accordingly, in an \emph{annealed} form which the
present data indicate is both sufficient and true, while the
orbit-by-orbit form stated at the end of \S\ref{sec:kps} is provably
too strong.  We also record the two supporting lemmas of the
programme, the supply lemma (\S\ref{sec:supply}) and the conditional
expander (\S\ref{sec:condexp}), whose proofs are short.

Throughout, $u=$ row $1$ and $v=$ row $2$; $\pi_P$ is the pair
permutation of a column pair $P$ on the row set (\S\ref{sec:marked}),
whose cycles are the column cycles of $P$; and the marked bit of an
instance is $[\,u\sim v \text{ in } \pi_P\,]$, i.e.\ the indicator of
$B$-status.

%----------------------------------------------------------------------
\subsection{Chords, interlacement, and the Cohn--Lempel theorem}

Let $\gamma$ be a cyclic permutation of a finite set $V$ (``the
circle'').  A \emph{chord} is an unordered pair $\{a,b\}\subseteq V$;
a set of chords is \emph{disjoint} if the pairs are pairwise disjoint.
Chords $\{a,b\}$ and $\{c,d\}$ \emph{interlace} (on $\gamma$) if $c,d$
lie in different arcs of the circle determined by $a,b$; a chord
$\{a,b\}$ \emph{separates} points $u,v\notin\{a,b\}$ if $u,v$ lie in
different arcs.  For a set $T$ of disjoint chords, identified with the
product $\tau_T=\prod_{\{a,b\}\in T}(a\,b)$ of the corresponding
commuting transpositions, the \emph{interlacement matrix}
$A_T\in\mathbb{F}_2^{T\times T}$ has $(A_T)_{ij}=1$ iff chords $i\neq
j$ interlace, and zero diagonal.

\begin{theorem}[Cohn--Lempel \cite{CL72}]\label{thm:CL}
Let $\gamma$ be a single cycle on $V$ and $T$ a set of disjoint
chords.  Then the number of cycles of $\gamma\,\tau_T$ equals
$\operatorname{null}_{\mathbb{F}_2}(A_T)+1$.
\end{theorem}

(We re-verified Theorem~\ref{thm:CL} instance-by-instance over random
circles and all activation subsets; \texttt{pf\_cl.py}.)

The matrix $A_T$ is alternating (symmetric, zero diagonal), so
$z^{\mathsf T}A_Tz=0$ for every $z$; this is what drives the next
lemma.

\begin{lemma}[bordered rank]\label{lem:border}
Let $A\in\mathbb{F}_2^{m\times m}$ be alternating and
$x\in\mathbb{F}_2^m$, and let
$B=\begin{psmallmatrix}A&x\\x^{\mathsf T}&0\end{psmallmatrix}$.  If
$x\in\operatorname{col}(A)$ then
$\operatorname{null}(B)=\operatorname{null}(A)+1$; otherwise
$\operatorname{null}(B)=\operatorname{null}(A)-1$.
\end{lemma}

\begin{proof}
If $x=Ay$, add the $y$-combination of the first $m$ columns to the
last column and the $y$-combination of the first $m$ rows to the last
row: the border becomes $(0,\dots,0,\;y^{\mathsf T}Ay)=(0,\dots,0,0)$
since $A$ is alternating; hence
$\operatorname{rank}(B)=\operatorname{rank}(A)$ and
$\operatorname{null}(B)=(m{+}1)-\operatorname{rank}A
=\operatorname{null}(A)+1$.  If $x\notin\operatorname{col}(A)$, then
$\operatorname{rank}[A\;x]=\operatorname{rank}(A)+1$, and the last row
$(x^{\mathsf T},0)$ is independent of the rows of $[A\;x]$: a relation
$x^{\mathsf T}=z^{\mathsf T}A$ would give $x=Az$ by symmetry.  So
$\operatorname{rank}(B)=\operatorname{rank}(A)+2$ and
$\operatorname{null}(B)=\operatorname{null}(A)-1$.
\end{proof}

\begin{theorem}[the bit as a rank condition; single cycle]
\label{thm:bitrank}
Let $\gamma$ be a single cycle on $V$, let $u,v\in V$, and let $T$ be
a set of disjoint chords avoiding $u,v$.  Let $x\in\mathbb{F}_2^{T}$
be the \emph{separation vector}, $x_i=[\,\text{chord }i\text{
separates }u,v\,]$.  Then
\[
u\sim v\ \text{in}\ \gamma\,\tau_T
\qquad\Longleftrightarrow\qquad
x\in\operatorname{col}(A_T).
\]
\end{theorem}

\begin{proof}
Let $e$ denote the chord $\{u,v\}$; since $T$ avoids $u,v$, the set
$T\cup\{e\}$ is disjoint, its interlacement matrix is the bordered
matrix of Lemma~\ref{lem:border} (chord $i$ interlaces $e$ iff it
separates $u,v$), and multiplication of any permutation $\rho$ by
$(u\,v)$ splits the cycle through $u,v$ if $u\sim v$ in $\rho$ and
joins two cycles otherwise.  Thus, writing $\rho=\gamma\tau_T$, we
have $u\sim v$ in $\rho$ iff
$c(\rho\,(u\,v))=c(\rho)+1$, iff (Theorem~\ref{thm:CL} twice)
$\operatorname{null}(A_{T\cup e})=\operatorname{null}(A_T)+1$, iff
(Lemma~\ref{lem:border}) $x\in\operatorname{col}(A_T)$.
\end{proof}

Both directions of the split/join dichotomy for $(u\,v)$ are the
switch-toggle mechanism of Lemma~\ref{lem:switchtoggle} in disguise:
adjoining the virtual chord $e$ is the CGW switch at the marked pair.

%----------------------------------------------------------------------
\subsection{Arbitrary base permutations and two-sided products}

Two reductions extend Theorem~\ref{thm:bitrank} to the situation
actually produced by an orbit: the base permutation $\pi_0=\pi_P$ has
many cycles, and the switch actions multiply on both sides.

\begin{lemma}[passenger construction]\label{lem:passenger}
Let $\pi_0$ have cycles $C_1,\dots,C_p$ (written as sequences), let
$s_1,t_1,\dots,s_{p-1},t_{p-1}$ be $2(p-1)$ fresh points, let
\[
\gamma=(\,C_1\;s_1\;C_2\;t_1\;s_2\;C_3\;t_2\;\cdots\;
s_{p-1}\;C_p\;t_{p-1}\,)
\]
be the single cycle obtained by the indicated concatenation, and let
$B=\{\{s_i,t_i\}:1\le i\le p-1\}$ (the \emph{bridges}; pairwise
disjoint).  Then $\gamma\,\tau_B$ has cycles
$C_1\cup\{s_1,\dots,s_{p-1}\}$ (with the $s_i$ inserted consecutively
after the last element of $C_1$) and $C_i\cup\{t_{i-1}\}$ for $2\le
i\le p$; in particular, deleting the fresh points (first-return map)
recovers $\pi_0$ exactly, and for any set $T$ of chords on
$V$ avoiding the fresh points,
$u\sim v$ in $\pi_0\tau_T$ iff $u\sim v$ in $\gamma\,\tau_{B\cup T}$.
\end{lemma}

\begin{proof}
Direct computation of successors in $\gamma\tau_B$: for $x\in C_i$ not
last, $\gamma\tau_B(x)=\gamma(x)$ is the successor in $C_i$; for the
last element $x$ of $C_i$ ($i\ge2$), $\gamma(x)=t_{i-1}$ and
$\gamma\tau_B(t_{i-1})=\gamma(s_{i-1})$ is the first element of $C_i$,
closing the cycle $C_i\cup\{t_{i-1}\}$; for the last element of $C_1$,
$\gamma(x)=s_1$, $\gamma\tau_B(s_i)=\gamma(t_i)=s_{i+1}$ for $i<p-1$,
and $\gamma\tau_B(s_{p-1})=\gamma(t_{p-1})=$ first element of $C_1$
(the wrap), closing $C_1\cup\{s_1,\dots,s_{p-1}\}$.  Since first
return commutes with right multiplication by transpositions avoiding
the deleted points, the last claim follows.
\end{proof}

\begin{lemma}[point splitting for two-sided products]\label{lem:split}
Let $\pi_0$ be a permutation of $V$, and let $\rho=\tau_R$,
$\mu=\tau_M$ be products of transpositions where the chords within $R$
are pairwise disjoint and likewise within $M$, all avoiding $u,v$, but
chords of $R$ may meet chords of $M$.  Form the \emph{doubled circle}:
replace each point $x$ by two consecutive copies $(x,0),(x,1)$ in each
cycle of $\pi_0$; lift each chord $\{a,b\}\in M$ to
$\{(a,0),(b,0)\}$ and each $\{a,b\}\in R$ to $\{(a,1),(b,1)\}$.  Then
all lifted chords are pairwise disjoint, and the permutation
$\widetilde P=\widetilde{\pi_0}\,\tau_{\widetilde R\cup\widetilde M}$
of the doubled point set (where $\widetilde{\pi_0}$ is the doubled
cyclic structure) satisfies: the first-return map of $\widetilde P$ on
the copies $\{(x,0)\}$ is $(x,0)\mapsto(\pi_0\rho\mu(x),0)$.  In
particular $u\sim v$ in $\pi_0\rho\mu$ iff $(u,0)\sim(v,0)$ in
$\widetilde P$.
\end{lemma}

\begin{proof}
Compute $\widetilde P$ on the two kinds of points.  On $(x,0)$: the
only lifted chord that can contain $(x,0)$ is the lift of the
$M$-chord through $x$ (if any), so
$\tau(\,(x,0)\,)=(\mu(x),0)$ in every case, and
$\widetilde{\pi_0}((\mu(x),0))=(\mu(x),1)$: hence $\widetilde
P((x,0))=(\mu(x),1)$.  On $(x,1)$: the only chord containing $(x,1)$
is the lift of the $R$-chord through $x$ (if any), so
$\tau((x,1))=(\rho(x),1)$ and
$\widetilde{\pi_0}((\rho(x),1))=(\pi_0\rho(x),0)$: hence $\widetilde
P((x,1))=(\pi_0\rho(x),0)$.  Composing,
$\widetilde P^2\!\restriction_{\text{copies }0}:
(x,0)\mapsto(\pi_0\rho\mu(x),0)$, and every second point of each
$\widetilde P$-cycle is a copy-$0$ point, so the first-return map is
as claimed and cycle membership transfers.
\end{proof}

%----------------------------------------------------------------------
\subsection{The master formula for the orbit bit}

Fix an orbit: a base square $L_0\in S(\lambda)$ (any point of the
orbit), a marked pair $P=\{j,j'\}$, the stable collection
$\mathcal I_0$ of disjoint stable intercalates in rows $3..n$ (in the
sense of \cite[\S7]{KPS25}, with respect to a fixed template), and the
uniform configuration $\omega\in\{0,1\}^{\mathcal I_0}$.  Partition
the \emph{useful} coordinates by column incidence with $P$:
$\mathcal I_j$ (meets column $j$ only), $\mathcal I_{j'}$ (column $j'$
only), $\mathcal I_{jj'}$ (both).  As computed in \S\ref{sec:kps}: a
switch in $\mathcal I_j$ with rows $\{r,r'\}$ precomposes $\pi_P$ with
$(r\,r')$; one in $\mathcal I_{j'}$ postcomposes (left-multiplies); one
in $\mathcal I_{jj'}$ conjugates; all other coordinates fix $\pi_P$.
Two stable intercalates meeting a common column cannot share a row
(they would share a cell, contradicting isolation), so the chords
within each of the families $R:=\mathcal I_j\cup\mathcal I_{jj'}$
(right side) and $M:=\mathcal I_{j'}\cup\mathcal I_{jj'}$ (left side)
are pairwise disjoint; between the two sides rows may be shared.
Writing $\pi(\omega)=M_\omega\,\pi_0\,R_\omega$ and conjugating by
$M_\omega^{-1}$ (which fixes $u,v$: all supports avoid rows $1,2$),
\[
\text{bit}(\omega)
=[\,u\sim v\ \text{in}\ \pi_0\,R_\omega\,M_\omega\,].
\]
Now apply Lemma~\ref{lem:split} (doubling; $M$-chords to copies $0$,
$R$-chords to copies $1$), then Lemma~\ref{lem:passenger} (bridges
$B$) to the doubled multi-cycle structure, then
Theorem~\ref{thm:bitrank} on the resulting single circle:

\begin{theorem}[master formula]\label{thm:master}
For each orbit there is an explicitly computable alternating matrix
$\bar A$ over $\mathbb{F}_2$, indexed by the lifted chords of the
useful coordinates together with the bridge set $B$
($|B|=\#\text{cycles}(\pi_0)-1$), and a vector $\bar x$ (separation of
$(u,0),(v,0)$ on the augmented circle), both determined by the
cyclic-order positions of the intercalate rows in the cycles of
$\pi_0=\pi_P(L_0)$, such that for every configuration $\omega$,
\[
\mathrm{bit}(\omega)=
\bigl[\,\bar x_{\,\varphi(\omega)\cup B}
\in\operatorname{col}\bigl(\bar A_{\,\varphi(\omega)\cup B}\bigr)\bigr],
\]
where $\varphi(\omega)$ is the set of lifted chords of the active
useful coordinates (a coordinate in $\mathcal I_{jj'}$ contributes its
two lifts together).  Consequently the conditional law of the marked
bit given the orbit is the law of a rank event of a uniformly random
principal submatrix (containing the frozen block $B$) of the fixed
matrix $\bar A$.
\end{theorem}

\emph{Status: PROVED (Theorem~\ref{thm:CL} cited; all other
ingredients proved above), and VERIFIED instance-by-instance:
\texttt{pf\_cl.py} checks (CL), (BIT), the passenger construction,
and randomized two-sided configurations with shared rows and
conjugating coordinates ($1262$ cases, all activation subsets,
exhaustively) against direct computation.}

The non-linearity of the bit is thus precisely the non-linearity of
$S\mapsto\operatorname{null}(\bar A_S)$; the $\mathbb{F}_2$-linear
parity functionals of \cite{KPS25} correspond to the special case
where the relevant matrix acts by a fixed linear map on the
configuration.  The coagulation--fragmentation dynamics
(\cite{Sch05}) has been integrated out: no walk analysis is needed at
fixed orbit.

%----------------------------------------------------------------------
\subsection{The exact bias identity and extremal configurations}

Let $k$ be the number of useful coordinates and, for $S$ a uniform
random subset, let
$\mathrm{bias}=\Pr_S[\mathrm{bit}]-\tfrac12$ (conditional on the
orbit).

\begin{proposition}[bias identity]\label{prop:bias}
With $e$ the virtual chord $\{(u,0),(v,0)\}$ adjoined to the ground
set,
\[
\mathrm{bias}
=\tfrac12\,\E_S\bigl[\operatorname{null}\bar A_{S\cup B\cup e}
-\operatorname{null}\bar A_{S\cup B}\bigr]
\]
--- the expected nullity increment of the virtual chord.  Moreover,
for any single coordinate $i$, the toggle status of $i$ (whether
flipping $\omega_i$ flips the bit) is a function of
$\omega_{-i}$ alone, and
$|\mathrm{bias}|\le\tfrac12\Pr[\,i\text{ does not toggle}\,]$.
\end{proposition}

\begin{proof}
The first claim is Lemma~\ref{lem:border} applied at each $S$ (the
increment is $+1$ exactly on the event $\mathrm{bit}=1$).  For the
second: multiplying by the transposition of coordinate $i$ splits or
joins according to the current configuration of the others, and the
toggle relation is symmetric between $\omega_i=0,1$; conditionally on
$\omega_{-i}$ with $i$ toggling, the bit is exactly Bernoulli$(1/2)$.
\end{proof}

Three exactly solvable families calibrate what balance requires
(all verified in \texttt{pf\_bias.py}):

\begin{enumerate}[leftmargin=2em]
\item \textbf{Ladders.}  $r$ separating chords, pairwise interlacing
($\bar A=J+I$ on the ladder, $x=\mathbf 1$), no other chords: the
kernel computation gives $\mathrm{bit}(S)=[\,|S|\ \text{even}\,]$
--- the bit collapses to an $\mathbb{F}_2$-linear parity and
$\mathrm{bias}=0$ exactly, for every $r\ge1$.
\item \textbf{Nested families.}  $r$ separating chords, pairwise
non-interlacing ($\bar A=0$, $x=\mathbf 1$):
$\mathrm{bit}(S)=[S=\emptyset]$, $\mathrm{bias}=2^{-r}-\tfrac12$.
This is the worst case, and is the exact analogue of the failure
mode that the expander hypothesis exists to exclude.
\item \textbf{Cut chords.}  A single separating chord $c$ interlacing
no other chord (an ``uncrossed cut'') forces $\mathrm{bit}=0$ whenever
$c$ is active: $e_c$ is then a kernel vector with $x_c=1$, killing the
rank event.  Hence, exactly,
$\mathrm{bias}=\tfrac12\,\mathrm{bias}'-\tfrac14\in[-\tfrac12,0]$,
where $\mathrm{bias}'$ is the conditional bias given $c$ inactive.  In
particular good structure elsewhere cannot restore balance: with a
ladder in the remainder ($\mathrm{bias}'=0$), $\mathrm{bias}=-\tfrac14$
exactly (verified).  Balance is therefore \emph{not} implied by the
presence of good structure; it requires the \emph{absence} of bad
structure --- every separating chord must itself be crossed,
recursively.
\end{enumerate}

Item 3 has a consequence for the shape of the mixing lemma: no
argument that only \emph{finds} favourable switches (a ladder) can
close the programme; the statement must exploit the measure on
orbits.

%----------------------------------------------------------------------
\subsection{Generic position: \texorpdfstring{$k^{-1/2}$}{k**-1/2}
decay, and the annealed reformulation}\label{sec:annealed}

The handoff formulation of the mixing lemma (``conditional bit-law
$=\frac12+o(1/\log n)$ for all but $o(1/\log n)$ of orbits'') is
\emph{too strong}: if rows $1,2$ are at distance $O(1)$ along their
common $\pi_0$-cycle and no intercalate row falls in the short arc
(an event of probability $\Theta(1/k)$ under any spread-out measure,
and of constant conditional probability on some $S(\lambda)$-positive
events), then no available chord separates $u,v$, $x=0$, and
$\mathrm{bit}\equiv1$: $\mathrm{bias}=+\tfrac12$ on an orbit set of
measure $\asymp 1/k$, not $o(1/\log n)$ small in general.  The correct
target is the \emph{annealed} (orbit-averaged) statement, which is
also all that Theorem~\ref{thm:marked} requires:

\begin{problem}[annealed orbit-mixing --- the revised open core]
\label{prob:annealed}
For $L$ uniform on $S(\lambda)$ restricted to the intercalate-expander
event, $j$ a supplied column (Lemma~\ref{lem:supply}), show
\[
\E_L\,\bigl|\mathrm{bias}(\mathrm{orbit}(L),P)\bigr|
\;=\;o(1/\log n)
\qquad\text{(averaged over marks as in \S\ref{sec:marked})}.
\]
\end{problem}

Computationally, generic position gives far more than is needed.  For
the ensemble ``$u,v$ and $k$ disjoint chords in uniformly random
position on a circle'' (all activation subsets exact for $k\le20$,
Monte Carlo beyond; \texttt{bias\_scale.c}):
\[
\begin{array}{lccccc}
\toprule
k & 8 & 16 & 32 & 64 & 128\\
\midrule
\E|\mathrm{bias}| & 0.124 & 0.099 & 0.075 & 0.044 & 0.032\\
\text{median} & 0.102 & 0.082 & 0.053 & 0.028 & 0.018\\
\text{mean (signed)} & +0.015 & +0.000 & -0.000 & +0.000 & +0.007\\
\bottomrule
\end{array}
\]
i.e.\ $\E|\mathrm{bias}|\approx0.35\,k^{-1/2}$, with the signed mean
consistent with $O(1/k)$.  The dependence on the $u$--$v$ arc length
is flat except when the arc contains $O(1)$ chord endpoints (measured:
at $k=64$, mean bias $+0.12$ with $\approx1$ endpoint-slot in the
arc, indistinguishable from generic already at $\approx3$); the
short-arc contribution to the annealed average is therefore
$O(1/k)$, self-regularizing --- no separate ``join-then-toggle chain
argument'' is required in the annealed formulation.  With
$k\ge6\log^{11}n$ (the supply lemma), a bound of quality
$\E|\mathrm{bias}|\le k^{-c}$ for \emph{any} constant $c>1/11$ closes
Problem~\ref{prob:annealed}; the generic-position value $c=1/2$ has a
margin of a factor $5$ in the exponent.

\emph{Status: the reduction of the weak conjecture to
Problem~\ref{prob:annealed} together with
Lemmas~\ref{lem:supply}--\ref{lem:condexp} is PROVED modulo the two
bookkeeping items recorded in Remark~\ref{rem:bookkeeping} [session 13:
item (i) is closed by Proposition~\ref{prop:nodecouple}]; the
$k^{-1/2}$ scaling is COMPUTED (not proved) and refers to the
iid-position ensemble, not yet to the true orbit ensemble.}

What would a proof of Problem~\ref{prob:annealed} need?  By
Proposition~\ref{prop:bias} it suffices to show that under the true
orbit measure the chord configuration is \emph{generic}: e.g.,
(i) with probability $1-o(1/\log n)$ every separating chord is
interlaced by some other available chord (excluding cut chords and
nested families), and (ii) a quantitative rank-balance for the
resulting random interlacement matrices.  Both are statements about
the joint law of the positions, in the cycles of $\pi_P$, of the
stable-intercalate rows of a uniform $L\in S(\lambda)$ --- a
concentration/anti-concentration problem, not a worst-case one.  The
exchangeability of rows $3..n$ (which acts on everything by
conjugation, though not template-equivariantly) and the
Harer--Zagier-type computation of expected nullity for random chord
systems are the natural tools.  [Session-7 update: conditions of
type (i) are provably insufficient at any fixed crossing
multiplicity (Proposition~\ref{prop:pendant} and the exhaustive scan
of \S\ref{sec:pendant}); the productive formalisation of (ii) is the
second-moment factorization of \S\ref{sec:factor}, which reduces the
iid model to Problem~\ref{prob:mixing}.]

%----------------------------------------------------------------------
\subsection{The supply lemma}\label{sec:supply}

\begin{lemma}[supply]\label{lem:supply}
Let $H$ be a partial Latin square that is an
$(\ell,\beta)$-intercalate-expander (Definition~8.1 of
\cite{KPS25}) on $n$ symbols, $n$ large --- in the application,
$H=T\cap L$; we avoid the letter $P$, which denotes the marked pair
--- and let $u,v$ be two rows,
$j'$ a column, and $t\ge1$ with $4t\ell+2t+\ell+2\le\beta n/2$ and
$\beta\le\tfrac12$ (for general $\beta$ replace the hypothesis by
$2t\ell+2\le(1-\beta)n$, $t\ell+\ell+1\le(1-\beta)n$; the application
has $\beta=1/10$).  Then
for all but fewer than $t\ell$ columns $c\neq j'$ there exist $t$
stable intercalates of $H$ through $c$, pairwise disjoint in rows,
partner columns and symbols, all avoiding rows $\{u,v\}$ and column
$j'$.  Moreover \emph{all} stable intercalates through a common column
are automatically pairwise row-disjoint (two of them sharing a row
would share a cell in that column, contradicting isolation).
With $\ell=\log^{11}n$ and $t=6\ell$: all but $6\log^{22}n$ columns
carry at least $6\log^{11}n$ such intercalates.
\end{lemma}

\begin{proof}
Process rounds $\tau=1,\dots,t$; every column starts alive with an
empty family, and in round $\tau$ each alive column attempts to add
one intercalate to its family subject to the constraints; a column
that fails is declared dead and never retried.  Suppose $\ell$ columns
die in some round $\tau$.  Each dead column has a current family of
$\tau-1\le t$ intercalates.  Apply the expander property with:
$C^*=$ the $\ell$ dead columns; $R=R^*=$ all rows except $u,v$ and the
$\le 2t\ell$ rows used by the dead columns' families (size $\ge
n-2-2t\ell\ge\beta n$); $C=$ all columns except $C^*$, $j'$, and the
$\le t\ell$ used partner columns (size $\ge n-\ell-1-t\ell\ge\beta
n$); $S=S^*=$ all symbols except the $\le2t\ell$ used ones.  The
guaranteed stable intercalate has one column $c_0\in C^*$ and avoids
every constraint of $c_0$'s attempted extension --- contradicting the
death of $c_0$.  So fewer than $\ell$ columns die per round, hence
fewer than $t\ell$ die overall, and every surviving column finishes
with a full family of size $t$.
\end{proof}

Note that Definition~8.1 is monotone in the five $\beta n$-sets (a
larger set contains a $\beta n$-subset), which the proof uses freely.

\begin{remark}[bookkeeping for exceptional columns; template
averaging]\label{rem:bookkeeping}
Two quantifier items remain to make the assembly fully rigorous, both
of which we record with proposed fixes rather than proofs.
(i) \emph{Exceptional marks.}  The $t\ell$ exceptional columns of
Lemma~\ref{lem:supply} are square-dependent; for merges whose
$m$-cycles occupy few columns ($\lambda_m m=n^{o(1)}$, say) the
per-square fraction of bad marks is not automatically $o(1/\log n)$.
[Session-6 text proposed column-relabelling (template) averaging as
a fix.  \emph{Withdrawn}: Proposition~\ref{prop:vacuity} shows that
averaging over relabelled templates $g(T)$ leaves the $X$-averaged
bad-mark fraction exactly unchanged, by joint equivariance.  The
genuine requirement is the supply--mark decoupling stated as
Problem~\ref{prob:decouple}.]
(ii) \emph{Rows-$1,2$ conditioning.}  The expander Lemma 8.3 of
\cite{KPS25} is stated for uniform $L$; conditioning on $S(\lambda)$
is handled by Lemma~\ref{lem:condexp} below, but the supplied
intercalates must additionally avoid rows $1,2$, which
Lemma~\ref{lem:supply} already builds in through the row sets $R,R^*$.
\end{remark}

%----------------------------------------------------------------------
\subsection{The conditional expander}\label{sec:condexp}

\begin{lemma}[expander conditioned on the type]\label{lem:condexp}
Let $T$ be the template of \cite[Lem.~8.3]{KPS25} and
$\ell=\log^{11}n$.  For every type $\lambda$,
\[
\Pr\bigl[\,T\cap L\ \text{is an }(\ell,\beta)\text{-intercalate-expander}
\ \big|\ L\in S(\lambda)\bigr]
\;\ge\;1-e^{-\omega(n\log^2 n)}.
\]
\end{lemma}

\begin{proof}
$\Pr[\lnot\text{expander}\mid S(\lambda)]\le
\Pr[\lnot\text{expander}]/\Pr[S(\lambda)]$, and
$\Pr[\lnot\text{expander}]\le e^{-\omega(n\log^2n)}$ by
\cite[Lem.~8.3]{KPS25}, so it suffices that
$\Pr[S(\lambda)]\ge e^{-O(n\log n)}$ for every $\lambda$.  Writing
$\Pr[S(\lambda)]=D_\lambda N(\lambda)/\sum_\nu D_\nu N(\nu)$:
every $D_\nu/D_\lambda\le n!/D_\lambda\le e^{O(n\log n)}$ (the worst
type, all $2$-cycles, has $D_\lambda=n!/(2^{n/2}(n/2)!)\ge
n!\,e^{-O(n\log n)}$), the number of types is $e^{O(\sqrt n)}$, and
$N(\nu)/N(\lambda)\le3^{n}$ for all $\nu,\lambda$ by composing the
two-sided merge bounds ($\frac12\le \CC_n(\text{merge})/
\CC_n(\text{split})\le\frac32$: the upper bound is
\cite[Thm.~3.13]{CGW08}, the lower is that theorem or
Theorem~\ref{thm:marked}) along a path of at most $n$ merges through
the one-part type.  Combining, $\Pr[S(\lambda)]\ge
e^{-O(n\log n)}$.
\end{proof}

In particular no transfer through the comparison machinery
(\cite[Lem.~4.7]{KPS25}) is needed for the conditioning step: direct
division of probabilities suffices, for every type simultaneously.

%----------------------------------------------------------------------
\subsection{Summary of the revised programme}

\begin{enumerate}[leftmargin=2em]
\item Conditioning is free; orbits re-randomise exactly
(\S\ref{sec:kps}, unchanged).
\item Supply: for all but $o(1/\log n)$ of marks, both marked
columns carry $\ge6\log^{11}n$ disjoint stable switches avoiding
rows $1,2$ (Lemma~\ref{lem:supply} + Lemma~\ref{lem:condexp}).
[PROVED for marks in classes of size $\lambda_m m\ge\log^{24}n$,
where the deterministic bound $|E|\le6\log^{22}n$ gives bad fraction
$O(1/\log^2 n)$; for smaller classes OPEN as
Problem~\ref{prob:decouple} --- template averaging does not close
it, Proposition~\ref{prop:vacuity}.]
\item The conditional bit given the orbit is an explicit
$\mathbb{F}_2$ rank event in a fixed interlacement matrix
(Theorem~\ref{thm:master}).  [PROVED.]
\item Open core: the annealed bias bound
(Problem~\ref{prob:annealed}).  Generic position gives
$\E|\mathrm{bias}|\approx0.35k^{-1/2}$, a factor-$5$ exponent margin
over the budget; extremal counterexamples (nested families, cut
chords) show the bound must come from the measure, not from
worst-case structure.  [OPEN; scaling COMPUTED.]
\item $\mathrm{EQ}(o(1/\log n))$ then follows by
Theorem~\ref{thm:marked}, and the weak conjecture by
Theorem~\ref{thm:reduction}.
\end{enumerate}


%======================================================================
\section{The annealed bound: factorization of the second moment,
refutation of local genericity, and real-ensemble data}
\label{sec:anneal}

This section records the session-7 progress on
Problem~\ref{prob:annealed}.  The positive result is a
\emph{two-point factorization} of the annealed second moment
(Theorem~\ref{thm:factor}), which reduces the iid model of the
annealed bias bound to a single one-parameter mixing statement ---
polynomial $L^2$ decay under a ``fresh-chord averaging'' Markov
operator (Problem~\ref{prob:mixing}) --- verified numerically with a
large margin.  The negative results sharpen the target: every local
genericity condition of the kind contemplated at the end of
\S\ref{sec:annealed} is refuted by an explicit family
(Proposition~\ref{prop:pendant} and an exhaustive scan), and the
template-averaging fix proposed in Remark~\ref{rem:bookkeeping}(i) is
shown to be \emph{exactly neutral} (Proposition~\ref{prop:vacuity}),
which re-opens the exceptional-mark bookkeeping as a genuine (if
presumably easier) probabilistic sub-problem, stated as
Problem~\ref{prob:decouple}.  Finally, the master formula and the
bias statistics are grounded on \emph{real} random squares at
$n=30$--$100$ via Jacobson--Matthews sampling.

%----------------------------------------------------------------------
\subsection{The two-point factorization of the annealed second
moment}\label{sec:factor}

Work in the \emph{iid ensemble}: $u,v$ fixed on the circle, and
chords $Z_1,\dots,Z_k$ with iid endpoint pairs, uniform on the circle
(a.s.\ disjoint); this is the continuum version of the
generic-position ensemble of \S\ref{sec:annealed}, and the argument
below uses nothing about the position law except that the $Z_i$ are
iid given $(u,v)$.  For a finite set $W$ of disjoint chords let
\[
\varepsilon(W)\;=\;c\bigl(\gamma_W\,\tau_W\,(u\,v)\bigr)
-c\bigl(\gamma_W\,\tau_W\bigr)\;\in\;\{\pm1\},
\]
the cycle increment of the virtual chord: $\varepsilon(W)=+1$ iff
$u\sim v$ in $\gamma_W\tau_W$, i.e.\ $\varepsilon=2\,\mathrm{bit}-1$.
For $S\subseteq[k]$ write $\varepsilon_S=\varepsilon(\{Z_i:i\in S\})$.
Then $\mathrm{bias}=\tfrac12\,\E_S[\varepsilon_S]$ conditionally on
the positions, so with $S,T$ independent uniform subsets of $[k]$,
\begin{equation}\label{eq:bias2}
4\,\E\bigl[\mathrm{bias}^2\bigr]
=\E_{S,T}\,\E_{\mathrm{pos}}\bigl[\varepsilon_S\,\varepsilon_T\bigr].
\end{equation}

\begin{theorem}[two-point factorization]\label{thm:factor}
Fix $S,T\subseteq[k]$ and let $j=|S\cap T|$, $d_1=|S\setminus T|$,
$d_2=|T\setminus S|$.  Let $C=(Z_i)_{i\in S\cap T}$ and define the
\emph{fresh-averaged sign}
\[
M_d(C)\;:=\;\E\bigl[\varepsilon(C\cup F_d)\,\big|\,C\bigr],
\qquad F_d=d\text{ fresh iid chords}.
\]
Then
\[
\E_{\mathrm{pos}}\bigl[\varepsilon_S\varepsilon_T\bigr]
=\E_C\bigl[M_{d_1}(C)\,M_{d_2}(C)\bigr].
\]
Consequently, with
$\psi(j,d):=\E_{C_j}\bigl[M_d(C_j)^2\bigr]$ ($C_j=j$ iid chords),
\[
4\,\E[\mathrm{bias}^2]
\;\le\;\max_{\,j\le k,\ d\ge k/8}\psi(j,d)\;+\;2e^{-k/32}.
\]
\end{theorem}

\begin{proof}
Given $C$, the position blocks $(Z_i)_{i\in S\setminus T}$ and
$(Z_i)_{i\in T\setminus S}$ are independent, and
$\varepsilon_S$ is a function of $C$ and the first block only,
$\varepsilon_T$ of $C$ and the second only: so
$\E[\varepsilon_S\varepsilon_T\mid C]=M_{d_1}(C)M_{d_2}(C)$; take
expectations.  For the second claim: by Cauchy--Schwarz,
$\E_C[M_{d_1}M_{d_2}]\le\psi(j,d_1)^{1/2}\psi(j,d_2)^{1/2}
\le\max_{d\in\{d_1,d_2\}}\psi(j,d)$.  In \eqref{eq:bias2},
$(j,d_1,d_2)$ is multinomial with cell probabilities $1/4$, so
$d_1,d_2\ge k/8$ outside probability $2e^{-k/32}$ (Chernoff), where
$|\varepsilon_S\varepsilon_T|\le1$ absorbs the exceptional mass.
\end{proof}

\begin{corollary}\label{cor:psienough}
If $\psi(j,d)\le K\,d^{-2c}$ for all $j\le k$, $d\ge k/8$, then
$\E|\mathrm{bias}|\le\E[\mathrm{bias}^2]^{1/2}\le K'k^{-c}$.  With
the supply $k\ge6\log^{11}n$ of Lemma~\ref{lem:supply}, \emph{any}
$c>1/11$ closes the iid model of Problem~\ref{prob:annealed}.
\end{corollary}

The quantity $M_d$ has a Markov-semigroup description: if $P$
denotes the operator that adds \emph{one} fresh active chord and
averages, i.e.\ $(Pf)(C)=\E_z[f(C\cup z)]$, then by exchangeability
of the fresh chords
\[
M_d\;=\;P^{\,d}\,\varepsilon,
\qquad
\psi(j,d)\;=\;\bigl\|P^{\,d}\varepsilon\bigr\|_{L^2(\mu_j)}^2,
\]
$\mu_j$ the law of $j$ iid chords.  Adding a fresh chord is one step
of a coagulation--fragmentation dynamics on the cycle structure of
$\gamma\tau$ (cf.\ \cite{Sch05}); the one-point function
$\E[\varepsilon(F_m)]$ is a Harer--Zagier-type quantity
\cite{HZ86}.  The open core of the iid model is thus:

The array $\psi$ is monotone in a useful way:

\begin{lemma}[freezing dominates averaging]\label{lem:freeze}
$\psi(j,d+1)\le\psi(j+1,d)$ for all $j,d\ge0$; hence
$\psi(j,d)\le\psi(j+d-d',d')$ for all $0\le d'\le d$, and
$\bar\psi(d):=\sup_j\psi(j,d)$ is non-increasing in $d$.
\end{lemma}

\begin{proof}
Conditional Jensen: $(Pf)(C)^2=\E_z[f(C\cup z)]^2
\le\E_z[f(C\cup z)^2]=P(f^2)(C)$.  Hence
$\psi(j,d+1)=\E_{\mu_j}[(P(P^{d}\varepsilon))^2]
\le\E_{\mu_j}[P((P^{d}\varepsilon)^2)]
=\E_{\mu_{j+1}}[(P^{d}\varepsilon)^2]=\psi(j+1,d)$, since
integrating one fresh chord into $C\sim\mu_j$ produces exactly
$\mu_{j+1}$.  Iterate.
\end{proof}

\begin{problem}[polynomial $L^2$ mixing of fresh-chord
averaging]\label{prob:mixing}
Show that $\bar\psi(d)\le K d^{-2c}$ for some
$c>1/11$ (numerically true with $c\approx0.52$: the measured
$\psi(256,d)$ fits $d^{-1.05}$, and the $j$-dependence saturates;
see below).
\end{problem}

\emph{Status: Theorem~\ref{thm:factor} and
Corollary~\ref{cor:psienough} are PROVED (three-line conditioning
argument; exact in the iid ensemble).  Problem~\ref{prob:mixing} is
OPEN; the numerics below (COMPUTED, \texttt{s7\_twopoint.c}) give
$\psi\approx0.8/d$ with at most slow growth in $j$:}
\[
\begin{array}{lcccc}
\toprule
\psi(j,d) & d=4 & d=16 & d=64 & d=256\\
\midrule
j=16  & 0.124 & 0.019 & 0.0017 & 0.0004\\
j=64  & 0.174 & 0.041 & 0.0084 & 0.0010\\
j=256 & 0.197 & 0.050 & 0.0114 & 0.0027\\
j=1024&  --   & 0.064 & 0.0128 &  --  \\
\bottomrule
\end{array}
\]
\emph{Directly measured (cycle route, independent of the rank
route): $4\E[\mathrm{bias}^2]\cdot k=0.95,1.08,1.22,1.28,1.27,1.44$
at $k=8,16,32,64,128,256$, i.e.\
$\E[\mathrm{bias}^2]\approx0.33/k$, consistent with
$\E|\mathrm{bias}|\approx0.35k^{-1/2}$; the signed one-point
function $\E[\varepsilon(F_m)]$ is $+0.024,+0.005,+0.001$ at
$m=4,8,16$ and $\le10^{-3}$ (statistically $0$) for $m\ge16$.}

Two further remarks.  (i) The factorization argument uses exactly
one probabilistic fact: the positions attached to disjoint index
sets are independent.  Transferring it to the true orbit measure ---
where the chord positions are the stable-intercalate row positions
in the cycles of $\pi_P$, exchangeable-in-rows but not independent
--- therefore requires either a conditional-resampling argument
(given the shared intercalates, re-randomise the rest of the square)
or a quantitative correlation bound; this is the transfer problem
7a(i), not an additional obstacle discovered here.  (ii) The
short-arc orbits are visible in the factorization as the atom of
$C$-configurations with no fresh-chord separation available; in the
antipodal continuum ensemble this atom is absent, and in the true
ensemble it contributes the $O(1/k)$ term already accounted for in
\S\ref{sec:annealed}.

%----------------------------------------------------------------------
\subsection{No local genericity condition suffices}\label{sec:pendant}

The discussion at the end of \S\ref{sec:annealed} contemplated a
genericity condition of the form ``with high probability every
separating chord is interlaced by some other available chord''.
This is \emph{qualitatively insufficient}, in the strongest possible
sense: crossing multiplicity of any fixed order does not help.

\begin{proposition}[pendant stars]\label{prop:pendant}
For every $m\ge1$ and $r\ge1$ there is a configuration of
$k=r(m+1)$ disjoint chords in which every separating chord is
interlaced by exactly $m$ other chords, and
\[
\mathrm{bias}
=\bigl(1-2^{-(m+1)}\bigr)^{r}-\tfrac12
\;\xrightarrow[r\to\infty]{}\;-\tfrac12 .
\]
\end{proposition}

\begin{proof}
Take $r$ pairwise non-interlacing (nested) separating chords
$c_1,\dots,c_r$, and for each $i$ a cluster of $m$ pairwise nested
``pendant'' chords straddling one endpoint of $c_i$, short enough to
interlace $c_i$ and nothing else, and to separate nothing.  The
interlacement graph is a disjoint union of $r$ stars $K_{1,m}$ with
separation vector $e_{\mathrm{center}}$ on each; $\bar A_S$ is block
diagonal, so the bit is the product of the block bits.  In a block,
$x_S\in\operatorname{col}(A_S)$ fails iff the center is active and
all $m$ pendants are inactive (probability $2^{-(m+1)}$): the block
bit has $\Pr=1-2^{-(m+1)}$, and the blocks are independent.
\end{proof}

(Verified exactly for $m\le3$, $r\le4$: \texttt{s7\_bias\_lab.py
pendant}.)  To keep $|\mathrm{bias}|\le k^{-c}$ in this family one
needs $m\gtrsim\log r$: no condition ``every separating chord is
crossed by $\ge m$ others'' with $m$ fixed (or growing slower than
logarithmically) can imply the annealed bound.  An exhaustive scan of
\emph{all} chord configurations with $k\le6$ chords
(\texttt{s7\_bias\_lab.py exhaust}; $10{,}395$ matchings $\times$
all $13$ $v$-positions at $k=6$) confirms the moral.  Within the
classes cut by the natural local predicates --- number $r$ of
separating chords, minimum crossing number of a separating chord (by
all chords, or by separating chords only), connectivity or
$\mathbb{F}_2$-rank of the separating-chord interlacement subgraph
--- $|\mathrm{bias}|$ vanishes identically only on ladders and on
the class ``$r=1$ with the separating chord crossing \emph{nothing}'',
where the cut-chord identity of \S\ref{sec:annealed} gives
$\mathrm{bias}=\tfrac12\cdot(+\tfrac12)-\tfrac14=0$ exactly (the
remainder has empty separation vector, hence conditional bias
$+\tfrac12$).  At $k=6$, every other class with $r\ge1$ contains configurations
with $|\mathrm{bias}|\ge1/16$; the only classes whose maxima stay
below $15/64$ are the all-separating ($r=k$) high-crossing classes,
and the class maxima grow with $k$ (e.g.\ the class ``$r=2$, minimum
crossing $1$'' has maxima $0.188,0.219,0.234$ at $k=4,5,6$).  The bound must come from the measure; the moment
route of \S\ref{sec:factor} is the appropriate formalisation.

Relatedly, the single-coordinate toggle bound of
Proposition~\ref{prop:bias} cannot by itself produce decay: in the
generic ensemble the measured $\E[\min_i\Pr(i\text{ does not
toggle})]$ is $\approx0.35$--$0.40$ and \emph{non-decreasing} over
$k=4,\dots,12$ (\texttt{s7\_bias\_lab.py toggle}), while
$\E|\mathrm{bias}|$ already decays over the same range.

%----------------------------------------------------------------------
\subsection{Template averaging is exactly neutral}
\label{sec:vacuity}

Remark~\ref{rem:bookkeeping}(i) proposed to handle the
square-dependence of the exceptional columns of
Lemma~\ref{lem:supply} by relabelling the template by a uniform
column permutation $g$ and averaging.  This does not work:

\begin{proposition}[vacuity of template averaging]
\label{prop:vacuity}
Fix $n,\lambda,\alpha$ and a template $T$.  For a mark
$(L,P)$ and a template $T'$, let
$\mathrm{bad}(T';L,P)\in\{0,1\}$ indicate that a column of $P$ fails
the supply property of Lemma~\ref{lem:supply} for the partial square
$T'\cap L$ (with the exceptional set defined intrinsically: the set
of columns $c$ that do not carry $t$ disjoint stable intercalates
avoiding rows $1,2$ and the partner column).  Then for
\emph{every} column permutation $h$,
\[
\E_{(L,P)\sim\mathrm{Unif}(X)}\bigl[\mathrm{bad}(h(T);L,P)\bigr]
=\E_{(L,P)\sim\mathrm{Unif}(X)}\bigl[\mathrm{bad}(T;L,P)\bigr].
\]
In particular averaging over uniform $g$ leaves the $X$-averaged
bad-mark fraction unchanged, and the same holds verbatim when $T$ is
itself random with a column-exchangeable law (e.g.\ the iid
Bernoulli templates behind \cite[Lem.~8.3]{KPS25}).
\end{proposition}

\begin{proof}
Column permutations act on squares by $h(L)(r,h(c))=L(r,c)$.  This
action fixes $\sigma$ as a permutation of symbols (the pair
$(L(1,c),L(2,c))$ moves with the column), hence preserves the type
and maps $S(\lambda)$ bijectively onto itself; marked pairs move
covariantly, $P\mapsto h(P)$, so $(L,P)\mapsto(h(L),h(P))$ is a
bijection of $X$.  Moreover $h(T)\cap h(L)=h(T\cap L)$, and the
supply property is intrinsic (rows are untouched), so
$\mathrm{bad}(h(T);h(L),h(P))=\mathrm{bad}(T;L,P)$.  Substituting
$(L,P)=(h(L_0),h(P_0))$ in the left-hand average gives the claim.
\end{proof}

The remark's proposed fix is therefore withdrawn.  What the assembly
actually needs is a \emph{decoupling} between the supply structure
and the mark distribution:

\begin{problem}[supply--mark decoupling]\label{prob:decouple}
Show that
$\Pr_{(L,P)\sim\mathrm{Unif}(X(n,\lambda;\alpha,\beta))}
\bigl[\mathrm{bad}(T;L,P)\bigr]=o(1/\log n)$, uniformly over the
admissible $(\lambda,\alpha,\beta)$ --- including types whose
$m$-cycle class occupies as few as $m=O(1)$ columns, where the
deterministic bound $|E|\le6\log^{22}n$ of Lemma~\ref{lem:supply} is
vacuous.
\end{problem}

Heuristically Problem~\ref{prob:decouple} should hold with much
room: the supply structure of $T\cap L$ is (almost) a function of
rows $3..n$, while the class membership of a column is a function of
rows $1,2$.  A concrete framing: condition on the bottom
$(n-2)\times n$ rectangle $L_\downarrow$.  This fixes the
``leftover'' $2$-regular bipartite graph on columns whose cycles are
completed to rows $1,2$ by independent binary orientation choices
(tilted by the $S(\lambda)$ conditioning), so the $m$-class columns
are a function of $L_\downarrow$ plus these orientations, while the
exceptional set is a function of $L_\downarrow$ \emph{up to} the
weak dependence of stability (isolation and non-criticality in
\cite[\S7]{KPS25}) on intercalates of $T\cap L$ meeting rows $1,2$.
Quantifying either the orientation-randomness or a per-column
failure probability is open; note that a per-column version of
Lemma~\ref{lem:supply} cannot be extracted from the expander
property alone, which bottoms out at exceptional sets of size
$\ell$, nor from conditional division as in
Lemma~\ref{lem:condexp}, since the per-tuple failure rates inside
\cite[Lem.~8.3]{KPS25} are far too weak to survive division by
$\Pr[S(\lambda)]\ge e^{-O(n\log n)}$.

\emph{Status: Proposition~\ref{prop:vacuity} PROVED;
Problem~\ref{prob:decouple} OPEN (new, replaces
Remark~\ref{rem:bookkeeping}(i)'s proposed fix).  Session-8 empirical
support (\texttt{s8\_real\_bias.py}, \texttt{s8\_dec\{30,50\}.csv}):
on JM-uniform squares the supply $k$ of a marked instance is
statistically independent of the $\sigma$-cycle length $m$ of its
marked pair --- $\mathrm{corr}(k,m)=+0.006\pm0.018$ at $n=30$
($3200$ instances), $-0.034\pm0.022$ at $n=50$ ($2000$), with
$\E[k\mid m]$ flat across the full range of $m$ and no supply-poor
instances at all ($\Pr[k\le4]=0$ in-sample) --- consistent with the
heuristic that supply is a rows-$3..n$ property and class membership
a rows-$1,2$ property.  The programme
summary of \S\ref{sec:orbit} now carries two open items:
Problem~\ref{prob:annealed} (main) and Problem~\ref{prob:decouple}
(bookkeeping).}

%----------------------------------------------------------------------
\subsection{Real-ensemble grounding at \texorpdfstring{$n=30$--$100$}
{n=30-100}}\label{sec:realbias}

Finally, the master formula and the bias statistics were taken to
\emph{real} random squares (\texttt{s7\_real\_bias.py}; uniform
sampling by Jacobson--Matthews \cite{JM96}, marked pairs
$P=\{j,\sigma^\alpha(j)\}$ sampled uniformly among admissible
instances, greedy maximal cell-disjoint intercalate families through
the marked columns in rows $3..n$, orbit data built exactly as in
\texttt{pf\_master.py}).

\[
\begin{array}{lccccccc}
\toprule
n & \text{inst.} & \bar k & \E|\mathrm{bias}| &
\E|\mathrm{bias}|\,(r{>}0) & \Pr[r{=}0] & \text{signed mean} &
\E[\mathrm{bias}^2]\cdot\bar k\\
\midrule
30  & 4200 & 12.9 & 0.168 & 0.145 & 0.065 & +0.038 & 0.67\\
50  & 3300 & 22.6 & 0.123 & 0.108 & 0.039 & +0.023 & 0.75\\
100 & 1800 & 26\ (\text{capped}) & 0.114 & 0.101 & 0.029 & +0.013
& 0.73\\
\bottomrule
\end{array}
\]
(Families capped at $26$ coordinates; $k>14$ instances by MC over
$1500$--$2000$ activation subsets.)  Findings:
\begin{enumerate}[leftmargin=2em]
\item \textbf{The master formula holds on real squares at
$n\le100$}: $3250$ random activation subsets across $n=30,50,100$
checked against the actually-switched square --- $0$ mismatches.
\item \textbf{The annealed decay survives on the real ensemble}:
$\E[\mathrm{bias}^2]\cdot\bar k\approx0.7$ across $n$, i.e.\ the
$1/k$ second-moment law of the iid ensemble holds with a constant
$\approx2.2\times$ larger; equivalently
$\E|\mathrm{bias}|\approx1.5$--$1.6\times$ the iid value at the same
$k$.  At fixed $k=26$ the value is the same at $n=50$ and $n=100$:
the elevation is a property of the orbit position law at given
supply, not an $n$-effect.
\item \textbf{The short-arc atom behaves exactly as predicted in
\S\ref{sec:annealed}}: the orbits with no separating chord ($r=0$,
$\mathrm{bias}=+\tfrac12$ exactly) have measure $\approx0.85/\bar k$
($6.5\%,3.9\%,2.9\%$), and the \emph{signed} mean bias is accounted
for by this atom alone ($+0.038$ vs.\ atom$\times\tfrac12=+0.032$;
$+0.023$ vs.\ $+0.020$; $+0.013$ vs.\ $+0.015$): the annealed
average is self-regularizing, no separate treatment needed.
\item Cut-chord configurations occur at frequency
$2.1\%,1.6\%,1.1\%$ (decreasing in supply), and no new pathology
class appeared: the per-instance bias histogram is consistent with
the iid shape plus the $r=0$ atom.
\end{enumerate}

\emph{Status: COMPUTED (tables above; per-instance data in
\texttt{s7\_inst\{30,50,100\}.csv}); the ground-truth re-verification
of the master formula at $n=30,50,100$ is VERIFIED (0/3250).
Session-8 extension to $n=200$ (\texttt{s8\_inst200.csv}, $144$+
instances at $k=26$ capped, JM thinning $2n^3$): $\E|\mathrm{bias}|
=0.118$, $\E[\mathrm{bias}^2]\cdot k=0.87\pm0.15$, signed mean
$-0.006\pm0.014$ --- the fixed-$k$ values agree with $n=50,100$
($0.107$--$0.116$): the elevation is confirmed to be a $k$-level
property out to $n=200$, and the $1/k$ law shows no $n$-drift.}


%======================================================================
\section{The mixing problem as a marked split--merge chain: symmetry,
factorization, and the exact rate from below}
\label{sec:mixing}

This section records the session-8 progress on
Problem~\ref{prob:mixing}.  The fresh-chord averaging operator $P$ of
\S\ref{sec:factor} is identified \emph{exactly} with one step of the
uniform split--merge (coagulation--fragmentation) dynamics on measured
partitions of the circle, carrying two marked points
(Lemma~\ref{lem:splitmerge}); the chain admits an exact sign-reversing
involution (Theorem~\ref{thm:iota}) which factors the fresh-averaged
sign as $M_d=\varepsilon_0\cdot\Phi_d(\text{geometry})$
(Corollary~\ref{cor:quotient}) and removes the label from the problem;
the stationary law of the relevant ``gap'' coordinate is computed
exactly (linear density; Proposition~\ref{prop:gaplawstat}); and these
combine into a \emph{proved} lower bound
$\bar\psi(d)\ge c/d$ (Proposition~\ref{prop:lower}): the numerically
measured rate $\psi\approx1.0/d$ is sharp, and $c=1/2$ is the true
exponent of Problem~\ref{prob:mixing} if the matching upper bound
holds.  The upper bound is reduced to two sharply scoped sub-problems
(Problems~\ref{prob:gaplaw} and~\ref{prob:survival}), each with the
mechanism confirmed numerically, plus a spectral route via
reversibility (\S\ref{sec:spectral}).  \emph{Session 9
(\S\ref{sec:s9}) shows the first of the two is unnecessary over the
horizon $j\le8d$ relevant to Theorem~\ref{thm:factor}: any
$\alpha>2/9$ in Problem~\ref{prob:survival} closes the iid model.
Session 10 (\S\ref{sec:s10}) shows the environment must enter that
problem through recharge, proves that the block count on the horizon is
$O(\log d)$ (the Harer--Zagier law of the number of curves), so the
environment is healthy where it is needed, and reduces
Problem~\ref{prob:survival} to two conditions (R1)--(R2) on the
log-macroscopic post-flip chain via an exact Markov-renewal identity.}

Throughout, the ensemble is that of \S\ref{sec:factor}: $u,v$ fixed on
the circle of unit circumference ($u$ at $0$, $v$ at $1/2$ in the
code), chords $=$ iid uniform endpoint pairs.

%----------------------------------------------------------------------
\subsection{Identification with the uniform split--merge chain}
\label{sec:smident}

For a finite chord set $W$, the cycles of $\gamma\tau_W$ trace out a
partition of the circle (minus the endpoint set) into finite unions of
arcs (``curves''); each curve carries its Lebesgue measure, and its
internal cyclic order (the order in which $\gamma\tau_W$ traverses its
arcs) transports Lebesgue measure to a uniform measure along the
curve.  The \emph{marked state} is
\[
X \;=\;\Bigl(\{\text{curve masses}\}, \ \text{curves of }u,v,\
\begin{cases}
(x,y) & u\sim v:\ \text{masses of the two }u\to v\text{ curve arcs},\\
(p,q) & u\not\sim v:\ \text{masses of the curves through }u,\ v,
\end{cases}\Bigr)
\]
and $\varepsilon(W)=+1$ iff $u\sim v$.

\begin{lemma}[fresh chord $=$ one split--merge move]
\label{lem:splitmerge}
Adding one fresh chord (iid uniform endpoints $c_1,c_2$) maps the
marked state by the following rules, which depend on the current state
only.  Write $B(c)$ for the curve containing $c$; positions of $c_i$
along their curves are independent uniform.
\begin{enumerate}[leftmargin=2em]
\item $B(c_1)=B(c_2)=B$ (probability $|B|^2$): $B$ splits into the two
arcs of $B$ between the cuts.  Sub-cases:
 (a) $B$ unmarked, or marked by exactly one of $u,v$ (mass $p$): the
 mark keeps the piece of mass $p-|t_1-t_2|$, $t_i\sim U(0,p)$ iid; the
 complementary piece of mass $|t_1-t_2|$ joins the environment.
 (b) $B$ the common block, sides $(x,y)$: with the cycle parametrised
 $u\to v$ ($x$-side) then $v\to u$ ($y$-side), cuts uniform on
 $(0,x{+}y)$: both cuts on one side shrinks that side as in (a);
 one cut per side (probability $2xy/(x{+}y)^2$ given the split)
 \emph{separates} $u,v$: the new blocks have masses
 $q'=t_2-t_1$ ($v$'s) and $p'=x+y-q'$ ($u$'s), i.e.\
 $(p',q')\overset{d}{=}(\alpha+y-\beta,\ x-\alpha+\beta)$,
 $\alpha\sim U(0,x)$, $\beta\sim U(0,y)$.
\item $B(c_1)\ne B(c_2)$: the two curves merge (cut-and-join at
$c_1,c_2$).  Sub-cases: (a) at most one is marked, or one holds $u$
and the other neither: masses add, marks carried along.  (b) $u$'s and
$v$'s blocks (masses $p,q$) merge: $u,v$ become joined with sides
$(x,y)=(\alpha+\beta,\ p+q-\alpha-\beta)$, $\alpha\sim U(0,p)$,
$\beta\sim U(0,q)$.  (c) the common block (sides $x,y$) merges with an
environment block of mass $z$: the insertion point is on the $x$-side
with probability $x/(x{+}y)$, and that side gains $z$.
\end{enumerate}
In particular $P$ (add one fresh chord and average) is the transition
operator of this \emph{marked uniform split--merge chain}, the
continuum limit of the cycle structure of the random-transposition
walk (cf.\ \cite{Sch05}), and
$M_d=P^d\varepsilon$ is its $d$-step averaged together-indicator sign.
\end{lemma}

\begin{proof}
Adding the chord $z=(c_1,c_2)$ replaces $\gamma\tau_W$ (computed on
the endpoint set enlarged by $c_1,c_2$; refining $\gamma$ at
non-endpoint positions does not change the curve partition) by
$\gamma\tau_W\,(c_1c_2)$.  Right multiplication by a transposition
whose entries lie on one cycle splits it into the two arcs between
them, and merges two cycles otherwise --- with the traversal orders
spliced at the cuts; this is the standard surgery computation.  The
probabilistic statements: $c_1,c_2$ are iid Lebesgue, the curves
partition the circle, and Lebesgue restricted to a curve is uniform
along it; all sub-case laws then read off from the traversal order,
e.g.\ in 1(b) the cut positions along the cycle are iid $U(0,x{+}y)$
and $u,v$ separate iff the cuts straddle both sides.
\end{proof}

\emph{Status: PROVED.  VERIFIED end-to-end
(\texttt{s8\_splitmerge.c}): the continuum chain reproduces the
discrete-surgery values of $\psi(j,d)$ (fs$=1$ unbiased runs;
$\psi(16,4)$: $0.1341(22)$ vs $0.1305(22)$;
$\psi(64,16)$: $0.0391(26)$ vs $0.0395(26)$), the one-point function
($f(4)$: $+0.0232(4)$ vs $+0.0237(5)$), and the
Poisson--Dirichlet sanity values at stationarity
($\E\sum p_i^2=0.4979$ vs $\mathrm{PD}(1)$'s $1/2$;
$\E[\text{common-block mass}\mid u\sim v]=0.6685$ vs $2/3$, see
Proposition~\ref{prop:gaplawstat}).  NOTE: repeat runs show the
session-7 $\psi$ table of \S\ref{sec:factor} carries sampling scatter
$\approx\pm0.005$ (its inner-average standard errors were optimistic);
its fitted exponent is unaffected.}

%----------------------------------------------------------------------
\subsection{The sign-reversing involution and the quotient chain}
\label{sec:iota}

\begin{theorem}[label-flip symmetry]\label{thm:iota}
Let $\iota$ be the involution of marked states that swaps the label
$\{u\sim v\}\leftrightarrow\{u\not\sim v\}$ while keeping the
environment and the ordered geometry pair fixed:
\[
\iota:\ \bigl(\text{together},(x,y),\mathrm{env}\bigr)
\;\longleftrightarrow\;
\bigl(\text{apart},(p,q)=(x,y),\mathrm{env}\bigr).
\]
Then $\iota$ commutes with the transition kernel of the marked
split--merge chain, and $\varepsilon\circ\iota=-\varepsilon$.
Consequently $M_d\circ\iota=-M_d$ for every $d$.
\end{theorem}

\begin{proof}
Match the transition laws of Lemma~\ref{lem:splitmerge} case by case
under the correspondence $x\leftrightarrow p$ (the $u$-coordinate),
$y\leftrightarrow q$ ($v$-coordinate).  The total marked mass is
$x+y=p+q$ in both states, so the probabilities of the events ``both
cuts in marked material'', ``one cut in marked material, one in a
given environment block'', ``both in environment'' agree, as does the
finer classification by $u$- vs $v$-material ($x^2/(x{+}y)^2$ etc.,
against $p^2/(p{+}q)^2$ etc.).  Given the event: both cuts in
$u$-material shrinks $x$ (resp.\ $p$) by $|t_1-t_2|$ with the same
law; a cut pair straddling $u$- and $v$-material is the
\emph{separation} in the together state and the \emph{join} in the
apart state --- i.e.\ it flips the label in both --- and the outgoing
geometry laws agree:
$(\alpha+y-\beta,\,x-\alpha+\beta)
\overset{d}{=}(\alpha+\beta',\,x+y-\alpha-\beta')$ jointly, by
$\beta\mapsto y-\beta'$; a single cut in $u$-material merges an
environment block $z$ into the $u$-coordinate in both states; and
environment-only moves are identical.  This exhausts all cases, so
$K(\iota s,\iota\,\cdot)=\iota_*K(s,\cdot)$; $\varepsilon$ flips by
definition of $\iota$, and $M_d=P^d\varepsilon$ inherits oddness.
\end{proof}

\emph{Status: PROVED (case check as above).  VERIFIED
(\texttt{s8\_splitmerge iota}): from mirrored starts
$(\mathrm{tog},(x,y),z)$ and $(\mathrm{apart},(x,y),z)$, measured
$P_A(u\sim v)+P_B(u\sim v)=1$ within Monte Carlo error at all tested
$t\le32$ and geometries, with identical $u$-coordinate marginals.}

\begin{corollary}[quotient factorization; the label carries all the
sign]\label{cor:quotient}
Let $\Gamma_t=(a_t,b_t;\mathrm{env}_t)$ be the quotient of the marked
chain by $\iota$ (the geometry with the label forgotten;
$a\leftrightarrow x$ or $p$, $b\leftrightarrow y$ or $q$).  Then:
\begin{enumerate}[leftmargin=2em]
\item $\Gamma_t$ is Markov, with the transition laws of
Lemma~\ref{lem:splitmerge} read on $(a,b;\mathrm{env})$; the
\emph{flip moves} (cut pair straddling $a$- and $b$-material) occur
with probability $2a_tb_t$ per step and update
$(a,b)\to(\alpha+\beta,\,a+b-\alpha-\beta)$.
\item The label evolves deterministically along the quotient path:
$\varepsilon_d=\varepsilon_0\cdot(-1)^{N_d}$ with $N_d$ the number of
flip moves, a functional of the quotient path alone.
\item Hence, with
$\Phi_d(\Gamma):=\E_\Gamma\bigl[(-1)^{N_d}\bigr]$,
\[
M_d(C)\;=\;\varepsilon(C)\cdot\Phi_d\bigl(\Gamma(C)\bigr),
\qquad
\psi(j,d)\;=\;\E_{\mu_j}\bigl[\Phi_d(\Gamma_0)^2\bigr].
\]
\item Every stationary law of the marked chain is $\iota$-invariant
wherever uniqueness holds on its support; in particular the label is
a fair coin given the geometry at stationarity.
\end{enumerate}
\end{corollary}

\begin{proof}
(1),(2): the flip moves are exactly the label-changing moves in both
sectors (Theorem~\ref{thm:iota}), and their effect on the geometry has
the stated common law; all other moves preserve the label.  (Flip
moves are a.s.\ identifiable on the quotient path: they are the only
moves that change both marked coordinates.)  (3):
condition on the quotient path.  (4): $\iota$ pushes any stationary
law to a stationary law with the same quotient marginal.
\end{proof}

\emph{Status: PROVED.  Problem~\ref{prob:mixing} is thereby freed of
the sign structure: it is now a statement about the \emph{parity}
functional $\Phi_d=\E[(-1)^{N_d}]$ of an autonomous
positive-recurrent chain.  Consistency check: at stationarity the
label-fairness (4) predicts $P(u\sim v)=1/2$; measured $0.5065(29)$ at
burn-in $8192$.}

\begin{remark}[renewal structure at flip times]\label{rem:renewal}
Let $\tau$ be the first flip move.  Strong Markov gives the exact
renewal identity
\[
\Phi_d(\Gamma)\;=\;\Pr_\Gamma[\tau>d]\;-\;
\E_\Gamma\bigl[\Phi_{d-\tau}(\Gamma_\tau);\ \tau\le d\bigr].
\]
The flip clock runs at rate $2a_tb_t$; after a flip the $u$-coordinate
is \emph{recharged} ($a'=\alpha+\beta$ is macroscopic unless both
$a,b$ were small).  The correlation-survival mechanism is therefore
gap-trapping: $\Pr[\tau>d]$ is dominated by trajectories on which
$g_t=\min(a_t,b_t)$ stays $O(1/d)$, since an untouched small side
freezes ($\Pr[\text{the }g\text{-material is hit in }d\text{ steps}]
\le 4gd$), while a touched one is typically recharged to macroscopic
size by a size-biased merge.
\end{remark}

%----------------------------------------------------------------------
\subsection{The stationary gap law, and the exact rate from below}
\label{sec:gaplaw}

\begin{proposition}[stationary gap law]\label{prop:gaplawstat}
Let $\pi$ be the stationary law of the marked chain (the $n\to\infty$
limit of the cycle structure of a uniform permutation of $n$ circle
points with two marked points; top of the size spectrum
$\mathrm{PD}(1)$).  Under $\pi$: $\Pr[u\sim v]=\tfrac12$; given
$u\sim v$ the common-block mass $s$ has density $2s$ on $(0,1)$ (so
$\E[s\mid u\sim v]=2/3$) and $x\mid s\sim U(0,s)$; given $u\not\sim v$
the pair $(p,q)$ is uniform on the triangle $\{p,q>0,\ p+q<1\}$
(so the joint density of $(\text{apart},p,q)$ is $1$ there, mass
$\tfrac12$).  Consequently, with
$g=\min(x,y)$ resp.\ $\min(p,q)$,
\[
\Pr_\pi[g\le\delta]\;=\;4\delta+O(\delta^2)\qquad(\delta\to0),
\]
and the density of $g$ at $0^+$ is $4$.
\end{proposition}

\begin{proof}
At $n$ points: for uniform $\sigma$, $\Pr[u,v$ in a common $m$-cycle$]
=(m-1)/\bigl(n(n-1)\bigr)$ (choose and order the cycle), summing to
$1/2$; given the cycle, the $u\to v$ distance is uniform on
$\{1,\dots,m-1\}$; if apart,
$\Pr[|C_u|=m_1,v\notin C_u]=\tfrac1n\cdot\tfrac{n-m_1}{n-1}$ and
given this $|C_v|$ is uniform on $\{1,\dots,n-m_1\}$, so the joint
law of $(p,q)=(m_1,m_2)/n$ converges to
$(1-p)\,dp\cdot\tfrac{dq}{1-p}=dp\,dq$ on the triangle.  Pass to the
limit ($m/n\to s$).  The
gap: given together,
$\Pr[\min(x,y)\le\delta]=\int_0^1 2s\min(1,2\delta/s)\,ds
=4\delta-4\delta^2$; given apart,
$\Pr[\min(p,q)\le\delta]=4\delta+O(\delta^2)$ likewise; average the
two sectors.
\end{proof}

\emph{Status: PROVED at finite $n$ with a routine limit exchange
(the marked chain is the lumping of the transposition walk, see
\S\ref{sec:spectral}; over any fixed horizon the kernels converge).
VERIFIED: measured dyadic masses at stationarity
($\Pr[g\in(2^{-6},2^{-5})]$: together $0.0297$ vs predicted $0.0312$,
apart $0.0288$ vs $0.0312$; lower bins consistent with density $4$
throughout), and the label is a fair coin in every gap bin.}

\begin{proposition}[the mixing rate is at most $1/d$: lower bound]
\label{prop:lower}
There are $c_0>0$, $d_0$ such that the stationary-start mixing
functional obeys
\[
\psi_\pi(d)\;:=\;\E_\pi\bigl[M_d^2\bigr]\;\ge\;\frac{c_0}{d},
\qquad d\ge d_0 .
\]
Hence $\bar\psi(d)\ge c_0/d$ along any initial family
$\{\mu_j\}$ whose gap law dominates a multiple of the stationary one
(numerically true for the iid-chord $\mu_j$ already at $j=16$), and
no exponent $c>1/2$ is attainable in Problem~\ref{prob:mixing}.
\end{proposition}

\begin{proof}
Fix $K$ and let
$G=\{g_0\le1/(Kd)\}$.  On $G$, the $g$-material (the short side if
together, the small block if apart) is hit by no cut during $d$ steps
except with probability $\le 4g_0d\le4/K$: each step's two cuts avoid
it with probability $\ge1-4g_0$ while it remains unhit, and while
unhit its mass is frozen.  If it is never hit there is no flip move
(a flip requires a cut in both $a$- and $b$-material), so
$\varepsilon_d=\varepsilon_0$ and, for the conditional expectation,
$|M_d|=|\E[\varepsilon_d\mid X_0]|\ge1-8/K$ on $G$.  Therefore
$\psi_\pi(d)\ge\E[M_d^2;G]\ge(1-8/K)^2\,\pi(G)$, and
$\pi(G)\ge c_1/(Kd)$ by Proposition~\ref{prop:gaplawstat}.  Take
$K=16$.
\end{proof}

\emph{Status: PROVED (given Proposition~\ref{prop:gaplawstat}).  This
settles the sharpness question left open in \S\ref{sec:factor}: the
measured $\psi\approx0.8/d$ is the true order, the budget margin of
Corollary~\ref{cor:psienough} is exactly $c=1/2$ vs $c>1/11$
(factor $5.5$ in the exponent), and no route can do better than
$1/d$ --- in particular any proof of Problem~\ref{prob:mixing} must
tolerate the gap-trapped trajectories, they are really there.}

%----------------------------------------------------------------------
\subsection{What remains for the upper bound}
\label{sec:upper}

The decomposition
$\psi(j,d)=\E[\Phi_d^2;\,g_0<\varepsilon]+\E[\Phi_d^2;\,g_0\ge
\varepsilon]$ with $\varepsilon\asymp1/d$, together with
$|\Phi_d|\le1$, reduces the upper bound to two independent
statements.

\begin{problem}[uniform gap law]\label{prob:gaplaw}
Show $\Pr_{\mu_j}[g_0\le\varepsilon]\le C\varepsilon^{\kappa}$ (some
$\kappa\in(0,1]$, uniformly in $j$) for the iid-chord initial laws
$\mu_j$ --- equivalently for the chain run $j$ steps from the
one-block start.
\end{problem}

\emph{Partial progress: (i) at stationarity this is
Proposition~\ref{prop:gaplawstat} with $\kappa=1$; (ii) the per-step
\emph{creation} flux into $\{g\le\varepsilon\}$ has the universal
bound $\Pr[\text{step creates }g'\in d\varepsilon]\le
C\varepsilon\,d\varepsilon$ for every current state (each channel ---
shrink of a marked coordinate, flip recharge landing small --- has
outgoing density $\le2\varepsilon$ at $\varepsilon$: e.g.\ a split of
the $u$-block of mass $p$ leaves $u$ with $p-|t_1-t_2|$, of density
$2\varepsilon'/p^2\cdot p^2$ at $\varepsilon'$); (iii) escape from
$\{g\le\varepsilon\}$ happens at rate $\ge2g(1-g)\cdot
\Pr[\text{merge partner macroscopic}]$, so a renewal bound gives
$\kappa=1$ \emph{given} a uniform-in-time ``no-dust'' estimate (the
environment retains a macroscopic block); it is the no-dust step that
is open, and only it.  Numerically the dyadic gap masses under
$\mu_j$ match the stationary density already at $j=16$.}

\begin{lemma}[recharge estimates]\label{lem:recharge}
For the quotient chain: (i) flip moves preserve the marked total
$s=a+b$; (ii) a flip from $(a,b)$ lands at gap
$g'=\min(\alpha+\beta,\,s-\alpha-\beta)$ with
$\Pr[g'\le\eta]\le\eta^2/(ab)$ for $\eta\le\min(a,b)$; (iii) from
\emph{any} state, each move channel sends a marked coordinate into
$(0,\eta]$ from above with probability $\le\eta^2$ per step (e.g.\
the shrink channel: both cuts in the $a$-material has probability
$a^2$, and $|t_1-t_2|\in[a-\eta,a)$ has conditional probability
$\le\eta^2/a^2$); so the total creation flux into
$\{g\le\eta\}$ is $\le4\eta^2$ per step.
\end{lemma}

\begin{proof}
(i) is immediate from the flip law.  (ii): $\alpha+\beta\le\eta$ has
probability $\eta^2/(2ab)$ ($\alpha\sim U(0,a)$, $\beta\sim U(0,b)$
independent), likewise the other end.  (iii): direct integration of
the stated densities, channel by channel (shrink of $a$, shrink of
$b$, and the two flip ends).
\end{proof}

\emph{Status: PROVED.  Together with Remark~\ref{rem:renewal} this
pins the intended endgame: from a good start the first-flip time
$\tau$ has tail $\asymp1/t$ (trapping is the only way to stall, and
Lemma~\ref{lem:recharge}(iii) prices its creation), each flip
recharges the gap with quadratic safety (ii), so the flips form an
approximate renewal process with infinite-mean, regularly-varying
inter-arrival tail $\sim c/t$; for such renewals the alternating sum
$\E[(-1)^{N_d}]$ collapses to $\Theta(\text{tail}(d))=\Theta(1/d)$
(the signed renewal measure $\sum_k(-1)^kG^{*k}$ has total mass
$\tfrac12$).  Two genuine gaps remain: the inter-flip times are
state-dependent rather than iid, and the escape-rate lower bound
$\ge c\,g$ behind the $1/t$ tail needs the environment to offer a
macroscopic merge partner --- a ``no-dust'' estimate; at stationarity
dust configurations have probability $\approx$ the Dickman mass
$\rho(8)\sim10^{-8}$, small but not $o(1)$ in $d$, so the assembly
must let dust episodes delay rather than veto escape (dust coarsens
under merges).  Note also the second trap coordinate: the marked total
$s$ itself can be small (then flips are rare no matter the gap); but
$\Pr_\pi[s\le\delta]=O(\delta^2)$ in both sectors (the together mass
$s$ has density $2s$, and $(p,q)$ is uniform on the triangle), so the
$s$-trap contributes only $O(1/d^2)$ and rides along harmlessly in
the renewal bookkeeping.}

\begin{problem}[parity decay off small gaps]\label{prob:survival}
Show $\E_{\mu_j}\bigl[\Phi_d(\Gamma_0)^2;\ g_0\ge\varepsilon\bigr]\le
C(\varepsilon d)^{-2\alpha}$ for some $\alpha>0$ (uniformly in $j$),
i.e.\ trajectories started with both marked coordinates
$\ge\varepsilon$ lose the parity memory at a polynomial rate in
$\varepsilon d$.
\end{problem}

Any pair $(\kappa,\alpha)$ solves Problem~\ref{prob:mixing}: from
$\psi\le C\bigl(\varepsilon^\kappa+(\varepsilon d)^{-2\alpha}\bigr)$,
optimising at $\varepsilon=d^{-2\alpha/(\kappa+2\alpha)}$ gives
$\psi\le Cd^{-2c'}$ with $c'=\alpha\kappa/(\kappa+2\alpha)$; the
budget needs only $c'>1/11$, e.g.\ $(\kappa,\alpha)=(1,1/2)$ gives
$c'=1/4$, and any $\kappa=\alpha>3/10$ still closes
($c'=\kappa/3>1/11$ there).  \emph{The mechanism behind
Problem~\ref{prob:survival} is confirmed numerically: binning
$\psi$-contributions by $g_0$ (\texttt{s8\_splitmerge psibar}) shows
$\E[\Phi_d^2\mid g_0]$ is a scaling function of $g_0d$ alone ---
$\approx1$ for $g_0d\lesssim1/4$, $\lesssim0.05$ for $g_0d\gtrsim4$
--- so the conditional decay is fast; what is missing is a proof
device for the parity cancellation (the renewal identity of
Remark~\ref{rem:renewal} plus recharge is the intended route: after
the first flip the state is macroscopic w.h.p., flips then recur at
rate $\Theta(1)$, and the alternating sum cancels to the heavy-tail
term $\Pr[\tau>d]$).}

%----------------------------------------------------------------------
\subsection{The spectral route}\label{sec:spectral}

The discrete precursor of the marked chain --- right multiplication
of $\sigma_t\in S_n$ by uniform transpositions, marked cycle
structure observed --- is a strongly lumpable function of a chain
that is reversible with respect to the uniform measure; the lumped
chain is therefore reversible with respect to its stationary law, and
this persists in the $n\to\infty$ limit chainwise (the continuum
statement deserves its own write-up; nothing below is used elsewhere).
Granting it: $\psi_\pi(d)=\int\lambda^{2d}\,d\nu_\varepsilon(\lambda)$
for the spectral measure $\nu_\varepsilon$ of $\varepsilon$ under the
self-adjoint $P$, and Problem~\ref{prob:mixing} (stationary start)
becomes: \emph{$\nu_\varepsilon$ has at most polynomial mass near
$|\lambda|=1$}, with Proposition~\ref{prop:lower} asserting at least
linear mass near $\lambda=1$.  Two structural facts make this
attractive.  (i) By Theorem~\ref{thm:iota} the sign lives in the
$\iota$-odd sector, where the Dirichlet form has a \emph{massive}
term on flip moves: for $f=\varepsilon\cdot F(\Gamma)$,
\[
\mathcal{E}(f,f)=\tfrac12\,\E_\pi\Bigl[(F-F')^2\,;\,\text{non-flip}
\Bigr]+\tfrac12\,\E_\pi\Bigl[(F+F')^2\,;\,\text{flip}\Bigr],
\]
so a low-energy odd function must be small wherever the flip rate
$2ab$ is not, i.e.\ must localise on $\{g\lesssim\delta\}$, whose
mass is $\asymp\delta$ (Proposition~\ref{prop:gaplawstat}): this is
precisely the geometric picture of a linear spectral density, and
trial functions $F=\mathbf{1}\{g\le\delta\}$ realise it (Rayleigh
quotient $\asymp\delta$ at mass $\asymp\delta$), reproving
Proposition~\ref{prop:lower} spectrally.  (ii) The converse direction
--- \emph{every} odd $F$ with
$\mathcal{E}\le\delta\|F\|^2$ has
$|\langle F,\varepsilon\rangle|^2\le C\delta\|F\|^2$ --- is a
weighted Hardy-type inequality in the gap coordinate and is the
cleanest formulation of the missing upper bound.  A caveat for
whoever takes this up: the chain is $2$-periodic
($(-1)^{\#\mathrm{blocks}}$ flips deterministically, an exact
$\lambda=-1$ eigenfunction, and it is $\iota$-odd since $\iota$
changes the block count by one), so the top-of-spectrum analysis must
be run on $P^2$ or per parity sector; $\psi$ decays, so
$\langle\varepsilon,(-1)^{\#\mathrm{blocks}}\rangle_\pi=0$, which is
consistent with phase symmetry of $\pi$, but intermediate
inequalities must not silently assume aperiodicity.

%----------------------------------------------------------------------
\subsection{Transfer to the orbit measure: the mixture route}
\label{sec:transfer}

The factorization proof of Theorem~\ref{thm:factor} uses exactly one
probabilistic input --- chords attached to disjoint index sets are
conditionally independent given the shared ones --- and the session-7
caution stands: this fails per-instance.  The split--merge picture
suggests the correct weakening.  The orbit chord positions
$(Z_1,\dots,Z_k)$ are exchangeable given the frame (the cycles of
$\pi_P$, the points $u,v$); a finitely exchangeable vector drawn from
a population of $N$ position slots is a mixture of iid$(\theta)$
vectors up to total variation $O(k^2/N)$ (Diaconis--Freedman), with
$\theta$ the (random) empirical position law of the instance, and the
cell-disjointness of the intercalate family is a pairwise hard-core
correction of the same $O(k^2/N)$ order.  Two observations make the
mixture usable:
\begin{enumerate}[leftmargin=2em]
\item Conditionally on $\theta$, Theorem~\ref{thm:factor} applies
verbatim (its proof needs only iid given $(u,v,\theta)$), giving
$4\,\E[\mathrm{bias}^2]\le
\E_\theta\bigl[\max_{j,d\ge k/8}\psi^\theta(j,d)\bigr]
+2e^{-k/32}+O(k^2/N)$.
\item The entire analysis of
\S\S\ref{sec:smident}--\ref{sec:gaplaw} is \emph{measure-covariant}:
for chords with iid $\theta$-distributed endpoints, the induced chain
is the split--merge chain in which every rate is read off the
$\theta$-masses of curves and sides (cuts land in a curve with
probability its $\theta$-mass, uniformly in its
$\theta$-parametrisation).  In particular Lemma~\ref{lem:splitmerge},
Theorem~\ref{thm:iota}, Corollary~\ref{cor:quotient} and
Lemma~\ref{lem:recharge} hold verbatim with ``mass'' $=$
$\theta$-mass, provided $\theta$ is atomless; atoms of mass $m$ act
as permanent traps of that scale and cap the decay at
$\psi^\theta\gtrsim\max_i\theta(\{i\})$, which for orbit frames
($N\asymp n$ near-uniform slots) is $O(1/n)$ --- far below the
$k^{-2c}$ budget, $k=\Theta(\log^{11}n)$.
\end{enumerate}
So the transfer problem 7a(i) is reduced to (a) a quantitative
exchangeable-to-mixture step for the greedy maximal family under the
orbit measure (the hard-core correction; KSS-type concentration
should control the conflict density), and (b) the
$\theta$-uniform version of the mixing bound, i.e.\
Problem~\ref{prob:survival} with constants depending only on
$\max_i\theta(\{i\})$ and not otherwise on $\theta$
(Proposition~\ref{prop:shorthorizon} is already $\theta$-covariant:
Lemma~\ref{lem:recharge}(iii) reads verbatim in $\theta$-mass for
atomless $\theta$).  Empirically the conclusion already holds on real squares:
binning the session-7 per-instance data
(\texttt{s7\_inst\{30,50,100\}.csv}) by $k$ gives
$\E[\mathrm{bias}^2\mid k]\cdot k=0.6$--$0.75$, flat across
$k=8$--$26$ and $n=30$--$100$ (and \texttt{s8\_inst200.csv} extends
this, see \S\ref{sec:mixnum}): the $1/k$ law transfers with a
constant $\approx2.2\times$ the iid one, no drift.

\emph{Status: reduction formulated; both steps OPEN.  This replaces
the vaguer options ($\alpha$)/($\beta$) of the session-7 list.}

%----------------------------------------------------------------------
\subsection{Numerics at large $d$}\label{sec:mixnum}

All session-8 runs, continuum chain unless stated
(\texttt{s8\_psibar.log}, \texttt{s8\_4eb2\_big.log},
\texttt{s8\_onept.log}); COMPUTED.

\emph{(a) The stationary rate is $1/d$ on the nose.}  Stationary-start
(burn-in $16384$) values of $\psi_\pi(d)\cdot d$:
$0.96$ ($d{=}64$), $1.06$ ($d{=}128$), $1.11$ ($d{=}256$),
$0.93$ ($d{=}512$), $0.87$ ($d{=}1024$; errors grow from $\pm0.05$
to $\pm0.25$ along the list) --- constant $\approx1.0$, fitted
exponent $1.03\pm0.09$ over $d=64$--$1024$ (and $1.05$ on the
session-7 $j$-table route), no logarithmic growth, matching
Proposition~\ref{prop:lower} from below and pinning the conjectured
upper bound at $c=1/2$ exactly.

\emph{(b) The scaling collapse.}  Binning the $\psi$-contribution by
the initial gap: $\E[\Phi_d^2\mid g_0]$ is a function of $u=g_0d$
alone across $d=64,128,512$ --- the dyadic profile
$\varphi(u)\approx0.06,\,0.23,\,0.49,\,0.69,\,0.82,\,0.9,\,{\to}1$ at
$u\approx0.7,\,0.35,\,0.18,\,0.09,\,0.04,\,0.02,\dots$, and
$\varphi\le0.02$ for $u\ge1.4$ (consistent with fast decay); the
transition sits at $g_0\approx1/d$ as the trap picture requires, the
together and apart rows agree bin by bin (the $\iota$-symmetry made
visible), and summing bins reproduces $\psi\approx1.0/d$ with the
$\{g_0\lesssim1/d\}$ bins carrying essentially all of it.  Small-$u$
behaviour $1-\varphi(u)\approx u^{1/2\text{--}0.6}$ (not the naive
linear touch-count) --- unexplained, low stakes.

\emph{(c) The iid-ensemble constant flattens.}  Extending the
session-7 drift $4\E[\mathrm{bias}^2]\cdot k=0.95\to1.44$
($k=8..256$): $k=512$: $1.448\pm0.229$; $k=1024$: $1.171\pm0.324$
(discrete code, $20$M/$40$M trials) --- constant $\approx1.3$--$1.45$,
\emph{no} $\log k$; the budget question of session 7 (8d(b)) is
settled: no log factor.

\emph{(d) The one-point exponent is exactly $2$.}  From the clean
start (single curve, $x=y=\tfrac12$):
$f(m)=+0.006280(50),\ +0.002951(41),\ +0.001713(35),\
+0.000814(29),\ +0.000458(25)$ at $m=8,12,16,24,32$;
$f(m)\cdot m^2=0.402,0.425,0.438,0.469,0.469$ and the local exponent
reaches $2.00$ on $m=24\to32$: $f(m)\to0.47\,m^{-2}$.  This confirms
the trap-creation mechanism: a clean start must \emph{create} its
trap (Lemma~\ref{lem:recharge}(iii) prices a depth-$1/m$ trap at
$\asymp1/m^2$), whereas the stationary start carries a ready-made
linear atom --- hence $1/m^2$ against the stationary $1/d$; it also
resolves session 7's 8d(c) (the value was below MC resolution
there).  \emph{Session-9 correction:} the same data at odd $m$ shows
$f(m)$ \emph{alternates} in sign, so the leading $1/m^2$ term is the
alternating one and not the frozen-label trap term; see
Lemma~\ref{lem:parity}(a) and Theorem~\ref{thm:exactonepoint}.

%----------------------------------------------------------------------
\subsection{Session 9: the gap law is not needed; parity; an exact
one-point formula}\label{sec:s9}

\subsubsection{Short horizon: Problem~\ref{prob:gaplaw} is
unnecessary}\label{sec:shorthorizon}

The reduction of \S\ref{sec:upper} asked for a gap law uniform in $j$.
But in Theorem~\ref{thm:factor} the initial law $\mu_j$ enters only for
$j\le k$ while $d\ge k/8$, i.e.\ $j\le8d$: the chain has run for at
most $8d$ steps before the $d$ fresh steps begin, and over such a short
horizon the state-independent creation flux of
Lemma~\ref{lem:recharge}(iii) alone controls the gap.

\begin{proposition}[short-horizon gap bound]\label{prop:shorthorizon}
Let $g_0^{\mathrm{ini}}=\min(|u-v|,1-|u-v|)$ be the initial separation
($=\tfrac12$ in the antipodal ensemble) and let $g_j$ be the gap of the
marked state after $j$ iid chords.  Then for every $\varepsilon>0$ and
every $j\ge0$,
\[
\Pr_{\mu_j}\bigl[g_j\le\varepsilon\bigr]
\;\le\;\mathbf 1\{g_0^{\mathrm{ini}}\le\varepsilon\}\;+\;4\varepsilon^2 j .
\]
\end{proposition}

\begin{proof}
If $g_j\le\varepsilon<g_0^{\mathrm{ini}}$ there is a first step
$t\le j$ at which the gap enters $(0,\varepsilon]$ from above; by
Lemma~\ref{lem:recharge}(iii) each step does so with probability at
most $4\varepsilon^2$ whatever the current state (the four channels
--- shrink of either coordinate, either end of a flip --- each
contribute at most $\varepsilon^2$; merges only increase coordinates
and environment-only moves leave them unchanged).  Sum over $t\le j$.
\end{proof}

\emph{Status: PROVED.  VERIFIED (\texttt{s9\_survival flux}): at
$j=8$ the ratio $\Pr[g_j\le\varepsilon]/(4\varepsilon^2j)$ is
$1.00\pm0.05$ for $\varepsilon\le2^{-8}$ --- the bound is
\emph{sharp} at small $\varepsilon$ (each channel's entry probability
is exactly $\varepsilon^2$ while all coordinates exceed
$\varepsilon$, and nothing has had time to escape); at $j=64$ and
$j=512$ the ratio is $\le0.93$ resp.\ $\le0.92$ in every bin and
decreases to the stationary linear law $4\varepsilon$ at moderate
$\varepsilon$.}

\begin{corollary}[$\kappa$-free budget]\label{cor:kappafree}
Suppose Problem~\ref{prob:survival} holds with exponent $\alpha$:
$\E_{\mu_j}[\Phi_d^2;\,g_j\ge\varepsilon]\le C(\varepsilon d)^{-2\alpha}$
for all $j\le8d$.  Then, in the antipodal ensemble,
\[
\psi(j,d)\;\le\;32\,\varepsilon^2d+C(\varepsilon d)^{-2\alpha}
\;\le\;C'\,d^{-\alpha/(1+\alpha)}\qquad(j\le8d),
\]
i.e.\ Problem~\ref{prob:mixing} holds with
$c'=\alpha/(2(1+\alpha))$, and the budget $c'>1/11$ of
Corollary~\ref{cor:psienough} is met as soon as $\alpha>2/9$.  In
particular the expected value $\alpha=1/2$ gives $c'=1/6$ (against
$c'=1/4$ with a linear gap law), and Problem~\ref{prob:gaplaw} is
\emph{not needed} for the iid model.
\end{corollary}

\begin{proof}
Combine Proposition~\ref{prop:shorthorizon} ($j\le8d$) with
$|\Phi_d|\le1$ in the decomposition of \S\ref{sec:upper} and optimise
$\varepsilon=d^{-(1+2\alpha)/(2+2\alpha)}$; $c'>1/11$ iff
$11\alpha>2+2\alpha$.
\end{proof}

\emph{Status: PROVED.  The whole ``no-dust''/escape discussion of
\S\ref{sec:upper} is thereby demoted to a route for the \emph{sharp}
exponent $c=1/2$ (or for the stationary-start statement); the iid
model needs only Problem~\ref{prob:survival} with any $\alpha>2/9$.
In the true ensemble the indicator term is the short-arc atom already
accounted for as the $O(1/k)$ term of \S\ref{sec:annealed}.}

\begin{lemma}[$s$-trap flux]\label{lem:strap}
From any state, the marked total $s=a+b$ enters $(0,\sigma]$ from above
with probability at most $2\sigma^2$ per step; hence
$\Pr_{\mu_j}[s_j\le\sigma]\le2\sigma^2j$ in the antipodal ensemble.
\end{lemma}
\begin{proof}
Flips and merges do not decrease $s$; a shrink of $a$ lands at
$s'=a'+b\le\sigma$ only if $b\le\sigma$ and $a'\le\sigma-b$, which has
probability $\le(\sigma-b)^2\le\sigma^2$ (the shrink channel of
Lemma~\ref{lem:recharge}(iii)); likewise for $b$.
\end{proof}
\emph{Status: PROVED.  Consequently in Corollary~\ref{cor:kappafree}
one may also assume $s_0\ge\sigma$ at cost $16\sigma^2d$, so the
target estimate of Problem~\ref{prob:survival} need only be proved
for starts with $a_0,b_0\ge\varepsilon$ and $s_0\ge\sigma\asymp\varepsilon$.}

\subsubsection{Parity: the label is a function of time and the
environment block count}\label{sec:parity}

\begin{lemma}[block-count parity]\label{lem:parity}
Let $B_t$ be the total number of curves and $E_t$ the number of
environment curves (those through neither $u$ nor $v$) at time $t$.
Every step is a split or a merge, so $B_t\equiv B_0+t\pmod 2$; since
$B_t=E_t+1+\mathbf 1\{u\not\sim v\}$,
\[
\varepsilon_t\;=\;-(-1)^{B_0+t}\,(-1)^{E_t},
\qquad
(-1)^{N_t}\;=\;(-1)^{t}\,(-1)^{E_t-E_0},
\]
with $N_t$ the number of flip moves.  Hence
$\Phi_d(\Gamma)=(-1)^d\,\E_\Gamma[(-1)^{E_d-E_0}]$.
\end{lemma}

\begin{proof}
Splits add a curve, merges remove one; a flip (separate or join)
changes $B$ by one and leaves $E$ unchanged, every other move leaves
the label and changes $E$ by one.
\end{proof}

\emph{Status: PROVED (trivial, but structurally decisive).  Three
consequences.}

\emph{(a) The sign of the one-point function alternates.}  From the
clean start, $f(m)=\E[\varepsilon(F_m)]$ obeys
$f(m)\cdot(-1)^m>0$: measured $f(7)=-0.00753(22)$, $f(8)=+0.00599(22)$,
$f(9)=-0.00526(22)$, $f(11)=-0.00353$, $f(12)=+0.00295$,
$f(13)=-0.00275$, $f(15)=-0.00192$, $f(16)=+0.00213$,
$f(17)=-0.00149$ ($2\times10^7$ trials each).  Session 8 sampled only
even $m$ and read $f(m)\approx0.47/m^2$ as the frozen-label trap
term; that reading is \emph{incomplete}: the leading term is the
alternating one, $f(m)\approx(-1)^m\,c_1/m^2$ with $m^2|f(m)|$
increasing through $0.40,0.43,0.44,0.47,0.47$ ($m=8..32$) and
plausibly $c_1=\tfrac12$ (Theorem~\ref{thm:exactonepoint} below gives
exactly $\tfrac12$ for the $(u,v)$-averaged function), while the
non-alternating (trap) part, $\tfrac12(f(m{-}1)+f(m{+}1))+f(m)$ at
even $m$, is $\lesssim2\times10^{-4}$ at $m=8$.  Nothing downstream
is affected ($\psi=\E\Phi^2$ is sign-blind), but the smooth object is
$(-1)^d\Phi_d$, not $\Phi_d$, and Remark~\ref{rem:renewal}'s heuristic
``$\Phi_d\approx\Pr[\tau>d]$'' must be read with this sign.

\emph{(b) An obstruction for coupling proofs.}  Any two words of the
same length whose final environments have the same block-count parity
have the same flip parity; in particular two trajectories from the
same start that agree in marked geometry \emph{and} environment at
time $d$ have the same label.  A coupling proof of
Problem~\ref{prob:survival} must therefore produce, at time $d$,
environments of \emph{different} block-count parity with the same
law --- e.g.\ one carrying an extra tiny block.  The natural ``swap a
flip for a shrink--re-merge pair and let the environment take a tiny
extra split'' coupling has success probability $O(1/d)$ per step
and is repeated $d$ times, so it yields only $|\Phi_d|\le1-c$: no
decay.  Same-time $\iota$-mirroring is impossible ($\iota$ changes
$B$ by one); a one-step shifted mirror gives $\Phi_d\approx-\Phi_{d+1}$,
which is the alternation (a), not decay.  The parity of $N_d$ is thus
a genuinely ``continuous-state'' functional and the proof of
Problem~\ref{prob:survival} has to go through the renewal or spectral
route.

\emph{(c) First-flip tail from good starts is $1/t^2$, not $1/t$.}
From a stationary environment conditioned on $g_0\in[1/16,1/8)$
(\texttt{s9\_survival phic}), $t^2\Pr[\tau>t]\to11.1$ ($t=128,256$;
$16.1$ at $t=64$, $29$ at $t=32$); from the clean start with
$g_0=1/4$, $t\Pr[\tau>t]$ halves per doubling of $t$ from $t=32$ on.
The renewal heuristic explains this: a trap at depth $\eta\le\varepsilon$
is created at rate $8\eta\,d\eta$ (Lemma~\ref{lem:recharge}(iii)) and
survives $m$ further steps without flipping with probability
$\approx e^{-2\eta sm}$, so the kernel $K(m)=\int_0^\varepsilon
8\eta e^{-2\eta sm}d\eta\approx2/(s^2m^2)$ for $m\gg1/(\varepsilon s)$,
with total mass $\approx4\varepsilon/s<1$: the first-flip time is a
\emph{subcritical} renewal sum with a $1/t^2$ kernel, hence
$\Pr[\tau>t]\asymp t^{-2}$ and, if the alternating-renewal
cancellation holds, $|\Phi_d|\lesssim(\varepsilon d)^{-2}$ off small
gaps --- i.e.\ the true $\alpha$ in Problem~\ref{prob:survival} is
plausibly $2$, far above the $2/9$ needed.  Consistently,
$\E[\Phi_d^2\mid g_0=u/d]$ is at the noise floor ($\lesssim10^{-3}$,
often negative estimates) for every bin $u\ge3$ in the session-8
\texttt{psibar} logs at $d=128$--$1024$.  COMPUTED, not proved.

\subsubsection{An exact one-point formula}\label{sec:exactonept}

Averaging the marked points over the circle turns the one-point
function into a class function of the discrete precursor, and hook
characters compute it exactly.

\begin{theorem}[exact $(u,v)$-averaged one-point function]
\label{thm:exactonepoint}
Let $\sigma_m=\gamma\,\tau_1\cdots\tau_m$ with $\gamma$ a uniform
$n$-cycle and $\tau_i$ iid uniform transpositions, and let $c_i$ be
its cycle lengths.  Then
\[
\E\Bigl[\sum_i c_i(\sigma_m)^2\Bigr]
=\frac{n(n+1)}2+\sum_{k=1}^{n-1}(n-k)\Bigl(1-\frac{2k}{n-1}\Bigr)^{m},
\]
so that for uniform distinct $u,v$,
$\Pr[u\sim v\text{ in }\sigma_m]=\bigl(\E\sum c_i^2-n\bigr)/(n(n-1))$.
Consequently, for the continuum uniform split--merge chain
$(p_i(m))$ started from a single block,
\[
\E\Bigl[\sum_i p_i(m)^2\Bigr]
=\frac12+\frac14\Bigl[\frac{1+(-1)^m}{m+1}+\frac{1-(-1)^m}{m+2}\Bigr]
=\begin{cases}\tfrac12+\tfrac1{2(m+1)},& m\text{ even},\\[2pt]
\tfrac12+\tfrac1{2(m+2)},& m\text{ odd},\end{cases}
\]
and the $(u,v)$-averaged one-point function
$\bar f(m)=\E_{u,v}\E[\varepsilon(F_m)]=2\E\sum p_i^2-1$ equals
$1/(m+1)$ for even and $1/(m+2)$ for odd $m$:
$\bar f(m)=\dfrac{2m+3}{2(m+1)(m+2)}+(-1)^m\dfrac{1}{2(m+1)(m+2)}$.
\end{theorem}

\begin{proof}
Right multiplication by a uniform transposition acts on the
$\lambda$-isotypic component of $L^2(S_n)$ by
$r_\lambda=\chi_\lambda(\tau)/d_\lambda$, and the law of $\gamma$ is
the class measure of $n$-cycles, so for any $F$,
$\E F(\sigma_m)=\sum_\lambda \rho_\lambda\,r_\lambda^{m}\,
\langle\chi_\lambda,F\rangle$ with
$\rho_\lambda=\chi_\lambda(n\text{-cycle})$ and
$\langle\chi,F\rangle=\frac1{n!}\sum_\sigma\chi(\sigma)F(\sigma)$
(Fourier inversion for class measures).  Only hooks
$\lambda_k=(n-k,1^k)$ have $\rho_{\lambda}\ne0$, with
$\rho_{\lambda_k}=(-1)^k$, and their content sum gives
$r_{\lambda_k}=\bigl[(n-k)(n-k-1)-k(k+1)\bigr]/(n(n-1))=1-2k/(n-1)$.
For $F=\sum_ic_i^2$ use the hook generating function
$\sum_k\chi_{\lambda_k}(\mu)t^k=(1+t)^{-1}\prod_i(1-(-t)^{\mu_i})$ and
the exponential formula with one marked part:
\[
\sum_{n}x^n\sum_k t^k\langle\chi_{\lambda_k},F\rangle
=\frac1{1+t}\Bigl[\sum_{j\ge1}j\,(1-(-t)^j)x^j\Bigr]
\exp\Bigl(\sum_{j\ge1}\frac{(1-(-t)^j)x^j}{j}\Bigr)
=\frac1{1+t}\Bigl[\frac{x}{(1-x)^2}+\frac{tx}{(1+tx)^2}\Bigr]
\frac{1+tx}{1-x}.
\]
Extracting $[x^n]$: the first term gives
$\bigl[\binom{n+1}2+t\binom n2\bigr]/(1+t)$, the second
$t(1-(-t)^n)/(1+t)^2$; the coefficient of $t^k$ ($1\le k\le n-1$) is
$(-1)^kn+(-1)^{k-1}k=(-1)^k(n-k)$, and $\binom{n+1}2$ at $k=0$.
Multiply by $\rho_{\lambda_k}=(-1)^k$ and $r^m$.  The continuum limit
is the Riemann sum $\int_0^1(1-x)(1-2x)^m dx$, evaluated by
$y=1-2x$; the identification of $\sigma_m/n$ with the split--merge
chain from one block is Lemma~\ref{lem:splitmerge} (the $n$ points
carry arcs of length $1/n$).
\end{proof}

\emph{Status: PROVED.  VERIFIED: the finite-$n$ formula
(\texttt{s9\_hooks.py}, exact rationals; the inner products
$(-1)^k(n-k)$ were also checked against the full character sums for
$n\le12$) gives at $n=40$: $0.6007,\,0.5567,\,0.5314$ at $m=4,8,16$
against the continuum simulator's $0.6005,\,0.5561,\,0.5297$
(\texttt{s8\_splitmerge traj}), and the continuum values
$m=1:\ 2/3$, $m=2:\ 2/3$ agree with direct hand integration.
Remarks: (i) this is the $(u,v)$-averaged version of the one-point
function of \S\ref{sec:mixnum}(d); its non-alternating part
$\approx1/m$ is the frozen-label trap term fed by the linear initial
gap density (2 on $(0,\tfrac12)$) of a random separation, absent for
antipodal $u,v$, and its alternating part $(-1)^m/(2(m+1)(m+2))$ has
exactly the coefficient $\tfrac12$ towards which the antipodal
$m^2|f(m)|$ appears to converge; (ii) in spectral terms the averaged
start has spectral density $(1+y)/4$ on $y=1-2x\in[-1,1]$ --- linear
at \emph{both} ends $y=\pm1$, the two ends being the trap ($y=+1$)
and the alternating ($y=-1$, the $(-1)^{B}$ eigenfunction) edges;
(iii) the same technique cannot reach $\psi$ (a two-time quantity is
not a class function of the increment), but it does prove the
averaged one-point decay $\bar f(m)\le1/(m+1)$ outright, and gives
the exact law of $\E\sum p_i^2$ along the whole approach to
$\mathrm{PD}(1)$.}

\subsubsection{The spectral route on the quotient: twisted kernel}
\label{sec:twisted}

Lemma~\ref{lem:parity} and Corollary~\ref{cor:quotient} make the
spectral formulation of \S\ref{sec:spectral} concrete on the quotient
chain.  Write $P=P_{\mathrm{nf}}+P_{\mathrm f}$ for the non-flip and
flip parts of the quotient kernel and let
$P_\pi:=P_{\mathrm{nf}}-P_{\mathrm f}$ (the kernel twisted by
$(-1)^{F}$).  Then $\Phi_d=P_\pi^{\,d}\mathbf 1$, and since flips are
reversed by flips, $P_\pi$ is self-adjoint on $L^2(\tilde\pi)$
($\tilde\pi$ the quotient stationary law), so
\[
\psi_\pi(d)=\langle\mathbf 1,P_\pi^{2d}\mathbf 1\rangle
=\int_{[-1,1]}\lambda^{2d}\,\nu(d\lambda),\qquad
\nu=\text{spectral measure of }\mathbf 1\text{ under }P_\pi,\ \nu([-1,1])=1,
\]
and Problem~\ref{prob:mixing} at stationarity reads
$\nu(\{|\lambda|\ge1-\delta\})\le C\delta^{\beta}$
($\psi_\pi(d)\le Cd^{-\beta}$).  Two facts, one numerical and one a
proof sketch:

\emph{(a) No mass at the $-1$ end at stationarity (COMPUTED).}
$(-1)^{E}$ is an exact $\lambda=-1$ eigenfunction of $P_\pi$
(Lemma~\ref{lem:parity}), so the $2$-periodicity caveat of
\S\ref{sec:spectral} is a priori relevant.  But
$\langle P^d\varepsilon,P^{d+1}\varepsilon\rangle_\pi
=\int\lambda^{2d+1}d\nu$ versus $\int\lambda^{2d}d\nu$
(\texttt{s9\_survival oddeven}): $0.0559/0.0575=0.971$ at $d=16$ and
$0.0146/0.0150=0.974$ at $d=64$ (both $\pm0.03$), so the relative
mass of $\nu$ near $-1$ is $\lesssim3\%$ and compatible with zero; in
the continuum $E=\infty$ under $\tilde\pi$ and the eigenfunction
$(-1)^E$ is not in $L^2(\tilde\pi)$.  The sign alternation of
\S\ref{sec:parity}(a) is a finite-block (clean-start) phenomenon,
invisible at stationarity.

\emph{(b) A Cauchy--Schwarz/smoothing route to $\beta=1/2$
(SKETCH).}  The Dirichlet form of $P_\pi$ is
$\mathcal E(F)=\tfrac12\E[(F-F')^2;\mathrm{nf}]+\tfrac12\E[(F+F')^2;\mathrm f]$.
Let $k(x)=2ab$ be the flip rate and $Q(x,\cdot)=P_{\mathrm f}(x,\cdot)/k(x)$
the normalised flip kernel, which acts on each segment
$\{(a',s-a';\mathrm{env})\}$ separately by the trapezoidal density
$a'\mapsto|\{\alpha\in(0,a):a'-\alpha\in(0,s-a)\}|/(a(s-a))$.  By
Cauchy--Schwarz in the flip term,
$\mathcal E(F)\ge\tfrac12\int\tilde\pi(dx)\,k(x)\,(F+QF)^2(x)$, so
$\mathcal E(F)\le\delta\|F\|^2$ forces $F\approx-QF$ in
$L^2(\tilde\pi k)$, hence $F\approx Q^2F$: $F$ is nearly constant,
then nearly equal to minus itself, on every segment away from its
ends --- i.e.\ $\int_{\{ab\ge\eta\}}F^2d\tilde\pi\le C\delta\|F\|^2/\eta$
once the one-dimensional smoothing inequality for $Q^2$ on a segment
is quantified (the trapezoid density is $\ge c/s$ away from the
ends, $\asymp a'/s^2$ near them).  With $\tilde\pi(ab<\eta)\le C\eta$
(Proposition~\ref{prop:gaplawstat} plus the quadratic $s$-law),
$\langle F,\mathbf 1\rangle^2\le C\|F\|^2(\delta/\eta+\eta)$, and
$\eta=\sqrt\delta$ gives $\nu([1-\delta,1])\le C\sqrt\delta$, i.e.\
$\psi_\pi(d)\le Cd^{-1/2}$: $c=1/4>1/11$ \emph{at stationarity}.
Gaps: (i) the segment smoothing inequality (a weighted Poincar\'e
for $Q^2$ with weight $a'(s-a')$); (ii) the $-1$ end, to be handled
either by (a) made rigorous or by running the argument for
$\langle F,(I+P_\pi)F\rangle$; (iii) the passage from $\pi$ to
$\mu_j$, $j\le8d$ --- the short-horizon Proposition~\ref{prop:shorthorizon}
handles the gap coordinate but not the rest of the state, and
Lemma~\ref{lem:freeze} only trades $d$ for $j$.  (iii) is the real
obstacle for the iid model; a dust-insensitive comparison of
$\mu_{j}$ with $\pi$ for $j\ge d/2$ (a block of mass $w$ influences
$d$ steps with probability $\le2wd$) is the natural tool.

%----------------------------------------------------------------------
\subsection{Session 10: the block count is Harer--Zagier, the environment
is healthy on the short horizon, the exact size-biased law, and the
Markov-renewal form of the survival problem}\label{sec:s10}

Session 10 was to fix the exact statement behind
Problem~\ref{prob:survival} (TODO~10a, step~0) and choose a route.  The
first finding is negative and structural: \emph{the environment cannot
be avoided} --- without merges the parity of the flip count from a start
with $\min(a_0,b_0)=\varepsilon$ decays on the time scale
$\varepsilon^{-2}$, not $\varepsilon^{-1}$, and the budget of
Corollary~\ref{cor:kappafree} forces $\varepsilon\ll d^{-1/2}$, so a
proof must use the recharge of small marked coordinates by
environment blocks (\S\ref{sec:noenv}).  The second finding removes the
obstacle this creates.  The number of curves $B_t$ after $t$ chords has
an \emph{exact} law --- it is the Harer--Zagier vertex-count
distribution of a random gluing of a $2t$-gon --- with
$\E[3^{B_t}]=3+4t+2t^2$; a Chernoff bound then gives $B_t\le C\log d$
simultaneously for all $t\le9d$ except with probability $d^{-\gamma}$
(\S\ref{sec:blockcount}).  On that event every state on the horizon
has $\sum p_i^2\ge1/(C\log d)$ and an environment block of mass
$\ge(1-s)/(C\log d)$: the ``no-dust'' question is closed, up to
logarithms, in exactly the form the short horizon needs.  The same hook
computation gives the exact law of the mass of the curve through a
uniform point at every time (\S\ref{sec:sizebiased}; this contains
Theorem~\ref{thm:exactonepoint} and closes TODO~10b(b)).  Finally the
survival problem is put in Markov-renewal form
(\S\ref{sec:markovrenewal}): an exact operator identity for the
generating function $\sum_d\Phi_dz^d$ separates the two remaining
contents --- moments of the first-flip time (traps) and invertibility
of $I+K_z$ (cancellation) --- and a conditional decay proposition
states precisely what is left.

\subsubsection{The environment cannot be avoided}\label{sec:noenv}

Consider the marked pair alone, with merges suppressed (an adversarial
``dust'' environment): per step, $a\to a\sqrt U$ with probability
$a^2$, $b\to b\sqrt U$ with probability $b^2$, a flip
$(a,b)\to(\alpha+\beta,s-\alpha-\beta)$ with probability $2ab$, and
nothing otherwise.  Since merges only increase marked coordinates and
are the only moves that do, this is the chain that any argument using
the environment solely through upper bounds would have to control.
Its scaling is elementary: all rates are quadratic in the coordinates,
so from $(a_0,b_0)=(\varepsilon,\varepsilon)$ nothing happens for
$\asymp\varepsilon^{-2}$ steps, and from $(c_0,\varepsilon)$ the large
coordinate decays like $t^{-1/2}$ before a flip (rate $2\varepsilon b$)
can occur, the first flip arriving only at $t\asymp\varepsilon^{-2}$.
Measured (\texttt{s10\_nomerge}, $\Phi_t=\E[(-1)^{N_t}]$):
from $(\tfrac14,\tfrac14)$, $\Phi_t=0.130,\,0.031,\,0.0070,\,0.0024$
at $t=8,16,32,64$ (then noise floor $10^{-3}$; $t^2\Phi_t\approx8$--$10$
on $16\le t\le64$) --- a healthy $1/t^2$-like decay without any
environment; but from $(\tfrac14,0.02)$,
$\Phi_t=0.53,\,0.25,\,0.088,\,0.040$ at $t=64,256,1024,4096$, and from
$(0.02,0.02)$, $\Phi_t=0.90,\,0.68,\,0.30,\,0.13$ at
$t=64,256,1024,8192$: decay only on the scale $\varepsilon^2t$, and
slow even there.  In the budget of Corollary~\ref{cor:kappafree} the
trap-entry flux forces $\varepsilon\ll d^{-1/2}$, hence
$\varepsilon^2d\ll1$: a bound depending on the environment only
through upper bounds is worthless.  \emph{Status: COMPUTED (the scaling
statement is trivial).  Consequence: any proof of
Problem~\ref{prob:survival} must lower-bound the masses of environment
blocks over the horizon --- the recharge mechanism of
Remark~\ref{rem:renewal}.  The next subsection supplies exactly this.}

\subsubsection{The exact law of the number of curves}\label{sec:blockcount}

Let $B_t$ be the number of curves after $t$ iid chords (the total
block count of the marked chain from the one-block start; it does not
depend on the marks).

\begin{theorem}[the block count is Harer--Zagier]\label{thm:blockcount}
For every $x$,
\[
\sum_{t\ge0}\E\bigl[x^{B_t}\bigr]z^t
=\frac1{2z}\Bigl[\Bigl(\frac{1+z}{1-z}\Bigr)^{x}-1\Bigr],
\qquad\text{i.e.}\qquad
\Pr[B_t=m]=[z^{t+1}]\,\frac{(2\operatorname{artanh}z)^m}{2\,m!},
\]
so $B_t\equiv t+1\pmod2$, $\Pr[B_t=t+1-2g]=\varepsilon_g(t)/(2t-1)!!$
with $\varepsilon_g(t)$ the Harer--Zagier numbers \cite{HZ86}, and for
integer $x\ge1$ the expectation $\E[x^{B_t}]$ is a polynomial in $t$ of
degree $x-1$:
\[
\E[2^{B_t}]=2+2t,\qquad \E[3^{B_t}]=3+4t+2t^2,\qquad
\E[4^{B_t}]=4+\tfrac{20}3t+4t^2+\tfrac43t^3=\tfrac43(t+1)(t^2+2t+3),\qquad
\Pr[B_t=1]=\frac{1+(-1)^t}{2(t+1)}.
\]
\end{theorem}

\begin{proof}
The $2t$ endpoints are iid uniform on the circle, a.s.\ distinct; by
exchangeability, conditionally on their positions the pairing into
chords is uniform among the $(2t-1)!!$ pairings.  Refining $\gamma$ at
the endpoints does not change the curve partition
(Lemma~\ref{lem:splitmerge}), so $B_t$ is the number of cycles of
$\gamma_{2t}\tau$, with $\gamma_{2t}$ the rotation of the $2t$
endpoints and $\tau$ the pairing involution --- the number of vertices
of the surface obtained by gluing the sides of a $2t$-gon along a
uniform pairing.  Harer--Zagier's formula \cite{HZ86} is
$1+2\sum_{t\ge0}\bigl[\sum_g\varepsilon_g(t)N^{t+1-2g}\bigr]
z^{t+1}/(2t-1)!!=\bigl((1+z)/(1-z)\bigr)^N$; dividing by $(2t-1)!!$ turns
the bracket into $\E[N^{B_t}]$, which is the displayed identity for
$x=N$, and both sides are polynomials in $N$ for each $t$.  Expanding
$((1+z)/(1-z))^x=\exp(2x\operatorname{artanh}z)$ in powers of $x$ gives
the law; the polynomials are the coefficients of
$(1-z)^{-x}\sum_k\binom{x}{2k+1}z^{2k}$.
\end{proof}

\emph{Status: PROVED (the statement is classical once the chord
diagram is identified; the identification is the content).  A second
proof runs through the finite-$n$ hook machinery of
Theorem~\ref{thm:exactonepoint}: $x^{c(\sigma)}=\sum_\lambda
s_\lambda(1^x)\chi_\lambda(\sigma)$, so
$\E[x^{c(\sigma_t)}]=\sum_{k}(-1)^k r_k^{\,t}\,s_{(n-k,1^k)}(1^x)$ with
$s_{(n-k,1^k)}(1^x)=x(x{+}1)\cdots(x{+}n{-}k{-}1)\cdot(x{-}1)\cdots(x{-}k)
/\bigl(n\,(n{-}k{-}1)!\,k!\bigr)$ by hook-content; for integer $x$ only
$k\le x-1$ survive, e.g.\ $\E[2^{c(\sigma_t)}]=(n+1)-(n-1)r_1^{\,t}$,
$\E[3^{c(\sigma_t)}]=\tfrac12(n+1)(n+2)-(n^2-1)r_1^{\,t}+\tfrac12(n-1)(n-2)r_2^{\,t}$,
whose $n\to\infty$ limits are $2+2t$ and $3+4t+2t^2$
(\texttt{s10\_blocks.py}, exact rationals, Lagrange-extrapolated in
$1/(n-1)$, gives the polynomials for $x\le5$).  VERIFIED: the full law
against a direct simulation of the split--merge chain
(\texttt{s10\_blocksim.py}, $2\times10^5$ trials, $t\le6$): every
probability agrees to $\pm0.001$, e.g.\ $\Pr[B_6=1,3,5,7]=
0.1429,0.6222,0.2222,0.0127$ predicted vs.\ $0.1431,0.6230,0.2212,0.0127$
measured.}

\begin{corollary}[block count over the horizon]\label{cor:horizon}
For all $T\ge1$ and $K\ge1$,
\[
\Pr\Bigl[\max_{0\le t\le T}B_t\ge K\Bigr]\;\le\;3^{\,1-K}(T+1)^3 .
\]
In particular, for $\gamma>0$ and
$K_\gamma(T):=\lceil(3+\gamma)\log_3(T+1)\rceil+1$, the event
$H_T=\{B_t<K_\gamma(T)\ \forall t\le T\}$ has probability
$\ge1-(T+1)^{-\gamma}$, and on $H_T$, for every $t\le T$: the total number of
blocks is $<K_\gamma(T)$, $\sum_ip_i(t)^2\ge1/K_\gamma(T)$, the
environment has fewer than $K_\gamma(T)$ blocks and hence a block of
mass $\ge(1-s_t)/K_\gamma(T)$.
\end{corollary}

\begin{proof}
Markov with $x=3$: $\Pr[B_t\ge K]\le3^{-K}(3+4t+2t^2)\le3^{1-K}(t+1)^2$;
sum over $t\le T$.  With $K=K_\gamma(T)$,
$3^{1-K}\le(T+1)^{-3-\gamma}$.  The consequences are Cauchy--Schwarz
($\sum p_i^2\ge1/\#\text{blocks}$) and pigeonhole.
\end{proof}

\emph{Status: PROVED.  In Problem~\ref{prob:survival} the start
$\Gamma_0\sim\mu_j$ ($j\le8d$) followed by $d$ steps is the chain from
one block over at most $9d$ steps, so $H_{9d}$ has probability
$\ge1-(9d+1)^{-\gamma}$ for any fixed $\gamma$ at the price
$K=O_\gamma(\log d)$.  This is the ``no-dust'' estimate that
\S\ref{sec:upper} and the session-8/9 handoffs listed as the open
ingredient of every route: over the short horizon it holds for free,
up to a logarithm, because the chain simply has not had time to create
many blocks --- $\E[B_t]=[z^t]\,(1+z)\operatorname{artanh}(z)/(z(1-z))=\log t+O(1)$
and the pgf makes the tail geometric.  (At stationarity the count is infinite and the analogous
statement is false as stated; the short horizon of
Theorem~\ref{thm:factor} is again what makes the iid model tractable.)
Since the required exponent in Corollary~\ref{cor:kappafree} is only
$\alpha>2/9$, losing logarithms is harmless.}

\subsubsection{The exact law of the block of a uniform point}
\label{sec:sizebiased}

\begin{theorem}[size-biased block mass at all times]\label{thm:sizebiased}
Let $U$ be uniform on the circle, independent of the chords, and let
$\zeta_m$ be the mass of the curve through $U$ after $m$ chords (chain
from one block).  Then
\[
\Pr[\zeta_m\in dx]=\bigl(1-(1-2x)^m\bigr)\,dx\ \text{ on }(0,1),
\qquad
\Pr[\zeta_m=1]=\frac{1+(-1)^m}{2(m+1)} ,
\]
and consequently, for every $r\ge1$,
\[
\E\Bigl[\sum_ip_i(m)^r\Bigr]=\E[\zeta_m^{\,r-1}]
=\frac1r+\int_0^1(1-x^{r-1})(1-2x)^m\,dx .
\]
By rotation invariance the same law governs the mass of the curve
through any fixed point, in particular through $u$ (and through $v$).
At finite $n$: $\Pr[|C_{\sigma_m}(1)|=j]=(1-r_j^{\,m})/n$ for $j<n$ and
$\Pr[|C_{\sigma_m}(1)|=n]=\frac1n\sum_{k=0}^{n-1}r_k^{\,m}$, with
$r_k=1-2k/(n-1)$; equivalently
$(-1)^k\langle\chi_{(n-k,1^k)},\sum_ic_i^{\,r}\rangle=n^{r-1}-k^{r-1}$
for $1\le k\le n-1$ and $\sum_{j\le n}j^{r-1}$ for $k=0$.
\end{theorem}

\begin{proof}
For $F$ on $\{1,\dots,n\}$ put $G_F(\sigma)=\sum_{u}F(|C_\sigma(u)|)
=\sum_ic_iF(c_i)$.  The exponential formula with one marked cycle, as
in the proof of Theorem~\ref{thm:exactonepoint} (weight $jF(j)$ for a
marked cycle of length $j$ replaces $j\cdot j$ there), gives
\[
\sum_nx^n\sum_kt^k\langle\chi_{(n-k,1^k)},G_F\rangle
=\frac1{1+t}\Bigl[\sum_{j\ge1}F(j)\bigl(1-(-t)^j\bigr)x^j\Bigr]\frac{1+tx}{1-x}.
\]
Extracting $[x^n]$: $\sum_{j<n}F(j)(1-(-t)^j)(1+t)+F(n)(1-(-t)^n)$;
dividing by $1+t$ and extracting $[t^k]$ gives, for $1\le k\le n-1$,
$-(-1)^kF(k)+(-1)^kF(n)$, and $\sum_{j\le n}F(j)$ for $k=0$.  Hence,
by the Fourier inversion of Theorem~\ref{thm:exactonepoint},
$\E[G_F(\sigma_m)]=\sum_{j\le n}F(j)+\sum_{k=1}^{n-1}r_k^{\,m}(F(n)-F(k))
=\sum_{j<n}F(j)(1-r_j^{\,m})+F(n)\sum_{k=0}^{n-1}r_k^{\,m}$.  Since
$G_F(\sigma)/n=\E_U[F(|C_\sigma(U)|)]$ for a uniform point $U$, this is
the finite-$n$ law; $F(j)=j^{r-1}$ gives the hook inner products, and
$j=nx$, $r_j^{\,m}\to(1-2x)^m$ gives the continuum statements
(Lemma~\ref{lem:splitmerge} for the identification).
\end{proof}

\emph{Status: PROVED.  VERIFIED: the hook identity
$(-1)^k\langle\chi_k,\sum c^r\rangle=n^{r-1}-k^{r-1}$ by brute-force
character sums for $r\le6$, $n\le9$ (\texttt{s10\_hooks34.py}); the
continuum law and the moments $r=2,3,4$ against simulation
(\texttt{s10\_sizebiased.py}, $2\times10^5$ trials): at $m=8$ the ten
decile masses agree to $\pm0.001$, the atom $0.1115$ vs $1/9$, and
$\E\sum p^{2,3,4}=0.5561,0.3947,0.3140$ vs $0.5556,0.3939,0.3131$.
Remarks.  (i) $r=2$ is Theorem~\ref{thm:exactonepoint}; the atom is
$\Pr[B_m=1]$ of Theorem~\ref{thm:blockcount}, as it must be.  (ii) The
drift identity $\E[\Delta L]=2L^2-\tfrac73\sum p_i^4$ ($L=\sum p_i^2$)
is now closed in expectation along the whole approach to
$\mathrm{PD}(1)$: $\E\sum p_i^4(m)=\tfrac14+\int_0^1(1-x^3)(1-2x)^mdx$.
(iii) Small marked blocks: $\Pr[\text{mass of }u\text{'s curve}\le\eta]
=\int_0^\eta(1-(1-2x)^m)dx\le\min(\eta,\,m\eta^2)$ --- the
$m\eta^2$ regime is the creation flux of Lemma~\ref{lem:recharge}(iii)
made exact for the block through a fixed point, and the crossover at
$\eta\asymp1/m$ is the trap scale.  (iv) The density $1-(1-2x)^m$ is
$\ge1-e^{-2mx}$ on $(0,\tfrac12]$: the size-biased pick is macroscopic with probability
bounded below uniformly in $m\ge1$ --- the quantitative handle for
recharge that TODO~10a(A)(iii) asked for, now superseded by
Corollary~\ref{cor:horizon}.}

\subsubsection{Recharge under the block bound}\label{sec:recharge}

Fix $K\ge3$ and let $H=H_T$ be the event of Corollary~\ref{cor:horizon}
with $K_\gamma(T)$ replaced by $K$.

\begin{lemma}[recharge of the marked total]\label{lem:srecharge}
Let $K\ge3$, $s_0\ge\varepsilon$ with $\varepsilon\le1/(2K)$.  Then for every
$T\ge1$,
\[
\Pr\Bigl[\max_{t\le T}s_t<\tfrac1{2K}\Bigr]
\;\le\;e^{-\varepsilon T/K}+2\varepsilon^2T+\Pr[H^c].
\]
The same holds for each marked coordinate separately: if
$a_0\ge\varepsilon$ then
$\Pr[\max_{t\le T}a_t<1/(2K)]\le e^{-\varepsilon T/(3K)}+2\varepsilon^2T+\Pr[H^c]$,
and likewise for $b$.
\end{lemma}

\begin{proof}
On $H\cap\{s_t<1/(2K)\}$ the environment has mass $>1-1/(2K)\ge\tfrac12$
in fewer than $K$ blocks, so it has a block of mass $z\ge1/(2K)$; the
chord with one cut in the marked material and one in that block has
probability $2s_tz\ge\varepsilon/K$ as long as $s_t\ge\varepsilon$, and
it raises $s$ by $z\ge1/(2K)$.  Flips and merges never decrease $s$, and
$s$ enters $(0,\varepsilon]$ from above with probability $\le2\varepsilon^2$
per step (Lemma~\ref{lem:strap}); so the event $\{\max_{t\le T}s_t<1/(2K)\}$
is contained in $\{s_t<\varepsilon\text{ for some }t\le T\}\cup H^c\cup
\{T\text{ consecutive failures of an event of probability }\ge\varepsilon/K\}$.
For a single coordinate $a$: if $s_t\le\tfrac12$ the same merge
argument applies to $a$ (probability $\ge2a\cdot1/(2K)\ge\varepsilon/K$);
if $s_t>\tfrac12$ and $a_t<1/(2K)$ then $b_t\ge\tfrac13$, a flip has
probability $2a_tb_t\ge2\varepsilon/3$ and lands at
$a'=\alpha+\beta\ge\beta\sim U(0,b_t)$, so $a'\ge1/(2K)$ with conditional
probability $\ge1-3/(2K)\ge\tfrac12$: success probability
$\ge\varepsilon/3\ge\varepsilon/(3K)$ per step in both regimes; the
entry flux into $(0,\varepsilon]$ is $\le2\varepsilon^2$ per step
(shrink of $a$, flip landing small; Lemma~\ref{lem:recharge}(iii)).
\end{proof}

\emph{Status: PROVED.  With $T=(3K/\varepsilon)\log(1/\delta)$ the
failure probability is $\delta+6\varepsilon K\log(1/\delta)+\Pr[H^c]$,
small precisely when $\varepsilon\ll1/K$ --- the regime in which
recharge is needed at all (for $\varepsilon\ge1/(2K)$ the start is
already log-macroscopic).  At the optimising
$\varepsilon=d^{-(1+2\alpha)/(2+2\alpha)}\in[d^{-3/4},d^{-1/2})$ of
Corollary~\ref{cor:kappafree} ($\alpha<1$) and $K=O(\log d)$,
$T=O(\varepsilon^{-1}\log^2d)\ll d$ and the failure probability is
$O(d^{-\gamma})+O(\varepsilon\log^2d)$, which is
$o\bigl((\varepsilon d)^{-2\alpha}\bigr)=o(d^{-\alpha/(1+\alpha)})$ since
$\varepsilon^{1+2\alpha}d^{2\alpha}=d^{-1/(2+2\alpha)}$.}

\begin{lemma}[log-drift of the marked total]\label{lem:logdrift}
On $H$, for a state with $s\le\tfrac12$,
\[
\E[\log s_{t+1}-\log s_t\mid\Gamma_t]\;\ge\;\frac{s}{K}\log\Bigl(1+\frac1{2Ks}\Bigr)-\frac{s^2}2 ,
\]
which is $\ge s\,(\log2-\tfrac14)/K>0$ for $s\le1/(2K)$; the negative
jumps of $\log s$ are those of $\tfrac12\log U$ (exponential tails),
the positive ones are unbounded above only through merges.
\end{lemma}

\begin{proof}
Merges: with probability $\ge2sz$ the total gains $z$, and the largest
environment block has $z\ge(1-s)/K\ge1/(2K)$; $\log(1+z/s)$ is
increasing in $z$.  Shrinks: $a\to a\sqrt U$ with probability $a^2$
changes $\log s$ by $\log(1-(a/s)(1-\sqrt U))\ge\log\sqrt U$, of mean
$-\tfrac12$; likewise $b$; $a^2+b^2\le s^2$.  Flips preserve $s$;
environment-only moves do not touch it.
\end{proof}

\emph{Status: PROVED (a calculation).  It says that, once the block
count is bounded, the marked total is a one-dimensional process with
positive log-drift below the level $\asymp1/K$ and light-tailed
down-jumps; standard drift arguments (Lyapunov function
$s^{-\theta}$, small $\theta$) then bound the occupation time of
$\{s<c/K\}$ over any window by a small fraction with polynomially
small failure probability.  This is the ``marked pair health''
input for the tail of the first-flip time in
\S\ref{sec:markovrenewal}; it is NOT written out.}

\subsubsection{The Markov-renewal form of Problem~\ref{prob:survival}}
\label{sec:markovrenewal}

Let $\tau$ be the first flip time of the quotient chain and, for
$|z|\le1$ and bounded measurable $F$,
\[
(K_zF)(x):=\E_x\bigl[z^{\tau}F(\Gamma_\tau)\bigr],
\qquad \|K_z\|_{L^\infty\to L^\infty}\le\sup_x\E_x|z|^{\tau}\le|z| ,
\]
since $\tau\ge1$.  $K_1$ is the post-flip kernel (Markov whenever
$\tau<\infty$ a.s.).

\begin{proposition}[generating function of the parity]\label{prop:genfun}
For $|z|<1$, as bounded functions of the initial state,
\[
\sum_{d\ge0}\Phi_d\,z^d\;=\;\frac1{1-z}\,(I+K_z)^{-1}(I-K_z)\mathbf 1 .
\]
\end{proposition}

\begin{proof}
With flip times $\tau_1<\tau_2<\cdots$, $\Phi_d=\sum_{k\ge0}(-1)^k
\Pr_x[\tau_k\le d<\tau_{k+1}]$ and
$\sum_dz^d\mathbf 1\{\tau_k\le d<\tau_{k+1}\}=(z^{\tau_k}-z^{\tau_{k+1}})/(1-z)$;
by the strong Markov property at the flip times
$\E_x[z^{\tau_k}]=(K_z^k\mathbf 1)(x)$.  Sum the Neumann series, which
converges since $\|K_z\|\le|z|<1$.
\end{proof}

\emph{Status: PROVED.  It is exact and separates the problem: the
factor $(I-K_z)\mathbf 1=\E_\cdot[1-z^{\tau}]$ carries the first-flip
tail (traps), the resolvent $(I+K_z)^{-1}$ carries the cancellation
(the parity of the number of flips per unit time).  Note
$(I-K_1)\mathbf 1=0$ wherever $\tau<\infty$ a.s.\ (on $\{ab>0\}$, since
$\sum_t2a_tb_t=\infty$ a.s.\ there --- taken for granted below), so the
right side is a difference quotient at $z=1$: the decay of $\Phi_d$ is the regularity of
$G(z):=(I+K_z)^{-1}(I-K_z)\mathbf 1$ on the closed disc.}

\begin{proposition}[conditional decay]\label{prop:condecay}
Let $\mathcal B$ be a Banach space of functions on the state space
containing $\mathbf 1$, on which the evaluations $F\mapsto F(x)$,
$x\in S'$, are uniformly bounded, and let $\gamma\in(0,1)$.  Suppose
\begin{enumerate}[leftmargin=2em]
\item[(R1)] $z\mapsto K_z\in\mathcal L(\mathcal B)$ is $C^{1,\gamma}$
on $\{|z|\le1\}$ (for a weighted sup-norm space
$\mathcal B=\{F:\sup|F|/V<\infty\}$ it suffices that
$\sup_x\E_x[\tau^{1+\gamma}V(\Gamma_\tau)]/V(x)<\infty$, by
$|\tau z_1^{\tau-1}-\tau z_2^{\tau-1}|\le2\tau^{1+\gamma}|z_1-z_2|^\gamma$);
\item[(R2)] $I+K_z$ is invertible on $\mathcal B$ for every $|z|\le1$,
with $\sup_{|z|\le1}\|(I+K_z)^{-1}\|<\infty$.
\end{enumerate}
Then $\sup_{x\in S'}|\Phi_d(x)|\le C\,d^{-\gamma}$.
\end{proposition}

\begin{proof}
Under (R1)--(R2), $G$ is $C^{1,\gamma}$ on the closed disc with
$G(1)=0$, so $\hat\Phi(z)=(G(z)-G(1))/(1-z)$ extends to a
$C^{0,\gamma}$ function on the closed disc; it is analytic inside, so
its boundary Fourier coefficients are the $\Phi_d$, and the Fourier
coefficients of a H\"older-$\gamma$ function are $O(d^{-\gamma})$.
Evaluate at $x\in S'$.
\end{proof}

\emph{\textbf{SESSION-11 CORRECTION.}  Condition (R2) is FALSE at
$z=\pm1$ for every function space containing the bounded functions:
$K_1\varepsilon=-\varepsilon$, because exactly one flip has occurred at
time $\tau$ (Proposition~\ref{prop:r2false}).  Proposition~\ref{prop:condecay}
is therefore vacuous as stated.  The same holds in the
label-free bookkeeping, with $\kappa_{-1}[(-1)^E]=-(-1)^E$.  It is
repaired in \S\ref{sec:s11}: (R2) is replaced by the two conditions
(C1) and (C2) of \S\ref{sec:s11left}, both statements about the single
Markov kernel $\kappa_1$, and everything on the closed disc away from
$z=\pm1$ is proved unconditionally by
Theorem~\ref{thm:idleloop}.  The discussion of (R1) below is
unaffected, and so is the reduction paragraph at the end of this
subsection.  Read the rest of this paragraph with that substitution.}

\emph{Status: PROVED as a reduction; (R1) and (R2) are OPEN.  What each
needs, and what is in hand:}

\emph{(R1) --- traps.  A uniform bound $\sup_{x\in S'}\Pr_x[\tau>t]\le
C_Kt^{-2}$ (measured: $t^2\Pr[\tau>t]\to11$ from stationary
environments, \S\ref{sec:parity}(c)) gives every moment $\tau^{1+\gamma}$,
$\gamma<1$, hence $\gamma$ arbitrarily close to $1$ in
Proposition~\ref{prop:condecay}, far above the $4/9$ that
Corollary~\ref{cor:kappafree} needs.  Ingredients: the state-independent
trap-creation flux (Lemma~\ref{lem:recharge}(iii)); the recharge of
Lemma~\ref{lem:srecharge} and the drift of Lemma~\ref{lem:logdrift},
both available on $H$ by Corollary~\ref{cor:horizon}; and the fact
that a trap of depth $g$ is left by a flip alone, at rate
$2g(s-g)$, needing no environment.  The proof of the $1/t^2$ tail is
a renewal computation with these inputs (the subcritical kernel
$\approx2/(s^2m^2)$ of \S\ref{sec:parity}(c)), which we have not
written.}

\emph{(R2) --- cancellation.  This is the parity content and cannot be
had from trap accounting.  It is needed on the whole closed disc: the
trivial bound $\|K_z\|\le|z|$ gives nothing on $|z|=1$.  At $z=1$ it
is the statement that $-1$ is not in the spectrum of the post-flip
chain $K_1$ on $\mathcal B$, with a uniform inverse; near $z=1$ it
follows by perturbation under (R1).  At other points of the circle it
is the non-lattice condition of Markov renewal theory --- no $F\ne0$
with $\E_x[z^{\tau}F(\Gamma_\tau)]=-F(x)$ --- and at $z=-1$, where
$K_{-1}F=\E_\cdot[(-1)^{\tau}F(\Gamma_\tau)]$, it says the parity of the
inter-flip times is not asymptotically deterministic; this is where
the $(-1)^d$ alternation of \S\ref{sec:parity}(a) lives and it needs
its own statement.  The Doeblin structure is on the marked pair: two consecutive
flips from any $(a,b)$ with total $s$ produce $a''$ whose law dominates
$\tfrac14\cdot\mathrm{Unif}[s/4,3s/4]$ (the trapezoid of
Lemma~\ref{lem:recharge}(ii) is symmetric and unimodal, so the first
flip lands in $[s/4,3s/4]$ with probability $\ge\tfrac12$, and from
there the second has density $\ge1/s$ on that interval).  This is a
minorisation on the fibre $\{a'+b'=s,\ \mathrm{env}\}$ only; the base
$(s,\mathrm{env})$ is not regenerated, so $\mathcal B$ must be a
weighted space in which the post-flip chain is $V$-geometrically
ergodic in the base coordinates as well (Meyn--Tweedie), with the
weight matched to the moment condition in (R1).  On $H$ the base is
finite-dimensional ($<K$ blocks), which is what makes this plausible.
Because the label parity is determined by the environment block-count
parity (Lemma~\ref{lem:parity}), no argument can avoid the environment
in (R2) either; but the environment enters only through the ergodicity
of the base chain between flips, not through any fine property of
$\mathrm{PD}(1)$.}

\emph{Reduction to Problem~\ref{prob:survival}.  Take
$S'=S'_K:=\{a,b\ge1/(16K)\}$ and enforce $H_{9d}$ (so $C=C(K)$ in
Proposition~\ref{prop:condecay} may depend on the block bound, which
holds up to time $d$).  Granting Proposition~\ref{prop:condecay} on
$S'_K$ with $C(K)$ polynomial in $K$, and granting the occupation-time
consequence of Lemma~\ref{lem:logdrift} (that after $s$ first reaches
$1/(2K)$ it stays $\ge1/(4K)$ for at least half of the next $T_2$
steps, except with polynomially small probability; NOT written out),
the chain from $\Gamma_0\sim\mu_j$ with $a_0,b_0\ge\varepsilon$ enters
$S'_K$ at a stopping time $T'$: first $s$ reaches $1/(2K)$ within
$T_1=O(\varepsilon^{-1}K\log d)$ steps (Lemma~\ref{lem:srecharge});
then, on the times with $s_t\ge1/(4K)$ and $\min(a_t,b_t)\ge\varepsilon$
(the latter failing with probability $\le4\varepsilon^2(T_1+T_2)$ by
the flux bound), a flip occurs with probability $\ge2\varepsilon(s_t-\varepsilon)
\ge\varepsilon/(4K)$ per step and lands with $\min(a',b')\ge s_t/4\ge1/(16K)$
with conditional probability $\ge\tfrac12$ (the trapezoid is
symmetric and unimodal), so $T'\le T_1+T_2$ with
$T_2=O(\varepsilon^{-1}K\log d)$ except with probability
$O(d^{-\gamma})+O(\varepsilon K\log d)$.  Note that a single
application of Lemma~\ref{lem:srecharge} to each coordinate does not
suffice: it makes $a$ and $b$ large at different times, and a
coordinate of size $1/(2K)$ shrinks on the scale $K^2\ll\varepsilon^{-1}K$.
By the strong Markov property
$|\Phi_d(\Gamma_0)|\le\E|\Phi_{d-T'}(\Gamma_{T'})|+\Pr[T'>T_1+T_2]
\le C(K)(d/2)^{-\gamma}+O(d^{-\gamma})+O(\varepsilon K\log d)$, and
with $K=O(\log d)$ and $\varepsilon$ as in Corollary~\ref{cor:kappafree}
this is $\le C(\varepsilon d)^{-2\alpha}$ for $\alpha<\gamma/2$ (for
$\varepsilon\ge1/(16K)$ no recharge is needed).  Thus
\emph{Problem~\ref{prob:survival}, hence the iid model, is reduced to
(R1) and (R2) on the log-macroscopic set $S'_K$ plus the
occupation-time step}, with the environment controlled throughout by
Corollary~\ref{cor:horizon}.}

%----------------------------------------------------------------------
\subsection{Session 11: (R2) is false; the idle loop proves its repair
on the whole disc except two points; and the residual condition is one
scalar statement about the parity of the environment block
count}\label{sec:s11}

Session 11 was directed at condition (R2) of
Proposition~\ref{prop:condecay} --- uniform invertibility of $I+K_z$ on
the closed unit disc --- under the instruction that a fourth
reformulation of the cancellation core would not count as progress.
There are four findings, in logical order.

\begin{enumerate}[leftmargin=2.2em]
\item[(1)] \textbf{The label is a passenger.}  The quotient chain
$\bar\Gamma=(p_1,p_2;\mathrm{env})$ --- the two marked masses and the
environment, with the label $\varepsilon$ forgotten --- is
\emph{autonomous}: its transition law does not depend on
$\varepsilon$, and the involution $J$ that swaps the two sectors
keeping $\bar\Gamma$ fixed commutes with the chain
(Lemma~\ref{lem:autonomy}).  Hence $K_z=\kappa_z\oplus(-\kappa_z)$ for
the quotient first-flip kernel $\kappa_z$, and
$\spec K_z=\pm\spec\kappa_z$.  One line of this gives
$\Pr_\pi[u\sim v]=\tfrac12$.
\item[(2)] \textbf{(R2) is false.}  With the label kept,
$K_1\varepsilon=-\varepsilon$ --- exactly one flip has occurred at
time $\tau$ --- so $I+K_1$ has a bounded null vector.  On the quotient,
where that artefact is absent, $I+\kappa_{-1}$ still has the bounded
null vector $(-1)^{E}$ ($E=$ number of environment curves), by
Lemma~\ref{lem:parity}.  Either way
$\sup_{|z|\le1}\|(I+K_z)^{-1}\|=\infty$ and
Proposition~\ref{prop:condecay} is vacuous as stated
(Proposition~\ref{prop:r2false}).
\item[(3)] \textbf{An unconditional contraction off $z=\pm1$.}  The
chain can waste exactly two steps by splitting a curve and merging the
two pieces it just created, returning to the same state without
flipping.  This ``idle loop'' has probability
$\lambda(\Gamma)=\tfrac13\sum_jr_j^4$ --- an exact formula, $r_j$ the
arc-refined mass list --- and it forces
$\|\kappa_z\|_\infty\le(1-\lambda)/|1-\lambda z^2|
\le1/(1+\lambda(1-\Re z^2))$ for all $|z|\le1$
(Theorem~\ref{thm:idleloop}).  On the block-bounded state space this is
uniform, with $\lambda\ge\tfrac13K^{-3}$.  The non-lattice condition of
Markov renewal theory on the \emph{whole} circle --- Pitfall~45,
TODO~11b(v), and the third name of the cancellation core --- is
therefore \emph{proved}, unconditionally and with explicit constants.
\item[(4)] \textbf{What is left is two scalar conditions on one Markov
kernel.}  After (3), $\widehat\Phi(z)=(1-z)^{-1}(I+\kappa_z)^{-1}
(I-\kappa_z)\mathbf 1$ can fail to be regular only at $z=1$ and
$z=-1$.  At $z=-1$ the $\bar M$-conjugation of Lemma~\ref{lem:conj}
turns the problem into the ordinary renewal pole at $+1$ applied to the
test function $(-1)^{E}$, and the condition is
\[
\text{(C2)}\qquad
\bigl\|\kappa_1^{\,k}(-1)^{E}\bigr\|_\infty\le C\rho^k
\quad\text{for some }\rho<1 ,
\]
i.e.\ \emph{the parity of the environment block count decorrelates
geometrically along the post-flip chain}.  At $z=1$ the condition is
$-1\notin\spec\kappa_1$ with a uniform inverse (C1), and we prove that
a single Doeblin minorisation shared by the $n$-th and $(n{+}1)$-st
post-flip laws implies it, with
$\|(I+\kappa_1)^{-1}\|\le(2n+1)/(2\delta)$
(Proposition~\ref{prop:c1}).  Neither condition contains a $z$, a
phase, or an alternating sign that is not the sign of a state function.
(C2) is exactly what \texttt{s11\_kflip} measures, and it holds
numerically with $\rho\approx0.3$ from both a macroscopic and a deeply
trapped start (\S\ref{sec:s11num}).
\end{enumerate}

\subsubsection{The label is a passenger}\label{sec:s11quotient}

Recall the state: a partition of the circle into curves with masses,
two marks $u,v$, and the pair $(p_1,p_2)$ --- the two arcs $(x,y)$ of
the marked curve when $u\sim v$, the two marked curve masses $(a,b)$
when $u\not\sim v$ --- with $s=p_1+p_2$, and $\varepsilon=\pm1$ the
label.  Let $\bar\Gamma=(p_1,p_2;\mathrm{env})$ be the state with the
label forgotten and $\pi_\ast:\Gamma\mapsto\bar\Gamma$ the projection.

\begin{lemma}[label autonomy]\label{lem:autonomy}
The transition kernel of $\Gamma$ projects to a kernel on
$\bar\Gamma$, and in both sectors the moves and their probabilities are
\[
\text{flip: }2p_1p_2,\qquad
p_i\to p_i\sqrt U:\ p_i^2,\qquad
p_i\to p_i+q_j:\ 2p_iq_j,\qquad
\text{environment moves: unchanged,}
\]
a flip replacing $(p_1,p_2)$ by $(V,s-V)$ with
$V\sim\mathrm{U}(0,p_1)\star\mathrm{U}(0,p_2)$ and toggling
$\varepsilon$.  Consequently:
\begin{enumerate}[leftmargin=2em]
\item[(i)] the involution $J$ of the state space that swaps the two
sectors keeping $\bar\Gamma$ fixed commutes with the transition kernel
and fixes every flip time;
\item[(ii)] in the decomposition
$F=g\circ\pi_\ast+\varepsilon\cdot(f\circ\pi_\ast)$ of bounded
functions, $K_z=\kappa_z\oplus(-\kappa_z)$, where
$(\kappa_zg)(\bar x)=\E_{\bar x}[z^\tau g(\bar\Gamma_\tau)]$; hence
$\spec K_z=\spec\kappa_z\cup(-\spec\kappa_z)$ and
$\widehat\Phi(z)=(1-z)^{-1}(I+\kappa_z)^{-1}(I-\kappa_z)\mathbf 1$;
\item[(iii)] if $\bar\pi$ is a stationary law of $\bar\Gamma$ then the
stationary label is uniform: $\Pr[u\sim v]=\tfrac12$.
\end{enumerate}
\end{lemma}

\begin{proof}
A step picks two uniform points of the circle; if they lie in the same
curve it is a split, otherwise a merge.  In the label-$+$ sector the
marked curve has mass $s$ and is picked twice with probability $s^2$;
given that, the two cuts lie on different arcs with probability
$2xy/s^2$ (a flip, separating $u$ from $v$) and on the same arc $i$
with probability $p_i^2/s^2$, in which case the arc loses a factor
$\sqrt U$ in law ($1-|U_1-U_2|\sim\sqrt U$) and the removed piece
becomes an unmarked curve.  A merge of the marked curve with an
environment curve $q_j$ has probability $2sq_j$, and the mass $q_j$ is
attached to arc $i$ with probability $p_i/s$, i.e.\ rate $2p_iq_j$.  In
the label-$-$ sector the same four rates hold with $(a,b)$ in place of
$(x,y)$: picking $u$'s and $v$'s curves is a merge --- a flip, joining
them --- with probability $2ab$; picking $u$'s curve twice splits it,
$a\to a\sqrt U$; picking $u$'s curve and $q_j$ merges them,
$a\to a+q_j$.  The post-flip law is the same trapezoid
$\mathrm{U}(0,p_1)\star\mathrm{U}(0,p_2)$ in both directions
(Lemma~\ref{lem:recharge}(ii)).  This is the displayed list, and it
does not mention $\varepsilon$.  (i) is immediate; (ii) follows because
$\varepsilon(\Gamma_\tau)=-\varepsilon(x)$ while the law of
$\bar\Gamma_\tau$ depends only on $\bar x$; (iii) because $J$ preserves
the stationary law --- it commutes with the kernel and fixes
$\bar\Gamma$ --- while exchanging the two sectors.
\end{proof}

\emph{Status: PROVED.  Implicit in the phrase ``quotient chain'' used
since session 8 but never stated; it is worth stating because it is
what makes the next two subsections possible.  Two consequences.  (a)
The first-flip time $\tau$ has the same law from $x$ and from $Jx$:
the alternating renewal governing $\Phi_d$ is \emph{symmetric}, which
is why $\Phi_d\to0$ rather than to the generic limit
$(\E_+\tau-\E_-\tau)/(\E_+\tau+\E_-\tau)$ of an asymmetric alternating
renewal.  This is the structural reason behind
$\Pr_\pi[\mathrm{tog}]=\tfrac12$, and it is exactly the residue that
has to vanish at $z=1$ below.  (b) All the parity content therefore
lives in $\kappa_z$, a kernel on a space where $\varepsilon$ does not
exist.}

\subsubsection{(R2) is false, in both bookkeepings}\label{sec:s11false}

Write $E(\Gamma)$ for the number of environment curves and let $\bar M$
be multiplication by $(-1)^{E}$, an isometric involution of every
weighted sup-norm space; $N$ denotes multiplication by $\varepsilon$.

\begin{lemma}[sign conjugations]\label{lem:conj}
For every $z\in\mathbb C$,
\[
\text{(i)}\quad N\,K_z\,N=-K_z,
\qquad\qquad
\text{(ii)}\quad \bar M\,\kappa_z\,\bar M=-\kappa_{-z} .
\]
Hence $\spec K_z=-\spec K_z$ and $\spec\kappa_{-z}=-\spec\kappa_z$, and
the resolvents $(I\pm\kappa_{\pm z})^{-1}$ have equal norms in pairs.
\end{lemma}

\begin{proof}
(i) Exactly one flip has occurred at time $\tau$, so
$\varepsilon(\Gamma_\tau)=-\varepsilon(x)$ a.s.
(ii) By Lemma~\ref{lem:parity}, $(-1)^{N_t}=(-1)^t(-1)^{E_t-E_0}$; at
$t=\tau$ we have $N_\tau=1$, so $(-1)^\tau=-(-1)^{E_\tau-E_0}$.
Therefore
$(\bar M\kappa_z\bar Mg)(\bar x)=(-1)^{E(\bar x)}
\E_{\bar x}\bigl[z^\tau(-1)^{E(\bar\Gamma_\tau)}g(\bar\Gamma_\tau)\bigr]
=-\E_{\bar x}\bigl[(-z)^\tau g(\bar\Gamma_\tau)\bigr]$.
\end{proof}

\emph{Remark (the label-carrying form of (ii)).  Since $B$ changes by
exactly $\pm1$ at every step, $(-1)^\tau=(-1)^{B(\Gamma_\tau)-B(x)}$
directly, so with $M$ = multiplication by $(-1)^{B}$ one has the
cleaner-looking $M K_zM=K_{-z}$ on the label-carrying space, and
$K_{-1}[(-1)^B]=+(-1)^B$.  That eigenvalue is $+1$ and is harmless;
the obstruction to (R2) is the eigenvalue $-1$ of
$\kappa_{-1}$ at $(-1)^E$ in Proposition~\ref{prop:r2false}.  The two
are consistent: $(-1)^B=-\varepsilon\,(-1)^E$, and
$N K_zN=-K_z$ supplies the sign.  Keeping these apart matters: in the
label-carrying picture $\E_x[(-1)^{\tau_k}]$ has no a-priori reason to
decay, while in the label-free picture its decay is exactly (C2).}

\begin{proposition}[failure of (R2)]\label{prop:r2false}
Assume $\tau<\infty$ a.s.\ (as is taken for granted after
Proposition~\ref{prop:genfun}), so that $\kappa_1$ is Markov.  Then
\[
K_1\varepsilon=-\varepsilon
\qquad\text{and}\qquad
\kappa_{-1}\bigl[(-1)^{E}\bigr]=-(-1)^{E} .
\]
Hence $I+K_z$ is not injective at $z=1$ in the label-carrying
bookkeeping and $I+\kappa_z$ is not injective at $z=-1$ in the
label-free one, on every function space containing the bounded
functions.  In particular condition (R2) of
Proposition~\ref{prop:condecay} never holds, and
Proposition~\ref{prop:condecay} is vacuous as stated.
\end{proposition}

\begin{proof}
The first identity is Lemma~\ref{lem:conj}(i) applied to $\mathbf 1$
together with $K_1\mathbf 1=\mathbf 1$.  For the second, by
Lemma~\ref{lem:conj}(ii) with $z=1$,
$\kappa_{-1}\bar M\mathbf 1=-\bar M\kappa_1\mathbf 1=-\bar M\mathbf 1$.
\end{proof}

\emph{Status: PROVED.  The failure is structural, not technical: the
null vectors $\varepsilon$ and $(-1)^E$ are bounded by $1$ and lie in
every space in sight, so no choice of weight, state space or test class
removes them.  Session 10 did not see it because
Proposition~\ref{prop:genfun} applies $(I+K_z)^{-1}$ only to
$(I-K_z)\mathbf 1$, which degenerates at the same points; the whole
question is whether the vanishing of the numerator beats the blow-up of
the resolvent, and that is not what (R2) says.  The two failures are
the same failure in the two bookkeepings, since
$\spec K_z=\pm\spec\kappa_z$ by Lemma~\ref{lem:autonomy}(ii).}

\begin{proposition}[the two singular points, exactly]\label{prop:repair}
For $|z|<1$,
$\widehat\Phi(z)=(1-z)^{-1}(I+\kappa_z)^{-1}(I-\kappa_z)\mathbf 1$, and
for every $w$,
\[
(I+\kappa_{-w})^{-1}(I-\kappa_{-w})\mathbf 1
=\bar M\,(I-\kappa_{w})^{-1}(I+\kappa_{w})\,\bar M\mathbf 1 .
\]
Thus near $z=1$ the singular factor is $(I+\kappa_z)^{-1}$ acting on
the numerator $(I-\kappa_z)\mathbf 1=\E_\cdot[1-z^\tau]=O(|1-z|)$,
while near $z=-1$ it is the ordinary renewal resolvent
$(I-\kappa_w)^{-1}$, $w\to1$, acting on
$(I+\kappa_w)(-1)^{E}\to2\,(-1)^{E}$.
\end{proposition}

\begin{proof}
The first statement is Lemma~\ref{lem:autonomy}(ii) substituted in
Proposition~\ref{prop:genfun}; the second is Lemma~\ref{lem:conj}(ii)
together with $\bar M^2=I$.
\end{proof}

\emph{Status: PROVED (algebra).  This is the sharp form of the
correction.  At $z=-1$ the singularity is the classical simple pole of
a Markov renewal resolvent at $w=1$, whose residue is the stationary
projection $\langle\bar\pi,\cdot\rangle\mathbf 1$; the residue of
$\widehat\Phi$ at $z=-1$ is therefore proportional to
$\langle\bar\pi,(-1)^{E}\rangle$, and $\Phi_d$ carries a non-decaying
alternating part $c\,(-1)^d$ unless that pairing vanishes.  That is
precisely the alternation found in Lemma~\ref{lem:parity}(a) --- where
it is measured to decay like $(-1)^dc_1/d^2$, i.e.\ the residue does
vanish.  At $z=1$ the numerator vanishes and the residue to be beaten
is the residue of a \emph{symmetric} alternating renewal, zero by
Lemma~\ref{lem:autonomy}(a).  Both residues therefore vanish for
structural reasons; what is unproved is the quantitative, uniform
version, and that is (C1) and (C2) of \S\ref{sec:s11left}.}

\subsubsection{The idle loop: an unconditional contraction off
\texorpdfstring{$z=\pm1$}{z=+-1}}\label{sec:idleloop}

Write the \emph{arc-refined mass list} $r(\Gamma)=(r_j)_j$ for the
masses of the unmarked curves together with $p_1$ and $p_2$; so
$\sum_jr_j=1$ and $r$ has at most $B(\Gamma)+1$ entries, $B$ the total
number of curves.

\begin{lemma}[the idle loop]\label{lem:idleloop}
Let $L$ be the event that step~1 splits some curve and step~2 merges
the two pieces that step~1 created, the returned piece being
re-attached to the arc it came from if the split curve was the marked
one in the label-$+$ sector.  On $L$ no flip occurs and
$\bar\Gamma_2=\bar\Gamma_0$; up to a null set
$L=\{\bar\Gamma_2=\bar\Gamma_0\}$; and
\[
\lambda(\Gamma):=\Pr_\Gamma[L]\;=\;\frac13\sum_j r_j(\Gamma)^4
\;\ge\;\frac1{3\,(B(\Gamma)+1)^3}.
\]
\end{lemma}

\begin{proof}
Fix an entry $r$ of the arc-refined list; the events for distinct
entries are disjoint.

\emph{$r=q$, an unmarked curve.}  Step~1 picks it twice with
probability $q^2$ and cuts it at two independent uniform points,
giving pieces $qV$ and $q(1-V)$ with $V=|U_1-U_2|$, density $2(1-v)$.
Step~2 picks those two curves with probability $2\,qV\cdot q(1-V)$ and
merges them back into a curve of mass $q$.  Total
$q^2\cdot2q^2\,\E[V(1-V)]=q^4/3$, since
$\E[V(1-V)]=\int_0^1v(1-v)\,2(1-v)\,dv=\tfrac16$.

\emph{$r=a=p_1$ in the label-$-$ sector.}  Identical, the mark staying
in the piece of mass $a(1-V)$: total $a^4/3$.

\emph{$r=x$, one arc in the label-$+$ sector.}  Both cuts must fall in
that arc, probability $x^2$ --- and this is a split, not a flip,
precisely because a flip needs the cuts on \emph{different} arcs.  The
piece has mass $w=xV$; the marked curve keeps $s-w$ and the arc becomes
$x-w$.  Step~2 merges the piece back with probability $2w(s-w)$, and
the returned mass is re-attached to the arc it came from with
conditional probability $(x-w)/(s-w)$.  The factor $s-w$ cancels:
$x^2\,\E[2\,xV\cdot x(1-V)]=2x^4\E[V(1-V)]=x^4/3$.

Summing gives $\lambda=\tfrac13\sum_jr_j^4$.  Conversely a merge at
step~1 followed by a split at step~2 reproduces the two merged masses
with probability $0$, and a split followed by the merge of two other
curves reproduces the mass list only on a null set.  Finally, with
$p\le B+1$ entries, Cauchy--Schwarz twice gives
$\sum r_j^4\ge(\sum r_j^2)^2/p\ge p^{-3}$.
\end{proof}

\emph{Status: PROVED.  VERIFIED (\texttt{s11\_ztau}, $2\times10^6$
two-step trials per start, exact state comparison at tolerance
$10^{-12}$): one curve, arcs $(\tfrac12,\tfrac12)$: $\lambda$ measured
$0.041530$ vs.\ $\tfrac13\cdot2\cdot2^{-4}=0.041667$; $s_0=0.6$, arcs
$(0.1,0.5)$, environment $0.4$: $0.029459$ vs.\ $0.029400$;
$s_0=0.5$, arcs $(0.02,0.48)$, environment $0.5$: $0.038556$ vs.\
$0.038528$.  The formula is an identity, not an estimate, and it is the
\emph{arc-refined} list that appears: splitting the marked curve across
its two arcs is a flip, not an idle loop.  The quantity $\sum_jr_j^4$
is the one whose expectation Theorem~\ref{thm:sizebiased} computes
exactly, $\E\sum p^4(m)=\tfrac14+\int_0^1(1-x^3)(1-2x)^mdx\to\tfrac14$,
so $\E\lambda\to\tfrac1{12}$; the block bound is needed only to turn
that average into a uniform bound.}

\begin{theorem}[idle-loop contraction]\label{thm:idleloop}
Let $\lambda_*=\inf_{\bar x}\lambda(\bar x)$ over the state space on
which $\kappa_z$ acts.  For all $|z|\le1$,
\[
\|\kappa_z\|_{\infty\to\infty}\;\le\;
\sup_{\bar x}\frac{1-\lambda(\bar x)}{|1-\lambda(\bar x)z^2|}
\;\le\;\frac1{1+\lambda_*\,(1-\Re z^2)} .
\]
Consequently if $1-\Re z^2\ge t_0$ then $I\pm\kappa_z$ and
$I-\kappa_z^2$ are invertible with
$\|(I\pm\kappa_z)^{-1}\|\le1+(\lambda_*t_0)^{-1}$.  On the state space
$\Omega_K=\{B\le K\}$ --- the chain killed at the first split that
would make $B=K+1$, whose $\Phi_d$ differs from the true one by at most
$\Pr[H^c]\le(9d+1)^{-\gamma}$ (Corollary~\ref{cor:horizon}) --- one has
uniformly
\[
\|\kappa_z\|_\infty\;\le\;1-\min\Bigl(\frac{t_0}{6K^3},\ \frac1{2K}\Bigr),
\qquad\text{so}\qquad
\|(I\pm\kappa_z)^{-1}\|\le\max\Bigl(\frac{6K^3}{t_0},\,2K\Bigr).
\]
\end{theorem}

\begin{proof}
$L$ is measurable with respect to the first two steps; on $L$ no flip
occurs and $\bar\Gamma_2=\bar x$.  By the Markov property at time $2$,
$\E_{\bar x}[z^\tau g(\bar\Gamma_\tau)\mathbf 1_L]
=z^2\lambda(\bar x)(\kappa_zg)(\bar x)$, whence
\[
(1-\lambda z^2)(\kappa_zg)(\bar x)
=\E_{\bar x}\bigl[z^\tau g(\bar\Gamma_\tau)\mathbf 1_{L^c}\bigr],
\qquad
|(\kappa_zg)(\bar x)|\le\frac{(1-\lambda)\|g\|_\infty}{|1-\lambda z^2|} .
\]
For $|z|\le1$, $|1-\lambda z^2|\ge\Re(1-\lambda z^2)
=(1-\lambda)+\lambda t$ with $t=1-\Re z^2$, and
$u\mapsto(1-u)/(1-u+ut)$ has derivative $-t/(\,\cdot\,)^2<0$, so the
supremum over $\bar x$ is attained at $\lambda_*$ and equals
$1/(1+\lambda_*t/(1-\lambda_*))\le1/(1+\lambda_*t)$.  The resolvent
bounds are Neumann series.

For the killed chain, split the states.  If $B(\bar x)\le K-1$ the idle
loop raises $B$ only to $B+1\le K$, so it stays inside $\Omega_K$ and
is not killed; Lemma~\ref{lem:idleloop} gives
$\lambda\ge\tfrac13K^{-3}$, and $1/(1+u)\le1-u/2$ for $u\le1$ gives the
first branch.  If $B(\bar x)=K$ then \emph{every} split kills, and a
split occurs at the next step with probability
$\sum_i(\text{curve mass})^2\ge1/K$, so
$|\kappa_zg(\bar x)|\le\|g\|_\infty\Pr_{\bar x}[\tau<\text{kill}]
\le(1-1/K)\|g\|_\infty$.  Take the worse of the two branches.
\end{proof}

\emph{Status: PROVED.  VERIFIED (\texttt{s11\_ztau}): from each of the
three deterministic starts above,
$|\E_{\bar x}[z^\tau]|=|\kappa_z\mathbf 1(\bar x)|$ was measured at
$\theta/\pi=0.1,\dots,0.9$ ($2\times10^5$ trials each) and compared
with $(1-\lambda)/|1-\lambda z^2|$: the bound holds at every point,
with margin (at $\theta=\pi/2$ from the $(\tfrac12,\tfrac12)$ start,
measured $0.446$ against the bound $0.920$).  The bound is crude --- it
uses one two-step return and nothing else --- but it is \emph{uniform
in the state}, which is the entire difficulty; and it is the right
mechanism, since $1-\lambda z^2$ is exactly the Euler factor of the
geometric component that the idle loop puts into the law of $\tau$ at
lag $2$, which is what an aperiodicity or non-lattice hypothesis is
normally assumed to supply.}

\emph{Three items on the open list are settled by
Theorem~\ref{thm:idleloop} and Lemma~\ref{lem:conj}.  (i) Pitfall~45
and TODO~11b(v) --- the non-lattice condition ``on the whole circle,
not only near $z=1$'' --- are PROVED, with rate $1-\Re z^2$ and
constant $3K^3$.  (ii) TODO~11b(iv), ``$-1\notin\spec K_1$'', is
DISPROVED and was the wrong question
(Proposition~\ref{prop:r2false}).  (iii) The point $z=-1$, treated as a
separate problem in every handoff since session 8, is by
Proposition~\ref{prop:repair} the ordinary renewal point $w=1$ with the
test function $(-1)^E$: not a separate problem, and not a place where
anything new has to be invented.}

\subsubsection{What is left: (C1) and (C2)}\label{sec:s11left}

By Theorem~\ref{thm:idleloop}, $\widehat\Phi$ is regular on
$\{|z|\le1\}\setminus\{\pm1\}$ with bounds polynomial in $K$.  Two
conditions remain, both about the single Markov kernel $\kappa_1$ ---
the post-flip quotient chain --- on $\Omega_K$:

\medskip
\noindent\textbf{(C1)}\quad$-1\notin\spec\kappa_1$, with
$\|(I+\kappa_1)^{-1}\|\le C(K)$;\\
\noindent\textbf{(C2)}\quad$\bigl\|\kappa_1^{\,k}(-1)^{E}\bigr\|_\infty
\le C(K)\rho^k$ for some $\rho=\rho(K)<1$.
\medskip

\noindent With (R1) and these two, the proof of
Proposition~\ref{prop:condecay} goes through: (C1) gives regularity at
$z=1$, where the numerator has its zero; (C2) gives, through
Proposition~\ref{prop:repair}, both the vanishing of the residue at
$z=-1$ --- since $\langle\bar\pi,(-1)^E\rangle=\lim_k\kappa_1^k(-1)^E=0$
--- and a Hölder modulus there.  \emph{Neither contains a $z$, a phase,
or an alternating sign that is not the sign of a state function: the
parity content of the problem has been reduced to the single statement
that the environment block count has no asymptotic parity bias along
the post-flip chain.}

\emph{Neither (C1) nor (C2) is false for the reason (R2) was.  The
only deterministically-toggled state function available on the quotient
is $(-1)^E$ (flips leave $E$ fixed; every non-flip step changes it by
$\pm1$), and by Lemma~\ref{lem:conj}(ii)
$\kappa_1[(-1)^E]=-(-1)^{E}\,\E_\cdot[(-1)^\tau]$, so $(-1)^E$ is an
eigenfunction of $\kappa_1$ if and only if $\E_x[(-1)^\tau]$ is
constant in $x$ --- which it measurably is not
($-0.1306\pm0.003$ from the macroscopic start against
$-0.0009\pm0.005$ from the trapped start, \S\ref{sec:s11num}).}

\begin{proposition}[a minorisation suffices for (C1)]\label{prop:c1}
Suppose there are $n\ge1$, $\delta>0$ and a probability measure $\nu$
with
\[
\Pr_{\bar x}\bigl[\bar\Gamma_{\tau_n}\in\cdot\,\bigr]\ge\delta\nu(\cdot)
\qquad\text{and}\qquad
\Pr_{\bar x}\bigl[\bar\Gamma_{\tau_{n+1}}\in\cdot\,\bigr]\ge\delta\nu(\cdot)
\qquad\text{for every }\bar x .
\]
Then $I+\kappa_1$ is invertible on $L^\infty$ with
$\|(I+\kappa_1)^{-1}\|\le(2n+1)/(2\delta)$.
\end{proposition}

\begin{proof}
Let $(I+\kappa_1)F=G$ and $m=\|F\|_\infty$.  Iterating
$F=G-\kappa_1F$ gives $F=H_j+(-1)^j\kappa_1^{\,j}F$ with
$H_j=\sum_{i<j}(-1)^i\kappa_1^{\,i}G$, $\|H_j\|\le j\|G\|$.  Write
$\kappa_1^{\,n}F=\delta c+(1-\delta)\rho^1_{\bar x}(F)$ and
$\kappa_1^{\,n+1}F=\delta c+(1-\delta)\rho^2_{\bar x}(F)$ with the same
$c=\nu(F)$ and probability measures $\rho^i_{\bar x}$.  Multiplying the
$j=n$ identity by $(-1)^n$ and the $j=n+1$ identity by $(-1)^{n+1}$ and
subtracting the results eliminates the $\delta c$ terms' difference and
leaves
\[
2F(\bar x)=\pm\bigl(H_n+H_{n+1}\bigr)(\bar x)
+(1-\delta)\bigl(\rho^2_{\bar x}-\rho^1_{\bar x}\bigr)(F),
\]
so $2m\le(2n+1)\|G\|+2(1-\delta)m$, i.e.\
$2\delta m\le(2n+1)\|G\|$.
\end{proof}

\emph{The proof uses no property of $\kappa_1$ beyond
$\|\kappa_1\|\le1$: it is valid verbatim for the \emph{sub}-Markov
$\kappa_1$ of the chain killed on $\Omega_K$, the residual measures
$\rho^i_{\bar x}$ then having mass $\le1-\delta$ rather than exactly
$1-\delta$.  This matters because $\Omega_K$ is where
Theorem~\ref{thm:idleloop} lives.}

\emph{\textbf{SESSION-12 CORRECTION.}  The hypothesis of this proposition
is never satisfied: for every $n$ there is no common minorant of the
$n$-th post-flip laws over $\Omega_K$, because a curve of mass
$\varepsilon$ survives the first $n$ flips untouched with probability
$\to1$ (Proposition~\ref{prop:nomino}).  The proposition stays correct
and vacuous; TODO~12b is closed.}

\emph{Status: PROVED as a reduction; the minorisation is OPEN.  It is
the honest remaining difficulty and it is the classical one: the
post-flip chain regenerates on the fibre --- two consecutive flips
produce a $p_1''$ whose law dominates
$\tfrac14\mathrm{Unif}[s/4,3s/4]$ (\S\ref{sec:markovrenewal}) --- but
not in the base $(s,\mathrm{env})$, which is carried along.  On
$\Omega_K$ the base has at most $K$ curves, and the numerics of
\S\ref{sec:s11num} show it equilibrating in about four flips from a
deliberately bad start, so a minorisation with $n=O(1)$ and a
$\delta$ polynomial in $1/K$ is the natural target; the $K$-dependence
of $\delta$ may be as bad as one likes, polynomially.}

\emph{Routes to (C2).  (a) Uniform ergodicity plus a parity symmetry:
the same minorisation gives
$\|\kappa_1^k(F-\bar\pi F)\|\le C\rho^k$ for all bounded $F$, reducing
(C2) to $\langle\bar\pi,(-1)^E\rangle=0$ --- the statement that the
post-flip stationary law on $\Omega_K$ has no environment block-count
parity bias.  On $\Omega_K$ that is a statement about a
finite-dimensional chain and is the cleanest target for session 12.
(b) Directly: by Lemma~\ref{lem:conj}(ii),
$\kappa_1^{\,k}(-1)^E=(-1)^{E}(-1)^k\kappa_{-1}^{\,k}\mathbf 1$ and
$\kappa_{-1}^{\,k}\mathbf 1(\bar x)=\E_{\bar x}[(-1)^{\tau_k}]$, so
(C2) says exactly that the parity of the $k$-th flip time decorrelates
geometrically --- which is what \texttt{s11\_kflip} measures.  A proof
along (b) needs a state-uniform lower bound on the probability that
two consecutive inter-flip times have opposite parity; the idle loop
supplies a uniform event at lag $2$, but it changes $\tau$ by an
\emph{even} amount, so a genuinely odd perturbation is still needed.
That is the one place where an odd/even asymmetry must be produced
rather than conjugated away, and it is the sharpest statement of the
residue of the whole problem.}

\emph{A Lyapunov remark.}  The proof of Theorem~\ref{thm:idleloop}
degrades only through $\lambda^{-1}=3(\sum_jr_j^4)^{-1}$, whose
expectation along the chain is known exactly
(Theorem~\ref{thm:sizebiased}).  A drift inequality
$\kappa_1V\le\varrho V+b$ with $V=(\sum_jr_j^4)^{-\theta}$, $\theta>0$
small, would replace the killing at $B=K$ by a weighted space and
supply the ``base'' half of the minorisation of
Proposition~\ref{prop:c1} at the same time.  It is a finite
computation of the same kind as Lemma~\ref{lem:logdrift} and is the
recommended first task of session 12.

\subsubsection{Numerics along the post-flip chain}\label{sec:s11num}

TODO~11d asked for $\E_{\bar x}[(-1)^{\tau_k}]$ along the post-flip
chain, predicting geometric decay under (R2).  The prediction needed
correcting before the run could be interpreted --- under (R2) as stated
there is no such prediction, (R2) being false --- but the quantity is,
by route (b) above, \emph{exactly} the residual condition (C2).
Measured (\texttt{s11\_kflip}: stationary environment burnt in over
$4096$ steps, accepted on the initial gap, $k\le16$, no censoring):

\emph{Macroscopic start, gap in $[\tfrac18,\tfrac38)$,
$1.2\times10^5$ samples (s.e.\ $0.0029$):
$\E[(-1)^{\tau_k}]=-0.1306,\ +0.0263,\ -0.0100,\ +0.0011$, and at the
noise floor for every $k\ge4$: geometric decay with ratio
$\approx0.2$--$0.4$ and no persistent $(-1)^k$ component.  Trap start,
gap in $[10^{-3},10^{-2})$, $4\times10^4$ samples (s.e.\ $0.0050$):
$-0.0009,\,+0.0018,\,+0.0024,\dots$, at the noise floor throughout ---
one long inter-flip time randomises the parity immediately.
\textbf{This is (C2), measured, with $\rho\approx0.3$.}}

\emph{\textbf{SESSION-12 CORRECTION.}  The ``$\rho\approx0.3$'' is the
transient.  With $2\times10^7$ samples from fixed low-block-count
starts, $a_k=(-1)^k(2.5\pm0.1)k^{-3}$ for $7\le k\le12$: the decay is
polynomial, not geometric, and the stationary-start average above
hides it because it averages the sign $(-1)^{B(\bar x)}$ over the
start (\S\ref{sec:s12num}).  (C2) is replaced by the polynomial
(C2$''$) of \S\ref{sec:s12left}, which suffices for a H\"older exponent
$(p-1)/(p+1)>4/9$ at $z=-1$ when $p>13/5$.}

\emph{Inter-flip parity $\E[(-1)^{\tau_k-\tau_{k-1}}]$ settles at
$-0.170\pm0.003$ for $k\ge3$ from both starts: bounded away from
$\pm1$, as (C2) needs, and strikingly stable.}

\emph{The post-flip chain equilibrates in about four flips.  From the
trap start, $\E[\mathrm{gap}]$ at $\tau_k$ runs
$0.1995,0.2375,0.2468,0.2484,0.2491,\dots$ and $\E[s]$ runs
$0.715,0.774,0.790,0.797,0.798$, matching the macroscopic start's
$0.2498$ and $0.799$ by $k=4$; the correlations
$\E[\mathrm{sgn}(g_0-\tfrac14)\,\mathrm{sgn}(g_k-\tfrac14)]$ decay
$0.317,0.090,0.022,0.012,\dots$, geometrically with ratio
$\approx0.3$.  This is direct evidence for the minorisation of
Proposition~\ref{prop:c1} with $n$ of order $4$.}

\emph{First-flip tails along the chain: $t^2\Pr[\tau_k-\tau_{k-1}>t]$
lies in $15$--$20$ on $16\le t\le128$ for $k=2,4,8$ from both starts
(and $\approx10^4$ for $k=1$ from the trap start, consistent with the
$g^{-2}$ prefactor), so (R1) holds uniformly \emph{along} the post-flip
chain and not merely at the start: TODO~11d(iii), answered in the
affirmative.}

\emph{Finally, \texttt{s11\_kflip} asserts
$B(\Gamma_t)-B(\Gamma_0)\equiv t\pmod2$ at every flip of every one of
the $1.6\times10^5$ trajectories --- the hypothesis of
Lemma~\ref{lem:conj}(ii), via Lemma~\ref{lem:parity} --- and it holds.}

%----------------------------------------------------------------------
\subsection{Session 12: the minorisation route is closed by inert dust;
(C2) is polynomial, not geometric}\label{sec:s12}

Session 12 was directed at (C1)+(C2) of \S\ref{sec:s11left}, through
TODO~12a (a Lyapunov drift for $\lambda$), 12b (the minorisation of
Proposition~\ref{prop:c1}) and 12c (an odd perturbation), with the
instruction to return one of (A1)--(A3) or a clean negative (B).  The
outcome is (B), in a form sharper than the one the brief anticipated.
Four findings, in logical order.

\begin{enumerate}[leftmargin=2.2em]
\item[(1)] \textbf{The environment carries inert dust, and this closes
every Doeblin-type route.}  A curve of mass $\varepsilon$ is touched by
a chord with probability $\le2\varepsilon$ per step, so it survives
\emph{unchanged} through the first $n$ flips with probability
$\to1$ as $\varepsilon\to0$ (Lemma~\ref{lem:inert}).  Consequently the
post-flip laws from the states $x_\varepsilon$ (a fixed macroscopic
marked pair, plus one dust curve of mass $\varepsilon$) are supported,
up to $o(1)$, on the pairwise disjoint sets ``some curve has mass
exactly $\varepsilon$''.  No common minorant exists: the hypothesis of
Proposition~\ref{prop:c1} fails on $\Omega_K$ for \emph{every} $n$
(Proposition~\ref{prop:nomino}), and so does uniform ergodicity of
$\kappa_1$ in $L^\infty(\Omega_K)$ and $V$-uniform ergodicity for
\emph{every} weight $V$ (Corollary~\ref{cor:noerg}).  Route (a) of
\S\ref{sec:s11left} and TODO~12b are therefore closed, not
deferred.  The proof uses only $\tau_n<\infty$ a.s.
\item[(2)] \textbf{TODO~12a fails on the full state space.}  From a
state whose non-dust material has total mass $m$ and whose remaining
mass sits in $N\gg m^{-3}$ equal pieces, one has
$\sum_jr_j^4\le17m^4$ deterministically for a long time, so
$\kappa_1^kV\ge\tfrac34(17m^4)^{-\theta}$ while $V\le m^{-4\theta}$:
no inequality $\kappa_1V\le\varrho V+b$ with $\varrho<1$ can hold
(Proposition~\ref{prop:nodrift}).  As instructed, the fallback is the
killed chain on $\Omega_K$, where $V$ is bounded anyway.
\item[(3)] \textbf{(C2) is true but polynomial.}  With $2\times10^7$
samples from the one-curve start,
$a_k(x)=\E_x[(-1)^{\tau_k}]=(-1)^k\,(2.5\pm0.1)\,k^{-3}$ for
$7\le k\le13$, and with $10^8$ samples the local exponent over
$8\le k\le20$ is $2.7\pm0.3$; a tail of the same form,
$k^3|a_k|\to2.0$--$2.4$, from the three-curve start $(0.3,0.3;\{0.4\})$,
and the same tail \emph{emerging} after a transient of $\approx8$
flips from a nine-curve start.  Session~11's ``$\rho\approx0.3$''
was the transient above a noise floor of $0.003$.  Since
$a_k=\pm\E_x[(-1)^{B(\bar\Gamma_{\tau_k})}]$ is the parity bias of the
block count at the $k$-th flip, the mechanism is the slow ($\log$)
growth of $B$: the block-count parity is a slow mode of $\kappa_1$ at
eigenvalue $+1$, not a mode near $-1$ (\S\ref{sec:s12num}).  The
polynomial decay is insensitive to inserting dust into the start
(identical to two standard errors) and is \emph{weaker} from starts
with $9$ or $17$ curves.
\item[(4)] \textbf{What the reduction actually needs is polynomial.}
At $z=-1$ Proposition~\ref{prop:condecay} needs only
$G\in C^{0,\gamma}$, $\gamma>4/9$, and a polynomial decay of the
twisted parity $\E_x[w^{\tau_k}(-1)^{E_{\tau_k}}]$ with exponent $p$,
uniformly for $w$ near $1$, gives $\gamma=(p-1)/(p+1)$ by a crude
bound (Proposition~\ref{prop:holder}), so the budget needs $p>13/5$;
if the parity memory factorises from the time weighting, as it does
numerically, the exponent improves to $\min(1,p-1)$ and the budget
needs only $p>13/9$.  The measured $p=2.7\pm0.3$ meets the second
comfortably and the first marginally.  (C2) is accordingly replaced by
(C2$''$) in \S\ref{sec:s12left}.  Any proof must be \emph{dust-blind},
i.e.\ work in a topology in which $x$ and $x_\varepsilon$ are close;
we record why the natural transport metric is not contracted by the
dynamics and why an adversarial environment cannot be allowed
(\S\ref{sec:s12left}).
\end{enumerate}

\subsubsection{Inert dust}\label{sec:s12dust}

Throughout, the chain is realised on the circle: the state is a
partition of the circle into curves (each a cyclically ordered union of
arcs, the cyclic order being that of the chord diagram), a step picks
two independent uniform points, and the move is determined by the
curves the points lie in and their positions along those curves.

\begin{lemma}[inert dust coupling]\label{lem:inert}
Let $y$ be a state, $D$ a sub-arc of the circle of mass $\varepsilon$
lying inside one curve of $y$ and containing neither mark, and let
$y_\varepsilon$ be the state obtained from $y$ by declaring $D$ a
separate (unmarked) curve.  Drive the chains from $y$ and from
$y_\varepsilon$ by the same uniform points.  Then, up to the first step
at which a point falls in $D$, the state of the $y_\varepsilon$-chain
is the state of the $y$-chain with $D$ declared a separate curve; in
particular the two chains have the same flip times, and the
$y_\varepsilon$-chain contains a curve of mass exactly $\varepsilon$.
The first point in $D$ falls at a step of probability at most
$2\varepsilon$, so for every $T$,
\[
\Pr\bigl[\text{the coupling holds through time }T\bigr]\ge1-2\varepsilon T .
\]
\end{lemma}

\begin{proof}
$D$ is a contiguous segment of the cyclic order of the curve $C\supset D$
containing it.  A split of $C$ at two points outside $D$ cuts the cycle
into two arcs, one of which contains $D$; in the $y_\varepsilon$-chain
the same cuts split $C\setminus D$ into the same two arcs with $D$
excised.  A merge of $C$ with another curve at points outside $D$
splices the two cycles at those points; in the $y_\varepsilon$-chain the
same splice is applied to $C\setminus D$.  Moves not involving $C$ are
identical.  A flip is a split or merge decided by which marked material
the two points lie in; the marked material of the two chains differs
only by $D$, which contains no point.  Induction on the steps.
\end{proof}

\begin{proposition}[no state-uniform minorisation]\label{prop:nomino}
Let $K\ge4$ and assume $\tau_n<\infty$ a.s.\ from the state
$x_0=(\tfrac14,\tfrac14;\{\tfrac12\})$ (apart sector).  Then there are
no $n\ge1$, $\delta>0$ and probability measure $\nu$ on $\Omega_K$ with
\[
\Pr^{\Omega_K}_{\bar x}\bigl[\bar\Gamma_{\tau_n}\in\cdot\,\bigr]\ge\delta\,\nu(\cdot)
\qquad\text{for every }\bar x\in\Omega_K\cap S'_K ,
\]
for the chain killed on $\Omega_K$ or for the unkilled chain.  A
fortiori the hypothesis of Proposition~\ref{prop:c1} is never
satisfied.
\end{proposition}

\begin{proof}
Suppose such $n,\delta,\nu$ exist.  For $i\ge3$ let
$\varepsilon_i=2^{-i}$, let $x_i$ be $x_0$ with a sub-arc $D_i$ of mass
$\varepsilon_i$ of its environment curve declared a separate curve
(so $x_i=(\tfrac14,\tfrac14;\{\tfrac12-\varepsilon_i,\varepsilon_i\})$,
$B(x_i)=4$, $x_i\in\Omega_K\cap S'_K$), and let
$S_i=\{\bar x\in\Omega_K:\ \text{some curve of }\bar x\text{ has mass
exactly }\varepsilon_i\}$.  Apply the minorisation at $x_i$ to the set
$\Omega_K\setminus S_i$.  By Lemma~\ref{lem:inert} with $T=T_i$, on the
event that the coupling holds through $T_i$ and $\tau_n\le T_i$ for the
unkilled $x_0$-chain, the $x_i$-chain performs its $n$-th flip at the
same time and, if it is still alive, its state then contains $D_i$
intact; if it has been killed, its state is the cemetery, which is not
in $\Omega_K\setminus S_i$.  Hence
\[
\delta\,\nu(\Omega_K\setminus S_i)\;\le\;2\varepsilon_iT_i+\Pr_{x_0}[\tau_n>T_i]
\qquad\text{for every }T_i ,
\]
and $T_i=\varepsilon_i^{-1/2}$ gives $\nu(S_i)\to1$ as $i\to\infty$.
But a configuration in $\Omega_K$ has at most $K$ curves, hence at most
$K$ distinct masses, hence lies in at most $K$ of the sets $S_i$; so
$\sum_i\nu(S_i)=\int\#\{i:\bar x\in S_i\}\,\nu(d\bar x)\le K<\infty$,
contradicting $\nu(S_i)\to1$.
\end{proof}

\emph{Status: PROVED.  Note what the proof does not use: no moment of
$\tau$, no property of the environment beyond the trivial one that a
curve of mass $\varepsilon$ is rarely hit, and no bound on $K$ beyond
$K\ge4$.  It also shows that the difficulty is not the one anticipated
in TODO~12b (``$K^{-2K}\to K^{-O(1)}$''): the obstruction is not the
\emph{cost} of driving the environment to a reference configuration but
the fact that no finite sequence of steps can remove a curve that is
never hit.  The same applies with the $(n,n+1)$ pair of
Proposition~\ref{prop:c1}, to any weighted space, and to any
``small set'' containing a state with a dust curve --- which every set
of positive $\bar\pi$-measure does, since dust is created at every
scale at rate $\asymp2\eta\,d\eta$ per step.}

\begin{corollary}[no uniform or $V$-uniform ergodicity]\label{cor:noerg}
Let $K\ge4$ and let $\kappa_1$ be the post-flip kernel of the unkilled
chain.
\begin{enumerate}[leftmargin=2em]
\item[(i)] There are no probability measure $\bar\pi$, $C<\infty$ and
$\rho<1$ with $\sup_{\bar x\in\Omega_K}|\kappa_1^kF(\bar x)-\bar\pi F|
\le C\rho^k\|F\|_\infty$ for all bounded $F$ and all $k$.
\item[(ii)] Let $\mathcal Y=\{(a',\tfrac12-a';\{\tfrac12\}):
a'\in[\tfrac18,\tfrac38]\}$ and assume
$(\mathrm M_p)$: $\sup_{y\in\mathcal Y}\E_y[\tau_k]\le Ck^p$ for some
$C,p<\infty$.  Then there is no $V\ge1$ for which $\kappa_1$ is
$V$-uniformly ergodic: no $\bar\pi$, $C'$, $\rho<1$ with
$|\kappa_1^kF(\bar x)-\bar\pi F|\le C'\rho^kV(\bar x)\|F\|_V$ for all
$F$ with $\|F\|_V=\sup|F|/V<\infty$ and all $\bar x$ with $B(\bar x)\le4$.
\end{enumerate}
\end{corollary}

\begin{proof}
(i) With $F_i=\mathbf 1_{S_i}$ (notation of the previous proof) and
$x_i$ as there, Lemma~\ref{lem:inert} gives
$\kappa_1^kF_i(x_i)\ge1-2\varepsilon_iT-\Pr_{x_0}[\tau_k>T]$, while
$\kappa_1^kF_i(x_0)=0$ because from $x_0$ every curve mass at every time
is either a sum of some of $\tfrac14,\tfrac14,\tfrac12$ or has an
atomless law, so no curve ever has mass exactly $\varepsilon_i\le\tfrac18$ a.s.  Let
$i\to\infty$: the oscillation of $\kappa_1^kF_i$ over $\Omega_K$ tends
to $1$ for each fixed $k$, contradicting $2C\rho^k<1$.

(ii) Suppose such $V,\bar\pi,C',\rho$ exist.  First, $\Pi=\mathbf 1\otimes\bar\pi$
is bounded on the $V$-space, so $\bar\pi V<\infty$ and
$\kappa_1V\le(C'\rho+\bar\pi V)\,V$; in particular $\kappa_1V(x_0)<\infty$.
Second, for $y\in\mathcal Y$ and $y_\varepsilon$ obtained by carving a
sub-arc of mass $\varepsilon$ from its environment curve, the argument
of (i) at $y$ gives $\kappa_1^kF_i(y_\varepsilon)\ge1-2\varepsilon T-\Pr_y[\tau_k>T]$
and $\bar\pi F_i=\lim_k\kappa_1^kF_i(x_0)=0$, so with $T=T_k:=4Ck^p$
and $\varepsilon\le\varepsilon_k:=1/(8T_k)$ (Markov and $(\mathrm M_p)$),
\[
V(y_\varepsilon)\;\ge\;\frac{1-\tfrac14-\tfrac14}{C'\rho^k}\;=\;\frac1{2C'\rho^k}
\qquad(y\in\mathcal Y,\ \varepsilon\le\varepsilon_k).
\]
Third, from $x_0$ the first step carves a piece of mass in
$d\varepsilon$ from the environment curve with probability
$\ge\tfrac12\,d\varepsilon$ (pick that curve twice, $\tfrac14$;
$|t_1-t_2|/\tfrac12$ has density $\ge2$ near $0$) and the second step
is a flip landing in $\mathcal Y$ with probability $\ge\tfrac18\cdot\tfrac34$;
the landing state is then some $y_\varepsilon$, $y\in\mathcal Y$.  So
\[
\kappa_1V(x_0)\;\ge\;\frac{3}{64}\sum_{k\ge1}\frac{\varepsilon_k-\varepsilon_{k+1}}{2C'\rho^k}
\;\ge\;c\sum_k\frac{k^{-p-1}}{\rho^{k}}\;=\;\infty ,
\]
contradicting the first step.
\end{proof}

\emph{Status: PROVED ((ii) under $(\mathrm M_p)$, which (R1) along the
chain implies with $p=1$ and which is measured: $\E[\tau_k]\approx6k$).
For the chain killed on $\Omega_K$ the same argument runs up to
$k\le c\,3^{K/(3p)}$, so any $V$-uniform ergodicity of the killed chain
has $(1-\rho)^{-1}\log(C'V(x_0)\bar\pi V)\ge c\,3^{K/(3p)}$: exponential
in $K$, i.e.\ fatal in the budget of the session-11 handoff.  This
settles, negatively, the sentence of \S\ref{sec:markovrenewal} that
asked for a weighted space ``in which the post-flip chain is
$V$-geometrically ergodic in the base coordinates'': there is none.
It does \emph{not} touch (C1) or (C2) themselves, whose null vectors
would have to sit near $-1$, whereas the dust modes sit near $+1$.}

\subsubsection{The drift for \texorpdfstring{$\lambda$}{lambda} fails
on the full state space (TODO~12a)}\label{sec:s12drift}

\begin{proposition}[no Lyapunov drift for $(\sum r_j^4)^{-\theta}$]
\label{prop:nodrift}
Let $V=(\sum_jr_j^4)^{-\theta}$ with $0<\theta\le\tfrac12$, on the full
(unkilled) state space, and assume $\tau_k<\infty$ a.s.\ from every
state.  There are no $\varrho<1$ and $b<\infty$ with
$\kappa_1V\le\varrho V+b$.
\end{proposition}

\begin{proof}
Suppose the inequality holds.  Iterating,
$\kappa_1^kV\le\varrho^kV+b/(1-\varrho)$ for all $k$.  Fix $k$ with
$\varrho^k\le\tfrac38\,17^{-\theta}$ and $m\in(0,\tfrac14)$.

\emph{The isolated process.}  Let $\bar x$ be the apart-sector state
with $a=m$, $b=m/10$ and one environment curve of mass $1-1.1m$, and
let $\mathrm P^{\rm iso}$ be the law of the chain from $\bar x$ modified
so that a chord with one point in the original environment curve (or in
any curve descended from it by splits and merges among such curves)
and one point elsewhere does nothing.  Under $\mathrm P^{\rm iso}$ the
``non-dust'' material of total mass $1.1m$ evolves as a split--merge
chain in its own right, slowed by the factor $(1.1m)^2$, with the
marked pair inside it; its flip times are a.s.\ finite, so there is
$T=T(m,k)$ with $\mathrm P^{\rm iso}[\tau_k>T]\le\tfrac18$.

\emph{The states $x_{m,N}$.}  Let $x_{m,N}$ be $\bar x$ with the
environment curve cut into $N$ equal pieces,
$N\ge16(T+1)^3\max(m^{-1},m^{-4/3})$.  Drive the chain from $x_{m,N}$
and the isolated process from $\bar x$ by the same points (the dust
pieces of $x_{m,N}$ are sub-arcs of the environment curve of $\bar x$;
dust--dust chords split or merge dust material in both and never
touch the non-dust material).  The two processes make the same non-dust moves
except at chords joining dust to non-dust, of which there are at most
$T$ up to time $T$; each such chord appends to one non-dust curve of the
$x_{m,N}$-chain a dust curve, i.e.\ a union of at most $T+1$ original pieces or
fragments of them, of mass $\le(T+1)/N$.  As in Lemma~\ref{lem:inert} (with the
appended material in the role of $D$, now inside a curve of one chain
and absent from the other), the non-dust configurations of the two
processes then agree up to the appended material, and their flip times
agree, until a point falls in appended material, which has probability
$\le2T(T+1)/N$ per step.  Hence
$\Pr_{x_{m,N}}[\tau_k>T]\le\tfrac18+2T^2(T+1)/N\le\tfrac14$.

\emph{The deterministic bound.}  Along the $x_{m,N}$-chain up to time
$T$, the non-dust mass is at most $1.1m+T(T+1)/N\le2m$ and every dust
curve has mass $\le(T+1)/N\le m^{4/3}$, so
\[
\sum_jr_j(\bar\Gamma_t)^4\;\le\;(2m)^4+\Bigl(\max_{\text{dust}}r_j\Bigr)^3
\sum_{\text{dust}}r_j\;\le\;16m^4+m^4=17m^4\qquad(t\le T),
\]
whence $\kappa_1^kV(x_{m,N})\ge\Pr_{x_{m,N}}[\tau_k\le T]\,(17m^4)^{-\theta}
\ge\tfrac34\,17^{-\theta}m^{-4\theta}$, while $V(x_{m,N})\le m^{-4\theta}$
because the $u$-curve alone contributes $m^4$.  The drift bound gives
$\tfrac34\,17^{-\theta}m^{-4\theta}\le\tfrac38\,17^{-\theta}m^{-4\theta}+b/(1-\varrho)$,
i.e.\ $m^{-4\theta}\le\tfrac83\,17^{\theta}b/(1-\varrho)$ for every
$m\in(0,\tfrac14)$: false.
\end{proof}

\emph{Status: PROVED.  This is the fallback the brief provided for: the
obstruction is the family ``all non-dust mass $\asymp m$, the rest in
ultra-fine mass-carrying dust'', which is of vanishing probability
under the chain from one curve (Corollary~\ref{cor:horizon}) but is a
legitimate state, and on it $\lambda$ cannot recover within one
inter-flip period because there is no macroscopic partner to merge
with.  On $\Omega_K$, $V\le(2K)^{4\theta}$ is bounded and the drift is
vacuous.  The killing at $B=K$ in Theorem~\ref{thm:idleloop} therefore
stays, and it is harmless: it costs $(9d+1)^{-\gamma}$ over the horizon.
The second purpose of TODO~12a --- supplying the base half of the
minorisation --- is void by Proposition~\ref{prop:nomino}.  A quick
numerical check of the mechanism (\texttt{s12\_driftfail}): from
$(0.3,0.03;\ 3\times10^4\text{ equal pieces})$, $\E[(\sum r^4(\bar\Gamma_\tau)/\sum r^4(x))^{-\theta}]=1.19$
$(\theta=0.1)$, $1.71$ $(\theta=0.3)$, against $0.91$, $0.77$ with a
single environment curve of the same mass.}

\subsubsection{TODO~12c: the stated criterion does not give (C2)}
\label{sec:s12odd}

The session-11 handoff proposed, as route 12c, a state-uniform lower
bound $c(K)$ on the probability that two consecutive inter-flip times
have opposite parity, ``which gives (C2) by a standard two-block
argument''.  It does not: if the inter-flip times were the
deterministic sequence $1,2,1,2,\dots$ then consecutive parities are
always opposite ($c=1$) while $(-1)^{\tau_k}$ is periodic with period
$4$ and does not decay.  What the two-block argument needs is the
\emph{operator} statement $\|\kappa_{-1}^{\,2}\|_\infty<1$, and that is
false exactly: $\kappa_{-1}(-1)^E=-(-1)^E$ (Proposition~\ref{prop:r2false}),
so $\|\kappa_{-1}^{\,m}\|_\infty=1$ for every $m$.  Any odd perturbation
that is exact returns a state with a different block-count parity,
which is visible; the only odd perturbations that are invisible are
dust insertions, whose weight at scale $\eta$ is $\asymp2\eta$ per step
and whose lifetime is $\asymp1/(2\eta)$, and one checks in one line
that they give a contraction $1-O(\eta)$ per flip and nothing
geometric.  Route 12c is closed as stated; the honest form of (C2) is
in \S\ref{sec:s12num} and \S\ref{sec:s12left}.

\subsubsection{The decay of the parity along the post-flip chain is
polynomial}\label{sec:s12num}

By Lemma~\ref{lem:parity}, $(-1)^{\tau_k}=(-1)^{B(\bar\Gamma_{\tau_k})-B(\bar x)}$
exactly, so
\begin{equation}\label{eq:akB}
a_k(\bar x):=\E_{\bar x}[(-1)^{\tau_k}]=(-1)^{B(\bar x)}\sum_{b\ge1}(-1)^b\,
\Pr_{\bar x}\bigl[B(\bar\Gamma_{\tau_k})=b\bigr],
\end{equation}
the parity bias of the block count at the $k$-th flip, and
$\kappa_1^{\,k}(-1)^E=(-1)^{E}(-1)^ka_k$, so (C2) is the statement that
\eqref{eq:akB} decays geometrically, uniformly in $\bar x$.

\emph{Measurement (\texttt{s12\_akdecay}, \texttt{s12\_akB}; fixed
deterministic starts, $2\times10^7$ independent trajectories each,
standard error $2.2\times10^{-4}$, no censoring).}  From the one-curve
start (arcs $(\tfrac12,\tfrac12)$):
\[
a_k=-0.4003,\ +0.1530,\ -0.0674,\ +0.0335,\ -0.0186,\ +0.0111,\ -0.0073,\
+0.0048,\ -0.0034,\ +0.0022,\ -0.0019,\ +0.0015\quad(k=1,\dots,12),
\]
so that $k^3|a_k|=2.51,2.48,2.49,2.20,2.54,2.54$ for $k=7,\dots,12$:
\[
a_k(\text{one curve})\;=\;(-1)^k\,(2.5\pm0.1)\,k^{-3}\qquad(7\le k\le12),
\]
with the successive ratios $|a_{k+1}/a_k|$ rising through
$0.38,0.44,0.50,0.56,0.60,0.66,0.66,0.70$ --- no geometric rate.  A
$10^8$-sample run to $k=20$ (standard error $10^{-4}$) gives
$|a_k|=0.00495,0.00349,0.00271,0.00195,0.00142,0.00115,0.00107,0.00095,
0.00091,0.00079,0.00055,0.00045,0.00035$ for $k=8,\dots,20$, all with
sign $(-1)^k$; $k^3|a_k|$ stays in $2.4$--$2.7$ up to $k=13$ and
fluctuates in $2.8$--$3.9$ (two to three standard errors) for
$k=14$--$20$; the least-squares exponent over $8\le k\le20$ is
$2.7\pm0.3$.  From
the three-curve apart start $(0.3,0.3;\{0.4\})$: $a_1=-0.004$,
then $-0.0606,+0.0385,-0.0219,+0.0132,-0.0086,+0.0060,-0.0041,+0.0031,
-0.0020,+0.0018,-0.0014$ ($k=2,\dots,12$), with $k^3|a_k|\to2.0$--$2.4$:
the same tail, sign $(-1)^{k+B_0+1}$.  Inserting a dust curve of mass
$10^{-4}$ into either start (mass taken from an environment curve or
from an arc) reproduces every $a_k$ to within two standard errors, as
Lemma~\ref{lem:inert} predicts.  From a $9$-curve start
$(0.3,0.3;\ 7\times0.057)$, $10^8$ samples: $a_k$ passes through zero
at $k=3$ and then \emph{grows} in absolute value to $0.00145$ at
$k=7$--$8$ before decaying, with the sign $(-1)^{k+1}$ at every
$k=4,\dots,20$ (seventeen consecutive alternations); $k^3|a_k|$ rises
from $0.5$ at $k=7$ to $2$--$3$ at $k=15$--$20$.  So the tail is
universal --- $C k^{-3}$ with $C\approx2.5$--$3$ --- and is reached
after a start-dependent transient, of about eight flips when the
environment starts with nine curves.  From a $17$-curve start
($5\times10^6$ samples) $|a_k|\le7\times10^{-4}$ for $3\le k\le16$,
the transient not yet over.  From the
burnt-in ``stationary'' start of session~11 (gap window
$[\tfrac18,\tfrac38)$, $5\times10^6$ samples), $a_k(\mu)$ is at the noise
floor $4.5\times10^{-4}$ for $k\ge5$ --- but that average carries the
factor $(-1)^{B(\bar x)}$ of \eqref{eq:akB} inside the mixture, which
cancels the tail; the sign-blind $\E_\mu[a_k(\bar x)^2]$ is at
$(4\pm5)\times10^{-4}$ for $k=5..8$, consistent with the fixed-start
values.  The decomposition of \eqref{eq:akB} by $b$ from the one-curve
start (\texttt{s12\_akB}) shows the mechanism: $\Pr[B_{\tau_k}=1]$
decays like $c/k$ with $c\approx0.5$ (the chain merges back to one curve with
probability $\asymp1/t$, and there $\tau_k$ is even), the law of
$B_{\tau_k}$ spreads only like $\log k$, and the alternating sum of a
lattice law that spreads logarithmically decays polynomially.

\emph{Status: MEASURED.  Conclusions.  (i) (C2) in the geometric form
of \S\ref{sec:s11left} is false on the full state space, and on
$\Omega_K$ it holds only with $1-\rho(K)\le e^{-cK}$ (the rate the
killing provides), which is outside the budget; the session-11 value
$\rho\approx0.3$ described the transient $k\le4$.  (ii) The polynomial
decay is a $+1$ phenomenon, not a $-1$ phenomenon: the generating
function $\sum_ka_kw^k$ is singular at $w=-1$, i.e.\ at the eigenvalue
$-1$ of $\kappa_{-1}$, which is the eigenvalue $+1$ of $\kappa_1$ (the
trivial one), approached sub-geometrically by the slow variable $B$.
Nothing in the measurement bears on (C1).  (iii) The decay is
sign-alternating in $k$ from every start tested, with the sign
$(-1)^{k+B_0+1}$; this is the statement that $B_{\tau_k}\equiv k+1$ is
slightly favoured, i.e.\ that the block count at the $k$-th flip
remembers that it started odd --- a memory that $\log$-growth of $B$
erases only polynomially.}

\subsubsection{What is left, restated}\label{sec:s12left}

\begin{proposition}[polynomial parity decay suffices at $z=-1$]
\label{prop:holder}
Fix $K$ and $S'_K$.  Suppose
\begin{itemize}[leftmargin=2em]
\item[(C2$''$)] there are $C,p$ and $\delta_0>0$ with
$\sup_{\bar x\in S'_K}\bigl|\E_{\bar x}[w^{\tau_k}(-1)^{E_{\tau_k}}]\bigr|\le Ck^{-p}$
for all $k\ge1$ and all $|w|\le1$ with $|1-w|\le\delta_0$;
\item[(M$_1$)] $\sup_{\bar x\in S'_K}\E_{\bar x}[\tau_k]\le C_1k$ for all $k$.
\end{itemize}
Then $G(z)=(I+\kappa_z)^{-1}(I-\kappa_z)\mathbf 1$, evaluated on $S'_K$,
is H\"older continuous of exponent $(p-1)/(p+1)$ on
$\{|z|\le1,\ |1+z|\le\delta_0\}$, with constant
$\le C_2\,(C+C_1)$.  In particular $p>13/5$ gives an exponent $>4/9$.
\end{proposition}

\begin{proof}
By Proposition~\ref{prop:repair}, for $z=-w$,
$G(-w)=\bar M(I-\kappa_w)^{-1}(I+\kappa_w)\bar M\mathbf 1
=\mathbf 1+2\bar M\sum_{k\ge1}g_k(w)$ with
$g_k(w)(\bar x)=\E_{\bar x}[w^{\tau_k}(-1)^{E_{\tau_k}}]$
(Neumann series, $\|\kappa_w\|\le|w|<1$; the boundary values are limits).
For $|w|,|w'|\le1$, $|w^n-w'^n|\le n|w-w'|$, so
$|g_k(w)-g_k(w')|\le\min\bigl(2Ck^{-p},\ C_1k|w-w'|\bigr)$ by (C2$''$)
and (M$_1$).  Summing with the cut at $k_*=|w-w'|^{-1/(p+1)}$:
$\sum_{k\le k_*}C_1k|w-w'|+\sum_{k>k_*}2Ck^{-p}
\le C_2(C+C_1)\,|w-w'|^{(p-1)/(p+1)}$.
\end{proof}

\emph{Status: PROVED as a reduction; the hypotheses are OPEN.  Two
comments.  (a) The exponent $(p-1)/(p+1)$ is crude: it bounds
$|g_k(w)-g_k(w')|$ by the trivial Lipschitz bound and never uses that
$g_k(1)=(-1)^{E}(-1)^ka_k$ has one sign for all $k$.  If the parity
memory factorises from the time weighting,
\[
\text{(C2$'''$)}\qquad
\E_{\bar x}[w^{\tau_k}(-1)^{\tau_k}]=\E_{\bar x}[w^{\tau_k}]\,a_k(\bar x)\,(1+O(|1-w|k)),
\]
then $\sum_kg_k(w)\approx(-1)^E\sum_kCk^{-p}\E[w^{\tau_k}]$, and with
$\E[w^{\tau_k}]\approx w^{6k}$ this is $C\,\mathrm{Li}_p(w^6)$ up to
smoother terms, whose modulus of continuity at $w=1$ is
$|1-w|^{\min(1,p-1)}$ (up to a logarithm at integer $p$): the budget
then needs only $p>13/9$.  (C2$'''$) was measured (\texttt{s12\_twist}, one-curve start,
$3\times10^7$ samples, $w=0.99,0.95,0.90$ real): the exact
factorisation fails --- the twisted parity $\E[w^{\tau_k}(-1)^{\tau_k}]$
exceeds $\E[w^{\tau_k}]\,a_k$ by a factor $1.2$ at $|1-w|k=0.08$ and
$2$--$4$ at $|1-w|k\approx0.4$--$0.8$ (short paths carry more parity
memory) --- but the uniform bound (C2$''$) holds with the same
constant, $|\E[w^{\tau_k}(-1)^{\tau_k}]|\le|a_k|$ at every $k\le16$ and
every $w$ tested.  More usefully, the sum itself was measured: with
$\Sigma(w):=\sum_k(-1)^kE_x[w^{\tau_k}(-1)^{\tau_k}]$ (all terms of one
sign), $\Sigma(1)-\Sigma(w)=0.026,\ 0.085,\ 0.143$ at
$|1-w|=0.01,0.05,0.10$ (tail beyond $k=16$ estimated from $2.5k^{-3}$),
i.e.\ a modulus of continuity $\asymp|1-w|^{0.75}$ over that range ---
above the crude $(p-1)/(p+1)\approx0.5$, below the $\min(1,p-1)$ of the
factorised heuristic, and far above the $4/9$ the budget needs.  The
same on the unit circle near $w=1$ is the first numerical task of
session~13.  With the measured $p=2.7\pm0.3$ the crude route is marginal
($p>13/5$ required); the measured modulus is not.  (b) (C2$''$) at $w=1$ is exactly the polynomial
(C2$'$): $\sup_{S'_K}|a_k|\le Ck^{-p}$.  Away from $w=1$ on the circle
Theorem~\ref{thm:idleloop} gives a geometric rate; (C2$''$) asks that
the polynomial rate survive a small complex twist near $w=1$, which is
a damping of long paths and should not spoil a parity cancellation that
is produced by the block count.  This is measurable and is the first
numerical task for session~13.}

\medskip
\noindent The open conditions are now:

\medskip
\noindent\textbf{(C1)}\quad$\|(I+\kappa_1)^{-1}\|_{L^\infty(\Omega_K)}\le C(K)$ ---
equivalently $\|(I-\kappa_{-1})^{-1}\|\le C(K)$, the summability of the
signed renewal series $\sum_k\kappa_{-1}^{\,k}G$ for every bounded $G$;
(C2$'$) is the case $G=\mathbf 1$.  Unchanged, not contradicted, and no
longer reachable by any Doeblin/ergodicity argument
(Proposition~\ref{prop:nomino}, Corollary~\ref{cor:noerg}).\\
\noindent\textbf{(C2$''$)}\quad as in Proposition~\ref{prop:holder},
with $p>13/5$ and $C=C(K)$ polynomial; measured $p\approx3$ at $w=1$.

\medskip
\emph{What a proof must look like.  (i) It must be dust-blind: by
Lemma~\ref{lem:inert}, $a_k(y_\varepsilon)=a_k(y)+O(\varepsilon\,\E_y\tau_k)$
while $(-1)^{E}$ changes sign, so any argument that treats $(-1)^E$ as
a function to be contracted in a sup-type norm on configurations is
doomed, and any argument that works must use a topology in which
$y_\varepsilon\to y$.  (ii) The natural candidate --- the transport
distance between mass configurations, under the coupling of
Lemma~\ref{lem:inert} --- is not contracted: when the dust is finally
hit, the two coupled configurations differ by a macroscopic merge
(``$r+r_0$'' against ``$r+\varepsilon$, $r_0-\varepsilon$''), so a
discrepancy of size $\varepsilon$ is amplified to size $O(1)$ rather
than resolved.  A contracting metric would have to compare coarse
\emph{laws} rather than configurations.  (iii) The environment's own
randomness is needed: an adversary allowed to choose the masses gained
by the marked pair at the merge steps (whose timing it cannot control)
can bias the parity of each inter-flip time toward a target by a fixed
amount, since a gain changes the flip hazard $2p_1p_2$ by a bounded
amount at a time of its choosing among the merge steps; so no
``uniformly over environments'' statement can be true, and the proof
must use that the true gains are size-biased picks from a partition
whose count parity is invisible to them.  (iv) Given
\eqref{eq:akB}, the cleanest target is a statement about the block
count alone: the parity of $B$ at the $k$-th flip equidistributes
polynomially.  Since the total block count from one curve has the
exact law of Theorem~\ref{thm:blockcount} at fixed times, and the flip
times are a renewal-like sampling of the time axis, this may be
provable from the exact formulas by an averaging argument, at least
for the annealed environment --- which is the form the downstream
reduction (Problem~\ref{prob:survival}, whose start is
$\mu_j$) actually consumes.}

%----------------------------------------------------------------------
\subsection{Session 13: audit of the reduction chain}\label{sec:s13}

Session 13 opened with the audit recommended in the session-11 health
note: nothing in the chain Theorem~\ref{thm:marked} $\to$
Problem~\ref{prob:annealed} $\to$ Theorem~\ref{thm:factor} $\to$
Problem~\ref{prob:mixing} $\to$ Problem~\ref{prob:survival} $\to$
(R1)+(C1)+(C2$''$) had been re-read since it was laid, and two carried
conditions had just been found false in consecutive sessions.  The
audit produced three corrections to the \emph{form} of the bottom of
the chain, none of which changes what is proved above it:
\begin{enumerate}[leftmargin=2.2em]
\item[(1)] the quantity the programme consumes is an $L^2$ average of
the quenched parity $\Phi_m$ over the start, not a supremum over
states (Proposition~\ref{prop:S2red}); the state-uniform target (S)
of \S\ref{sec:s12left} was introduced by the strong-Markov step of
the session-10 reduction, not demanded by anything above it, and the
inert-dust obstruction of \S\ref{sec:s12dust} does not touch the
$L^2$ form (Remark~\ref{rem:dustblind});
\item[(2)] the exponent budget carried since session 10
($\gamma>4/9$ at $z=-1$, hence $p>13/5$ in
Proposition~\ref{prop:holder}) is too demanding by a factor of
about five: the correct budget is $\gamma>1/11$, hence $p>6/5$
(Corollary~\ref{cor:budget}).  The measured $p\approx2.7$ and
modulus $0.75$ are far inside it, not marginal;
\item[(3)] $\tau_k<\infty$ a.s.\ from every state with $ab>0$, taken
for granted since session 10, has a four-line proof
(Lemma~\ref{lem:flipsrecur}).
\end{enumerate}
The occupation-time step and (M$_1$) are stated exactly and left open
(\S\ref{sec:s13occ}, \S\ref{sec:s13m1}).

\subsubsection{The reduction consumes an \texorpdfstring{$L^2$}{L2} average, not a supremum}
\label{sec:s13ltwo}

Recall (Corollary~\ref{cor:quotient}) that
$\psi(j,d)=\E_{\mu_j}[\Phi_d(\Gamma_0)^2]$ with
$\Phi_m(\Gamma)=\E_\Gamma[(-1)^{N_m}]$ the parity of the number of
flips in $m$ steps of the quotient chain, and that
Theorem~\ref{thm:factor} needs $\psi(j,d)\le Kd^{-2c}$ for
$j\le8d$ only.  Fix $K\ge3$ and
$S'_K=\{a,b\ge1/(16K)\}$ as in \S\ref{sec:markovrenewal}.

\begin{proposition}[$L^2$ form of the reduction to $S'_K$]\label{prop:S2red}
Let $d\ge1$, $j\le8d$, $0<\varepsilon\le1/(2K)$, and let
$T^*=T^*(\varepsilon,K,d)<d/2$ be any deterministic time.  Let
$T'=\inf\{t\ge0:\Gamma_t\in S'_K\}$ and let
$\nu=\nu_{j,\varepsilon,K,T^*}$ be the joint law of
$(\Gamma_{T'},T')$ under $\Pr_{\mu_j}(\,\cdot\,;\,g_0\ge\varepsilon,\
T'\le T^*)$, a sub-probability measure on $S'_K\times\{0,\dots,T^*\}$;
write $\nu_t=\nu(\cdot\times\{t\})$ and, when $\nu_t\ne0$, let
$\bar\nu_t=\nu_t/\nu_t(S'_K)$ be the \emph{time-$t$ entrance law}, a
probability measure on $S'_K$.  Then
\[
\E_{\mu_j}\bigl[\Phi_d(\Gamma_0)^2;\ g_0\ge\varepsilon\bigr]
\;\le\;2\int\Phi_{d-t}(x)^2\,\nu(dx,dt)
\;+\;2\,\Pr_{\mu_j}\bigl[T'>T^*;\ g_0\ge\varepsilon\bigr]
\;\le\;2\sup_{t\le T^*}\E_{\bar\nu_t}\bigl[\Phi_{d-t}^2\bigr]
\;+\;2\,\Pr_{\mu_j}\bigl[T'>T^*;\ g_0\ge\varepsilon\bigr].
\]
\end{proposition}

\begin{proof}
$T'$ is a stopping time and $T'\wedge T^*\le T^*<d$.  By the strong
Markov property at $T'\wedge T^*$ and the additivity of the flip
count, $\E[(-1)^{N_d-N_{T'}}\mid\mathcal F_{T'}]=\Phi_{d-T'}(\Gamma_{T'})$
on $\{T'\le T^*\}$, so
\[
\Phi_d(\Gamma_0)
=\E_{\Gamma_0}\bigl[(-1)^{N_{T'}}\Phi_{d-T'}(\Gamma_{T'});\ T'\le T^*\bigr]
+\E_{\Gamma_0}\bigl[(-1)^{N_d};\ T'>T^*\bigr]
=:A(\Gamma_0)+B(\Gamma_0).
\]
Then $|B|\le\Pr_{\Gamma_0}[T'>T^*]\le1$, so $B^2\le|B|$, and by
Cauchy--Schwarz
$A^2\le\E_{\Gamma_0}[\Phi_{d-T'}(\Gamma_{T'})^2;\,T'\le T^*]\cdot
\Pr_{\Gamma_0}[T'\le T^*]\le\E_{\Gamma_0}[\Phi_{d-T'}(\Gamma_{T'})^2;\,T'\le T^*]$.
Use $(A+B)^2\le2A^2+2B^2$ and integrate over $\Gamma_0\sim\mu_j$ on
$\{g_0\ge\varepsilon\}$; the first integral is, by definition of
$\nu$, $\int\Phi_{d-t}(x)^2\nu(dx,dt)=\sum_{t\le T^*}\nu_t(S'_K)\,
\E_{\bar\nu_t}[\Phi_{d-t}^2]$, and $\sum_t\nu_t(S'_K)\le1$.  (The
supremum over $t$ must stay inside the sum in this way: a bound in
terms of the time-marginal $\sum_t\nu_t$ alone would cost a factor
$T^*$.)
\end{proof}

\emph{Status: PROVED.  Compare the session-10 reduction paragraph at
the end of \S\ref{sec:markovrenewal}, which at the same point bounds
$|\Phi_d(\Gamma_0)|\le\E|\Phi_{d-T'}(\Gamma_{T'})|+\Pr[T'>T^*]$ and
then takes the supremum of $|\Phi_m|$ over $S'_K\cap H$.  The
supremum is where the state-uniform target (S) entered the
programme; nothing above this point asks for it.  What the programme
needs on $S'_K$ is therefore:}

\medskip
\noindent\textbf{(S$_2$)}\quad
$\displaystyle\E_{\bar\nu_t}\bigl[\Phi_{m}^2\bigr]\le C(K)\,m^{-2\gamma}$
\emph{for every $t\le T^*$ and $m=d-t\in[d/2,d]$, for the time-$t$
entrance laws $\bar\nu_t=\bar\nu_{j,\varepsilon,K,T^*;t}$ of $S'_K$,
$j\le8d$, with $K=K_{\gamma'}(9d)=O(\log d)$ and $C(K)$ polynomial
in $K$,}

\medskip
\noindent together with the entrance estimate
$\Pr_{\mu_j}[T'>T^*;\,g_0\ge\varepsilon]\le O(d^{-\gamma'})+O(\varepsilon K\log d)$
for $T^*=O(\varepsilon^{-1}K\log d)$, which is the content of the
session-10 reduction paragraph (Lemma~\ref{lem:srecharge} plus the
occupation-time step of \S\ref{sec:s13occ}).  The block-bound event
$H_{9d}$ of Corollary~\ref{cor:horizon} is also handled in $L^2$
without any state-by-state argument: with $H=\{B_r<K,\ r\le m\}$
for the chain from $x$, writing $\Phi_m=\Phi_m^H+\Phi_m^{H^c}$ with
$|\Phi_m^{H^c}(x)|\le\Pr_x[H^c]$, one has
$\E_{\bar\nu_t}[\Phi_m^2]\le2\E_{\bar\nu_t}[(\Phi_m^H)^2]+2\E_{\bar\nu_t}[\Pr_x[H^c]]$,
and $\sum_t\nu_t(S'_K)\,\E_{\bar\nu_t}[\Pr_x[H^c]]\le\Pr_{\mu_j}[H^c_{9d}]\le(9d+1)^{-\gamma'}$
because $\nu$ is a sub-law of the path measure from $\mu_j$, whose
clock reaches at most $j+d\le9d$.

\begin{remark}[what the dust obstruction does and does not touch]
\label{rem:dustblind}
(a) Proposition~\ref{prop:nomino} and Corollary~\ref{cor:noerg} are
statements about the supremum norm on $\Omega_K$: they show that no
Doeblin-type argument can give a bound on $\sup_{x\in S'_K}|\Phi_m(x)|$
(or on $\|\kappa_1^kF-\bar\pi F\|_\infty$ for any weight).  They say
nothing against (S$_2$), which is an average over laws $\bar\nu_t$
that are themselves marginals of the chain run from one block.
(b) Under the coupling of Lemma~\ref{lem:inert}, $\Phi_m$ is
Lipschitz in the dust: $|\Phi_m(y_\varepsilon)-\Phi_m(y)|\le4\varepsilon m$,
since the flip counts agree on the coupling event.  So $\Phi_m$ is a
continuous function in the coarse topology at scale $1/m$ --- the
topology in which $y_\varepsilon\to y$ --- and the dust of
Proposition~\ref{prop:nomino}, which lives at scales
$\varepsilon_i\ll1/T_i$, is invisible to the functions that (S$_2$)
is about.  This is the precise sense in which \S\ref{sec:s12left}(i)
asked for a ``dust-blind'' statement, and (S$_2$) is one.
(c) (S$_2$) is \emph{not} an annealed statement in the naive sense:
it bounds the second moment of the quenched parity, so a proof must
still show that $\Phi_m(x)$ is small for $\bar\nu_t$-typical $x$, with
the environment of $x$ frozen; only the \emph{uniformity} over $x$
has been removed.  The quantity $\E_{\bar\nu_t}[\Phi_m]$ (parity
averaged over the environment of the start) is not what is needed and
is not sufficient.
(d) The natural function space for (S$_2$) is $L^2(\bar\pi)$, in
which the chain is self-adjoint (\S\ref{sec:spectral}); the spectral
route of session 8, shelved in favour of the Markov-renewal route,
is the dust-blind route.  What it additionally needs is a comparison
of the entrance laws $\bar\nu_t$ with $\bar\pi$ on the functions
$\Phi_m^2$; this is the first sub-problem of \S\ref{sec:s13spectral}.
\end{remark}

\subsubsection{The exponent budget, recomputed}\label{sec:s13budget}

\begin{corollary}[the budget is $\gamma>1/11$, not $\gamma>4/9$]
\label{cor:budget}
Suppose (S$_2$) holds with exponent $\gamma\in(0,\tfrac12)$ and the
entrance estimate holds as stated after Proposition~\ref{prop:S2red}.
Then $\psi(j,d)\le C'(\log d)^{O(1)}\,d^{-2\gamma}$ for all $j\le8d$,
i.e.\ Problem~\ref{prob:mixing} holds with $c=\gamma$, and
Corollary~\ref{cor:psienough} closes the iid model as soon as
$\gamma>1/11$.  The same holds if (S$_2$) is replaced by the
state-uniform (S) with exponent $\gamma$.  In particular, in
Proposition~\ref{prop:holder} the crude route needs only $p>6/5$,
and (R1) needs only $C^{1,\gamma}$ regularity with $\gamma>1/11$,
i.e.\ moments $\E_x[\tau^{1+\gamma}]$ with $\gamma>1/11$.
\end{corollary}

\begin{proof}
Set $\varepsilon:=d^{-1/2-\gamma}$ and $T^*:=C\varepsilon^{-1}K\log d$
with $K=K_{\gamma'}(9d)$, $\gamma'\ge2\gamma$; then
$T^*=O(d^{1/2+\gamma}\log^2d)<d/2$ for $d$ large because
$\gamma<\tfrac12$, and $\varepsilon\le1/(2K)$.  By
Proposition~\ref{prop:shorthorizon} ($j\le8d$),
$\Pr_{\mu_j}[g_0<\varepsilon]\le4\varepsilon^2j\le32\varepsilon^2d=32d^{-2\gamma}$.
By Proposition~\ref{prop:S2red} and (S$_2$) (with $m=d-t\ge d/2$),
$\E_{\mu_j}[\Phi_d^2;g_0\ge\varepsilon]\le2C(K)(d/2)^{-2\gamma}
+2\bigl(O(d^{-\gamma'})+O(\varepsilon K\log d)\bigr)$, and
$\varepsilon K\log d=O(d^{-1/2-\gamma}\log^2d)=o(d^{-2\gamma})$ since
$\gamma<\tfrac12$.  Add, and use $|\Phi_d|\le1$ on $\{g_0<\varepsilon\}$.
The budget: Corollary~\ref{cor:psienough} needs $c>1/11$.
For Proposition~\ref{prop:holder}: $(p-1)/(p+1)>1/11\iff p>6/5$.
\end{proof}

\emph{Status: PROVED (arithmetic).  Where the factor was lost.  The
session-10 reduction paragraph ends by bounding $|\Phi_d(\Gamma_0)|$,
not $\Phi_d(\Gamma_0)^2$, by $C(\varepsilon d)^{-2\alpha}$ --- the
form of Problem~\ref{prob:survival}, which is a bound on the
\emph{second moment} --- and then replaces $(\varepsilon d)^{-2\alpha}$
by $d^{-2\alpha}$.  The first step costs a factor $2$ in the exponent,
the second a further factor $2(1+\alpha)=22/9$; together they turn
the true requirement $\gamma>1/11$ into $\gamma>2\alpha>4/9$.
Problem~\ref{prob:survival} and Corollary~\ref{cor:kappafree} are
correct as stated and, with the correct accounting, lossless: from
$\Phi_d^2\le C(K)d^{-2\gamma}$ after the entrance one gets
Problem~\ref{prob:survival} with
$(\varepsilon d)^{-2\alpha}=d^{-\alpha/(1+\alpha)}$ at the optimising
$\varepsilon$ for $\alpha=2\gamma/(1-2\gamma)$, and then
$c'=\alpha/(2(1+\alpha))=\gamma$ exactly.  The loss was entirely in
the two steps just named.  Every downstream statement that quotes
``$\gamma>4/9$'' (\S\ref{sec:markovrenewal} (R1) discussion,
\S\ref{sec:s11left}, Proposition~\ref{prop:holder} and
\S\ref{sec:s12left}, the handoff documents of sessions 10--12) should
be read with $1/11$ in its place.  Consequences: the ``marginal''
verdict on the crude route of Proposition~\ref{prop:holder}
($p>13/5$ against a measured $2.7\pm0.3$) is withdrawn --- the crude
route needs $p>1.2$; (R1) needs barely more than a first moment of
$\tau$ along the chain (the measured tail is $t^{-2}$); and
Proposition~\ref{prop:lower}'s sharp rate $c=1/2$ exceeds the budget
by a factor $5.5$, as \S\ref{sec:factor} already noted for
$\bar\psi$.  The programme has a wide margin everywhere except in
the one place where it has no proof.}

The budget chain, for reference (each line is proved above it):
\[
\begin{array}{lll}
\toprule
\text{statement} & \text{exponent} & \text{needs}\\
\midrule
\E|\mathrm{bias}|=o(1/\log n),\ k\ge6\log^{11}n & k^{-c} & c>1/11\\
\psi(j,d)\le Kd^{-2c},\ j\le8d\ \text{(Thm.~\ref{thm:factor})} & c & c>1/11\\
\text{(S)}/\text{(S}_2)\ \text{on }S'_K\ \text{(Cor.~\ref{cor:budget})} & \gamma=c & \gamma>1/11\\
G\in C^{0,\gamma}\ \text{at }z=-1\ \text{(Prop.~\ref{prop:condecay})} & \gamma & \gamma>1/11\\
\text{(C2$''$) via Prop.~\ref{prop:holder}, crude} & (p-1)/(p+1) & p>6/5\\
\text{(C2$''$) via Prop.~\ref{prop:holder}, factorised} & \min(1,p-1) & p>12/11\\
\text{measured (\S\ref{sec:s12num})} & p=2.7\pm0.3,\ \text{modulus }0.75 & \\
\bottomrule
\end{array}
\]

\subsubsection{Flips recur: \texorpdfstring{$\tau_k<\infty$}{tau\_k<infty} a.s.}\label{sec:s13tau}

\begin{lemma}[flips recur]\label{lem:flipsrecur}
From every state of the quotient chain with $a_0b_0>0$,
$\sum_t2a_tb_t=\infty$ a.s.; consequently flips occur infinitely
often a.s., and $\tau_k<\infty$ a.s.\ for every $k$.
\end{lemma}

\begin{proof}
Write $m_t=\min(a_t,b_t)$, $M_t=\max(a_t,b_t)$ and
$p_t=2a_tb_t=2m_tM_t$, the conditional probability of a flip at step
$t+1$.  By Lemma~\ref{lem:splitmerge} (read on the quotient,
Corollary~\ref{cor:quotient}), at a non-flip step a marked coordinate
either stays fixed, increases (a merge), or is multiplied by
$\sqrt U$, $U\sim U(0,1)$ (a shrink: both cuts in that coordinate's
material, probability equal to the square of the coordinate).  Hence
$m_t$ is non-decreasing between the steps at which it strictly
decreases, and at a non-flip step $m$ can decrease in exactly two
ways: a shrink of the smaller coordinate, probability $m_t^2$; or a
shrink of the larger coordinate below $m_t$, probability
$M_t^2\cdot\Pr[M_t\sqrt U<m_t]=M_t^2(m_t/M_t)^2=m_t^2$.  So the event
$E_{t+1}=\{\text{step }t{+}1\text{ is not a flip and }m_{t+1}<m_t\}$
has $\Pr[E_{t+1}\mid\mathcal F_t]\le2m_t^2\le2m_tM_t=p_t$.

Let $\Omega_0=\{\sum_tp_t<\infty\}$.  By L\'evy's conditional
Borel--Cantelli lemma, on $\Omega_0$ almost surely only finitely many
flips occur and only finitely many $E_t$ occur.  All marked
coordinates stay strictly positive a.s.\ ($U>0$, flips land in the
open interval), so after the last flip and the last $E_t$ the process
$m_t$ is non-decreasing and bounded below by some $m_*>0$; then
$p_t\ge2m_*^2$ for all large $t$ and $\sum_tp_t=\infty$,
contradicting $\Omega_0$.  Hence $\Pr[\Omega_0]=0$, i.e.\
$\sum_tp_t=\infty$ a.s., and L\'evy's lemma in the other direction
gives flips infinitely often a.s.
\end{proof}

\emph{Status: PROVED.  This is the hypothesis of
Propositions~\ref{prop:r2false}, \ref{prop:nomino}, \ref{prop:nodrift}
and of the remark after Proposition~\ref{prop:genfun}; it needs
nothing about the environment --- in particular it holds with an
all-dust environment --- because the only mechanism that decreases the
smaller marked coordinate has probability at most the flip
probability.  It gives no rate; rates are (R1).}

\subsubsection{The occupation-time step, stated}\label{sec:s13occ}

The session-10 reduction (end of \S\ref{sec:markovrenewal}) and the
entrance estimate after Proposition~\ref{prop:S2red} use, without
proof, the following statement.

\begin{lemma}[occupation time of the marked total; OPEN]\label{lem:occupation}
Let $K\ge3$ and let $H=H_T$ be the block-bound event with bound $K$.
There are $c_K,C_K>0$ such that for every start with
$s_0\ge1/(2K)$, every $\varepsilon\le1/(8K)$ and every $T\ge1$,
\[
\Pr\Bigl[\#\{t\le T:\ s_t<\tfrac1{4K}\}>\tfrac T2\Bigr]
\;\le\;C_K\,e^{-c_K\varepsilon T/K}\;+\;2\varepsilon^2T\;+\;\Pr[H^c].
\]
\end{lemma}

\emph{Status: OPEN; believed standard.  What it is used for: in the
entrance estimate, after $s$ first reaches $1/(2K)$ (within
$T_1=O(\varepsilon^{-1}K\log d)$ steps by Lemma~\ref{lem:srecharge}),
one needs, on at least half of the next $T_2=O(\varepsilon^{-1}K\log d)$
steps, $s_t\ge1/(4K)$ and $\min(a_t,b_t)\ge\varepsilon$ (the second
from the flux bound, cost $4\varepsilon^2T_2$); on those steps a flip
has probability $\ge2\varepsilon(s_t-\varepsilon)\ge\varepsilon/(4K)$
and lands in $S'_K$ with conditional probability $\ge\tfrac12$
(the trapezoid of Lemma~\ref{lem:recharge}(ii) is symmetric and
unimodal), so $T'\le T_1+T_2$ except with probability
$O(d^{-\gamma'})+O(\varepsilon K\log d)$.  What fails if it is false:
the entrance estimate, hence the reduction of (S$_2$) to the
log-macroscopic set $S'_K$; (S$_2$) would then have to be proved for
the laws $\mu_j$ restricted to $\{g_0\ge\varepsilon\}$ with
$\varepsilon=d^{-1/2-\gamma}$, i.e.\ with the trap accounting inside
the $L^2$ statement.  Nothing above Problem~\ref{prob:survival}
depends on it.  Intended proof: Lemma~\ref{lem:logdrift} gives
$\E[\Delta\log s\mid\Gamma_t]\ge c\,s_t/K$ on $H\cap\{s_t\le1/(2K)\}$,
with down-jumps those of $\tfrac12\log U$; so for small $\theta>0$,
$V=s^{-\theta}$ is a supermartingale on $\{s\le1/(2K)\}\cap H$ with
drift $\le-\theta c's^{1-\theta}/K$, and the occupation time of
$\{\varepsilon\le s<1/(4K)\}$ over $T$ steps is at most
$\varepsilon^{-(1-\theta)}$ times the accumulated drift; the time
below $\varepsilon$ is priced by Lemma~\ref{lem:strap}; excursions
above $1/(2K)$ are handled by stopping $V$ at $1/(2K)$.  The missing
piece is only the bookkeeping of the re-entries into
$\{s<1/(2K)\}$ from above, each of which starts $V$ at
$\le(2K)^\theta\cdot2^\theta$ (a single shrink from $\ge1/(2K)$
lands at $\ge s\sqrt U$, and $\Pr[\sqrt U\le\tfrac12]=\tfrac14$) and
of which there are at most $O(T/K^2)$ in expectation (shrinks have
probability $\le s^2$).  We have not written it out.}

\subsubsection{What (M\texorpdfstring{$_1$}{1}) needs}\label{sec:s13m1}

Proposition~\ref{prop:holder} assumes (M$_1$):
$\sup_{\bar x\in S'_K}\E_{\bar x}[\tau_k]\le C_1k$.  By the strong
Markov property at the flip times,
$\E_{\bar x}[\tau_k]=\sum_{i<k}\E_{\bar x}\bigl[\E_{\bar\Gamma_{\tau_i}}[\tau]\bigr]$,
so (M$_1$) follows from a bound on the mean first-flip time that is
uniform \emph{on average over post-flip states}.  The post-flip state
has marked total $s$ (preserved by the flip) and landing gap $g'$
with $\Pr[g'\le\eta]\le\eta^2/(ab)$ (Lemma~\ref{lem:recharge}(ii)),
so $\E[1/g']\le\int_0^{s}\eta^{-2}\Pr[g'\le\eta]\,d\eta\cdot
\text{(const)}<\infty$ is automatic from a macroscopic pre-flip
state; what is needed is therefore
\[
\text{(R1$'$)}\qquad
\E_{\bar x}[\tau]\;\le\;C(K)\,\bigl(1+1/g(\bar x)\bigr)\quad
\text{on }H,\ \text{uniformly over the environment},
\]
together with control of the pre-flip state at each flip (so that
$1/(ab)$ in the landing bound is $O(K^2)$ on average), which is the
recurrence of the marked pair to $S'_K$ --- the same recharge
machinery as the entrance estimate.  (R1$'$) is the first-moment
part of TODO~11a (which asks for the tail $\Pr_x[\tau>t]\le C_Kt^{-2}$,
measured in \S\ref{sec:parity}(c)); by Corollary~\ref{cor:budget},
(R1) itself now needs only a $(1+\gamma)$-th moment with
$\gamma>1/11$, so a tail bound $t^{-1-\delta}$ for any $\delta>1/11$
would suffice for both.  (M$_p$) in Corollary~\ref{cor:noerg}(ii) is
used only in a negative result and needs no further attention.

\subsubsection{Where this leaves the programme}\label{sec:s13spectral}

Items (I)--(III) of the overview are unchanged in substance, but (I)
is now: prove (S$_2$) with any $\gamma>1/11$, for the time-$t$
entrance laws $\bar\nu_t$ of $S'_K$, with $C(K)$ polynomial.  Two routes remain, neither of
which is Doeblin-type:
\begin{itemize}[leftmargin=2em]
\item[(i)] \emph{The spectral route} (\S\ref{sec:spectral}): the
quotient chain is reversible for $\bar\pi$; for the odd function
$f=\varepsilon F$ the Dirichlet form has the massive flip term, and
polynomial decay of $\psi_{\bar\pi}(d)=\|P^d\varepsilon\|^2_{L^2(\pi)}$
with exponent $\beta$ is equivalent to the spectral measure of
$\varepsilon$ having mass $O(\delta^\beta)$ within $\delta$ of
$\pm1$.  A sufficient functional inequality is
\[
\text{(H$_\beta$)}\qquad
\bigl|\E_{\bar\pi}[F]\bigr|^2\;\le\;C\,\mathcal E_{\mathrm{odd}}(F)^{\beta}\,\|F\|_{L^2(\bar\pi)}^{2(1-\beta)}
\qquad\text{for all }F\in L^2(\bar\pi),
\]
with $\mathcal E_{\mathrm{odd}}(F)=\tfrac12\E_{\bar\pi}[(F-F')^2;\text{non-flip}]
+\tfrac12\E_{\bar\pi}[(F+F')^2;\text{flip}]$: applied to
$F=\Pi_\delta\varepsilon$ (spectral projection onto
$[1-\delta,1]$) it gives $\nu_\varepsilon([1-\delta,1])\le C\delta^\beta$.
By Proposition~\ref{prop:lower}, $\beta\le1$; the budget needs only
$\beta>2/11$.  (The limiting case $\beta=1$ is a Hardy inequality
in the gap coordinate and is false by a logarithm: test functions
$F=g^{-1}\log^{-2}(1/g)\mathbf 1\{g\le\delta\}$ have
$|\E F|^2/\mathcal E_{\mathrm{odd}}(F)\asymp\log(1/\delta)$ --- consistent
with $\psi_{\bar\pi}\asymp1/d$, since linear spectral density gives
$\int d\nu_\varepsilon/(1-\lambda)=\infty$.)  The periodicity
caveat of \S\ref{sec:spectral} applies: the same inequality is
needed for $(-1)^BF$, i.e.\ near $\lambda=-1$.  What the spectral
route does \emph{not} give by itself is (S$_2$) for the entrance
laws $\bar\nu_t$: these are singular with respect to $\bar\pi$
(finitely many blocks against infinitely many), and the comparison
must be made at a coarse scale --- the blocks of mass $\gtrsim1/(m\log m)$,
which are the only ones $\Phi_m$ sees (Remark~\ref{rem:dustblind}(b))
--- where $\bar\nu_t$ has had time of order $m\log m\gg$ the
relaxation time $1/(2\eta)$ of each scale $\eta$ in that range.
This start-comparison is the first concrete sub-problem.
\item[(ii)] \emph{The exact-law route} of \S\ref{sec:s12left}(iv):
polynomial decay of the block-count parity at flip times from the
exact fixed-time laws, from the one-curve start, extended to the
$\mu_j$ starts.  Its target is now $p>6/5$ in
Proposition~\ref{prop:holder}, i.e.\ $a_k=O(k^{-6/5})$, far weaker
than the measured $k^{-3}$.  Since $B(\bar\Gamma_{\tau_k})$ spreads
only like $\log k$, the decay cannot come from the width of its law;
it comes from the smoothness of the law of $\tau_k$ on the time axis
(the parity of $B$ at a \emph{fixed} time is deterministic,
Lemma~\ref{lem:parity}), and quantifying that smoothness is exactly
what (ii) must do.
\end{itemize}
Both routes consume the same two inputs that the audit has left open
--- the occupation-time lemma and (R1$'$) --- and both are now aimed at
a target with a fivefold margin.



%======================================================================
\section{Open problems}\label{sec:open}

\begin{problem}
Prove Hypothesis~\ref{hyp:eq} in any nontrivial regime, e.g.\
$\delta(n)=o(1)$ ($\delta(n)=o(1/\log n)$ gives
Conjecture~\ref{conj:cgw} by Theorem~\ref{thm:reduction}).  By
Theorem~\ref{thm:marked} this is exactly the statement that a marked
pair $\{j,\sigma^\alpha(j)\}$ of a uniform $L\in S(\lambda)$ is a
$B$-pair with probability $\tfrac12+O(\delta)$; the programme at the
end of \S\ref{sec:kps}, as sharpened in \S\ref{sec:orbit}, reduces it
further to the annealed bias bound of Problem~\ref{prob:annealed}
--- a genericity statement for the rank event of
Theorem~\ref{thm:master} under the orbit measure, for which generic
position gives $\E|\mathrm{bias}|\asymp k^{-1/2}$ against a budget of
$k^{-1/11}$.  After \S\ref{sec:anneal} the outstanding pieces are:
the $L^2$ mixing bound of Problem~\ref{prob:mixing} (which closes the
iid model via the factorization Theorem~\ref{thm:factor}; after
\S\ref{sec:mixing} this is reduced to the parity-decay statement
Problem~\ref{prob:survival} alone with any exponent $\alpha>2/9$
--- the gap law Problem~\ref{prob:gaplaw} is not needed over the
short horizon $j\le8d$ of Theorem~\ref{thm:factor}
(Corollary~\ref{cor:kappafree}) --- and the rate $1/d$ is proved
sharp; after \S\ref{sec:s10} the environment is controlled over the
horizon by the Harer--Zagier block count, Corollary~\ref{cor:horizon},
and Problem~\ref{prob:survival} is reduced to the conditions
(R1)--(R2) of Proposition~\ref{prop:condecay}), the
transfer from iid to orbit positions, and the supply--mark decoupling
of Problem~\ref{prob:decouple}.
\end{problem}

\begin{problem}
Determine the layer profile $\gamma$ of \S\ref{sec:layers} and the constant
$c_2=\int_0^1\gamma$ ($\gamma(0)=1$ is Corollary~\ref{cor:thirdrow};
Theorems~\ref{thm:layer4} and~\ref{thm:layer5} pin the fourth and fifth
layers, and Remark~\ref{rem:layerk} every fixed layer, so $\gamma$ is
flat at the origin; see Remark~\ref{rem:c2caveat} for
why the $n\le11$ value $1.30$ of $n^3h_n(2)$ does not yet determine the
limit).  Even a heuristic derivation via permanents of complements of
random rectangles with a planted rows-$(1,2)$ intercalate would be
valuable; the remaining unknown is the bulk $k\asymp tn$, $t>0$.
\end{problem}

\begin{problem}
Prove the third-order intercalate expectation
$\E[\mathbf N]=n^2/4-n/4+\Theta(1/n)$ (\S\ref{sec:consequences}) --- a
statement plausibly within reach of the switching methods of \cite{McKW99}.
\end{problem}

\begin{problem}
Prove the sharp parity asymptotics
$\Pr(\sigma_{1,2}\text{ even})-\tfrac12\sim(-1)^{n-1}(n-1)/(2D_n)$, in the
spirit of \cite{KPS25}.
\end{problem}

\begin{problem}
Give a corrected exact joining bookkeeping for $\alpha=\beta\ge3$
(Remark~\ref{rem:albe}), restoring the identity for all adjacent pairs.
\end{problem}

\begin{problem}
Extend the switch calculus of \S\ref{sec:switch} to the fifth layer
(Theorem~\ref{thm:layer4} settles the fourth), where the
analogue of Proposition~\ref{prop:pairIE} is an inclusion--exclusion over
the three pairwise coincidence patterns of a triple $(\tau,\eta,\theta)$.
Beyond that lies the bulk of the profile $\gamma$.  (Classical context:
Riordan's aggregate $K(3,n)=\sum_k\binom nkD_{n-k}D_kU_{n-2k}$
\cite[Ch.~8, eq.~(30a)]{Riordan58}; Moser's very reduced $4\times n$
rectangles \cite{Moser67}; the $k$-discordant literature
\cite{Whitehead79}.  Theorem~\ref{thm:poisson} resolves the
$2$-cycle-level joint statistics for three rows; the full
cycle-type-resolved joint law remains open.)
\end{problem}

%======================================================================
\begin{thebibliography}{99}

\bibitem{CamBlog15} P.~J. Cameron, \emph{A niggling problem}, blog post,
24~January 2015,
\url{https://cameroncounts.wordpress.com/2015/01/24/a-niggling-problem/};
see also \emph{Some challenges on Latin squares}, 6~October 2023.

\bibitem{Cam92} P.~J. Cameron, Almost all quasigroups have rank 2,
\emph{Discrete Math.} \textbf{106/107} (1992), 111--115.

\bibitem{CGW08} N.~J. Cavenagh, C.~Greenhill and I.~M. Wanless, The cycle
structure of two rows in a random Latin square, \emph{Random Structures
Algorithms} \textbf{33} (2008), 286--309.

\bibitem{CL72} M.~Cohn and A.~Lempel, Cycle decomposition by disjoint
transpositions, \emph{J.\ Combin.\ Theory Ser.\ A} \textbf{13} (1972),
83--89.

\bibitem{FS90} P.~Flajolet and M.~Soria, Gaussian limiting distributions for
the number of components in combinatorial structures, \emph{J. Combin.
Theory Ser.~A} \textbf{53} (1990), 165--182.

\bibitem{HJ96} R.~H\"aggkvist and J.~C.~M. Janssen, All-even Latin squares,
\emph{Discrete Math.} \textbf{157} (1996), 199--206.

\bibitem{HZ86} J.~Harer and D.~Zagier, The Euler characteristic of the
moduli space of curves, \emph{Invent.\ Math.} \textbf{85} (1986),
457--485.

\bibitem{JM96} M.~T. Jacobson and P.~Matthews, Generating uniformly
distributed random Latin squares, \emph{J. Combin. Des.} \textbf{4} (1996),
405--437.

\bibitem{KPS25} M.~Kwan, K.~Petrova and M.~Sawhney, Parities in random Latin
squares, arXiv:2509.13125 (2025).

\bibitem{KSSS22} M.~Kwan, A.~Sah, M.~Sawhney and M.~Simkin, Substructures in
Latin squares, \emph{Israel J. Math.} (2023); arXiv:2202.05088.  See also
\emph{Large deviations in random Latin squares}, \emph{Bull. Lond. Math.
Soc.} \textbf{54} (2022).

\bibitem{Sch05} O.~Schramm, Compositions of random transpositions,
\emph{Israel J. Math.} \textbf{147} (2005), 221--243.

\bibitem{LP93} T.~{\L}uczak and L.~Pyber, On random generation of the
symmetric group, \emph{Combin. Probab. Comput.} \textbf{2} (1993), 505--512.

\bibitem{McKW99} B.~D. McKay and I.~M. Wanless, Most Latin squares have many
subsquares, \emph{J. Combin. Theory Ser.~A} \textbf{86} (1999), 323--347.

\bibitem{McKW05} B.~D. McKay and I.~M. Wanless, On the number of Latin
squares, \emph{Ann. Comb.} \textbf{9} (2005), 335--344.

\bibitem{Moser67} W.~O.~J. Moser, The number of very reduced $4\times n$
Latin rectangles, \emph{Canad. J. Math.} \textbf{19} (1967), 1011--1017.

\bibitem{Riordan54} J.~Riordan, Discordant permutations,
\emph{Scripta Math.} \textbf{20} (1954), 14--23.

\bibitem{Riordan58} J.~Riordan, \emph{An Introduction to Combinatorial
Analysis}, Wiley, New York, 1958; Chapter 8 (\emph{Permutations with
restricted position II}), esp.\ \S3, pp.~201--204.

\bibitem{Touchard53} J.~Touchard, Permutations discordant with two given
permutations, \emph{Scripta Math.} \textbf{19} (1953), 109--119.

\bibitem{Whitehead79} E.~G. Whitehead, Jr., Four discordant permutations,
\emph{J. Austral. Math. Soc. Ser.~A} \textbf{28} (1979), 369--377.

\end{thebibliography}

\appendix

\section{Data and computation manifest}\label{app:data}

All computations were performed in this working session; code and exact data
are available alongside these notes.

\begin{enumerate}[leftmargin=2em]
\item \texttt{cgw\_data.py}: the exact $\CC_n(\lambda)$, $4\le n\le11$, from
\cite[Tables 1--2]{CGW08}, verified against the published
$\PP_n(\lambda)$ fractions and against $L_n$ for all $n\le11$.
\item \texttt{fit\_additive.py}: additive fits of
$\log \CC_n(\lambda)$, chain increments, $\dtv$, spreads.
\item \texttt{third\_row.py}, \texttt{third\_row\_asymp.py}: exact
$N^{(3)}$ via matching polynomials; fusion identity verified for all types,
$n\le25$; exact ratios to $n=400$ with Richardson extrapolation.
\item \texttt{layer4.c}: exact $N^{(4)}$ for $n\le11$ by summation of
Ryser--Gray permanents over all third rows ($\approx5\times10^6$ permanents
of order $11$).
\item \texttt{verify\_identity.py}: complete enumeration of Latin squares of
orders $5$ and $6$ (via first-row normalisation) implementing the CGW
switch/flip operations; exact rational match of
Theorem~\ref{thm:identity} in the cases reported, and the exact values
$\Pr(B_1)=\Pr(B_2)$.
\item \texttt{jm.c}, \texttt{jm\_stats.py}: Jacobson--Matthews sampler
($O(1)$ per move) with fixed-grid proper-state sampling, and the
measurement pipeline.
\item \texttt{puncture.py}: the punctured-board calculus
(\eqref{eq:punctured}), verified against exhaustive enumeration at
$n=6,7$.
\item \texttt{pair\_ie.py}, \texttt{layer4\_production.py},
\texttt{layer4\_big.txt}: the pair inclusion--exclusion identity
(verified exactly at $n=6,7$), exact $T_j$ sums for $j\le3$, truncated
layer-4 effects to $n=18$ with validation against the exact
$n\le11$ values.
\item \texttt{three\_row\_moments.py}: exact joint intercalate moments of
$3\times n$ rectangles (certified against brute force at $n=7$),
$n\le13$ tables, the tilt identity and covariance data.
\item \texttt{switch.py}: the switch calculus of \S\ref{sec:switch}:
board $B_0$, punctured counts on arbitrary path/cycle boards, the empty
bracket $=M_{n-4}$ for $n\le14$, brute-force verification of
Lemma~\ref{lem:switchIE} at $n=6,7$, and per-$M$ bracket ratio scans.
\item \texttt{fusion\_lemmas.py}: symbolic proof-checks of
Lemma~\ref{lem:wronskian} and exact integer verification of the master
formula \eqref{eq:generic} on random punctures at $n=9,11,13$.
\item \texttt{valuation\_scan.py}: exhaustive verification of
Lemma~\ref{lem:valuation} (valuation $\ge2$; the $[x^2]$
classification) over all $|M|\le2$ at $n=9,10$.
\item \texttt{switch\_j1.py}: stratified sums $S_j$ \eqref{eq:strata}
against direct enumeration ($n=8,9$), class-resolved contributions,
and the worry-class (rows $1\wedge2$) resummation scans.
\item \texttt{loop\_reduction\_check.py}: complete-enumeration proof of
Remark~\ref{rem:reduced} at $n=5,6$ (the loop second-row type law is
exactly $\PP_n$).
\item \texttt{anatomy.c}: exact $\mathrm{Br}_G$, $N^{(4)}_0$, masses
$m^\pm$ and conditional means $E^\pm[A]$ by full enumeration of $A_0$
with per-$\tau$ Ryser permanents, $n\le11$.
\item \texttt{triple\_ie.py}: the triple inclusion--exclusion identity
(Proposition~\ref{prop:tripleIE}), both forms, verified against
brute-force triple enumeration at $n=5,6,7$ (both cycle types).
\item \texttt{n5\_direct.c}, \texttt{n5\_direct2.c}: exact $N^{(5)}$ by
direct enumeration with 64-/128-bit cell masks, $n\le9$, both types;
$N^{(3)},N^{(4)}$ reproduced as certification.
\item \texttt{switch5.py}, \texttt{anatomy5.c}: the fifth-layer switch
decomposition (Lemma~\ref{lem:switchIE5}) verified by complete
enumeration at $n=6,7$ and against the direct $\Delta N^{(5)}$ at
$n=8$; exact $\mathrm{Br}^{(5)}_G$, $N^{(5)}_0$, $G_2^{(5)}$-,
$H^{(5)}$-pieces on $A_0$.
\item \texttt{layer5\_production.py}: the truncated triple-IE series
$S(Q)$, $Q\le3$, both types, $n\le11$+, validated against the exact
$N^{(5)}$ anchors.
\item \texttt{pf\_orbit.py}: the CGW cross-switch
(\cite[Def.~3.9]{CGW08}) implemented and verified exhaustively at
$n=5$ (Latin, type-preserving, involutive, marked status preserved,
shifted-pair toggle profile), and the splitting-instance statistics
$\Pr(B_i\mid\text{marked }A/B)$.
\item \texttt{pf\_bridge.py}: instance-by-instance verification of the
bijections behind Theorem~\ref{thm:marked} at $n=5,6$ (unflip
statuses, roundtrips, hit multiplicities, orphan structure) for
$(2,3)\to(5)$, $(2,4)\to(6)$, $(2,2,2)\to(2,4)$, $(3,3)\to(6)$.
\item \texttt{xside.c}, \texttt{xside2.c}: the $X$-side counts
$|X_A|,|X_B|$ for all types and merges --- by complete enumeration of
first-row-normalised squares ($n\le6$), and per fixed representative
$\sigma_0$ of each type ($n=7$; conjugation symmetry); exact match
with $2\CC_n(\lambda)/\CC_n(\mu)-1$ in all cases.
\item \texttt{pf\_cl.py}: instance-by-instance verification of
Theorem~\ref{thm:CL}, Theorem~\ref{thm:bitrank}, the passenger
construction, and randomized two-sided products with shared rows and
conjugating coordinates (all activation subsets, exhaustively).
\item \texttt{pf\_master.py}: end-to-end verification of the master
formula on real squares --- $n=5$ complete ($7920$ subset checks),
$n=6$ sampled ($1{,}706{,}949$ checks, all four column-incidence
classes): $0$ failures.
\item \texttt{pf\_bias.py}, \texttt{bias\_scale.c}: the exactly
solvable bias families of \S\ref{sec:orbit} and the $k$-scaling of
the generic-position bias (rank route).
\item \texttt{s7\_bias\_lab.py}: pendant-star family
(Proposition~\ref{prop:pendant}), exhaustive predicate scan over all
chord configurations $k\le6$, per-coordinate toggle statistics.
\item \texttt{s7\_twopoint.c}: cycle-route measurements of
$4\E[\mathrm{bias}^2]$, the one-point function, and $\psi(j,d)$
(Theorem~\ref{thm:factor}); independent of the rank route.
\item \texttt{s7\_real\_bias.py} (with \texttt{jm.c}): the
real-ensemble study of \S\ref{sec:realbias} at $n=30$--$100$,
including ground-truth re-verification of the master formula.
\item \texttt{s8\_splitmerge.c}: the continuum marked split--merge
chain of \S\ref{sec:mixing} (modes: \texttt{psi}/\texttt{psibar}
gap-binned $\psi$, \texttt{one}, \texttt{iota} symmetry test,
\texttt{gapdist} stationary gap law, \texttt{traj} invariants);
validated against \texttt{s7\_twopoint.c} at matched $(j,d)$ and
against the $\mathrm{PD}(1)$ invariants.
\item \texttt{s8\_real\_bias.py}: \texttt{s7\_real\_bias.py} plus the
$\sigma$-cycle length (class size) column and $k=0$ rows, for the
supply--mark decoupling measurements; \texttt{s8\_inst200.csv} extends
the real-ensemble study to $n=200$.
\end{enumerate}

\section{Exact \texorpdfstring{$\bar q$}{qbar} values from the identity}\label{app:qbar}

All $\alpha\neq\beta$ joins at $n=11$, with
$\bar q=2\CC_{11}(\lambda)/\CC_{11}(\mu)-1$ (exact values, displayed to six
decimals) and the normalised deficit $\tfrac12(1-\bar q)n^3$:

\begin{center}
\small
\begin{tabular}{llll r}
\toprule
$\mu$ & join & $\lambda$ & $\bar q$ & $\tfrac12(1-\bar q)n^3$\\
\midrule
$(2,9)$ & $2{+}9$ & $(11)$ & $0.998080$ & $+1.278$\\
$(2,2,7)$ & $2{+}7$ & $(2,9)$ & $0.998078$ & $+1.279$\\
$(2,3,6)$ & $2{+}6$ & $(3,8)$ & $0.998084$ & $+1.275$\\
$(2,4,5)$ & $2{+}5$ & $(4,7)$ & $0.998084$ & $+1.275$\\
$(2,2,2,5)$ & $2{+}5$ & $(2,2,7)$ & $0.998080$ & $+1.278$\\
$(2,4,5)$ & $2{+}4$ & $(5,6)$ & $0.998077$ & $+1.280$\\
$(2,2,3,4)$ & $2{+}4$ & $(2,3,6)$ & $0.998079$ & $+1.278$\\
$(2,3,6)$ & $2{+}3$ & $(5,6)$ & $0.998046$ & $+1.301$\\
$(2,2,3,4)$ & $2{+}3$ & $(2,4,5)$ & $0.998048$ & $+1.299$\\
$(2,3,3,3)$ & $2{+}3$ & $(3,3,5)$ & $0.998054$ & $+1.295$\\
$(2,2,2,2,3)$ & $2{+}3$ & $(2,2,2,5)$ & $0.998052$ & $+1.297$\\
\midrule
$(3,8)$ & $3{+}8$ & $(11)$ & $0.999966$ & $+0.023$\\
$(2,3,6)$ & $3{+}6$ & $(2,9)$ & $0.999970$ & $+0.020$\\
$(3,3,5)$ & $3{+}5$ & $(3,8)$ & $0.999970$ & $+0.020$\\
$(3,4,4)$ & $3{+}4$ & $(4,7)$ & $0.999963$ & $+0.024$\\
$(2,2,3,4)$ & $3{+}4$ & $(2,2,7)$ & $0.999971$ & $+0.019$\\
\midrule
$(4,7)$ & $4{+}7$ & $(11)$ & $0.999997$ & $+0.002$\\
$(2,4,5)$ & $4{+}5$ & $(2,9)$ & $1.000001$ & $-0.001$\\
$(5,6)$ & $5{+}6$ & $(11)$ & $1.000004$ & $-0.003$\\
\bottomrule
\end{tabular}
\end{center}

\noindent
The deficits stratify precisely by the smallest part destroyed by the join
--- $\approx c_2 n^{-3}$ when a $2$-cycle is destroyed, $\approx h_n(3)$
when a $3$-cycle is, and at noise level below --- and are essentially
independent of the ambient type $\mu$: the identity turns the refined law of
Conjecture~\ref{conj:refined} into an equidistribution statement, join by
join.

\end{document}
